\documentclass[10pt]{amsart}
\usepackage{amsmath,amssymb,amsthm,mathtools,mathrsfs}
\usepackage[a4paper,hmargin=0.76in,vmargin=0.60in]{geometry}
\usepackage{booktabs}
\usepackage{array}
\usepackage{longtable}
\usepackage{tabularx}
\usepackage{xr-hyper}
\usepackage[hypertexnames=false]{hyperref}
\externaldocument[A-]{column-packages-character-families}
\hypersetup{hidelinks}
\usepackage[final,expansion=false]{microtype}
\usepackage{url}
\emergencystretch=2.5em

\newtheorem{theorem}{Theorem}[section]
\newtheorem{proposition}[theorem]{Proposition}
\newtheorem{lemma}[theorem]{Lemma}
\newtheorem{corollary}[theorem]{Corollary}
\theoremstyle{definition}
\newtheorem{definition}[theorem]{Definition}
\newtheorem{problem}[theorem]{Problem}
\theoremstyle{remark}
\newtheorem{remark}[theorem]{Remark}

\newcommand{\C}{\mathbb C}
\newcommand{\R}{\mathbb R}
\newcommand{\N}{\mathbb N}
\newcommand{\Z}{\mathbb Z}
\newcommand{\Q}{\mathbb Q}
\newcommand{\Rel}{\operatorname{Re}}
\newcommand{\Iml}{\operatorname{Im}}
\newcommand{\supp}{\operatorname{supp}}
\newcommand{\TV}{\operatorname{TV}}
\newcommand{\Kp}{\mathcal K^{\mathrm{prim}}_{2T}}
\newcommand{\Kw}{\mathcal K_{2T}}
\newcommand{\PT}{\mathcal P_T}
\newcommand{\IT}{I_T}
\newcommand{\lamP}{\lambda^{\mathrm{prim}}_T}
\newcommand{\lamW}{\lambda^{\mathrm W}_T}
\newcommand{\lamA}{\lambda_{\mathrm A}(T)}
\newcommand{\sT}{\sigma_T}
\newcommand{\gh}[1]{\widehat g_{#1}}
\newcommand{\nuh}{\widehat\nu}
\newcommand{\wh}{\widehat w}
\newcommand{\Fh}{\widehat F}
\newcommand{\cF}{c_F}
\newcommand{\Lev}{\mathfrak E}
\newcommand{\Wcl}{\mathcal W_T}
\newcommand{\NT}{\mathcal N_T}
\newcommand{\Aconst}{\mathsf A_T}
\newcommand{\Ainf}{\mathcal A_\infty}
\newcommand{\PiT}{\Pi_T}
\newcommand{\mT}{\mathfrak m_T}
\newcommand{\laminf}{\lambda^{\infty}_{\min}}
\newcommand{\laminfP}{\lambda^{\infty,\mathrm P}_{\min}}
\newcolumntype{Y}{>{\raggedright\arraybackslash}X}

\begin{document}
\addtolength{\abovedisplayskip}{-3.4pt}
\addtolength{\belowdisplayskip}{-3.4pt}
\addtolength{\belowdisplayshortskip}{-2pt}
\title[The sign of the radial intersection form]{The sign of the radial intersection form on the idele class group}
\author{David Alejandro Trejo Pizzo}
\date{}
\begin{abstract}
The radial trace on the idele class group of $\mathbb Q$ defines a Hermitian
form on functions supported in $(-T,T)$. We determine which sign of this form
corresponds to the intersection inequality on the kernel of the two Tate
functionals, and we prove that the form has this sign: $W_T(F)\ge0$ on the
entire form domain for every $T>0$, and every nontrivial zero of $\zeta$ has
real part $1/2$. For the smoothed prime sum $G$ we show that, for every
$1/2<\theta<1$, a bound $G(x)\le C_\theta x^\theta$ on arbitrarily large
intervals $[Y,Y^{J_\theta}]$ with fixed $J_\theta>1$ excludes zeros in
$\Re s>\theta$, and we reduce these bounds to an upper estimate for a finite
mean difference, on the same fields and scales as three quadratic twist
estimates, with an exponent loss tending to zero. We prove this mean
estimate from a slope inequality for Hermitian lattices, given the existence
of column packages with uniform constants. The existence of the column
packages is due to L.~Alpoge and rests on his Lean formalisation; it has not
yet been written out on paper in full. A companion paper gives a first
written account of it: it proves the reductions, identifies which steps are
known theorems re-proved in the formalisation, and states exactly the new
statements taken from it. We also show that the sign on all windows is
equivalent to positive semidefiniteness in every degree of a single sequence
of explicit matrices, and on bounded windows we obtain quantitative
estimates: for $T\le\tfrac12\log5$ the lower bounds on unit vectors are
$10^{-18}$ for the complete form and $2.5\cdot10^{-12}$ for its primitive
restriction, unconditionally, by a computer-assisted proof in interval
arithmetic; and for $T\le2.8022$ we prove the lower bound
$e^{-120110026807}$ for the complete form, assuming the numerical
verification that the zeros of $\zeta$ up to height $3\cdot10^{12}$ lie on
the critical line and an explicit zero count, with the finite computations
done in exact integer and interval arithmetic.
\end{abstract}
\maketitle
\setcounter{tocdepth}{1}\tableofcontents
\section{Introduction}\label{sec:intro}
In Weil's argument for a curve over a finite field the two degree
directions pair hyperbolically. The index inequality concerns the
restriction of self-intersection to their common kernel. For the
radial correspondences on $\operatorname{Spec}\mathbb Z$, the
carrier is the colimit of the semilocal idele class groups over the
finite sets of places of $\mathbb Q$ containing infinity. The
involution, trace and Tate functionals are constructed in \cite{P43}.
The degree-$n$ diagonal contact contributes $\Lambda(n)/\sqrt n$
by \cite{P44}; the sum of contacts and the archimedean term has
the zero expansion with polar correction of \cite{Paper45}.

These three constructions are Rows~A, B and C, respectively.
This fourth paper treats Row~D, the index inequality for the resulting
radial intersection pairing on $\operatorname{Spec}\mathbb Z$.
The question is whether every radial class of bidegree zero has nonpositive
self-intersection, for every support radius $T>0$. From Rows~A--C
we use the pairing, its two degree maps and the trace identity
\eqref{comparison}. The remaining assertion is the sign of that
pairing on the common kernel of the degree maps.

The series follows the rows of Weil's proof for a curve over a finite
field, and this paper is its fourth. In \cite{P43} the correspondences
of a curve are modelled by graph bimodules of idele classes over a crossed
product by the positive $S$-units, whose Hochschild complexes are determined
up to homotopy; the sum over $S$ of Weil's local terms is a trace $T$ on
$C_c^\infty(C_S)$, the form $T(F\star G^\dagger)$ is nondegenerate, and the
two Tate characters $F\mapsto\int F|w|_S^{\mp1/2}$ take the place of the two
rulings. In \cite{P44} the graphs of the powers of Frobenius are replaced by
probability measures $\mathrm{Frob}_n$ on $C_S$ with
$\mathrm{Frob}_m\star\mathrm{Frob}_n=\mathrm{Frob}_{mn}$; on test functions
approximating $\mathrm{Frob}_n$ the finite local terms of $T$ sum to
$\Lambda(n)n^{-1/2}h(1/n)$, in the place of the count of points by
intersection with the diagonal. In \cite{Paper45} the limit of the trace
over $S$ is evaluated on radial functions: it is Weil's explicit formula
$T_{\lim}(\Gamma_h)=\Phi_h(0)+\Phi_h(1)-\sum_\rho\Phi_h(\rho)$, the Tate
functionals being the residues at $s=0$ and $s=1$, and its polarization is
\eqref{comparison}. What remains of Weil's argument is the analogue of the
index inequality, Row~D, and that is the subject of this paper.

The arithmetic criterion in this paper concerns the trivial character
of conductor one, whose $L$-function is $\zeta$. Its hypotheses
require simultaneous estimates on the same fields and scales.
If these estimates hold, the criterion excludes zeros to the right
of the critical line. The functional equation then puts every
nontrivial zero on that line, and the zero expansion of Row~C gives
a sum of absolute squares for every compactly supported smooth test.
The quantitative Fourier estimates on
fixed supports and the matrix criterion for arbitrary supports give
two further descriptions of the same sign. On the windows
$T\le\frac12\log5$ the lower bounds (Theorem~\ref{thm:log5W}) are
unconditional and computer-assisted, the finite computations being done
in interval arithmetic. On the windows $T\le2.8022$ we prove
$W_T(F)\ge e^{-120110026807}\|F\|_2^2$ (Theorem~\ref{thm:power-window});
this bound depends on three facts about the zeros of $\zeta$, among them
the numerical verification that the zeros up to height $3\cdot10^{12}$ lie
on the critical line, and its finite computations are done in
exact integer and interval arithmetic.

The estimates required by the criterion are supplied by a theorem of
L.~Alpoge on the existence of families of column packages with uniform
constants (Theorem~\ref{thm:alpoge-families}). Its proof rests on his Lean
formalisation \cite{CharacterFamily} and has not yet been written out on
paper in full; the companion paper \cite{Annex} gives a first written
account of it. With this theorem we obtain the Riemann hypothesis
(Theorem~\ref{thm:RH}) and the sign $W_T\ge0$ on the entire form domain for
every $T>0$ (Corollary~\ref{cor:all-windows}).

The paper is organised as follows. Sections~\ref{sec:trace}--\ref{sec:ladder}
study the radial trace and its form: the archimedean multiplier, duality and
certificates, the explicit formula on flat-top weights, the two Tate
directions, lower bounds on finite windows, the reduction to piecewise
exponentials and a common family of test matrices for all windows.
Section~\ref{sec:primesum} reduces the sign to an estimate for the smoothed
prime sum, and that estimate to a mean bound for column counts
(Proposition~\ref{prop:simultaneous-slope}). Section~\ref{sec:meanbound}
proves the mean bound for a column package (Theorem~\ref{thm:mean-bound}),
together with the frame slope inequality for Hermitian lattices.
Section~\ref{sec:inputs} states Alpoge's theorem on families, deduces the
column estimate from it (Corollary~\ref{cor:obligation}), and summarises the
arithmetic inputs from which the theorem is derived in \cite{Annex}, with
the status of each. Section~\ref{sec:conclusion} proves the Riemann
hypothesis and the sign of the form on every window, and states on what the
proof rests.

\section{The radial trace}\label{sec:trace}
Write $I_T=(-T,T)$ and extend functions by zero. Our two transforms are
\[
 \widehat F(\tau)=\int F(t)e^{i\tau t}\,dt,\qquad
 \Phi_F(s)=\int F(t)e^{(s-1/2)t}\,dt .
\]
Thus $\Phi_F(1)=A(F)$ and $\Phi_F(0)=B(F)$ are the Tate moments.
A class is primitive when both vanish. Put
$c_F(u)=\int F(t)\overline{F(t-u)}\,dt$.
We use the zero-constant normalization of the radial trace. A finite chart containing all primes up to
$\exp(\sup\{|u|:u\in\supp c_F\})$ contains every contact of a smooth
real test. The colimit comparison \cite[Theorem~4.1]{Paper45} gives
\begin{equation}\label{comparison}
 I_\infty(\Gamma_F,\Gamma_F)
 =2A(F)B(F)-\sum_\rho\Phi_F(\rho)\Phi_F(1-\rho).
\end{equation}
Repeated integration by parts in the strip and $N(H)=O(H\log H)$ give
absolute convergence with multiplicities. The primitive inequality is $I_\infty\le0$.

On real tests the moment map $M=(A,B)$ has pairing
$H((a,b),(a',b'))=ab'+ba'$. Indeed reflection interchanges $A$
and $B$, and the Mellin transform takes convolution to multiplication.
An even nonnegative bump of positive mass translated by $r$ has moment
vector proportional to $(e^{r/2},e^{-r/2})$, with proportionality
constant independent of $r$. Two distinct small translates have
independent moments. Their common kernel therefore has codimension
two in each window. Polarizing \eqref{comparison} gives
\[
 I_\infty(\Gamma_F,\Gamma_G)=H(M(F),M(G))-
 \tfrac12\sum_\rho\{\Phi_F(\rho)\Phi_G(1-\rho)+\Phi_G(\rho)\Phi_F(1-\rho)\}.
\]
For $F\in L^2(\IT)$ its autocorrelation is continuous, supported
in $[-2T,2T]$, and has Fourier transform $|\Fh|^2$. Put
\begin{equation}\label{eq:symbol}
m(\tau)=\Rel\psi\!\left(\tfrac14+\tfrac{i\tau}{2}\right)-\log\pi,\qquad
\sT(\tau)=m(\tau)-\sum_{n\in\NT}v_n\cos(\tau\log n),\qquad v_n=\frac{2\Lambda(n)}{\sqrt n},
\end{equation}
where $\psi$ is the digamma function, $\NT=\{n\ge2:\ \log n<2T\}$ is the set of \emph{contacts} and
$\Aconst:=\sum_{n\in\NT}v_n$ is the total mass of the prime comb. The truncated Weil form is
\begin{equation}\label{eq:form}
\mathfrak a_T(F)=\frac1{2\pi}\int_\R\sT(\tau)\,|\Fh(\tau)|^2\,d\tau
\end{equation}
on the form domain $\mathcal D_T=\{F\in L^2(\IT):\int\log(2+|\tau|)|\Fh|^2<\infty\}$. The two \emph{Tate
moments} are $A(F)=\int F(t)e^{t/2}dt$ and $B(F)=\int F(t)e^{-t/2}dt$, the \emph{primitive subspace} is
$\PT=\{F\in L^2(\IT):A(F)=B(F)=0\}$, and the \emph{complete Weil form} on the window is
\begin{equation}\label{eq:W}
W_T(F)=\mathfrak a_T(F)+2\,\Rel\big(A(F)\overline{B(F)}\big).
\end{equation}
In particular the arithmetic side gives $I_\infty(\Gamma_F,\Gamma_F)=-\mathfrak a_T(F)$ for real smooth tests. The explicit formula identifies $W_T(F)$ with the zero sum in
\eqref{comparison} for real smooth tests \cite{Weil52,Paper45};
Lemma~\ref{lem:roughEF} treats the extension. We write
\[
\lamA=\inf\{\mathfrak a_T(F):F\in\PT\cap\mathcal D_T,\ \|F\|=1\},\qquad
\lamW=\inf\{W_T(F):F\in\mathcal D_T,\ \|F\|=1\};
\]
$\lamW\le\lamA$ because the two forms agree on $\PT$. Window positivity asks whether
\begin{equation}\label{eq:target}
\lamA\ge0\qquad(T>0),
\end{equation}
and complete positivity asks the same of $\lamW$. By monotonicity
(Lemma~\ref{lem:monoT}), \eqref{eq:target} is equivalent to
$\lim_{T\to\infty}\lamA\ge0$.
\subsection{Conventions}\label{subsec:conv}
$\widehat h(\tau)=\int_\R h(u)e^{i\tau u}du$; for even $h$ this is $\int h\cos(\tau u)$. For a finite even
measure $\nu$, $\nuh(\tau)=\int\cos(\tau u)\,d\nu(u)$, and for a probability measure $\mu$ on $\R$,
$\check\mu(u)=\int\cos(\tau u)\,d\mu(\tau)$. The letter $\tau$ always denotes a frequency, and $u$ a
variable dual to it. The right shift is $(V_aF)(t)=F(t-a)$, the reflection is $(JF)(t)=F(-t)$, and $Q_T$
denotes restriction to $\IT$ followed by extension by zero. We write $\mT(F)=\frac1{2\pi}\int m|\Fh|^2$
for the archimedean form, $\PiT(F)=-\sum_{n\in\NT}v_n\Rel\cF(\log n)$ for the prime form, so that
$\mathfrak a_T=\mT+\PiT$, and $\Ainf=\mT+2\Rel(A\bar B)$ for the complete archimedean form. The functions
\begin{equation}\label{eq:g}
g_0(u)=\cosh\tfrac u2\,\mathbf 1_{|u|<2T},\qquad g_1(u)=u\sinh\tfrac u2\,\mathbf 1_{|u|<2T},
\end{equation}
have transforms $\gh0,\gh1$ (Lemma~\ref{lem:gbounds}). For a parity label
$\mathrm{par}\in\{\mathrm{ev},\mathrm{odd}\}$, $\chi_{\mathrm{ev}}=\cosh\frac t2\,\mathbf 1_{\IT}$ and
$\chi_{\mathrm{odd}}=\sinh\frac t2\,\mathbf 1_{\IT}$. 

The variational approach follows \cite{Yoshida1992,Bombieri2000,Weil52,ConnesConsani2021}.

\section{The archimedean multiplier}\label{sec:levy}

\begin{lemma}[the L\'evy density]\label{lem:K}
For $u>0$ put $K(u)=\dfrac{2e^{3u/2}}{e^{2u}-1}=\dfrac{1}{2\sinh(u/2)}+\dfrac{1}{2\cosh(u/2)}
=2\sum_{k\ge0}e^{-(2k+\frac12)u}$. Then $K$ is positive, completely monotone,
$K(u)=u^{-1}+O(1)$ as $u\downarrow0$, $K(u)=2e^{-u/2}(1+O(e^{-u}))$ as $u\to\infty$, and
\begin{equation}\label{eq:m}
\begin{split}
m(\tau)&=\Rel\psi\!\left(\tfrac14+\tfrac{i\tau}2\right)-\log\pi
= C_0+\int_0^\infty(1-\cos\tau u)\,K(u)\,du,\\
C_0&=\psi(\tfrac14)-\log\pi=-5.3721834\ldots
\end{split}
\end{equation}
\end{lemma}
\begin{proof}
Writing the hyperbolic functions as exponentials gives
$K(u)=2e^{-u/2}/(1-e^{-2u})$. Its geometric series gives complete
monotonicity by termwise differentiation locally on $(0,\infty)$;
expansion at zero and infinity gives the stated bounds. In the integral
$\psi(z)=-\gamma+\int_0^\infty(e^{-t}-e^{-zt})/(1-e^{-t})\,dt$
subtract the value at $1/4$, take real parts and put $t=2u$.
The resulting integrand is $(1-\cos\tau u)K(u)$, integrable by
the bounds just obtained.
\end{proof}

\begin{theorem}[complete Bernstein structure]\label{thm:S}
Put $a_k=k+\frac14$ and $\varphi(x)=\sum_{k\ge0}\dfrac{x}{a_k(x+4a_k^2)}$ for $x\ge0$. Then
\emph{(i)} $\Rel\psi(\frac14+\frac{i\tau}2)-\psi(\frac14)=\varphi(\tau^2)$ and $\varphi$ is a
complete Bernstein function; \emph{(ii)} the L\'evy measure of the negative-definite function
$\tau\mapsto\varphi(\tau^2)$ on $\R$ is $\frac12K(|u|)du$, which is completely monotone on each
half-line, so the associated L\'evy process is a pure-jump subordinate Brownian motion
$B_{S_t}$ with a special subordinator $S$ of Laplace exponent $\varphi$; in particular it has no
Gaussian component; \emph{(iii)} $\int_{2T}^\infty K(u)\,du=\log\coth\frac T2+\pi-2\arctan e^{T}$.
\end{theorem}
\begin{proof}
Subtract the digamma series at $1/4$ from its value at
$1/4+i\tau/2$ and take real parts. Each difference equals
$\tau^2/[a_k(4a_k^2+\tau^2)]$. Thus $\varphi$ has the Stieltjes
representation with masses $1/a_k$ at $4a_k^2$, locally convergent
by comparison with $\sum_{k\ge1}k^{-3}$. Lemma~\ref{lem:K} identifies
the jump measure, and $\varphi(x)/x\to0$ excludes a Gaussian part.
The subordination statement follows from \cite[Chs.~6,13]{SSV2012}.
The tail formula follows by evaluating the primitives
$\log\tanh(u/4)$ and $2\arctan e^{u/2}$ of the two hyperbolic kernels.
\end{proof}

\begin{theorem}[the form as a L\'evy energy]\label{thm:A}
For every $F\in L^2(\IT)$, with $\Lev(\Phi):=\int_0^\infty(\Phi(0)-\Rel\Phi(u))K(u)\,du\in[0,\infty]$,
\begin{equation}\label{eq:Aform}
\mathfrak a_T(F)=C_0\|F\|^2+\Lev(\cF)-\sum_{n\in\NT}v_n\Rel\cF(\log n),
\end{equation}
both sides being finite exactly when $\int\log(2+|\tau|)|\Fh|^2<\infty$. Consequently, with
$\Kw=\{\cF:F\in L^2(\IT),\|F\|=1\}$ and the primitive class
\begin{equation}\label{eq:Kprim}
\Kp=\Big\{\Phi\in\Kw:\ \int\Phi(u)\cosh\tfrac u2\,du=0,\ \int\Phi(u)\,u\sinh\tfrac u2\,du=0\Big\},
\end{equation}
one has
\begin{align*}
\lamP&:=\inf_{\Phi\in\Kp}\Big[C_0+\Lev(\Phi)-\sum_{n\in\NT}v_n\Rel\Phi(\log n)\Big]=C_0-\Lambda_T,\\
\Lambda_T&:=\sup_{\Phi\in\Kp}\Big[\sum_{n\in\NT}v_n\Rel\Phi(\log n)-\Lev(\Phi)\Big],
\end{align*}
and $\lamP\le\lamA$.
\end{theorem}
\begin{proof}
Plancherel gives $\int\cos(\tau u)|\Fh|^2d\tau/(2\pi)=\Rel\cF(u)$.
Apply Tonelli to $(1-\cos(\tau u))K(u)|\Fh|^2$ in \eqref{eq:m}.
The resulting integral is finite on the stated domain since
$\int_0^1(1-\cos(\tau u))\,du/u=\log|\tau|+O(1)$ at infinity and
$K(u)-1/u$ is bounded near zero. Expanding the hyperbolic functions
in \eqref{eq:Kprim} gives $\Rel(A\bar B)$ and
$\Rel(A'\bar B-A\bar B')$, where primes insert $t$ into the moment
integrals. Both vanish on $\PT$, giving $\lamP\le\lamA$. For primitive
tests $\widehat{\cF}(z)=\widehat F(z)\overline{\widehat F(\bar z)}$
has double zeros at $\pm i/2$.
\end{proof}

\begin{remark}\label{rem:Kprim}
$\Kp$ is slightly larger than the set of primitive autocorrelations: for real $F$ its two
conditions read $AB=0$ and $A'B-AB'=0$, i.e.\ $F\in\PT$, or $A=A'=0$, or $B=B'=0$. Every
certificate below bounds $\lamP$, hence $\lamA$; the equality $\lamP=\lamA$ is not needed
anywhere and is not asserted. Since $W_T(F)=\mathfrak a_T(F)$ whenever $A\bar B$ is purely
imaginary, also $\lamW\le\lamP$. Whether $\lamP=\lamW$, i.e.\ whether the optimal multipliers of the
primitive problem are the polar ones, is characterized in Theorem~\ref{thm:mult}.
\end{remark}

\begin{corollary}[monotonicity; upper bound]\label{cor:mono}
With the contact set held fixed, $T\mapsto\Lambda_T$ is nondecreasing and $T\mapsto\lamP$
nonincreasing. Moreover
$\lamA\le C_0+\Aconst+\varphi\big((3\pi/2T)^2\big)$, which equals $1.291$ at $T=\frac12\log3$ and
$2.327$ at $T=\log 2$.
\end{corollary}
\begin{proof}
Monotonicity: $\Kp$ grows with $T$ and $\Lev$ does not depend on $T$. Upper bound: the form
$\Lev(\cF)$ is the Dirichlet form of the L\'evy process of Theorem~\ref{thm:S} killed on
leaving $\IT$; its eigenvalues are bounded by those of the subordinate killed Brownian motion,
$\lambda_n\le\varphi((n\pi/2T)^2)$ \cite[Prop.~2.3]{SongVondracek2003}, the Brownian motion being
normalized with characteristic exponent $\xi^2$; $\PT$ has codimension
$2$, so the bottom of $\Lev$ on $\PT$ is at most $\lambda_3$; and $-\sum v_n\cF(\log n)\le\Aconst\|F\|^2$.
\end{proof}

\begin{lemma}[elementary bounds]\label{lem:mbounds}
\emph{(i)} $m$ is increasing on $[0,\infty)$. \emph{(ii)} For every $\tau>0$,
\[
\log\frac{\tau}{2\pi}-\frac1\tau-\frac1{2\tau^2}\ \le\ m(\tau)\ \le\ \log\frac{\tau}{2\pi}+\frac1\tau+\frac12\log\Big(1+\frac1{4\tau^2}\Big),
\]
and the lower bound, which we denote by $m_{\mathrm{lo}}(\tau)$, is increasing in $\tau$. \emph{(iii)}
\[
-\sum_{n\in\NT}v_n\cos(\tau\log n)\ge-\Aconst,\qquad\text{with}\quad
\sup_\tau\sum_{n\in\NT}v_n\cos(\tau\log n)=\Aconst.
\]
\end{lemma}
\begin{proof}
(i) In $\Rel\psi(\frac14+iy)=-\gamma+\sum_{n\ge0}\big[\frac1{n+1}-\frac{n+\frac14}{(n+\frac14)^2+y^2}\big]$
every summand increases with $|y|$. (ii) Use $\psi(z)=\log z-\frac1{2z}-\int_0^\infty g(t)e^{-zt}dt$
for $\Rel z>0$, $g(t)=\frac1{e^t-1}-\frac1t+\frac12$ \cite[5.9.13]{DLMF}, at $z=\frac14+iy$, $y=\tau/2$.
Here $g(0^+)=0$, $g'(t)=t^{-2}-\frac14\sinh^{-2}(t/2)\in(0,\infty)$ and $g(\infty)=\frac12$, so one
integration by parts gives $|\int_0^\infty g e^{-zt}dt|=|z|^{-1}|\int_0^\infty g'e^{-zt}dt|\le\frac1{2|z|}\le\frac1{2y}$.
Also $\Rel\log z=\frac12\log(\frac1{16}+y^2)\in[\log y,\log y+\frac12\log(1+\frac1{16y^2})]$ and
$\Rel\frac1{2z}=\frac{1/8}{1/16+y^2}\in[0,\frac1{8y^2}]$. Substituting $y=\tau/2$ and $-\log\pi$
gives (ii); the derivative of the lower bound is $\tau^{-1}+\tau^{-2}+\tau^{-3}>0$. (iii) Each cosine is at most one, and equality holds simultaneously at $\tau=0$.
\end{proof}

\begin{lemma}[density and continuity in form norm]\label{lem:form-density}
Let $T>0$ and put
$\|F\|_{\mathcal D}^2=\frac1{2\pi}\int_\R\log(2+|\tau|)\,|\Fh(\tau)|^2\,d\tau$
on $\mathcal D_T$.
\begin{enumerate}
\item[(i)] The complex combinations of real functions in $C_c^\infty(\IT)$
are dense in $\mathcal D_T$ for $\|\cdot\|_{\mathcal D}$.
\item[(ii)] There is $C_T$ with $|W_T(F)|\le C_T\|F\|_{\mathcal D}^2$ on
$\mathcal D_T$. Hence $W_T$ and its polarization are continuous for
$\|\cdot\|_{\mathcal D}$.
\end{enumerate}
\end{lemma}
\begin{proof}
(i) Real and imaginary parts of $F\in\mathcal D_T$ lie in $\mathcal D_T$,
since the weight is even; so let $F$ be real. First dilate $F(t)$ to
$F(t/\lambda)$ with $\lambda\uparrow1$; the support lies in
$[-\lambda T,\lambda T]$. Fourier dilation is strongly continuous in
$L^2(\log(2+|\tau|)d\tau)$ by density of compactly supported continuous
functions and doubling of the weight. Mollification by a real mollifier at
scale below $(1-\lambda)T$ preserves interior support and converges in the
same norm by dominated convergence.

(ii) The function $m$ is even and continuous, so Lemma~\ref{lem:mbounds}(ii)
gives a constant $C$ with $|m(\tau)|\le C+\log(2+|\tau|)$, and
Lemma~\ref{lem:mbounds}(iii) bounds the prime part of $\sT$ by $\Aconst$.
With Plancherel and $\log2\,\|F\|_2^2\le\|F\|_{\mathcal D}^2$ this gives
$|\mathfrak a_T(F)|\le(1+(C+\Aconst)/\log2)\|F\|_{\mathcal D}^2$. By
Cauchy--Schwarz, $|A(F)|,|B(F)|\le(2e^T)^{1/2}\|F\|_2$. This proves the bound
for $W_T$; continuity of the quadratic form and of its polarization follows
from the bound and the polarization identity.
\end{proof}

\begin{lemma}[monotonicity across contact thresholds]\label{lem:monoT}
The functions $T\mapsto\lamA$, $T\mapsto\lamP$ and $T\mapsto\lamW$ are nonincreasing on $(0,\infty)$,
including at the points $T=\frac12\log n$ where a contact enters.
\end{lemma}
\begin{proof}
Let $T<T'$ and extend $F\in L^2(\IT)$ by zero to $I_{T'}$. The norm, the Tate moments $A(F),B(F)$ and
the autocorrelation $\cF$ are unchanged, and $\cF$ vanishes on $|u|\ge2T$. A contact
$n\in\mathcal N_{T'}\setminus\NT$ has $\log n\ge2T$, so $\cF(\log n)=0$, and by \eqref{eq:Aform}
$\mathfrak a_{T'}(F)=\mathfrak a_T(F)$. Hence $\PT\subset\mathcal P_{T'}$ isometrically with the same
form values, which gives the statement for $\lamA$; the same computation with $W_T=\mathfrak a_T+2\Rel(A\bar B)$
gives it for $\lamW$, and for $\lamP$ the class $\Kp$ grows with $T$, $\Lev$ is $T$-free and the new
contacts contribute $\Phi(\log n)=0$.
\end{proof}

\begin{lemma}[compactness]\label{lem:compact}
For every $V\in\R$ the set $\{F\in\mathcal D_T:\|F\|=1,\ \mathfrak a_T(F)\le V\}$ is relatively compact in
$L^2(\IT)$. Consequently the form $\mathfrak a_T$ restricted to $\PT\cap\mathcal D_T$ is the form of a
self-adjoint operator $H_T$ on $\PT$ with compact resolvent, $\lamA=\min\operatorname{spec}H_T$ is attained,
and the bottoms of $\mathfrak a_T$ and $W_T$ on the unit sphere of $L^2(\IT)$, and on any closed subspace, are
attained.
\end{lemma}
\begin{proof}
By Lemma~\ref{lem:mbounds} the shifted form norm is equivalent to the
weighted Fourier norm defining $\mathcal D_T$. This domain is complete
by Plancherel and closedness of the support condition. Its unit ball
has uniformly vanishing Fourier tails since $m(R)\to\infty$; on
$[-R,R]$ the Fourier restriction has square-integrable kernel
$e^{i\tau t}$ on $[-R,R]\times\IT$. Compactness on each band and
the tail bound give compactness of the form-domain embedding.
Closed restriction to $\PT$ and addition of the bounded polar form
preserve this argument. Compact resolvent and attainment follow.
\end{proof}

On $\PT$ the operator $H_T$ commutes with $J$, since $\sT$ is even and $A(JF)=B(F)$; so its eigenspaces
split into even and odd parts.

\begin{lemma}[continuity in the window]\label{lem:contT}\label{lem:contA}
Both $\lamW$ and $\lamA$ are continuous in $T$, as is the bottom
of the complete archimedean form. For the complete Weil form this
also follows from the parity minima in \cite[Thm.~5]{Bombieri2000}.
\end{lemma}
\begin{proof}
Work on $I_{2T}$ with zero extension, so by Lemma~\ref{lem:monoT}
the form of a vector is independent of its supporting window.
For $T_j\downarrow T$, choose unit minimizers. Their energies are
bounded above by the minimum at $T$. The polar term is bounded on the
unit ball, so compactness gives an $L^2$-convergent subsequence.
Its limit is a unit vector supported in $[-T,T]$; the Tate conditions,
when imposed, pass to the limit. Lower semicontinuity on the common
window and monotonicity give right continuity.

For the other direction take a unit minimizer $F$ at $T$ and put
$F_r(t)=r^{-1/2}F(t/r)$, $r\uparrow1$. Then $F_r\to F$ in form norm:
on Fourier transforms the dilations $g(s)\mapsto r^{1/2}g(rs)$
are uniformly bounded on $L^2(\log(2+|s|)\,ds)$ for $1/2<r<1$,
and converge strongly on continuous compactly supported functions,
a dense subspace. In the primitive case correct the two moments
of $F_r$ by a linear combination of two fixed smooth bumps in
$I_{T/2}$ with independent moment vectors. Its coefficients tend
to zero by $L^2$ convergence. The corrected vectors, normalized,
are supported in $I_{rT}$ and their energies converge to the minimum
at $T$. Monotonicity concludes. With the bounded prime comb removed
the same compactness and approximation arguments apply.
\end{proof}

\section{Duality for the window form}\label{sec:duality}

Fix $T$ and the weights $v=(v_n)_{n\in\NT}$ (any nonnegative weights may be used; $v_n=2\Lambda(n)/\sqrt n$
is the case of $\zeta$).

\begin{definition}\label{def:PD}
\textbf{(P)} $\ \lamP=\inf\big\{\int_\R\sT\,d\mu:\ \mu$ an even probability measure on $\R$,
$\check\mu(u)=0$ for $|u|\ge2T$, $\int\gh0\,d\mu=\int\gh1\,d\mu=0\big\}$.

\textbf{(D)} $\ D_T=\sup\big\{\inf_{\tau\in\R}\big[\sT(\tau)+\nuh(\tau)+\beta_0\gh0(\tau)+\beta_1\gh1(\tau)\big]:\
\nu$ a finite even signed Borel measure with $\supp\nu\subseteq\{|u|\ge2T\}$, $\beta\in\R^2\big\}$.
A triple $(\nu,\beta_0,\beta_1)$ with $f:=\sT+\nuh+\beta\cdot\widehat g\ge c$ on $\R$ is a
\emph{certificate of value $c$}.
\end{definition}

\begin{proposition}[the primal is the eigenvalue problem; weak duality]\label{prop:weak}
The value of \textup{(P)} is the $\lamP$ of Theorem~\ref{thm:A}. Every certificate of value $c$
satisfies $c\le\lamP\le\lamA$.
\end{proposition}
\begin{proof}
For $\cF\in\Kp$ take the even part of $|\Fh|^2d\tau/(2\pi)$.
Its cosine transform is $\Rel\cF$; it preserves the two zero moments
and the integral of the even symbol $\sT$. Conversely a feasible
$\mu$ has continuous positive definite transform $\check\mu$ supported
on $[-2T,2T]$ with $\check\mu(0)=1$. Boas--Kac factorization
\cite{BoasKac1945} gives $\check\mu=\cF$ for a unit vector on $\IT$;
the two constraints put $\cF$ in $\Kp$. Finally
$\int\nuh\,d\mu=\int\check\mu\,d\nu=0$ by the support conditions,
so integrating a certificate gives $c\le\int\sT\,d\mu$.
\end{proof}

\begin{lemma}[bounds on the multipliers]\label{lem:gbounds}
The multipliers $\gh0(\tau)=2\int_0^{2T}\cosh\frac u2\cos\tau u\,du$ and $\gh1(\tau)=2\int_0^{2T}u\sinh\frac u2\cos\tau u\,du$
are even, real, bounded, tend to $0$ at infinity, and for $\tau>0$
\[
|\gh0(\tau)|\le k_0(\tau):=\frac{2\cosh T}{\tau}+\frac{\sinh T}{\tau^2},\qquad
|\gh1(\tau)|\le k_1(\tau):=\frac{8T\sinh T}{\tau},
\]
both bounds decreasing in $\tau$. Moreover $\gh0(\tau)=\frac{2\cosh T\sin2T\tau}{\tau}+O(\tau^{-2})$,
and for any $g=y_0g_0+y_1g_1$, $\widehat g(\tau)=\frac{2g(2T^-)\sin2T\tau}{\tau}+\frac{2g'(2T^-)\cos2T\tau}{\tau^2}+O(\tau^{-3})$.
\end{lemma}
\begin{proof}
Integrate by parts: $\gh0=\frac{2\cosh T\sin2T\tau}\tau-\frac1\tau\int_0^{2T}\sinh\frac u2\sin\tau u\,du$,
and once more,
\[\frac1\tau\int_0^{2T}\sinh\frac u2\sin\tau u\,du=-\frac{\sinh T\cos2T\tau}{\tau^2}+\frac{\gh0(\tau)}{4\tau^2}.\]
Solving gives $|\gh0|(1+\frac1{4\tau^2})\le\frac{2\cosh T}\tau+\frac{\sinh T}{\tau^2}$. For $\gh1$, one
integration by parts gives $\gh1=\frac{4T\sinh T\sin2T\tau}{\tau}-\frac2\tau\int_0^{2T}\big(\sinh\frac u2+\frac u2\cosh\frac u2\big)\sin\tau u\,du$,
and the integrand is the derivative of $u\sinh\frac u2$, whose integral over $[0,2T]$ is $2T\sinh T$.
The asymptotic expansion is two integrations by parts for a function real-analytic on $[0,2T]$
with $g'(0)=0$.
\end{proof}


\begin{theorem}[strong duality; finitely atomic dual sequences]\label{thm:D1}
$D_T=\lamP$. Moreover, for every $\varepsilon>0$ there are finitely many atoms
$u_1,\dots,u_N\ge2T$, reals $a_1,\dots,a_N$ and $\beta\in\R^2$ such that
$\sT(\tau)+\sum_j2a_j\cos(u_j\tau)+\beta_0\gh0(\tau)+\beta_1\gh1(\tau)>\lamP-\varepsilon$ for all $\tau\in\R$.
In particular, for every $T$ with $\lamP>0$, $\lamA>0$ is provable by a finite certificate.
\end{theorem}
\begin{proof}
\emph{(i) Finite relaxations.} For a finite set $U=\{u_1<\dots<u_N\}\subset[2T,\infty)$ let
$P(U)=\inf\{\int\sT d\mu:\ \mu\ \text{even probability},\ \check\mu(u_i)=0\ \forall i,\ \int\gh j d\mu=0\}$
and $D(U)$ the supremum in (D) restricted to $\nu$ supported on $\pm U$. Then $P(U)\le\lamP$
(the feasible set of (P) is contained in that of $P(U)$) and $D(U)\le D_T\le\lamP$.

\emph{(ii) $P(U)=D(U)$.} Write $f_0\equiv1$, $f_i(\tau)=\cos(u_i\tau)$ ($1\le i\le N$),
$f_{N+1}=\gh0$, $f_{N+2}=\gh1$, $b=(1,0,\dots,0)\in\R^{N+3}$, and let
$C=\{\sum_kt_kf(\tau_k):\ t_k\ge0\}$ be the convex cone generated by the moment curve
$\tau\mapsto f(\tau)=(f_0,\dots,f_{N+2})(\tau)$. The value function
$w(b')=\inf\{\int\sT d\mu:\ \mu\ge0,\ \int f\,d\mu=b'\}$ is convex, $>-\infty$ everywhere
($\sT\ge-\Aconst-|C_0|$ and $\int f_0d\mu=b'_0$ bounds the mass) and finite on $C$ (finite atomic
$\mu$ have finite value). Its Legendre biconjugate at $b$ is the Lagrangian dual
$w^{**}(b)=\sup\{y_0:\ \sT-\sum_ky_kf_k\ge0\text{ on }\R\}$, which is $D(U)$ after
$a_i=-y_i/2$, $\beta_j=-y_{N+1+j}$. A proper convex function agrees with its biconjugate at every
interior point of its domain \cite[Thm.~7.4, Cor.~12.1.1]{Rockafellar}; so it suffices to show
$b\in\operatorname{int}C$.

Otherwise separation from the interior supplies $y\ne0$ with
$y\cdot z\ge y\cdot b$ for $z\in C$ (if $C$ has empty interior, use
a normal to its span, with sign chosen at $b$). Taking $z=0$ and
$z=tf(\tau)$ gives $y_0\le0$ and $G=y\cdot f\ge0$. The Ces\`aro mean
of $G$ is $y_0$, since the distinct positive frequencies have mean zero
and $\gh j\to0$. Hence $y_0=0$, and the mean of $|G|$ is zero.
Testing against $\cos(u_i\tau)$ gives $y_i/2=0$. There remains
$\widehat g\ge0$ with $g=y_{N+1}g_0+y_{N+2}g_1$. By the expansion in Lemma~\ref{lem:gbounds},
$\widehat g\ge0$ forces $g(2T^-)=0$ and then $g'(2T^-)=0$; but
$g(2T^-)=y_{N+1}\cosh T+2T\,y_{N+2}\sinh T$ and $g'(2T^-)=\frac12y_{N+1}\sinh T+y_{N+2}(\sinh T+T\cosh T)$,
whose determinant is $\sinh T\cosh T+T>0$. So $y=0$, a contradiction; hence $b\in\operatorname{int}C$.

Choose increasing finite $U_n$ with dense union in $[2T,\infty)$
and $1/n$-minimizers $\mu_n$. Since $\sigma_T$ is bounded below and tends to
infinity, these probabilities are tight. In a weak limit every cosine
constraint on the union survives, hence by continuity every constraint
on $[2T,\infty)$; the bounded continuous $\gh j$ likewise preserve the
two moments. Lower semicontinuity gives $\lamP\le\lim P(U_n)\le\lamP$.
Thus $D_T\ge\sup_n D(U_n)=\lamP$, with the required finite certificates.
\end{proof}

Deleting the two $\widehat g$ constraints in the proof gives
$\lamW=\sup_\nu\inf_\tau[\sT+2\gh0+\nuh]$ by Proposition~\ref{prop:P}.


\begin{theorem}[non-attainment]\label{thm:O}
Let $\nu$ be any finite even signed measure with $\supp\nu\subseteq\{|u|\ge2T\}$ and let
$h_0,\dots,h_J$ be finitely many continuous functions on $[-2T,2T]$ that annihilate $\Kp$
(for instance $g_0,g_1$). Then for every $\beta\in\R^{J+1}$
\[
\inf_{\tau\in\R}\Big[\sT(\tau)+\nuh(\tau)+\sum_j\beta_j\widehat h_j(\tau)\Big]<\lamP\quad\text{strictly.}
\]
The statement holds for every finite weight family, including $v=0$.
\end{theorem}
\begin{proof}
A minimizer $\Phi_\star=c_{F_\star}\in\Kp$ of the primal exists: the sublevel sets of
$\mathfrak a_T$ on the unit sphere of $L^2(\IT)$ are compact (Lemma~\ref{lem:compact}),
and, in the language of measures, the feasible set of (P) is tight and weak-$*$ closed and $\mu\mapsto\int\sT d\mu$ is
weak-$*$ lower semicontinuous, exactly as in step (iii) of the proof of Theorem~\ref{thm:D1}. Put $\varrho=\sT+\nuh+\sum\beta_j\widehat h_j-\lamP$, a
continuous function ($\nu$ finite). By the annihilation properties,
$\frac1{2\pi}\int\varrho|\Fh_\star|^2=\mathfrak a_T(F_\star)-\lamP=0$. If $\varrho\ge0$ then $\varrho=0$ a.e.\ on
$\{\Fh_\star\ne0\}$, whose complement is the discrete zero set of a nonzero entire function of
exponential type; so $\varrho\equiv0$ on $\R$, and by Fourier inversion, as distributions in $u$,
\[
\check\sT+\nu+\sum_j\beta_jh_j=\lamP\,\delta_0 .
\]
Restrict to the open set $0<|u|<2T$, where $\nu=0$: by \eqref{eq:m}, $\check m=-\frac12K(|u|)$
there, and the prime part of $\check\sT$ is $-\sum_{n\in\NT}\frac{v_n}2(\delta_{\log n}+\delta_{-\log n})$;
hence $\sum_j\beta_jh_j(u)=\frac12K(|u|)+\sum_{n\in\NT}\frac{v_n}2(\delta_{\log n}+\delta_{-\log n})$
on $0<|u|<2T$. The left side is a bounded continuous function; the right side is unbounded as
$u\to0$ ($K(u)\sim1/u$) and carries atoms as soon as $\NT\neq\emptyset$. Contradiction.
\end{proof}

Every value below a positive $\lamP$ admits an atomic lower
certificate by Theorem~\ref{thm:D1}. Maximizing sequences have
unbounded total variation (Corollary~\ref{cor:T2}); the primal
infimum is attained \cite[Thm.~3]{Bombieri2000}.


\begin{theorem}[finite-interval Wiener--Hopf equation]\label{thm:W}
Let $\mathsf S_T$ denote the convolution operator on $\R$ with symbol $\sT$,
$\widehat{\mathsf S_TF}=\sT\Fh$, acting on $L^2(\IT)$ extended by zero and taking values in tempered
distributions. A unit vector $F\in\PT\cap\mathcal D_T$ is a critical point of $\mathfrak a_T$ on the unit
sphere of $\PT$ with critical value $\lambda$ if and only if there are $\alpha,\beta\in\C$ (real for real
$F$) with
\[
(\mathsf S_T-\lambda)F=\alpha\cosh\tfrac t2+\beta\sinh\tfrac t2\qquad\text{on }\IT,
\]
in the sense of distributions on $\IT$. The symbol $\sT-\lambda$ is real, even, tends to
$+\infty$, and changes sign for every $\lambda>C_0-\Aconst$; in that range it is not elliptic, and no
Wiener--Hopf index is defined for it.
\end{theorem}
\begin{proof}
For $G\in C_c^\infty(\IT)$, $\langle(\mathsf S_T-\lambda)F,G\rangle=\frac1{2\pi}\int(\sT-\lambda)\Fh\,\overline{\widehat G}$
is the polarization of $\mathfrak a_T-\lambda\|\cdot\|^2$. Criticality on the unit sphere of $\PT$
with value $\lambda$ means that this vanishes for all $G$ in $\PT\cap\mathcal D_T$; since
$C_c^\infty(\IT)\cap\PT$ is dense in $\PT\cap\mathcal D_T$ for the form norm (mollify and correct the two
Tate moments by a fixed pair of $C_c^\infty$ functions), it is equivalent to require it for
$G\in C_c^\infty(\IT)$ with $\int Ge^{t/2}=\int Ge^{-t/2}=0$. This is a subspace of codimension $2$
of $C_c^\infty(\IT)$, and a distribution on $\IT$ that annihilates the common kernel of two continuous
linear functionals is a linear combination of them; the functionals are integration against
$e^{\pm t/2}$, i.e.\ against $\cosh\frac t2\pm\sinh\frac t2$. Conversely, if $(\mathsf S_T-\lambda)F$
is such a combination on $\IT$, it annihilates every $G$ in that subspace, so $F$ is critical. Finally
$\sT(0)=C_0-\Aconst$ and $\sT\to+\infty$ (Lemma~\ref{lem:mbounds}).
\end{proof}


\begin{lemma}[moment interior]\label{lem:W}
For $F\in L^2(\IT)$ real put $Q_0(F)=\int\cF\cosh\frac u2=A B$ and
$Q_1(F)=\int\cF\,u\sinh\frac u2=A'B-AB'$. The image of the set of unit vectors under
$F\mapsto(Q_0,Q_1)$ contains a neighborhood of $0$ in $\R^2$, realised by $C_c^\infty(\IT)$
functions; equivalently, $0$ is an interior point of the (convex) moment image
$\mathfrak m(\mu)=(\int\gh0d\mu,\int\gh1d\mu)$ of the feasible set of (P) without the two
constraints, and the interior is realised by measures of finite value.
\end{lemma}
\begin{proof}
The formulas for $Q_0,Q_1$ follow from expanding $\cosh\frac{t-s}2$ and $(t-s)\sinh\frac{t-s}2$
in exponentials. The functions $e^{t/2},e^{-t/2},te^{t/2},te^{-t/2}$ are linearly independent
on $\IT$, and a nontrivial combination is real-analytic, so it cannot annihilate
$C_c^\infty(\IT)$: the linear map $F\mapsto(A,B,A',B')$ is onto $\R^4$ from $C_c^\infty(\IT)$.
If $0$ were not interior, some $(t_0,t_1)\ne0$ would satisfy $t_0Q_0(F)+t_1Q_1(F)\ge0$ for all
$F$ (normalization only rescales $Q_j$ by a positive factor). Take $(A,B,A',B')=(1,0,0,\mp1)$:
$Q_0=0$, $Q_1=\pm1$, so $t_1=0$; take $(1,\pm1,0,0)$: $Q_1=0$, $Q_0=\pm1$, so $t_0=0$. A finite
convex combination of the four witnesses covers a ball, and each has finite value because
$\Fh$ is Schwartz. For the image of the unit sphere itself: a function with
$(A,B,A',B')=(1,x,0,-y)$ has $(Q_0,Q_1)=(x,y)$; adding an element of the kernel of the moment map
does not change $(Q_0,Q_1)$ and makes the norm any number above the minimal one, and
$(Q_0,Q_1)(F/\|F\|)=(Q_0,Q_1)(F)/\|F\|^2$. Since the minimal norm is bounded on $x^2+y^2=1$, every
point $t(x,y)$ with $0<t\le t_0$, $t_0>0$ uniform, is attained on the unit sphere, and $(0,0)$ is
attained on $\PT$.
\end{proof}

\begin{proposition}[the polar multiplier]\label{prop:P}
For every $F\in L^2(\IT)$ in the form domain,
$\frac1{2\pi}\int\gh0|\Fh|^2=\Rel\big(A(F)\overline{B(F)}\big)$ and hence
\[
\frac1{2\pi}\int\big(\sT+2\gh0\big)|\Fh|^2\,d\tau=W_T(F);
\]
if moreover $F$ is real and $\sum_\rho|\Phi_F(\rho)\Phi_F(1-\rho)|<\infty$ (in particular if
$F\in C_c^\infty(\IT)$ is real), then $W_T(F)=\sum_\rho\Phi_F(\rho)\Phi_F(1-\rho)$
(Lemma~\ref{lem:roughEF} below).
Consequently a certificate with $(\beta_0,\beta_1)=(2,0)$ and value $c$ proves $W_T\ge c$ on all
of $L^2(\IT)$; and $\beta_0=2$ is the only value for which
$\frac1{2\pi}\int(\sT+\beta_0\gh0)|\Fh|^2=W_T(F)$ for every $F$ in the form domain.
\end{proposition}
\begin{proof}
$\frac1{2\pi}\int\gh0|\Fh|^2=\int g_0\cF=\iint\cosh\frac{t-s}2F(t)\overline{F(s)}$ and
$\cosh\frac{t-s}2=\frac12e^{t/2}e^{-s/2}+\frac12e^{-t/2}e^{s/2}$ gives $\frac12(A\bar B+B\bar A)$.
The identification of $W_T$ with the zero sum is the explicit formula for the autocorrelation
$\cF$, \cite{Weil52,Bombieri2000}, \cite[Theorem~3.3]{Paper45}, whose polar term is
$2\Phi_F(0)\Phi_F(1)=2\Rel(A\bar B)$ for the real forms considered here. Uniqueness: for
$\beta_0\ne2$, $\frac1{2\pi}\int(\sT+\beta_0\gh0)|\Fh|^2=W_T(F)+(\beta_0-2)\Rel(A\bar B)$, and
$\Rel(A\bar B)$ is not identically zero on the form domain (it takes both signs on the unit sphere,
Lemma~\ref{lem:W}, witnesses $(1,\pm1,0,0)$), so the identity fails; in the same way a term
$\beta_1\gh1$ adds a multiple of $\Rel(A'\bar B-A\bar B')$, which is not identically zero, so the identity
also forces $\beta_1=0$.
\end{proof}

\cite[Theorem~3.3]{Paper45} is stated for real $F\in C_c^\infty$. The first equality of
Proposition~\ref{prop:P} holds on the whole form domain and is the only one the certificates use;
wherever the zero sum is used for a rough test (Theorem~\ref{thm:margin}), it enters through the
following lemma.

\begin{lemma}[the explicit formula for rough tests]\label{lem:roughEF}
Let $F\in L^2(\IT)$ be real, with \emph{(a)} $\int\log(2+|\tau|)|\Fh|^2<\infty$ and \emph{(b)}
$\sum_\rho|\Phi_F(\rho)\Phi_F(1-\rho)|<\infty$. Then
\[
\sum_\rho\Phi_F(\rho)\Phi_F(1-\rho)=W_T(F)=2\Phi_F(0)\Phi_F(1)+\mathfrak a_T(F),
\]
the sum being absolutely convergent. In particular $\mathfrak a_T(F)=\sum_\rho\Phi_F(\rho)\Phi_F(1-\rho)$
for $F\in\PT$, and a zero $\rho=\frac12+i\gamma$ on the critical line contributes $|\Fh(\gamma)|^2$.
\end{lemma}
\begin{proof}
For real $F$, $\Phi_F(1)=A(F)$ and $\Phi_F(0)=B(F)$, so $2\Phi_F(0)\Phi_F(1)=2\Rel(A\bar B)$. Let
$\eta_\varepsilon\ge0$, $\eta_\varepsilon\in C_c^\infty(-\varepsilon,\varepsilon)$, $\int\eta_\varepsilon=1$,
and $F_\varepsilon=F*\eta_\varepsilon\in C_c^\infty$, supported in $[-T-\varepsilon,T+\varepsilon]$.
Apply \cite[Theorem~3.3]{Paper45} to $F_\varepsilon$ at the window $T+\varepsilon$. Zero side:
$\Phi_{F_\varepsilon}=\Phi_F\Phi_{\eta_\varepsilon}$, $|\Phi_{\eta_\varepsilon}(s)|\le e^{\varepsilon|\Rel s-\frac12|}\le e^{\varepsilon/2}$
in the critical strip and $\Phi_{\eta_\varepsilon}\to1$ pointwise, so by (b) and dominated convergence
the zero sum (which for $F_\varepsilon$ converges absolutely, so that the symmetric order is
immaterial) tends to $\sum_\rho\Phi_F(\rho)\Phi_F(1-\rho)$. Polar term:
$\Phi_{F_\varepsilon}(0)\Phi_{F_\varepsilon}(1)\to\Phi_F(0)\Phi_F(1)$. Archimedean term:
$|\Fh_\varepsilon|^2=|\Fh|^2|\widehat{\eta_\varepsilon}|^2\le|\Fh|^2$ and $|m(\tau)|\le|C_0|+\log(1+|\tau|)$
(Lemma~\ref{lem:mbounds}(i),(ii); $m\le0$ for $|\tau|\le\tau_0$), so by (a) and dominated convergence
$\frac1{2\pi}\int m|\Fh_\varepsilon|^2\to\frac1{2\pi}\int m|\Fh|^2$. Prime term: $c_{F_\varepsilon}\to\cF$
uniformly, the sum runs over the finitely many $n<e^{2T+2\varepsilon}$, and for
$e^{2T}\le n<e^{2T+2\varepsilon}$ one has $\cF(\log n)=0$; so it tends to $-\sum_{n\in\NT}v_n\cF(\log n)$.
\end{proof}


\section{Atomic certificates}\label{sec:tails}

A certificate is an inequality $f\ge c$ on all of $\R$. The lemmas of this section reduce it to a
bounded interval; they rest on Lemmas~\ref{lem:mbounds} and \ref{lem:gbounds}.

For a feasible probability $\mu$ of value at most $\lamP+\varepsilon$
put $G=\lamP+\varepsilon+\Aconst-C_0$. The pointwise inequalities
$0\le m-C_0\le\sT+\Aconst-C_0$ give
$\int(m-C_0)d\mu\le G$. By monotonicity the mass outside $[-R,R]$
is at most $G/(m(R)-C_0)$. This is at most $\eta\in(0,1)$ if
$R\ge\max(2,2\pi\exp(1+C_0+G/\eta))$ by Lemma~\ref{lem:mbounds}.
For a certificate $f\ge c$, weak duality gives
$\lamP-c\le\int(f-c)d\mu$. Outside the interval,
$f-c\le m-C_0+(E_R+C_0)^+$ with
$E_R=\Aconst+\|\nu\|_{\TV}+\sum_i|\beta_i|k_i(R)-c$.
The omitted integral therefore costs at most
$\int_{|\tau|>R}(m-C_0)d\mu+\eta(E_R+C_0)^+$.
If $f\ge c$ is known only on the interval, set
$V=\Aconst+\|\nu\|_{\TV}+\sum_i|\beta_i|\widehat g_i(0)-C_0$.
Since $g_i\ge0$, one has $f\ge-V$ everywhere. For $c\ge-V$,
integration yields $c(1-\eta)-V\eta\le\int f\,d\mu\le\lamP+\varepsilon$.
For $c<-V$ the same bound follows from $\int f\,d\mu\ge-V$.

\subsection{Lattice atoms: the graded periodic tail}

\begin{lemma}[graded periodic tail]\label{lem:F1}
Let $d=2T/M$, $u_j=jd$ for $M\le j\le J$ (so $u_M=2T$), $\nu=\sum_ja_j(\delta_{u_j}+\delta_{-u_j})$,
$P=2\pi/d$, so that $\nuh(\tau)=\sum_j2a_j\cos(u_j\tau)$ is $P$-periodic. Let
$f=\sT+\nuh+\beta_0\gh0+\beta_1\gh1$. Fix $\tau_1>0$, $N\ge1$, and for $x\in[\tau_1,\tau_1+P)$ put
\begin{align*}
Y_k(x)&=m(x+kP)-\sum_{n\in\NT}v_n\cos\big((x+kP)\log n\big)-|\beta_0|k_0(x+kP)-|\beta_1|k_1(x+kP)\ \ (k<N),\\
Y_N(x)&=m(x+NP)-\Aconst-|\beta_0|k_0(x+NP)-|\beta_1|k_1(x+NP).
\end{align*}
If $\nuh(x)+Y_k(x)\ge c$ for all $x\in[\tau_1,\tau_1+P)$ and all $0\le k\le N$, then $f(\tau)\ge c$
for all $\tau\ge\tau_1$.
\end{lemma}
\begin{proof}
Every $\tau\ge\tau_1$ is uniquely $x+kP$ with $x\in[\tau_1,\tau_1+P)$, $k\ge0$, and
$\nuh(\tau)=\nuh(x)$ since $u_jP=2\pi j$. For $k<N$ the $k$-th row is $f(x+kP)$ with the exact
prime term and the two multiplier terms replaced by the lower bounds of Lemma~\ref{lem:gbounds}.
For $k\ge N$: $m(x+kP)\ge m(x+NP)$, $k_i(x+kP)\le k_i(x+NP)$, and the prime term is $\ge-\Aconst$.
\end{proof}

The prime term is never required to be periodic: it is evaluated exactly on the first $N$
copies of one period of $\nuh$ and replaced by its infimum beyond. The lemma therefore holds
for any finite set of contacts; the incommensurability of $\log2$ and $\log3$ plays no role.
It is a sufficient criterion only (the rows are conservative in $\beta$ by
$|\beta_0|k_0+|\beta_1|k_1$).

\subsection{Exponential columns: the monotone tail}

\begin{proposition}[exponential columns]\label{prop:T0}
For $s>0$ and an integer $p\ge0$ let $d\nu_{s,p}(u)=\frac12(|u|-2T)^pe^{-s(|u|-2T)}\mathbf 1_{|u|\ge2T}du$,
a positive finite measure of mass $p!/s^{p+1}$ supported outside the window. Then
\[
\nuh_{s,p}(\tau)=\Rel\Big[e^{2iT\tau}\frac{p!}{(s-i\tau)^{p+1}}\Big],\qquad
|\nuh_{s,p}(\tau)|\le\operatorname{cap}_{s,p}(\tau):=\frac{p!}{(s^2+\tau^2)^{(p+1)/2}},
\]
and $\operatorname{cap}_{s,p}$ is decreasing in $\tau\ge0$.
\end{proposition}
\begin{proof}
$\nuh_{s,p}(\tau)=\int_{2T}^\infty(u-2T)^pe^{-s(u-2T)}\cos\tau u\,du=\Rel\big[e^{2iT\tau}\int_0^\infty t^pe^{-(s-i\tau)t}dt\big]$.
\end{proof}

\begin{lemma}[monotone tail]\label{lem:T}
Let $\nu=\sum_{k=1}^Kc_k\nu_{s_k,p_k}+\sum_{j=1}^Ja_j(\delta_{u_j}+\delta_{-u_j})$ with $u_j\ge2T$, and
$f=\sT+\nuh+\beta_0\gh0+\beta_1\gh1$. Define for $\tau>0$
\[
L(\tau)=\log\frac\tau{2\pi}-\frac1\tau-\frac1{2\tau^2}-\Aconst-\sum_k|c_k|\operatorname{cap}_{s_k,p_k}(\tau)
-2\sum_j|a_j|-|\beta_0|k_0(\tau)-|\beta_1|k_1(\tau).
\]
Then $f\ge L$ on $(0,\infty)$, $L$ is nondecreasing, and $L(\tau_1)\ge c$ for one $\tau_1$ implies
$f\ge c$ on $[\tau_1,\infty)$.
\end{lemma}
\begin{proof}
Lemma~\ref{lem:mbounds}(ii),(iii), Proposition~\ref{prop:T0}, the trivial bound on the atoms and
Lemma~\ref{lem:gbounds}; every term of $L$ is nondecreasing.
\end{proof}

\begin{corollary}[cost of the tail]\label{cor:T2}
For an atom-free $\nu$ with $A:=\sum_k|c_k|p_k!$ and $A':=|\beta_0|(2\cosh T+\sinh T)+|\beta_1|8T\sinh T$,
the tail inequality $L(\tau_1)\ge c$ holds for every $\tau_1\ge\max\{A+A'+2,\ 2\pi e^{c+\Aconst+1}\}$,
so the tail costs linearly in the $\ell^1$ size of the certificate. For atoms it costs
exponentially, $\tau_1\approx2\pi e^{\Aconst+2\sum|a_j|}$, which is why Lemma~\ref{lem:F1} and not
Lemma~\ref{lem:T} is used for lattice certificates. Neither cost is uniform over the dual class:
by Theorem~\ref{thm:O}, $\|\nu\|\to\infty$ along every maximizing sequence.
\end{corollary}
\begin{proof}
For $\tau\ge1$, $\operatorname{cap}_{s,p}(\tau)\le p!/\tau$ and $|\beta_i|k_i\le A'/\tau$, so
$L(\tau)\ge\log\frac\tau{2\pi}-\Aconst-\frac{A+A'+2}{\tau}$.
\end{proof}

\begin{lemma}[hybrid tail]\label{lem:hybrid}
Let $\nu=\nu_{\mathrm{lat}}+\sum_\ell c_\ell\nu_{s_\ell,p_\ell}$ with $\nu_{\mathrm{lat}}$ as in Lemma~\ref{lem:F1}.
For $x\in[\tau_1,\tau_1+P)$ define $H_k(x)=\nuh_{\mathrm{lat}}(x)+Y_k(x)+\sum_\ell c_\ell\nuh_{s_\ell,p_\ell}(x+kP)$
for $0\le k<N$, and $H_N(x)=\nuh_{\mathrm{lat}}(x)+Y_N(x)-\sum_\ell|c_\ell|\operatorname{cap}_{s_\ell,p_\ell}(x+NP)$.
If $H_k\ge c$ on $[\tau_1,\tau_1+P)$ for all $0\le k\le N$ then $f\ge c$ on $[\tau_1,\infty)$.
\end{lemma}
\begin{proof}
As in Lemma~\ref{lem:F1}; on the rows $k<N$ the column term is evaluated exactly, and for
$k\ge N$ it is bounded by $-\sum_\ell|c_\ell|\operatorname{cap}(x+kP)\ge-\sum_\ell|c_\ell|\operatorname{cap}(x+NP)$
since $\operatorname{cap}$ decreases.
\end{proof}



\begin{theorem}[exponential columns exhaust the dual]\label{thm:E}
Let $\mathcal E$ be the real span of $\{\nu_{s,p}:s>0,p\ge0\}$. Then
$\sup\{\inf_\tau[\sT+\nuh+\beta\cdot\widehat g]:\ \nu\in\mathcal E,\beta\in\R^2\}=\lamP$, and the same
holds with $\mathcal E$ replaced by the absolutely continuous measures with bounded density and
compact support in $\{|u|\ge2T\}$.
\end{theorem}
\begin{proof}
Weak duality gives the upper bound. By Theorem~\ref{thm:D1} choose an
atomic certificate within $\varepsilon$ of $\lamP$. Since $\sT\to\infty$,
the tail bound beyond some $R$ holds for every replacement of no larger
total variation. Spread each positive-location atom uniformly over an
interval of length $\eta$ to its right and reflect at zero. This preserves
the support condition and changes the transform on $[-R,R]$ by at most
$R\eta\|\nu\|_{\rm TV}$. The densities $t^pe^{-t}$ span a dense subspace
of $L^1(0,\infty)$: an annihilator in $L^\infty$ has Laplace transform
with every derivative zero at one, hence is zero by uniqueness.
Translate this approximation by $2T$ and reflect. Total variation error
bounds Fourier error uniformly, proving both assertions.
\end{proof}


\begin{proposition}[the canonical column; band-limitation]\label{prop:M}
Let $d\nu_K=\frac12K(|u|)\mathbf 1_{|u|\ge2T}du$, a positive measure of mass
$c_K=\int_{2T}^\infty K$ (Theorem~\ref{thm:S}(iii)). Then for all $\tau$
\[
\sT(\tau)+\nuh_K(\tau)=\big(C_0+c_K\big)+\int_0^{2T}(1-\cos\tau u)K(u)\,du-\sum_{n\in\NT}v_n\cos(\tau\log n),
\]
a function whose distributional Fourier transform is supported in $[-2T,2T]$ (it still grows
like $\log|\tau|$: the singularity of $K$ at $u=0$ lies inside the window). Moreover
$d\nu_K=2e^{-T}\big(1+O(e^{-2|u|})\big)\,d\nu_{\frac12,0}$:
the archimedean L\'evy density selects the scale $s=\frac12$ among the exponential columns.
\end{proposition}
\begin{proof}
Splitting \eqref{eq:m} at $2T$ gives the identity. The constant,
the prime cosines and each $1-\cos(\tau u)$ with $u\le2T$ have
Fourier support in $[-2T,2T]$. Integration against $K$ preserves
this support by Fubini on Schwartz tests, using
$|1-\cos(\tau u)|\le\min(2,\tau^2u^2)$. The density ratio follows
from $K(u)/2=e^{-u/2}/(1-e^{-2u})$ for $u>0$.
\end{proof}

At $T=\tfrac12\log3$ only the contact at two contributes. For spacing
$\log(3)/24$, Lemma~\ref{lem:F1} reduces the tail test to periodic rows.

\begin{proposition}[second-order cover]\label{prop:TAY}
Write $f=m+q+r$ with $q=-\sum_{n\in\NT}v_n\cos(\tau\log n)+\nuh$ ($\nu$ atomic) and
$r=\beta_0\gh0+\beta_1\gh1$, and put
$M_2=\sum_{n\in\NT}v_n\log^2n+2\sum_j|a_j|u_j^2+|\beta_0|H_0+|\beta_1|H_1$ with
$H_0=\int_{-2T}^{2T}u^2\cosh\frac u2\,du$, $H_1=\int_{-2T}^{2T}|u|^3\sinh\frac{|u|}2\,du$. Then
$\sup|q''+r''|\le M_2$, and on every cell $[x_0-\frac w2,x_0+\frac w2]\subset[0,\infty)$
\[
\min_{\text{cell}}f\ \ge\ m\!\big(x_0-\tfrac w2\big)+(q+r)(x_0)-|(q+r)'(x_0)|\,\tfrac w2-M_2\tfrac{w^2}8 .
\]
The same bound applies to each row $\nuh+Y_k$ of Lemmas~\ref{lem:F1}, \ref{lem:hybrid} with
$m(x+kP)$ and $-|\beta_i|k_i(x+kP)$ evaluated at the left endpoint. A cell is certified when the
right side exceeds the target and is bisected otherwise; near a contact of $f$ the admissible width
is $\approx\sqrt{8c/M_2}$, so the number of evaluations grows like $c^{-1/2}$ rather than the
$c^{-1}$ of a Lipschitz cover.
\end{proposition}
\begin{proof}
Taylor's formula with integral remainder for $q+r\in C^2$, and monotonicity of $m$ and of
$-|\beta_i|k_i$.
\end{proof}



If a feasible $\mu$ has mass $e$ outside $[-R,R]$, then
$\int\sT\,d\mu\ge C_0-\Aconst+e(m(R)-m(0))$. For the reverse scale
take any feasible measure within $\delta$ of $\lamP$ and mix its
frequency translates by $0,\pm R$ in proportions $1-e,e/2,e/2$.
This multiplies $\check\mu$ by $1-e+e\cos Ru$. The moment defect is
$v=e\,o(1)$ by bounded convergence since $\gh j$ vanish at infinity.
For $v\ne0$ mix in the bounded-cost measure of Lemma~\ref{lem:W}
with moments $-r_0v/|v|$ in proportion $|v|/(r_0+|v|)$.
The corrected measure has mass at least $e(1-o(1))$ outside
$[-R+\sqrt R,R-\sqrt R]$, since $\mu([-\sqrt R,\sqrt R])\to1$.
Its value is at most $\lamP+\delta+e(\log R+O(1))$: the translate
cost follows from $\log(2+|t\pm R|)\le\log(2+R)+\log(2+|t|)$
and the finite log moment of $\mu$.

The direct tail bound subtracts $\Aconst+\|\nu\|_{\rm TV}
+\sum_i|\beta_i|k_i$ from $m_{\rm lo}$. Partial summation in the prime
number theorem gives $\Aconst=(4+o(1))e^T$.

\section{The explicit formula on flat-top weights}\label{sec:Z}

\begin{definition}\label{def:W}
$\Wcl$ is the set of even measurable $w:\R\to\R$ with $w\equiv1$ on $[-2T,2T]$ and
$|w(u)|\le Ce^{-(\frac12+\delta)|u|}$ for $|u|\ge2T$, for some $C,\delta>0$. $\Wcl^{\mathrm{reg}}\subset\Wcl$
consists of the $w\in C^\infty(\R)$ with compact support. For $w\in\Wcl$ define the even measure
\begin{equation}\label{eq:nuw}
d\nu_w=\Big[\tfrac12\big(1-w(u)\big)K(|u|)+2w(u)\cosh\tfrac u2\Big]\mathbf 1_{|u|\ge2T}\,du
-\sum_{n\notin\NT}\tfrac{v_n}2\,w(\log n)\,\big(\delta_{\log n}+\delta_{-\log n}\big),
\end{equation}
and $f_w:=\sT+\nuh_w+2\gh0$.
\end{definition}
The measure $\nu_w$ is finite ($K$ is integrable on $[2T,\infty)$, $|w|\cosh\frac u2$ is
integrable by the decay condition, and $\sum_nv_n|w(\log n)|\le2C\sum\Lambda(n)n^{-1-\delta}$) and
supported in $\{|u|\ge2T\}$. Hence $(\nu_w,2,0)$ is a dual feasible point, and by
Propositions~\ref{prop:weak} and \ref{prop:P}, $W_T\ge\inf f_w$ on $L^2(\IT)$.

\begin{theorem}\label{thm:Z}
\emph{(i)} For every $w\in\Wcl$ and $\tau\in\R$,
\begin{equation}\label{eq:fw}
\begin{split}
f_w(\tau)=C_0+\int_0^\infty\big(1-w(u)\cos\tau u\big)K(u)\,du&-\sum_{n\ge2}v_n\,w(\log n)\cos(\tau\log n)\\
&+2\int_\R w(u)\cosh\tfrac u2\cos(\tau u)\,du,
\end{split}
\end{equation}
which is the right-hand side of the Weil explicit formula for the even test function
$h_\tau(u)=w(u)\cos(\tau u)$.
\emph{(ii)} For $w\in\Wcl^{\mathrm{reg}}$, with $\wh(t)=\int w(u)e^{itu}du$ (an even entire function
of exponential type),
\begin{equation}\label{eq:Z}
f_w(\tau)=\sum_\rho\wh(\tau-\gamma_\rho),\qquad\rho=\tfrac12+i\gamma_\rho,
\end{equation}
the sum running over the non-trivial zeros of $\zeta$ with multiplicity, and converging absolutely.
A zero on the critical line contributes $\wh(\tau-\gamma_\rho)$ with $\gamma_\rho$ real, so on
the zeros of the critical line $f_w$ is the smoothing of the zero-counting measure by the kernel $\wh$.
\end{theorem}
\begin{proof}
For (i), insert \eqref{eq:nuw} into $f_w$. On $[0,2T]$ one has $w=1$;
on its complement the density $(1-w)K/2$ supplies the correction to
\eqref{eq:m}. The atoms supply exactly the omitted prime powers and the
remaining density completes the polar integral, giving \eqref{eq:fw}.
For (ii), apply the explicit formula \cite{Weil52,IK}
to $h=w\cos(\tau\,\cdot)$. Its transform is
$(\wh(t-\tau)+\wh(t+\tau))/2$; evenness and $\rho\leftrightarrow1-\rho$
give \eqref{eq:Z}. Indeed Fourier inversion against $1$ and $\cos(tu)$
gives $1$ and $w(u)\cos(\tau u)$, respectively, in the archimedean
term. The interchange is justified by
$\int|\wh(t-\tau)|\int_0^\infty\min(2,t^2u^2)K(u)\,du\,dt<\infty$.
Uniform rapid decay on $|\Iml t|\le1/2$ and $N(t)\ll t\log t$
give absolute convergence over zeros.
\end{proof}


\begin{corollary}[certificates are statements about zeros]\label{cor:zeros}
For $w\in\Wcl^{\mathrm{reg}}$, a lower bound $f_w\ge c$ applies
to the full zero sum by \eqref{eq:Z}. Finite symmetric truncations
to real ordinates take negative values by Theorem~\ref{thm:finZ}.
\end{corollary}

\begin{theorem}[exhaustion by flat-top weights]\label{thm:Zexh}
Let $T\ge\frac12\log2$. \emph{(i)} $\sup_{w\in\Wcl^{\mathrm{reg}}}\inf_{\tau\in\R}f_w(\tau)=\lamW$.
\emph{(ii)} $\sup\{\inf_\tau[f_w+(\beta_0-2)\gh0+\beta_1\gh1]:\ w\in\Wcl^{\mathrm{reg}},\beta\in\R^2\}=\lamP$.
\end{theorem}
\begin{proof}
Weak duality gives the upper bounds. Theorems~\ref{thm:D1}
and~\ref{thm:E} give finite even atomic certificates $\nu$ approaching
the two suprema, with $\beta=(2,0)$ in (i). If $2T=\log n$, let
$B=-v_n(\delta_{2T}+\delta_{-2T})/2$; otherwise put $B=0$.
Every flat-top weight gives this same boundary atom since $w(2T)=1$.
Replace each atom of $\nu-B$ by a smooth bump of the same mass,
supported strictly outside the window, respecting reflection. The
resulting density $h$ has $\|h\|_1\le\|\nu-B\|_{\TV}$, and the
transform of $B+h\,du$ approximates $\widehat\nu$ uniformly on
any prescribed compact frequency interval by taking narrow bumps.
Choose that interval first: $m(\tau)\to\infty$ while the prime comb,
moment multipliers and transforms of all these measures are bounded
independently of the bump widths. Thus the lower bound for the
certificate loses at most $\varepsilon$.

To realize the density, set $w=(h-K/2)/(2\cosh(u/2)-K/2)$ on
$u>2T$ and reflect evenly. The denominator is positive there by
monotonicity and $2T\ge\log2$. Interpolate smoothly to $1$ at the
window edge, and cut off beyond a large radius, where $w=O(e^{-u})$.
These changes cost arbitrarily little in the $L^1$ norm of the density
in \eqref{eq:nuw}. Inside the cutoff force $w$ to vanish in small
disjoint neighborhoods of the finitely many $\log n>2T$ by smooth
cutoffs. Their total density cost tends to zero with their lengths.
The resulting atoms are exactly $B$, so the final density error
perturbs the transform uniformly by arbitrarily little.
\end{proof}

\begin{theorem}\label{thm:finZ}
Let $Z$ be a nonempty finite symmetric multiset of real ordinates.
If $w$ is even, continuous, real and integrable, with $w=1$ on
$[-2T,2T]$, $T>0$, then $\sum_{\gamma\in Z}\widehat w(t-\gamma)$
takes negative values.
\end{theorem}
\begin{proof}
Write $g(u)=\sum_{\gamma\in Z}e^{iu\gamma}$ and $f=\widehat{wg}$.
If $f\ge0$, Gaussian damping followed by monotone convergence gives
$\int f(t)\,dt=2\pi g(0)$ by continuity at zero. The measure
$d\nu=f(t)\,dt/(2\pi)$ therefore has Fourier transform $wg$ by
inversion and evenness. On the window this is also the transform
of $\mu=\sum_{\gamma\in Z}\delta_\gamma$.
Let $z_1,\ldots,z_l$ be the distinct ordinates. For $0<a<2T/l$,
the integral of $|\prod_{j=1}^l(e^{iat}-e^{iaz_j})|^2$ uses only
Fourier values at $ka$, $|k|\le l$. It vanishes for $\mu$, hence
for $\nu$. Along a sequence $a\downarrow0$, almost every $t$ for
$\nu$ satisfies $a(t-z_j)\in2\pi\mathbb Z$ for some recurrent $j$;
then $t=z_j$. An absolutely continuous measure supported on finitely
many points has mass zero, contradicting $\nu(\mathbb R)=|Z|$.
\end{proof}

For a smooth compactly supported flat-top weight on $[-U,U]$ one
can also control the truncation error directly. If $D^{k+2}w$ is a
measure of total variation $V$, integration by parts gives
$|\widehat w(z)|\le Ve^{U|\Im z|}|z|^{-k-2}$ for $z\ne0$.
For $|t|\le H/2$ and $H\ge2$, group the omitted ordinates
$\gamma_\rho=(\rho-1/2)/i$ in shells
$2^jH<|\gamma_\rho|\le2^{j+1}H$. Their distances from $t$ are
at least $2^{j-1}H$, and $|\Im\gamma_\rho|\le1/2$.
The zero count $N(y)=O(y\log y)$ gives a total error
$O_{U,k}(VH^{-k-1}\log H)$ after summation. For autocorrelations
the on-line zeros instead contribute the nonnegative terms
$|\widehat F(\gamma)|^2$ individually.

\subsection{The Kre\u{\i}n extension problem at one window}\label{subsec:KEP}
Section~\ref{sec:Z} prescribes the Fourier transform of a positive measure on an interval.
This is the distributional analogue of Kre\u{\i}n's extension problem
\cite{Krein1940,Sasvari1994}.

\begin{definition}[the window datum, and $\mathrm{KEP}(T)$]\label{def:KEP}
Let $g$ be the even distribution defined on even $h\in C_c^\infty(\R)$ by
\begin{equation}\label{eq:gdatum}
\langle g,h\rangle=C_0\,h(0)+\int_0^\infty\big(h(0)-h(u)\big)K(u)\,du-\sum_{n\ge2}v_n\,h(\log n)
+2\int_\R h(u)\cosh\tfrac u2\,du,
\end{equation}
with the normalization of Theorem~\ref{thm:Z}(ii). Let $\mathfrak S(T)$ consist of the positive
symmetric tempered measures $\nu$ satisfying $\langle\nuh-g,h\rangle=0$
for every even $h\in C_c^\infty(-2T,2T)$.
The \emph{Kre\u{\i}n extension problem} $\mathrm{KEP}(T)$ is \emph{solvable} if
$\mathfrak S(T)\ne\emptyset$, \emph{determinate} if $\#\mathfrak S(T)=1$, and \emph{indeterminate}
if $\#\mathfrak S(T)\ge2$.
\end{definition}
Even tests determine the symmetric part of a measure. By Schwartz--Bochner,
solvability means that the window datum extends to a positive definite tempered distribution.

\begin{lemma}[one negative square in the polar term]\label{lem:polar}
The polar kernel $2\cosh((s-t)/2)$ has one positive and one negative
square on any set of at least two distinct points.
\end{lemma}
\begin{proof}
Writing $a_i=e^{t_i/2}$, $b_i=e^{-t_i/2}$ gives the matrix $ab^t+ba^t$.
Its two nonzero eigenvalues $a\cdot b\pm|a||b|$ have opposite signs
by strict Cauchy--Schwarz.
\end{proof}

\begin{theorem}[a certificate is an absolutely continuous solution]\label{thm:KEP1}
Let $T\ge\frac12\log2$ and let $w\in\Wcl^{\mathrm{reg}}$ satisfy $f_w\ge c$ on $\R$.
\emph{(i)} If $c\ge0$ then $\nu_w:=f_w(\tau)\,d\tau/2\pi$ belongs to $\mathfrak S(T)$; in particular
$\mathrm{KEP}(T)$ is solvable, by an absolutely continuous solution.
\emph{(ii)} If $c>0$ then $\mathfrak S(T)$ contains a continuum of pairwise distinct absolutely
continuous measures; in particular $\mathrm{KEP}(T)$ is indeterminate.
\end{theorem}
\begin{proof}
(i) Theorem~\ref{thm:Z} gives $\widehat{wg}=f_w$. Its logarithmic growth
gives temperedness, evenness gives symmetry, and $wg=g$ on the window.
(ii) Take a nonzero even nonnegative smooth bump $\chi$ supported beyond
$\max(2T,\log2)$ in absolute value and away from prime atoms.
There $\chi g$ has density $\chi(2\cosh(u/2)-K(|u|)/2)$, strictly
positive wherever $\chi>0$ by the proof of Theorem~\ref{thm:Zexh}.
Its Fourier transform $r$ is bounded and nonzero. For
$|\varepsilon|<c/\|r\|_\infty$, the densities
$(f_w+\varepsilon r)/(2\pi)$ give distinct positive measures with
the same window datum, since $\chi g$ vanishes there.
\end{proof}

\begin{theorem}[solvability implies Weil positivity at that window]\label{thm:KEP2}
If $\mathfrak S(T)\ne\emptyset$ then $\lamW\ge0$.
\end{theorem}
\begin{proof}
For real $F\in C_c^\infty(\IT)$ its autocorrelation $h$ is even and
supported in $(-2T,2T)$. By \eqref{eq:gdatum} and \eqref{eq:W},
$W_T(F)=\langle g,h\rangle=\int|\Fh|^2d\nu\ge0$ for $\nu\in\mathfrak S(T)$.
Real and imaginary parts add for the real symmetric form, so $W_T\ge0$ on
complex combinations of real smooth tests. Lemma~\ref{lem:form-density}
extends the inequality to $\mathcal D_T$.
\end{proof}

\begin{theorem}[the dichotomy at one window]\label{thm:KEPdich}
For every $T\ge\frac12\log2$,
\[
\lamW>0\ \Longrightarrow\ \mathrm{KEP}(T)\ \text{is indeterminate}\ \Longrightarrow\ \lamW\ge0 ,
\]
so certificates and indeterminacy coincide except on the set $\{T:\lamW=0\}$.
\end{theorem}
\begin{proof}
If $\lamW>0$, Theorem~\ref{thm:Zexh}(i) gives $w\in\Wcl^{\mathrm{reg}}$ with $\inf f_w>0$, and
Theorem~\ref{thm:KEP1}(ii) applies; the second implication is Theorem~\ref{thm:KEP2}.
\end{proof}

\section{The archimedean form on primitive tests}\label{sec:archstr}

Put, for $F\in L^2(\IT)$,
\[
\Ainf(F)=2A(F)B(F)+\frac1{2\pi}\int_\R m(\tau)|\Fh(\tau)|^2\,d\tau
=2A(F)B(F)+C_0\|F\|^2+\Lev(\cF),
\]
the complete Weil form with the prime comb deleted (the second equality is Theorem~\ref{thm:A} with
$\NT=\emptyset$), and write $\PiT(F)=-\sum_{n\in\NT}v_n\Rel\cF(\log n)$ for the prime form (the real part is redundant for real $F$), so that
$W_T=\Ainf+\PiT$. Let $\laminf(T)$ be the bottom of $\Ainf$ on the unit sphere of $L^2(\IT)$ and
$\laminfP(T)$ its bottom on the unit sphere of $\PT$, where the pole term vanishes identically and
$\Ainf=C_0\|F\|^2+\Lev(\cF)$.

\begin{proposition}[two archimedean thresholds]\label{prop:arch}
(i) Both minima decrease weakly with $T$, and $\laminfP$ decreases strictly.
(ii) $C_0\le\laminfP(T)\le C_0+\varphi((3\pi/2T)^2)$ and
$\laminf(T)\le\laminfP(T)$. In particular $\laminfP(T)\downarrow C_0$.
(iii) For $2T<\log2$, $W_T=\Ainf$ and $\laminf(T)\ge0$.
(iv) The suprema $T^*_\infty,T^*_{\infty,\mathrm P}$ of the radii of
strict positivity for the two minima satisfy
$\tfrac12\log2\le T^*_\infty\le T^*_{\infty,\mathrm P}\le T_0(0)$,
where $\varphi((3\pi/(2T_0(0)))^2)=-C_0$ and $T_0(0)\simeq0.74921$.
\end{proposition}
\begin{proof}
Zero extension enlarges the feasible sets; strictness on primitive tests
is Lemma~\ref{lem:strict}. Inclusion orders the minima, $\Lev\ge0$,
and the three-dimensional trial space of Corollary~\ref{cor:mono}
gives (ii) with the prime comb omitted. Its upper bound tends to $C_0$.
There are no prime contacts below $\tfrac12\log2$; strict positivity
at that radius \cite{Yoshida1992} and (ii) give (iii)--(iv).
Connes and Consani \cite{ConnesConsani2021} prove the positivity of the
Weil functional on $g*g^*$ for test functions $g$ supported in
$[2^{-1/2},2^{1/2}]$ whose Fourier transform vanishes at $i/2$ and at $0$;
in the additive variable this is the window $T=\frac12\log2$.
See also \cite{Bombieri2000,Zhu2026}.
\end{proof}

\begin{proposition}[nodal structure of $\PT$]\label{prop:nodes}
The ground state of $\mT=C_0+\Lev$ on the full window is simple and can
be chosen strictly positive almost everywhere. Every nonzero real primitive
test, including a real primitive ground state, has at least two sign changes:
its signs cannot be ordered on two consecutive subintervals, even allowing an
empty subinterval. For an even test a change of sign on $(0,T)$ is required.
\end{proposition}
\begin{proof}
Zero extension and symmetry give the jump and killing terms
\[
 \Lev(\cF)=\tfrac14\iint_{\IT^2}K(|x-y|)|F(x)-F(y)|^2\,dx\,dy
 +\tfrac12\int_{\IT}|F(x)|^2\!\int_{|y|>T}K(|x-y|)\,dy\,dx.
\]
The Markov contraction decreases this form and $K>0$ connects every
pair of positive-measure sets, so its semigroup is positivity improving.
Compactness follows from fixed support and the uniform Fourier tail
bound $\int_{|\tau|>R}|\Fh|^2\le2\pi\Lev(\cF)/(m(R)-C_0)$;
Perron--Frobenius gives the assertion. If the signs of $F$ were
ordered about $t_0$, then $G=e^{t/2}-e^{t_0}e^{-t/2}$ or $-G$
would have the matching signs and $\int FG\ne0$, contradicting
both Tate conditions. For even $F$ the positive weight $\cosh(t/2)$
already forces both signs on $(0,T)$.
\end{proof}

For $j=\mathrm{ev},\mathrm{odd}$ let $\PT^j$ be the corresponding parity
subspace, $\lambda_j,n_j$ its spectral bottom and negative index
(counting multiplicities), and $T_j^*=\sup\{T:\lambda_j(T)>0\}$.
Its moment condition is orthogonality to $\cosh(t/2)$ or $\sinh(t/2)$,
respectively.

\begin{lemma}[parity splitting]\label{lem:parity}
On $\PT$ the negative indices add and $\laminfP=\min_j\lambda_j$.
\end{lemma}
\begin{proof}
Reflection exchanges the Tate moments and commutes with the real even
multiplier $m$. Its two eigenspaces thus split both the norm and the form
orthogonally, giving the assertions by compact resolvent. Complexification
preserves the eigenvalues and their multiplicities.
\end{proof}

\begin{lemma}[zero extension]\label{lem:zeroext}
In each sector $n_j$ increases weakly and $\lambda_j$ decreases weakly
with the radius.
\end{lemma}
\begin{proof}
Zero extension preserves parity, both moments, the norm and the
autocorrelation on $\mathbb R$, hence also $\mT$.
\end{proof}

\begin{lemma}[strict monotonicity]\label{lem:strict}
For every $T>0$ the bottoms $\lambda_{\mathrm{ev}}(T)$, $\lambda_{\mathrm{odd}}(T)$ and
$\laminfP(T)$ are attained, and all three are strictly decreasing on $(0,\infty)$. Consequently, for each
sector, $\lambda_{\mathrm{par}}(T)>0$ for $T<T^*_{\mathrm{par}}$ and $\lambda_{\mathrm{par}}(T)<0$ for
$T>T^*_{\mathrm{par}}$; $\lambda_{\mathrm{par}}(T)<0$ if and only if $n_{\mathrm{par}}(T)\ge1$;
$T^*_{\infty,\mathrm P}=\min(T^*_{\mathrm{ev}},T^*_{\mathrm{odd}})$; and $\laminfP$ is positive on
$(0,T^*_{\infty,\mathrm P})$ and negative on $(T^*_{\infty,\mathrm P},\infty)$.
\end{lemma}
\begin{proof}
The form domain embeds compactly by Proposition~\ref{prop:nodes}. Both
moment constraints and parity are closed in $L^2$, so the direct method
gives a minimizer; real and imaginary parts split and one can take it real.
If a sector bottom agrees at $T'<T$, its normalized minimizer $F^*$ at
$T'$ also minimizes at $T$ by zero extension. Test its eigenvalue equation
against $G_0+\epsilon JG_0$, where $\epsilon$ is its parity and
$G_0\in C_c^\infty(T',T)$ has both moments zero. Disjoint supports and
reflection give $2\mT(F^*,G_0)=0$. Polarizing the jump integral, this
equals $-\iint K(|y-x|)F^*(x)G_0(y)\,dx\,dy$. Hence
$\Phi(y):=\int_{I_{T'}}K(|x-y|)F^*(x)\,dx$ satisfies $\int\Phi G_0=0$ for every such $G_0$, so
$\Phi=\alpha e^{y/2}+\beta e^{-y/2}$ on $(T',T)$. For $y>T'\ge x$, Lemma~\ref{lem:K} and Tonelli give
$\Phi(y)=2\sum_{k\ge0}c_ke^{-(2k+\frac12)y}$ with $c_k=\int F^*(x)e^{(2k+\frac12)x}dx$,
$|c_k|\le\|F^*\|_1e^{(2k+\frac12)T'}$; so $e^{y/2}\Phi(y)=2\sum_kc_kw^k$, $w=e^{-2y}$, is a convergent
power series for $0<w<e^{-2T'}$, and the identity extends to all $y>T'$. As $y\to\infty$, $\Phi\to0$,
so $\alpha=0$; then $2\sum_kc_kw^k\equiv\beta$, so $c_k=0$ for $k\ge1$ (the value $c_0=\beta/2$ is not
needed). With $u=e^{2x}$, $c_k=\int f(u)u^k\,du$ for
$f(u)=\frac12F^*(\frac12\log u)u^{-3/4}\in L^1[e^{-2T'},e^{2T'}]$, so all moments of $uf(u)$ vanish and
$uf=0$ by Weierstrass, whence $F^*=0$, contradicting $\|F^*\|=1$. For $\laminfP$ use Lemma~\ref{lem:parity}.

Strict decrease gives the asserted signs on either side of the supremum.
The negative index is positive exactly when the minimum is negative,
and the primitive threshold is the smaller sector threshold.
\end{proof}
Testing against the same two-moment functions removes the polar term
also in the unrestricted eigenvalue equation, giving strict decrease
of the unrestricted bottoms with or without that bounded term.

\section{The two Tate directions}\label{sec:complete}

\begin{lemma}[non-attainment for the complete form]\label{lem:Onat}
Let $T\ge\frac12\log2$ and $w\in\Wcl$. Then $\inf_\R f_w<\lamW$.
\end{lemma}
\begin{proof}
Weak duality gives $c=\inf f_w\le\lamW$. Equality would imply
$\int(f_w-c)|\Fh_\star|^2=0$ for a normalized ground state
by Proposition~\ref{prop:P} and compactness (Theorem~\ref{thm:O}).
The nonzero entire transform and continuity then force $f_w\equiv c$.
On $0<u<\min(2T,\log2)$ Fourier inversion gives
$2\cosh(u/2)=K(u)/2$, impossible near zero.
\end{proof}

\begin{theorem}[flat-top certificates and the complete bottom]\label{thm:Al}
\emph{(a)} If $f_w\ge0$ on $\R$ for some $w\in\Wcl$, regular or not, then $W_T\ge0$ on the form domain;
hence $\lamW\ge0$ and $\lamA\ge0$.
\emph{(b)} Let $T\ge\frac12\log2$. Then there is $w\in\Wcl$ with $f_w\ge0$ on $\R$ if and only if
$\lamW>0$, and $w$ can then be taken in $\Wcl^{\mathrm{reg}}$ with $\inf f_w>0$. Likewise, if $\lamP>0$
there are $w\in\Wcl^{\mathrm{reg}}$ and $\beta\in\R^2$ with $\inf_\tau[f_w+(\beta_0-2)\gh0+\beta_1\gh1]>0$.
\end{theorem}
\begin{proof}
(a) $(\nu_w,2,0)$ is dual feasible and $W_T\ge\inf f_w$ on $L^2(\IT)$ (the paragraph after
Definition~\ref{def:W}), and $W_T=\mathfrak a_T$ on $\PT$. (b) ``$\Leftarrow$'' and the regular choice:
Theorem~\ref{thm:Zexh}(i). ``$\Rightarrow$'': $\lamW\ge\inf f_w\ge0$ by weak duality, and $\lamW>\inf f_w$ by
Lemma~\ref{lem:Onat}. The last sentence is Theorem~\ref{thm:Zexh}(ii).
\end{proof}

\begin{definition}[the alignment inequality]\label{def:Al}
For $w\in\Wcl$ let $\mathrm{Al}_T(w)$ be the statement: for every $\tau\in\R$,
\[
-2\Rel\sum_{n\ge2}\frac{\Lambda(n)w(\log n)}{n^{\frac12+i\tau}}\ \ge\
-\Big[C_0+\int_0^\infty\big(1-w(u)\cos\tau u\big)K(u)\,du\Big]-2\int_\R w(u)\cosh\tfrac u2\cos(\tau u)\,du .
\]
The left side is the smoothed prime sum on the critical line, $-\sum_nv_nw(\log n)\cos(\tau\log n)$; by
\eqref{eq:fw}, $\mathrm{Al}_T(w)$ is exactly the statement $f_w\ge0$ on $\R$ of Theorem~\ref{thm:Al}.
\end{definition}

\begin{lemma}[the alignment inequality lives on a band]\label{lem:Alband}
For $w\in\Wcl$ let $\tau_{\mathrm{env}}(w)$ be the unique positive root of
$m_{\mathrm{lo}}(\tau)-\Aconst-\|\nu_w\|_{\TV}-2k_0(\tau)=0$, with $m_{\mathrm{lo}}$ as in Lemma~\ref{lem:mbounds}(ii)
and $k_0$ as in Lemma~\ref{lem:gbounds}. Then $f_w>0$ on $(\tau_{\mathrm{env}}(w),\infty)$, so $\mathrm{Al}_T(w)$ is
equivalent to its restriction to $|\tau|\le\tau_{\mathrm{env}}(w)$.
\end{lemma}
\begin{proof}
$m_{\mathrm{lo}}$ increases from $-\infty$ to $+\infty$ and $k_0$ decreases (Lemmas~\ref{lem:mbounds}, \ref{lem:gbounds}),
so the left side is strictly increasing and the root is unique. The envelope for $(\nu_w,2,0)$ gives $f_w\ge m_{\mathrm{lo}}-\Aconst-\|\nu_w\|_{\TV}-2k_0$, and $f_w$ is even.
\end{proof}

For $T\ge\frac12\log2$ the optimal uniform margin is $\lamW$
by Theorem~\ref{thm:Zexh}(i). For regular weights,
Theorem~\ref{thm:Z}(ii) identifies $\mathrm{Al}_T(w)$ with
$\sum_\rho\wh(\tau-\gamma_\rho)\ge0$.

\begin{remark}[the prime side of a primitive test is the Chebyshev remainder]\label{rem:Cheb}
Let $\psi_\Lambda(x)=\sum_{n\le x}\Lambda(n)$ and $r_\Lambda(u):=\psi_\Lambda(e^u)-e^u$. For real $F\in\PT\cap\mathcal D_T$,
\begin{equation}\label{eq:Cheb}
\mathfrak a_T(F)=C_0\|F\|^2+\Lev(\cF)+2\int_0^{2T}\cF(u)e^{-u/2}\,du-2\int_{(0,2T)}e^{-u/2}\cF(u)\,dr_\Lambda(u).
\end{equation}
The prime sum integrates $2e^{-u/2}\cF$ against $d\psi_\Lambda(e^u)$,
with zero boundary contribution since $\cF(2T)=0$. Split off $e^u du$.
Evenness and $\int_{-2T}^{2T}\cF(u)e^{u/2}du=A(F)B(F)=0$
give \eqref{eq:Cheb}. If $\lamA\ge0$, unit tests approaching the
infimum therefore make the last integral, with its factor two, agree
with the preceding three terms up to any error $\eta>\lamA$.
Corollary~\ref{cor:margin1} gives $\lamA\le4\cdot10^9Te^{-15.449T}$ for $T\ge1$.
\end{remark}

To pass from primitive to complete positivity at a fixed window it
suffices to consider real tests, since real and imaginary parts add.
For $j=\mathrm{ev},\mathrm{odd}$ let $\mathcal D_T^j$ denote the
corresponding parity sector, $\ell_j(F)=\langle F,\chi_j\rangle$
as in \S\ref{subsec:conv}, and $\varepsilon_j=+1,-1$, respectively.
Then $A=\ell_{\rm ev}+\ell_{\rm odd}$, $B=\ell_{\rm ev}-\ell_{\rm odd}$,
and $\mathcal Q_T^j=\ker(\ell_j|\mathcal D_T^j)$.
Write $\lambda_{\rm A}^j(T)$ for the bottom of $\mathfrak a_T$ on
$\mathcal Q_T^j$ and $\lambda^{\rm W}_{T,j}$ for that of $W_T$
on $\mathcal D_T^j$, using unit vectors in both cases.

\begin{theorem}[the parity Schur complement of the complete form]\label{thm:parityschur}
Fix $T>0$. In one parity sector abbreviate its moment, moment kernel,
primitive bottom and complete bottom by $\ell,\mathcal Q,\lambda,\omega$.
\emph{(o)} The sectors are orthogonal for $L^2$ and both forms;
$W_T(F)=\mathfrak a_T(F)+2\varepsilon\ell(F)^2$ on each. The
bottom of either form is the minimum of its two sector bottoms.

\emph{(a)} If zero belongs to the resolvent of the primitive operator,
there is a unique $h$ of moment one form-orthogonal to $\mathcal Q$.
Write $g=\mathfrak g_{\rm par}(T)=W_T(h)$. Then
\emph{(a1)} $W_T(G+xh)=\mathfrak a_T(G)+gx^2$ for $G\in\mathcal Q$.
\emph{(a2)} The negative index increases by $[g<0]$, and the radical
is $\R h$ for $g=0$ and trivial otherwise.
\emph{(a3)} For $\lambda>0$ both forms attain their moment-one
minima at $h$. \emph{(a4)} Put $b=g/\|h\|^2$. Then $\omega\le b$
and $\|h\|\ge1/\|\chi_{\rm par}\|$. For $g,\lambda>0$ one has
$2\omega\ge\min(\lambda,b)$, hence $b\le2\omega$ if $2\omega<\lambda$.

\emph{(b)} If $\lambda=0$, let $N$ be the primitive kernel. For
$z\in N$ write $\mathfrak a_T(z,\cdot)=\alpha(z)\ell$ on the sector.
If $\alpha\ne0$, then $\omega<0$. If $\alpha=0$ (the range condition),
$N$ lies in the radical of $W_T$ and (a) applies on $N^\perp$.

\emph{(c)} Put $\mathfrak G_T=\operatorname{diag}(\mathfrak g_{\rm ev},
\mathfrak g_{\rm odd})$ whenever both entries are defined.
For invertible primitive operator the negative indices of the primitive
form and $\mathfrak G_T$ add to that of $W_T$, and the nullity of
$W_T$ equals that of $\mathfrak G_T$. At every window $\lamW>0$ is
equivalent to $\lamA>0$ and $\mathfrak G_T\succ0$; $\lamW\ge0$ is
equivalent to $\lamA\ge0$, the range condition in each zero sector,
and $\mathfrak G_T\succeq0$.
\end{theorem}
\begin{proof}
Reflection preserves the Fourier multiplier and interchanges $A,B$,
giving (o) since $AB=\ell_{\rm ev}^2-\ell_{\rm odd}^2$. Choose
a smooth $H_0$ in the sector with $\ell(H_0)=1$. Subtracting
$\ell(F)H_0$ from smooth approximants gives density on the moment
kernel. The closed restricted form has compact resolvent by
Lemma~\ref{lem:compact}. If its operator is invertible, its negative
space is finite dimensional and on its positive space the spectral gap
gives coercivity in the shifted form norm. Thus finite-dimensional
inversion and Lax--Milgram give a unique $Y\in\mathcal Q$ with
$\mathfrak a_T(Y,G)=\mathfrak a_T(H_0,G)$ for every $G\in\mathcal Q$.
Here the right side is continuous by shifted form Cauchy--Schwarz.
Now $h=H_0-Y$, and subtracting $\ell(F)h$ gives (a1). Reading the
two blocks gives (a2) and (a3). Testing $h$ gives the upper bound
in (a4); the lower bound follows from
$\|G+xh\|^2\le2(\|G\|^2+x^2\|h\|^2)$.

For (b), compact resolvent gives a finite-dimensional zero eigenspace
and a gap above it. For $z\in N$ the functional vanishes on $\ker\ell$,
so $\alpha(z)=\mathfrak a_T(z,H_0)$. The energy on $H_0+xz$ is
$W_T(H_0)+2x\alpha(z)$, which takes negative values if $\alpha(z)\ne0$.
Otherwise, if $\alpha$ vanishes identically, orthogonal projection off
$N$ preserves the moment and gives the asserted reduction. Adding the
sector decompositions gives (c), using (a4) for strict positivity.
\end{proof}

\begin{proposition}[$\mathfrak G_T$ is the polar capacity]\label{prop:polarcap}
Under $\mathrm{WP}_T$ and the range condition in every zero sector,
let $E_T(\beta,\alpha)$ be the infimum of $W_T$ at real tests in
$\mathcal D_T$ with $B=\beta$, $A=\alpha$. This infimum is attained and
\[
 E_T(\beta,\alpha)=\tfrac14\bigl[\mathfrak g_{\rm ev}(\alpha+\beta)^2+
 \mathfrak g_{\rm odd}(\alpha-\beta)^2\bigr].
\]
Thus $\lamW\ge0$ exactly when $E_T\ge0$ on $\R^2$. In parity
coordinates the polar correction to the primitive Schur complement
is $D_{\rm pol}=\operatorname{diag}(2,-2)$.
\end{proposition}
\begin{proof}
The parity moments are $c=(\alpha+\beta)/2$ and $s=(\alpha-\beta)/2$.
The decomposition of Theorem~\ref{thm:parityschur} adds
$\mathfrak a_T(G)$, $G\in\PT$, to the displayed value; its minimum
is zero at $G=0$. The polar correction is $2c^2-2s^2=2\alpha\beta$.
\end{proof}

To compare the primitive multipliers with the polar value $(2,0)$, put
$A'(F)=\int tF(t)e^{t/2}dt$, $B'(F)=\int tF(t)e^{-t/2}dt$, and set
\[
 \mathfrak m(F)=\big(\Rel(A\bar B),\Rel(A'\bar B-A\bar B')\big),\qquad
 V_T(\beta)=\inf_{F\in\mathcal D_T,\ \|F\|=1}
       [\mathfrak a_T(F)+\beta\cdot\mathfrak m(F)] .
\]
The two components of $\mathfrak m$ have multipliers $\gh0,\gh1$ by
Theorem~\ref{thm:A}, and $V_T(2,0)=\lamW$ by Proposition~\ref{prop:P}.
All three forms add on the real and imaginary parts since the symbols
are even and $\Iml(\widehat u\,\overline{\widehat v})$ is odd for real $u,v$.

\begin{theorem}[the optimal multipliers]\label{thm:mult}
Fix $T>0$.
\emph{(i)} $V_T$ is finite, concave and continuous on $\R^2$, $V_T(\beta)\to-\infty$ as $|\beta|\to\infty$, and
\[
\lamP=\max_{\beta\in\R^2}V_T(\beta).
\]
Let $\mathcal M_T\subset\R^2$ be the set of maximizers; it is nonempty, compact and convex.
\emph{(ii)} If $(\nu_k,\beta^{(k)})$ are certificates of values $c_k\to\lamP$, then
$\operatorname{dist}(\beta^{(k)},\mathcal M_T)\to0$. Hence every such sequence satisfies $\beta^{(k)}\to(2,0)$ if and
only if $\mathcal M_T=\{(2,0)\}$, and $\lamP=\lamW$ if and only if $(2,0)\in\mathcal M_T$.
\emph{(iii)} Let $\mathcal Z_T$ be the set of real unit vectors $F\in\mathcal D_T$ with $W_T(F)=\lamW$; it is nonempty
and compact in $L^2(\IT)$. Then $(2,0)\in\mathcal M_T$ if and only if
$0\in\operatorname{conv}\{\mathfrak m(F):F\in\mathcal Z_T\}$.
\emph{(iv)} Suppose $\lamW<\lamA$. Then every $F\in\mathcal Z_T$ of definite parity has
$\Rel(A\bar B)=\ell_{\mathrm{ev}}(F)^2>0$ if $F$ is even and $\Rel(A\bar B)=-\ell_{\mathrm{odd}}(F)^2<0$ if $F$ is odd.
Consequently, if moreover $\lambda^{\mathrm W}_{T,\mathrm{ev}}\ne\lambda^{\mathrm W}_{T,\mathrm{odd}}$, then $\lamW<\lamP$, and
the multipliers of every maximizing sequence of certificates stay at a positive distance from $(2,0)$.
\emph{(v)} If $\lamW=\lamA$, then $\lamP=\lamW$ and $(2,0)\in\mathcal M_T$.
\end{theorem}
\begin{proof}
(i) By Cauchy--Schwarz the two moments of a unit vector are bounded
by $2\sinh T$ and $4T\sinh T$. Thus the affine functions defining
$V_T$ have uniformly bounded slope; their infimum is finite, concave
and Lipschitz since $\mathfrak a_T$ is bounded below.
The four finite-energy measures of Lemma~\ref{lem:W} realize a ball
of moments of radius $r>0$ by convex combination and the factorization
of Proposition~\ref{prop:weak}. Their energies have a common upper
bound $V_0$, so the moment $-r\beta/|\beta|$ gives
$V_T(\beta)\le V_0-r|\beta|$ for $\beta\ne0$. Feasible measures
have zero moments and hence $V_T(\beta)\le\lamP$. Every certificate
of value $c$ with multiplier $\beta$ gives $c\le V_T(\beta)$:
its annihilating measure integrates to zero against each autocorrelation.
Strong duality (Theorem~\ref{thm:D1}) now gives $\sup V_T=\lamP$.
Continuity, decay at infinity and concavity give a nonempty compact
convex set of maximizers.

(ii) Since $c_k\le V_T(\beta^{(k)})\le\lamP$, the multipliers
remain bounded and all their limit points maximize $V_T$.
For the converse at a fixed maximizer $\beta$, apply the separation
argument of Theorem~\ref{thm:D1} to the objective
$\sT+\beta\cdot\widehat g$ with only the support constraints.
This gives certificates with multiplier $\beta$ and values approaching
$V_T(\beta)$: the added symbol is bounded and the moment-zero
constraints are omitted on both sides. The last equivalence follows
from $V_T(2,0)=\lamW$.
(iii) Unit sublevels of $W_T$ are compact in $L^2$: the polar
term is bounded there, so Lemma~\ref{lem:compact} gives relative
compactness and lower semicontinuity gives closedness. Hence the
bottom is attained. Additivity on real and imaginary parts gives
a real unit minimizer and expresses the moment vector of each
complex unit minimizer as a convex combination of real ones.
Put $\beta_*=(2,0)$. For each minimizer $F$,
$V_T(\beta)\le\lamW+(\beta-\beta_*)\cdot\mathfrak m(F)$.
If their moment vectors have zero in their convex hull, averaging
gives $V_T(\beta)\le V_T(\beta_*)$. Otherwise compactness and
strict separation give $d\cdot\mathfrak m(F)\ge\delta>0$ on all
minimizers. By compactness of sublevels and continuity of the moments,
$d\cdot\mathfrak m(F)\ge\delta/2$ on some unit sublevel
$W_T\le\lamW+\eta$, $\eta>0$. Outside it,
$W_T(F)+t d\cdot\mathfrak m(F)\ge\lamW+\eta-t C_T|d|$,
where $C_T$ bounds the moments in (i). For sufficiently small $t>0$
both regions give a uniform strict improvement over $\lamW$,
so $\beta_*$ cannot maximize $V_T$.

(iv) If $\lamW<\lamA$, a unit minimizer of definite parity
has nonzero parity moment, since otherwise it would be primitive.
Its first moment coordinate is $A B=A^2$ in the even sector and
$-A^2$ in the odd sector. By parity additivity any component
in a sector of strictly larger bottom must vanish in a minimizer.
Thus when the sector bottoms differ the first coordinates of all
real minimizers have one strict sign. Apply (iii) and (ii).
(v) $\lamW\le\lamP\le\lamA$ (Theorem~\ref{thm:A}, Remark~\ref{rem:Kprim}) and (ii).
\end{proof}

\section{Lower bounds on finite windows}\label{sec:log5}

Let $T_5=\tfrac12\log5$. For $\tfrac12\log4<T\le T_5$ the
contacts are $2,3,4$ and $\Aconst=2.94197\ldots$.
For related radii and margins see \cite[Thm.~1.2]{Zhu2026}
and \cite[Thms.~A, B]{Liu2026}.

\begin{theorem}[Lower bounds at $T_5$]\label{thm:log5W}\label{thm:log5A}
At $T_5$ the even and odd primitive minima are at least
$2.5\cdot10^{-12}$ and $3.0\cdot10^{-10}$ respectively.
The complete even and odd minima are at least $10^{-18}$ and
$10^{-16}$. For $0<T\le T_5$ one has $\lamA\ge2.5\cdot10^{-12}$,
$\lamW\ge10^{-18}$ and $\mathfrak G_T\succ0$. These bounds are
unconditional; their proof is computer-assisted, the finite computations
of \S\ref{subsec:encllog5} being done in interval (ball) arithmetic.
\end{theorem}
\begin{proof}
In Proposition~\ref{prop:lowercert} use $c=0$ on the primitive spaces
and $c=-\varepsilon_{\rm par}/2$ on the complete parity spaces.
The ball enclosures in \S\ref{subsec:encllog5} give the four bounds.
Take the smaller parity minimum by Theorem~\ref{thm:parityschur}(o)
and restrict the support by Lemma~\ref{lem:monoT}.
Theorem~\ref{thm:parityschur}(c) gives the matrix inequality.
\end{proof}

\begin{corollary}[indeterminacy for $T\le\frac12\log5$]\label{cor:KEPlog5}
$\mathrm{KEP}(T)$ is indeterminate for every $0<T\le\frac12\log5$.
\end{corollary}
\begin{proof}
Theorems~\ref{thm:log5W} and \ref{thm:KEPdich} give the claim for
$\frac12\log2\le T\le T_5$. Below $\frac12\log2$, use the same weight as at
that endpoint: flatness on the larger window implies flatness on the smaller,
\eqref{eq:fw} is independent of $T$, and the perturbation in
Theorem~\ref{thm:KEP1}(ii) is supported outside $(-\log2,\log2)$.
\end{proof}

\begin{lemma}[a Ky Fan ceiling for the time--band limiting operator]\label{lem:KyFan}
Let $S_c=Q_TB_sQ_T$ on $L^2(\IT)$, where $Q_T$ is restriction to $\IT$ and $B_s$ the band projection
to $|\tau|<s$, and put $c=Ts$. Let $j_n$ be the spherical Bessel functions and
$G_k(x)=\sum_{j<k}(2j+1)j_j(x)^2\in[0,1]$. Then, in the orthonormal Legendre basis
$e_j(t)=\sqrt{(2j+1)/2T}\,P_j(t/T)$,
\[
\langle S_ce_j,e_j\rangle=\frac{2(2j+1)}\pi\int_0^cj_j(x)^2dx,\qquad
\sum_{j\ge0}\langle S_ce_j,e_j\rangle=\frac{2c}\pi=\operatorname{tr}S_c ,
\]
and for every $k\ge0$
\[
\lambda_k(S_c)\ \le\ \Lambda_k(c):=\frac2\pi\int_0^c\big(1-G_k(x)\big)\,dx .
\]
In particular $\Lambda_0(c)=2c/\pi$ is the trace and $\Lambda_k(c)=O(c^{2k+1})$ as $c\to0$.
\end{lemma}
\begin{proof}
The Fourier--Legendre identity gives
$|\widehat e_j(\tau)|^2=2T(2j+1)j_j(\tau T)^2$. Plancherel gives
the diagonal entry and $\sum_j(2j+1)j_j(x)^2=1$ gives the trace.
Positivity and Ky Fan's principle give
$\lambda_k\le\sum_{j\ge k}\lambda_j\le\operatorname{tr}S_c-
\sum_{j<k}\langle S_ce_j,e_j\rangle=\Lambda_k(c)$.
\end{proof}

\begin{corollary}[quadrature-free form]\label{cor:P1}
$|j_n(x)|\le x^n/(2n+1)!!$ for $x\ge0$, hence
$\lambda_k(S_c)\le\Lambda_k(c)\le\Lambda^*_k(c):=\frac2\pi\sum_{j\ge k}c^{2j+1}/\big((2j+1)!!\big)^2$, a convergent
series of rationals in $c$: the bound is evaluated by arithmetic alone, with no quadrature.
\end{corollary}


\begin{lemma}[bordered inertia]\label{lem:border}
Let $q$ be a real quadratic form on a real vector space $V$, $f$ a linear functional on
$V$ and $v_0\in V$ with $f(v_0)=1$ and $q(\cdot,v_0)$ finite. On $V\oplus\R$ put
$b(v,\xi)=q[v]+2\xi f(v)$. Then $n_-(b)=n_-(q|_{\ker f})+1$, $n_-$ being the largest dimension of a
subspace on which the form is negative definite.
\end{lemma}
\begin{proof}
Write $v=w+tv_0$ with $w\in\ker f$. Then $b=q[w]+2tq(w,v_0)+t^2q[v_0]+2\xi t$, and the invertible
substitution $\eta=\xi+q(w,v_0)$ gives $b=q[w]+(t^2q[v_0]+2t\eta)$: a $b$-orthogonal sum of
$q|_{\ker f}$ and a binary form of determinant $-1$, whose inertia is $(1,1)$.
\end{proof}

\subsection{The certificates at \texorpdfstring{$T=\frac12\log5$}{T = (1/2) log 5}}\label{subsec:encllog5}
Use the Legendre basis and the band projection of Lemma~\ref{lem:KyFan} in each real parity space. Here the lower form is a frequency truncation of the full symbol $\sT$, and the operator subtracted from a
multiple of the identity is indefinite.

\emph{The lower form.} Fix $\tau_N>0$ and a number $s_\infty\le m(\tau_N)-\Aconst$. Since $m$ is increasing
with $m(0)=C_0$ and the comb is at least $-\Aconst$ (Lemma~\ref{lem:mbounds}(i),(iii)), the even function
$\underline s$ equal to $\sT$ on $|\tau|\le\tau_N$ and to $s_\infty$ on $|\tau|>\tau_N$ satisfies
$\underline s\le\sT$. Write $g=s_\infty-\underline s=(s_\infty-\sT)\mathbf 1_{|\tau|<\tau_N}=g_+-g_-$ with
$g_\pm\ge0$, and $\mathcal G_{(\pm)}=Q_Tg_{(\pm)}(D)Q_T$. Then
\[
\mathfrak a_T(F)\ \ge\ s_\infty\|F\|^2-\langle\mathcal GF,F\rangle,\qquad
0\le\mathcal G_\pm\le\omega_\pm S_{c_*},\quad c_*=T\tau_N,
\]
with $\omega_+=s_\infty-C_0+\Aconst$ and $\omega_-=m(\tau_N)+\Aconst-s_\infty$. For $\mu\in\R$ and
$c\in\{0,\frac12,-\frac12\}$ put, on $L^2_{\mathrm{par}}(\IT)\oplus\R$,
\[
\mathfrak b_{\mu,c}(F,\xi)=(s_\infty-\mu)\|F\|^2-\langle\mathcal GF,F\rangle+2\xi\,\ell_{\mathrm{par}}(F)+c\,\xi^2 .
\]

\begin{proposition}[the lower certificate]\label{prop:lowercert}
Let $k_1<k_2$, let $E$ be the span of the $e_j$, $j<k_1$, of parity $\mathrm{par}$, $Y$ its orthogonal
complement in $L^2_{\mathrm{par}}(\IT)$, and $P_Y$ the orthogonal projection onto $Y$.
\emph{(a)} On $Y$, $\mathfrak b_{\mu,c}\ge q_Y:=s_\infty-\mu-\omega_+\Lambda_{k_1}(c_*)$.
\emph{(b)} Suppose $q_Y>0$, and let $\mathsf R$ be the symmetric matrix on $E\oplus\R$ with
\[
\begin{gathered}
\mathsf R_{ij}=(s_\infty-\mu)\delta_{ij}-\langle\mathcal Ge_i,e_j\rangle
-\frac{\langle P_Y\mathcal Ge_i,P_Y\mathcal Ge_j\rangle}{q_Y},\\
\mathsf R_{i\xi}=\langle e_i,\chi_{\mathrm{par}}\rangle+\frac{\langle P_Y\mathcal Ge_i,P_Y\chi_{\mathrm{par}}\rangle}{q_Y},\qquad
\mathsf R_{\xi\xi}=c-\frac{\|P_Y\chi_{\mathrm{par}}\|^2}{q_Y}.
\end{gathered}
\]
Then $n_-(\mathfrak b_{\mu,c})\le n_-(\mathsf R)$, and the same holds for any symmetric $\mathsf R'\le\mathsf R$.
\emph{(c)} If $n_-(\mathsf R')\le1$ with $c=0$, then $\mathfrak a_T\ge\mu$ on the unit sphere of
$\PT^{\mathrm{par}}$. If $n_-(\mathsf R')=0$ with $c=\frac12$ and $\mathrm{par}=\mathrm{odd}$, or $n_-(\mathsf R')\le1$ with
$c=-\frac12$ and $\mathrm{par}=\mathrm{ev}$, then $W_T\ge\mu\|\cdot\|^2$ on $\mathcal D^{\mathrm{par}}_T$.
\end{proposition}
\begin{proof}\sloppy
(a) $-P_Y\mathcal GP_Y\ge-P_Y\mathcal G_+P_Y$ and
$\|P_Y\mathcal G_+P_Y\|\le\omega_+\operatorname{tr}(P^\perp S_{c_*}P^\perp)=\omega_+\Lambda_{k_1}(c_*)$.
(b) For $i<k_1$ and $y\in Y$, $\mathfrak b_{\mu,c}(e_i,y)=-\langle P_Y\mathcal Ge_i,y\rangle$ and the border pairs
with $y$ through $+\langle P_Y\chi_{\mathrm{par}},y\rangle$. Completing the square in the $Y$ variable gives $\mathfrak b_{\mu,c}\cong\mathfrak s\oplus\mathfrak b_{\mu,c}|_Y$ with
$\mathfrak b_{\mu,c}|_Y\ge q_Y>0$, so $n_-(\mathfrak b_{\mu,c})=n_-(\mathfrak s)$, and the subtracted term of the Schur
complement $\mathfrak s$ is at most $q_Y^{-1}$ times the Gram matrix of the $Y$-columns
$-P_Y\mathcal Ge_i$ and $P_Y\chi_{\mathrm{par}}$; hence $\mathfrak s\ge\mathsf R\ge\mathsf R'$.
(c) Put $\mathfrak q_\mu(F)=\mathfrak b_{\mu,0}(F,0)$, so that $\mathfrak a_T-\mu\|\cdot\|^2\ge\mathfrak q_\mu$. For
$c=0$, Lemma~\ref{lem:border} gives $n_-(\mathfrak q_\mu|_{\ker\ell_{\mathrm{par}}})=n_-(\mathfrak b_{\mu,0})-1=0$. For
$c=\frac12$, $\mathfrak b_{\mu,1/2}=\mathfrak q_\mu-2\ell_{\mathrm{par}}^2+\frac12(\xi+2\ell_{\mathrm{par}})^2$, so
$\mathfrak q_\mu-2\ell_{\mathrm{odd}}^2\ge0$, and $W_T=\mathfrak a_T-2\ell_{\mathrm{odd}}^2$ on odd $F$
(Theorem~\ref{thm:parityschur}(o)). For $c=-\frac12$,
$\mathfrak b_{\mu,-1/2}=\mathfrak q_\mu+2\ell_{\mathrm{par}}^2-\frac12(\xi-2\ell_{\mathrm{par}})^2$; the substitution
$\eta=\xi-2\ell_{\mathrm{par}}(F)$ is a congruence, so
$n_-(\mathfrak q_\mu+2\ell_{\mathrm{ev}}^2)=n_-(\mathfrak b_{\mu,-1/2})-1=0$, and $W_T=\mathfrak a_T+2\ell_{\mathrm{ev}}^2$ on even
$F$. The statements pass from $L^2_{\mathrm{par}}(\IT)$ to the form domain because $\mathfrak q_\mu$ is bounded.
\end{proof}

\emph{The computation.} For $i-l$ even,
\[
\langle\mathcal Ge_i,e_l\rangle=\frac{2T}\pi\sqrt{(2i+1)(2l+1)}\,(-1)^{(l-i)/2}\int_0^{\tau_N}g(\tau)j_i(\tau T)j_l(\tau T)\,d\tau
\]
(proof of Lemma~\ref{lem:KyFan}), and the entries with $i-l$ odd vanish. The integral is computed by composite
Gauss--Legendre quadrature with $46$ nodes on panels of length $\frac12$, with the remainder bound for
functions analytic in a Bernstein ellipse \cite[Thm.~19.3]{Trefethen2013}, $\rho=3.2$. The ellipse lies
within distance $0.439<\frac12$ of the panel, while the poles of $\psi(\frac14\pm\frac{i\tau}2)$ lie at
$|\Iml\tau|\ge\frac12$. The modulus on the ellipse is bounded with $|j_n(z)|\le e^{|\Iml z|}$,
$|\cos(z\log n)|\le\cosh(\Iml z\log n)$ and, for $\psi$, the recurrence $\psi(w)=\psi(w+8)-\sum_{k<8}(w+k)^{-1}$
together with Binet's bound $|\psi(u)|\le\log|u|+\frac\pi2+\frac1{2\Rel u}+\frac1{12(\Rel u)^2}$ for
$\Rel u\ge8$. The resulting remainder, $1.5\cdot10^{-43}$ per entry (below $2\cdot10^{-42}$ under either
point-count convention for the rule), is added as a ball to every entry. The inner products over $Y$ are
exact sums over $k_1\le l<k_2$ of the given parity plus tail boxes: by Cauchy--Schwarz for the positive
operators $\mathcal G_\pm$, $\|P^\perp_{k_2}\mathcal G_\pm e\|^2\le\|P^\perp_{k_2}\mathcal G_\pm P^\perp_{k_2}\|
\langle\mathcal G_\pm e,e\rangle\le\omega_\pm^2\Lambda_{k_2}(c_*)$, so
$\|P^\perp_{k_2}\mathcal Ge\|\le(\omega_++\omega_-)\Lambda_{k_2}(c_*)^{1/2}$, and
$\|P^\perp_{k_2}\chi_{\mathrm{par}}\|^2=\sinh T\pm T-\sum_{l<k_2}\langle e_l,\chi_{\mathrm{par}}\rangle^2$ with
$\langle e_l,\chi_{\mathrm{par}}\rangle=\sqrt{(2l+1)/2T}\,2T\,i_l(T/2)$. The tails $\Lambda_k$ are bounded by
$\Lambda^*_k$ of Corollary~\ref{cor:P1}. Everything is midpoint--radius ball arithmetic at $640$ bits
(Arb \cite{Johansson2017}, through \texttt{python-flint}), including the digamma and Bessel values and the
Gauss--Legendre nodes and weights; every Bessel ball is asserted to have radius below $10^{-60}$, and
comparisons are ball predicates.

For the inertia test, replacing $\|P_Y\chi_{\rm par}\|^2/q_Y$ by its
upper bound $\delta'$ gives $\mathsf R'\le\mathsf R$. Interval Cholesky
proves $n_-(\mathsf R')=0$ in the positive cases. Otherwise choose a
Householder matrix $H$ from the midpoint of the $\xi$ column and certify
$\|H^tH-I\|_\infty<1$. Congruence by $\operatorname{diag}(H,1)$ gives
a block $B$ on $(e_0,\xi)$ with certified $\det B<0$ and positive Schur
complement by interval Cholesky. Inertia additivity gives one negative
eigenvalue \cite{Haynsworth1968}.

\emph{Data.} $T=\frac12\log5$, $\tau_N=240$, $c_*=193.13$, $k_1=280$, $k_2=330$ ($\dim E=140$ in each parity).
From $m(\tau_N)=3.64276\ldots$: $s_\infty=0.700787608\ldots$, $\omega_+=9.0150$, $\omega_-=5.8840$,
$\Lambda^*_{k_1}(c_*)\le4.41\cdot10^{-20}$, $\Lambda^*_{k_2}(c_*)\le3.87\cdot10^{-70}$, $q_Y\ge0.70078$, tail
box $(\omega_++\omega_-)\Lambda^*_{k_2}(c_*)^{1/2}\le2.94\cdot10^{-34}$; the entry radii of $\mathsf R$ are at
most $4.4\cdot10^{-41}$. The verified cases are:
\begin{center}\small
\begin{tabular}{@{}llccc@{}}
\toprule
form & sector & certified $\mu$ & minimal pivot at that $\mu$ & fails at\\
\midrule
$\mathfrak a_T|\PT$ ($c=0$) & even & $2.5\cdot10^{-12}$ & $8.47\cdot10^{-3}$ & $3\cdot10^{-12}$\\
$\mathfrak a_T|\PT$ ($c=0$) & odd & $3.0\cdot10^{-10}$ & $2.03\cdot10^{-2}$ & $4\cdot10^{-10}$\\
$W_T$ ($c=-\frac12$) & even & $10^{-18}$ & $4.47\cdot10^{-4}$ & $10^{-16}$\\
$W_T$ ($c=\frac12$) & odd & $10^{-16}$ & $8.14\cdot10^{-14}$ & $10^{-14}$\\
\bottomrule
\end{tabular}
\end{center}
\subsection{The margin decays}\label{subsec:margin}
An upper bound related to Zhu's conjectural law $-\log\lamW\asymp e^{2T}$ \cite[Conj.~12.1]{Zhu2026}, namely that
the margin tends to $0$ at an explicit rate, can be proved from three facts about the zeros of $\zeta$.
These are also the zero inputs of Theorem~\ref{thm:power-window}.
\begin{itemize}
\item[(Z1)] $\rho_1=\frac12+i\gamma_1$, $\gamma_1=14.134725141734\ldots$, is a simple zero on the critical line,
and $\zeta$ has no non-trivial zero $\rho\ne\rho_1,\bar\rho_1$ with $|\Iml\rho|\le21$
\cite{RosserYoheSchoenfeld1968,PlattTrudgian2021}. Only $\gamma_1\in[14.13,14.14]$ and the gap $(\gamma_1,21]$
are used.
\item[(Z2)] Every non-trivial zero $\rho$ with $|\Iml\rho|\le H_{\mathrm{PT}}:=3\cdot10^{12}$ lies on the
critical line \cite{PlattTrudgian2021}.
\item[(Z3)] For $t\ge2$, $\big|N(t)-\frac t{2\pi}\log\frac t{2\pi e}-\frac78\big|\le0.137\log t+0.443\log\log t+1.588$,
where $N(t)$ counts the zeros with $0<\Iml\rho\le t$ with multiplicity \cite[Thm.~19]{Rosser1941}.
\end{itemize}
Besides these, we use that the non-trivial zeros lie in $0<\Rel s<1$, and Lemma~\ref{lem:roughEF}.
\begin{lemma}[a crude zero count]\label{lem:Ncount}
$N(t)\le\frac t{2\pi}\log t$ for $t\ge21$.
\end{lemma}
\begin{proof}
$\frac t{2\pi}\log\frac t{2\pi e}=\frac t{2\pi}\log t-\frac{1+\log2\pi}{2\pi}t$ and
$\frac{1+\log2\pi}{2\pi}=0.45166\ldots$, so by (Z3)
$\frac t{2\pi}\log t-N(t)\ge0.45166\,t-\frac78-1.588-0.137\log t-0.443\log\log t$, which is $6.11\ldots$ at
$t=21$ and has derivative $0.45166-0.137/t-0.443/(t\log t)>0$ for $t\ge21$.
\end{proof}

\begin{lemma}[zero sums by counting]\label{lem:zerosum}
For $p\ge2$ and $H\ge21$, with the sum over the ordinates $\gamma=\Iml\rho>0$ of the non-trivial zeros
counted with multiplicity,
\[
\Sigma_p(H):=\sum_{\gamma>H}\gamma^{-p}\ \le\ \frac p{2\pi}\,H^{1-p}\Big[\frac{\log H}{p-1}+\frac1{(p-1)^2}\Big].
\]
\end{lemma}
\begin{proof}
$\sum_{\gamma>H}\gamma^{-p}=\int_{H^+}^\infty t^{-p}dN(t)=\big[t^{-p}N(t)\big]_H^\infty+p\int_H^\infty t^{-p-1}N(t)\,dt$;
the term at $\infty$ vanishes and the one at $H$ is $-H^{-p}N(H)\le0$ by Lemma~\ref{lem:Ncount}, which
also bounds the integral by $\frac p{2\pi}\int_H^\infty t^{-p}\log t\,dt=\frac p{2\pi}H^{1-p}\big[\frac{\log H}{p-1}+\frac1{(p-1)^2}\big]$.
\end{proof}

\begin{lemma}[the test family]\label{lem:testfam}
Fix $T>0$ and an integer $k\ge5$; let $b_k=\frac k2\mathbf 1_{[-1/k,1/k]}$ and $\omega_k=b_k^{*k}$ (the
$k$-fold convolution).
\emph{(i)} $\omega_k\ge0$, $\int\omega_k=1$, $\supp\omega_k\subseteq[-1,1]$, $\omega_k\in C^{k-2}(\R)$ with
$\omega_k^{(j)}(\pm1)=0$ for $j\le k-2$, and $\|\omega_k\|_2^2\ge\frac12$.
\emph{(ii)} $\widehat\omega_k(x)=\big(\sin(x/k)/(x/k)\big)^k$; $|\widehat\omega_k(x)|\le\min\{1,(k/|x|)^k\}$ for real
$x$, and $|\widehat\omega_k(z)|\le(k/|\Rel z|)^ke^{|\Iml z|}$ for complex $z$.
\emph{(iii)} Put $G(t)=\omega_k(t/T)$, $\varpi(c)=(c^2-\frac14)(c^2+\gamma_1^2)$, $\partial=d/dt$, and
\[
F:=\varpi(\partial)G=G''''+(\gamma_1^2-\tfrac14)G''-\tfrac{\gamma_1^2}4G .
\]
Then $F$ is real, supported in $[-T,T]$, bounded, and $F\in\PT$: four integrations by parts, the
boundary terms carrying $G^{(j)}(\pm T)=0$ for $j\le3$, give $\int Fe^{ct}dt=\varpi(c)\int Ge^{ct}dt$
for every $c\in\C$, and $\varpi(\pm\frac12)=0$.
\emph{(iv)} $\Fh(\tau)=(\tau^2+\frac14)(\tau^2-\gamma_1^2)\,T\,\widehat\omega_k(T\tau)$,
$\Phi_F(s)=\varpi(s-\frac12)\Phi_G(s)$ and $\Phi_G(\frac12+\beta+i\gamma)=T\,\widehat\omega_k\big(T(\gamma-i\beta)\big)$;
in particular $\Phi_F(\rho_1)=\Phi_F(\bar\rho_1)=0$.
\emph{(v)} With $Y:=T\gamma_1/\sqrt2$ and $\varepsilon_k(Y):=\frac1\pi k^{2k}Y^{1-2k}/(2k-1)$,
$\|F\|^2\ge\frac{\gamma_1^4}{64}\,T\,\big(\frac12-\varepsilon_k(Y)\big)$ whenever $\varepsilon_k(Y)<\frac12$.
\end{lemma}
\begin{proof}
(i) $\omega_k$ is a spline of degree $k-1$ supported in $[-1,1]$ and vanishing there to order $k-1$;
$1=(\int_{-1}^1\omega_k)^2\le2\int\omega_k^2$. (ii) $\widehat b_k(x)=\sin(x/k)/(x/k)$; $|\sin y/y|\le\min(1,1/|y|)$
on $\R$, and $|\sin(a+ib)|\le\cosh b\le e^{|b|}$. (iii) $G\in C^3$ with $G'''$ absolutely continuous
and $G''''\in L^\infty$ ($k=5$), and $G^{(j)}(\pm T)=0$ for $j\le3$; $\int G''e^{ct}=c^2\int Ge^{ct}$.
(iv) is (iii) with $c=i\tau$, resp.\ $c=s-\frac12$, and $\varpi(\pm i\gamma_1)=0$. (v) On
$|\tau|\le\gamma_1/\sqrt2$ one has $(\tau^2+\frac14)^2\ge\frac1{16}$ and $(\tau^2-\gamma_1^2)^2\ge\gamma_1^4/4$,
so $\|F\|^2=\frac1{2\pi}\int|\Fh|^2\ge\frac{\gamma_1^4}{64}\,T\,\frac1{2\pi}\int_{|x|\le Y}|\widehat\omega_k|^2$, and
$\frac1{2\pi}\int_{|x|\le Y}|\widehat\omega_k|^2\ge\|\omega_k\|_2^2-\frac1\pi\int_Y^\infty(k/x)^{2k}dx\ge\frac12-\varepsilon_k(Y)$.
\end{proof}

\begin{theorem}[the margin decays]\label{thm:margin}
Let $T>0$, $k\ge5$ an integer, $p:=2k-8$, $Y:=T\gamma_1/\sqrt2$, and suppose $\varepsilon_k(Y)<\frac12$. Then
\begin{equation}\label{eq:Bmargin}
\lamW\ \le\ \lamP\ \le\ \lamA\ \le\ \mathcal B(T,k):=
\frac{2\cdot1.0012\cdot T^2(k/T)^{2k}\big[\Sigma_p(21)+e^T\,\Sigma_p(H_{\mathrm{PT}})\big]}
{(14.13^4\,T/64)\,\big(\frac12-\varepsilon_k(Y)\big)},
\end{equation}
with $\Sigma_p(H)$ bounded by Lemma~\ref{lem:zerosum}. Consequently $\limsup_{T\to\infty}\lamA\le0$, and the
same for $\lamP$ and $\lamW$.
\end{theorem}
\begin{proof}
Take $F$ of Lemma~\ref{lem:testfam}. (1) $F\in\PT$, so $\lamA\le\mathfrak a_T(F)/\|F\|^2$, and
$\lamW\le\lamP\le\lamA$ by Theorem~\ref{thm:A} and Remark~\ref{rem:Kprim}.
(2) For $|\tau|\ge21$, $(\tau^2+\frac14)^2\le1.0012\,\tau^4$ and $(\tau^2-\gamma_1^2)^2\le\tau^4$, so
$|\Fh(\tau)|^2\le1.0012\,\tau^8T^2(k/(T|\tau|))^{2k}=O(|\tau|^{8-2k})$ with $8-2k\le-2$: $F$ is in the
form domain, which is hypothesis (a) of Lemma~\ref{lem:roughEF}; hypothesis (b) is verified in step (6).
Hence $\mathfrak a_T(F)=\sum_\rho\Phi_F(\rho)\Phi_F(1-\rho)$.
(3) $\rho_1,\bar\rho_1$ contribute $0$ (Lemma~\ref{lem:testfam}(iv)); by (Z1) no other zero has $|\Iml\rho|\le21$.
(4) A zero with $21<|\gamma|\le H_{\mathrm{PT}}$ lies on the line by (Z2) and contributes
$|\Fh(\gamma)|^2\le1.0012\,T^2(k/T)^{2k}|\gamma|^{8-2k}$.
(5) A zero $\rho=\frac12+\beta+i\gamma$ with $|\gamma|>H_{\mathrm{PT}}$ has $|\beta|<\frac12$;
$|\varpi(\beta+i\gamma)|\le(\gamma^2+\beta^2+\frac14)(\gamma^2+\beta^2+\gamma_1^2)\le(1+10^{-22})\gamma^4$, $\varpi$ is even,
and $|\Phi_G(\rho)|,|\Phi_G(1-\rho)|\le T(k/(T|\gamma|))^ke^{T|\beta|}$ by Lemma~\ref{lem:testfam}(ii),(iv);
so $|\Phi_F(\rho)\Phi_F(1-\rho)|\le1.0012\,T^2(k/T)^{2k}|\gamma|^{8-2k}e^T$.
(6) Summing (4)--(5) over $\pm\gamma$ with Lemma~\ref{lem:zerosum} ($p=2k-8\ge2$) gives
$\sum_\rho|\Phi_F(\rho)\Phi_F(1-\rho)|\le2\cdot1.0012\,T^2(k/T)^{2k}\big[\Sigma_p(21)+e^T\Sigma_p(H_{\mathrm{PT}})\big]<\infty$,
which is hypothesis (b) of Lemma~\ref{lem:roughEF} and bounds $\mathfrak a_T(F)$.
(7) Divide by Lemma~\ref{lem:testfam}(v) with $\gamma_1^4\ge14.13^4$. The limit statement is
Corollary~\ref{cor:margin1}.
\end{proof}

\begin{corollary}[closed form]\label{cor:margin1}
For every $T\ge1$,
\begin{equation}\label{eq:closedmargin}
\lamW\ \le\ \lamP\ \le\ \lamA\ \le\ 4\cdot10^9\,T\,e^{-15.449\,T},\qquad 15.449<\tfrac{42}e=15.4494\ldots\footnotemark
\end{equation}
\footnotetext{The exponent $\frac{42}{e}$ belongs to the test family of Lemma~\ref{lem:testfam}: one factor removes only the first zero, so the cutoff is $21$, just below the second ordinate $\gamma_2=21.022\ldots$, and $\frac{42}e$ comes from optimizing $k$ in $(k/21T)^{2k}$. It is a property of that family, not a sharp rate for $\lamA$.}
\end{corollary}
\begin{proof}
Take $k:=\lceil21T/e\rceil\ge8$ and write $k=21T/e+\theta$, $\theta\in[0,1)$. Then
$(k/21T)^{2k}=\exp\big(2k[\log(1+\theta e/21T)-1]\big)\le\exp(-42T/e+2\theta^2e/21T)\le e^{-15.4494\,T+0.259}$.
Next $(k/T)^{2k}\Sigma_p(21)\le21^9(k/21T)^{2k}\cdot\frac1{2\pi}\frac{2k-8}{2k-9}\big[\log21+\frac1{2k-9}\big]\le21^9\cdot0.5797\,(k/21T)^{2k}$
for $k\ge8$. The bound for the $H_{\mathrm{PT}}$-term is at most $10\,e^T(21/H_{\mathrm{PT}})^{2k-9}$ times that for the first, with
$21/H_{\mathrm{PT}}=7\cdot10^{-12}$ and $2k-9\ge15.449\,T-9$, hence less than $10^{-70}$ of it for $T\ge1$.
Finally $\varepsilon_k(Y)\le0.01566$ for all $T\ge1$: for fixed $k$, $\varepsilon_k$ decreases in $T$, so it
suffices to evaluate it at the left end of each $k$-range, namely $T=1$ for $k=8$ ($\varepsilon_8=0.0060$)
and $T=(k-1)e/21$ for $k\ge9$, where the largest value is $\varepsilon_9=0.01566$ at $T=8e/21$, the values
decreasing thereafter. So $\frac12-\varepsilon_k\ge0.484$, and
$\mathcal B(T,k)\le\frac{2\cdot1.0012\cdot0.5797\cdot21^9e^{0.259}}{(14.13^4/64)\cdot0.484}\,T\,e^{-15.4494\,T}(1+10^{-70})
\le3.97\cdot10^9\,T\,e^{-15.4494\,T}$.
\end{proof}

For the complete minimum Mori reports the upper bound
$7.58\cdot10^{24}\mu^5(\log\mu)^3e^{-4\pi\mu}$ for $\mu=e^{2T}\ge50$,
using finite-height zero data \cite{MoriUpper}. Corollary~\ref{cor:margin1}
also bounds the primitive minimum.

By monotonicity (Lemma~\ref{lem:monoT}), the three minima have limits
$L^{\rm W}\le L^{\rm prim}\le L^{\rm A}$ in $[-\infty,0]$, since
\eqref{eq:closedmargin} tends to zero. In particular \eqref{eq:target}
is equivalent to $L^{\rm A}=0$.

\subsection{The window $T\le2.8022$}\label{subsec:power-window}
In this subsection we prove a lower bound on the window $T\le\frac{14011}{5000}$. Unlike
Theorem~\ref{thm:log5W} it uses information about the zeros of $\zeta$, namely (Z1)--(Z3) of
\S\ref{subsec:margin}, and in particular the numerical verification (Z2) that the zeros up to height
$H_{\mathrm{PT}}=3\cdot10^{12}$ lie on the critical line. The idea is to apply the arithmetic expression
of the form to a filtered copy of the test, whose transform is concentrated where the archimedean
multiplier dominates the prime operator, and the zero expression to the remainder. Throughout,
\[
\begin{gathered}
T_1=\tfrac{14011}{5000},\quad U=\tfrac{57}{20},\quad R=2014985964,\quad H=2673507086718,\quad b=13,\\
M=9115062736,\quad d=M/2,\quad n_\chi=120122483651,\quad D=120110026806 .
\end{gathered}
\]
The finite computations used below (the Schur comparison of Lemma~\ref{lem:pw-schur} and the scalar
inequalities marked \emph{(computation)}) are reproduced in the notebook
\texttt{positivity-radial-intersection-form.ipynb} accompanying the paper: the Schur comparison in exact
integer arithmetic, the scalar inequalities in $192$-bit Arb ball arithmetic
\cite{Johansson2017} with python-flint~0.9.0.

\begin{theorem}[the window $T\le2.8022$]\label{thm:power-window}
Assume (Z1)--(Z3). Then for every $0<T\le T_1=2.8022$ and every $F\in\mathcal D_T$,
\begin{equation}\label{eq:power-window}
 W_T(F)\ge e^{-120110026807}\|F\|_2^2 .
\end{equation}
Consequently $\lamW\ge e^{-120110026807}$, and $\lamA\ge e^{-120110026807}$, for $0<T\le2.8022$.
\end{theorem}

(Z2) enters through the vanishing of the corrector of Lemma~\ref{lem:pw-zero} below for zeros of height at
most $H_{\mathrm{PT}}$, and (Z1), (Z3) through the local count of Lemma~\ref{lem:pw-count}. The
statement for $\lamA$ follows from that for $\lamW$ since $\lamW\le\lamA$.

\begin{lemma}[a local zero count]\label{lem:pw-count}
Assume (Z1) and (Z3). For every integer $j$ the number of non-trivial zeros $\rho$, counted with
multiplicity, with $j\le\Iml\rho<j+1$ is at most $100\log(3+|j|)$.
\end{lemma}
\begin{proof}
Put $N(t)=0$ for $t\le0$. We show $N(y+1)-N(y-1)\le100\log(3+y)$ for $y\ge0$. If $y+1<14.13$ the left
side is $0$ by (Z1). Otherwise $y-1\ge12$, and (Z3) applies at $y\pm1$. With
$f(t)=\frac t{2\pi}\log\frac t{2\pi e}$, $f'(t)=\frac1{2\pi}\log\frac t{2\pi}$ is increasing, so
$f(y+1)-f(y-1)\le2f'(y+1)\le\frac1\pi\log(y+1)$, and the error term of (Z3) is increasing in $t\ge2$. Hence
$N(y+1)-N(y-1)\le(\frac1\pi+0.274)\log(y+1)+0.886\log\log(y+1)+3.176\le100\log(3+y)$. For $j\ge0$ the
interval $[j,j+1)$ lies in $(j-1,j+1]$; for $j<0$, conjugation maps the zeros with ordinate in $[j,j+1)$
to those with ordinate in $(|j|-1,|j|]\subset(|j|-1,|j|+1]$, with multiplicities; there is no zero with
ordinate $0$, by (Z1).
\end{proof}

\begin{lemma}[the zero side with a cutoff]\label{lem:pw-zero}
Assume (Z1)--(Z3). Let $\nu_h$ be the uniform probability measure on $[-h,h]$ and
$\chi=\mathbf 1_{[-8,8]}*\nu_{1/4}*\nu_{1/4}*\nu_{1/n_\chi}^{*n_\chi}$. For a non-trivial zero
$\rho=\frac12+\delta+i\gamma$ put
\[
k_\rho(\tau)=\int_\R\chi(t)\big(\cosh(\delta t)-1\big)e^{i(\gamma-\tau)t}\,dt,\qquad
\mathcal C(\tau)=\frac1{2\pi}\sum_\rho k_\rho(\tau).
\]
\emph{(i)} $k_\rho$ is real, $k_\rho=0$ when $\delta=0$, in particular when $|\gamma|\le H_{\mathrm{PT}}$, and the series
defining $\mathcal C$ converges absolutely with $|\mathcal C(\tau)|\le10^9\log(3+|\tau|)$.
\emph{(ii)} If $|\tau|\le H$, then $|\mathcal C(\tau)|<e^{-D-2}$.
\emph{(iii)} For every real $F\in C_c^\infty(I_U)$,
\[
W_U(F)=\sum_\rho|\Fh(\gamma_\rho)|^2+\int_\R|\Fh(\tau)|^2\,\mathcal C(\tau)\,d\tau ,
\]
where $\gamma_\rho=\Iml\rho$ and the zero sum converges absolutely.
\end{lemma}
\begin{proof}
The function $\chi$ is supported in $[-9.5,9.5]$, and $\chi=1$ on $[-6.5,6.5]\supset[-2U,2U]$. Its
transform $\widehat\chi(z)=\int\chi(t)e^{izt}dt$ is the entire function
\[
\widehat\chi(z)=\frac{2\sin 8z}{z}\Big(\frac{\sin(z/4)}{z/4}\Big)^2\Big(\frac{\sin(z/n_\chi)}{z/n_\chi}\Big)^{n_\chi},
\]
and $k_\rho(\tau)=\frac12\big(\widehat\chi(s-i\delta)+\widehat\chi(s+i\delta)\big)-\widehat\chi(s)$ with $s=\gamma-\tau$;
$k_\rho$ is real because $\chi$ and $\cosh(\delta t)-1$ are even. Since $0<\Rel\rho<1$, $|\delta|<\frac12$. For
$z=x+iy$ one has $|\sin(hz)/(hz)|=|\frac12\int_{-1}^1e^{ihzw}dw|\le e^{h|y|}$, and also
$|\sin(hz)/(hz)|\le e^{h|y|}/(h|x|)$; likewise $|2\sin(8z)/z|\le16e^{8|y|}$. For $|y|\le\frac12$ the
exponents add up to at most $9.5/2$, and $\min(1,4/|x|)^2\le25(1+|x|)^{-2}$, so
$|\widehat\chi(x+iy)|\le400\,e^{4.75}(1+|x|)^{-2}\min(1,n_\chi/|x|)^{n_\chi}$. Hence
\[
\frac{|k_\rho(\tau)|}{2\pi}\le10^5\,(1+|s|)^{-2}\min(1,n_\chi/|s|)^{n_\chi},\qquad s=\gamma-\tau,
\]
since $800e^{4.75}/(2\pi)<1.48\cdot10^4$. The zeros with $k\le|\gamma-\tau|<k+1$ ($k\ge0$ an integer) have
ordinates in two intervals of length one, each meeting at most two intervals $[j,j+1)$ with
$|j|\le|\tau|+k+1$; by Lemma~\ref{lem:pw-count} there are at most $400\log(4+|\tau|+k)$ of them. Therefore
\begin{equation}\label{eq:pw-Csum}
|\mathcal C(\tau)|\le\sum_\rho\frac{|k_\rho(\tau)|}{2\pi}\le4\cdot10^7\sum_{k\ge0}\frac{\log(4+|\tau|+k)}{(1+k)^2}\min\Big(1,\frac{n_\chi}k\Big)^{n_\chi}.
\end{equation}
(i) Bound the minimum by $1$ and use $\log(4+|\tau|+k)\le\log(3+|\tau|)+\log(2+k)$,
$\sum_k(1+k)^{-2}<1.645$, $\sum_k\log(2+k)(1+k)^{-2}<1.801$ and $\log(3+|\tau|)\ge\log3$: the right side of
\eqref{eq:pw-Csum} is at most $4\cdot10^7(1.645+1.801/\log3)\log(3+|\tau|)<1.4\cdot10^8\log(3+|\tau|)$. If
$\delta=0$ the integrand of $k_\rho$ vanishes; by (Z2) this is the case for $|\gamma|\le H_{\mathrm{PT}}$.

(ii) Let $|\tau|\le H$ and $\Delta=H_{\mathrm{PT}}-H$. Only zeros with $|\gamma|>H_{\mathrm{PT}}$ contribute, and
for them $|\gamma-\tau|>\Delta$, so only $k\ge\Delta-1$ occurs in \eqref{eq:pw-Csum}. \emph{(Computation:)}
$\Delta-1-2.718\,n_\chi>2717$. For $k\ge\Delta-1$ one has $4+H+k\le(3+H_{\mathrm{PT}})\,k/(\Delta-1)$ (equality at
$k=\Delta-1$, and the right side grows faster), and $\log(XY)\le Y\log X$ for $X\ge e$, $Y\ge1$; so
$\log(4+|\tau|+k)\le\log(3+H_{\mathrm{PT}})\,k/(\Delta-1)$. Also
$(n_\chi/k)^{n_\chi}\le2.718^{-n_\chi}((\Delta-1)/k)^{n_\chi}$ and $(k/(\Delta-1))((\Delta-1)/k)^{n_\chi}\le1$.
Summing $(1+k)^{-2}$ over $k\ge\Delta-1$ gives at most $1$, so
$|\mathcal C(\tau)|\le4\cdot10^7\log(3+H_{\mathrm{PT}})\,2.718^{-n_\chi}$, and \emph{(computation)}
$10^9\log(3+H_{\mathrm{PT}})\,2.718^{-n_\chi}<e^{-D-2}$.

(iii) By Proposition~\ref{prop:P} and Lemma~\ref{lem:roughEF}, $W_U(F)=\sum_\rho\Phi_F(\rho)\Phi_F(1-\rho)$,
absolutely convergent; for real $F$, $\Phi_F(\rho)\Phi_F(1-\rho)=\int\cF(t)e^{(\delta+i\gamma)t}dt$. The
multiset of zeros is invariant under $\rho\mapsto1-\bar\rho$ (functional equation and conjugation), which changes $\delta$ to $-\delta$ and fixes
$\gamma$; averaging the two absolutely convergent expressions gives
$W_U(F)=\sum_\rho\int\cF(t)e^{i\gamma t}\cosh(\delta t)\,dt$. As $\supp\cF\subset[-2U,2U]$ and $\chi=1$ there,
\[
\int\cF(t)e^{i\gamma t}\cosh(\delta t)\,dt=\int\cF(t)e^{i\gamma t}dt+\int\cF(t)e^{i\gamma t}\chi(t)\big(\cosh(\delta t)-1\big)dt
=|\Fh(\gamma)|^2+\frac1{2\pi}\int|\Fh(\tau)|^2k_\rho(\tau)\,d\tau ,
\]
by $\cF(t)=\frac1{2\pi}\int|\Fh(\tau)|^2e^{-i\tau t}d\tau$ and Fubini. The sum $\sum_\rho|\Fh(\gamma)|^2$ converges
since $\Fh$ decays rapidly on $\R$ and $N(t)=O(t\log t)$; by (i) and $\int\log(3+|\tau|)|\Fh|^2<\infty$, summation
over $\rho$ and integration in $\tau$ may be interchanged.
\end{proof}

\begin{lemma}[Fourier leakage]\label{lem:pw-leak}
Let $F\in L^2(\IT)$ and $\Omega>0$. Then
\[
\int_{|\tau|>\Omega}|\Fh|^2\,d\tau\ \ge\ \eta(T,\Omega)\int_\R|\Fh|^2\,d\tau,\qquad
\eta(T,\Omega)=\frac1{1+(1+T\Omega)e^{2T\Omega+2}} .
\]
\end{lemma}
\begin{proof}
Put $l=1/T$, $A=\Omega+l$, $E_{\rm in}=\int_{|\tau|\le\Omega}|\Fh|^2$, $E_{\rm out}=\int_{|\tau|>\Omega}|\Fh|^2$,
$\Fh^\#(z)=\overline{\Fh(\bar z)}$ and $K(z)=\int_{-l}^{l}\Fh(z+s)\Fh^\#(z+s)\,ds$. Since $F\in L^1$ is supported
in $[-T,T]$, $K$ is entire and $|K(z)|\le2l\|F\|_1^2e^{2T|\Iml z|}$. For real $x$,
$K(x)=\int_{x-l}^{x+l}|\Fh|^2\ge0$, and $K(x)\le E_{\rm out}$ if $|x|\ge A$. On the strip $|\Rel\theta|\le\frac\pi2$
put $\phi(\theta)=K(A\sin\theta)e^{-2TA\cos\theta}$. With $\theta=\alpha+i\beta$,
$\Rel\cos\theta=\cos\alpha\cosh\beta\ge|\cos\alpha\sinh\beta|=|\Iml\sin\theta|$, so $\phi$ is bounded; on the
boundary lines $\sin\theta=\pm\cosh\beta$ is real of modulus at least $1$ and $\cos\theta$ is imaginary, so
$|\phi|\le E_{\rm out}$ there. By the three lines theorem $|\phi|\le E_{\rm out}$ on the strip, hence
$K(x)\le e^{2TA}E_{\rm out}$ for $|x|<A$. Each $\tau$ with $|\tau|\le\Omega$ lies in $[x-l,x+l]$ for all $x$ of
an interval of length $2l$ inside $[-A,A]$, so
$2lE_{\rm in}\le\int_{-A}^AK(x)\,dx\le2Ae^{2TA}E_{\rm out}$, i.e.\ $E_{\rm in}\le(1+T\Omega)e^{2T\Omega+2}E_{\rm out}$.
\end{proof}

\emph{The filter.} Put $v=\tau/H$, $r=R/H$ and
\[
u(\tau)=1-\Big(\frac{\sin v}{v}\Big)^2\Big(\frac{\cosh(b\sqrt{1-v^2})}{\cosh b}\Big)^2,\qquad
a=r^2\Big(\frac13+\frac{2b}{1+\sqrt{1-r^2}}\Big),
\]
where $\cosh(b\sqrt{1-v^2})=\sum_k b^{2k}(1-v^2)^k/(2k)!$ is entire, so $u$ is an even entire function, real
on $\R$. \emph{(Computation:)} $0.0000075739<a<0.0000075740$. Let $C_d$ be the Chebyshev polynomial,
$C_d(\cos\theta)=\cos d\theta$, so $|C_d|\le1$ on $[-1,1]$ and $|C_d(x)|=\cosh(d\operatorname{arcosh}|x|)$ for
$|x|\ge1$ ($d$ is even). Put
\[
h(s)=\frac{C_d(2s/a-1)}{C_d(2/a-1)},\qquad \kappa=C_d(2/a-1)^{-2}.
\]
The zeros $s_1,\dots,s_d$ of $h$ lie in $(0,a)$ and $h(1)=1$, so $h(s)=\prod_i(s-s_i)/(1-s_i)$ and, for
$0\le y<1-\max_is_i$,
\[
\frac1{h(1-y)}=\prod_{i=1}^d\Big(1-\frac y{1-s_i}\Big)^{-1}=\sum_{j\ge0}c_jy^j,\qquad
0\le c_j\le\binom{d+j-1}j(1-a)^{-j}\le\Big(\frac d{1-a}\Big)^j\le M^j ,
\]
by comparison with $(1-y/(1-a))^{-d}$; $c_0=1$. Put
\[
p(s)=h(s)\sum_{j=0}^{24}c_j(1-s)^j,\qquad \Theta(\tau)=p(u(\tau))^2,\qquad \mathsf J=25(2M)^{24},\qquad
z=M\operatorname{sech}^2b .
\]

\begin{lemma}[the filter]\label{lem:pw-filter}
\emph{(i)} $0\le u\le1$ on $\R$; $u\le a$ for $|\tau|\le R$; $u\ge\frac72a$ for $|\tau|\ge2R$; and
$-\frac{b+2}{2H^2}\le u(\frac i2)\le0$.
\emph{(ii)} $0\le\Theta\le1$ on $\R$, and $\Theta\le\kappa\mathsf J^2<1$ for $|\tau|\le R$.
\emph{(iii)} $\Theta(\tau)\ge\frac\kappa4e^{M\operatorname{arcosh}6}$ for $|\tau|\ge2R$.
\emph{(iv)} $z<1$, and $0\le1-\Theta(\tau)\le\frac{2z^{25}}{1-z}(H/|\tau|)^{50}$ for $|\tau|\ge H$.
\emph{(v)} $p(u(\frac i2))^2\le\kappa\mathsf J^2\exp\big(M\operatorname{arcosh}\big(1+\frac{b+2}{H^2a}\big)\big)$.
\emph{(vi)} If $F\in L^2(I_{T_1})$ and $F_1\in L^2(\R)$ is defined by $\widehat{F_1}=p(u)\Fh$, then
$\widehat{F_1}(\tau)=p(u(\tau))\Fh(\tau)$, $\tau\in\C$, is entire and $\supp F_1\subset[-T_1-\sigma,T_1+\sigma]\subset I_U$, with
$\sigma=(b+1)(M+48)/H$.
\end{lemma}
\begin{proof}
(i) For $|v|\le1$ both factors of the product lie in $[0,1]$; for $|v|>1$,
$\cosh(b\sqrt{1-v^2})=\cos(b\sqrt{v^2-1})$, so the product lies in $[0,\operatorname{sech}^2b/v^2]$. Hence
$0\le u\le1$. For $X,Y\in[0,1]$, $1-XY\le(1-X)+(1-Y)$; here $1-(\sin v/v)^2\le v^2/3$ and, as
$|(\log\cosh)'|\le1$, $1-Y\le-\log Y\le2b(1-\sqrt{1-v^2})=2bv^2/(1+\sqrt{1-v^2})$. Both bounds increase with
$|v|$, so $u\le a$ for $|v|\le r$. On $2r\le|v|\le1$ the two unsquared factors are positive and decreasing in
$|v|$, so $u$ increases, and \emph{(computation)} $u(2R)>\frac72a$; for $|v|\ge1$,
$u\ge1-\operatorname{sech}^2b>\frac72a$. At $\tau=\frac i2$, $v=iy$ with $y=1/(2H)$, and
$u(\frac i2)=1-(\sinh y/y)^2(\cosh(b\sqrt{1+y^2})/\cosh b)^2\le0$; as $\sinh y/y\le e^{y^2/6}$ and
$\log\cosh(b\sqrt{1+y^2})-\log\cosh b\le b(\sqrt{1+y^2}-1)\le by^2/2$, $|u(\frac i2)|\le e^{(b+1)y^2}-1\le
2(b+1)y^2\le\frac{b+2}{2H^2}$.

(ii) For $0\le s\le a$, $|h(s)|\le\kappa^{1/2}$ and $|\sum_{j\le24}c_j(1-s)^j|\le\sum_{j\le24}M^j\le\mathsf J$, so
$p(s)^2\le\kappa\mathsf J^2$, and $\kappa\le4e^{-M\operatorname{arcosh}(2/a-1)}$ gives $\kappa\mathsf J^2<1$ by
(vii) below. For $a\le s\le1$, $y=1-s\in[0,1-a]$, $h(s)>0$, and $h(s)\sum_{j\ge0}c_jy^j=1$; as $c_j\ge0$ and
$c_0=1$, $h(s)\le p(s)\le1$. With (i) this gives $0\le\Theta\le1$, and $\Theta\le\kappa\mathsf J^2$ where
$u\le a$.

(iii) For $|\tau|\ge2R$, $x=2u/a-1\ge6$, so $\Theta\ge h(u)^2=\kappa\cosh^2(d\operatorname{arcosh}x)\ge
\frac\kappa4e^{2d\operatorname{arcosh}6}$.

(iv) \emph{(Computation:)} $z<1$. For $|\tau|\ge H$, $u\ge a$ and $y=1-u\le\operatorname{sech}^2b/v^2$, so
$My\le z/v^2<1$ and $1-p(u)=h(u)\sum_{j\ge25}c_jy^j\le\sum_{j\ge25}(My)^j\le\frac{z^{25}}{1-z}v^{-50}$; and
$1-p^2\le2(1-p)$ for $0\le p\le1$.

(v) Let $s_0=u(\frac i2)$. By (i), $2s_0/a-1\in[-1-\frac{b+2}{H^2a},-1]$, so
$h(s_0)^2\le\kappa\cosh^2\big(d\operatorname{arcosh}(1+\frac{b+2}{H^2a})\big)\le\kappa
e^{M\operatorname{arcosh}(1+\frac{b+2}{H^2a})}$; and $|1-s_0|\le2$ gives
$|\sum_{j\le24}c_j(1-s_0)^j|\le\sum_{j\le24}(2M)^j\le\mathsf J$.

(vi) For $\Iml v=y$, $|\sin v/v|\le e^{|y|}$, and if $w=\sqrt{1-v^2}$ then
$(\Rel w)^2=\frac12\big(\Rel(1-v^2)+|1-v^2|\big)\le\frac12\big(1-(\Rel v)^2+y^2+1+|v|^2\big)=1+y^2$,
so $|\cosh(bw)/\cosh b|\le2e^{b(\sqrt{1+y^2}-1)}\le2e^{b|y|}$. Hence
$|u(\tau)|\le5e^{2(b+1)|\Iml\tau|/H}$, and since $p$ is a polynomial of degree $d+24$,
$|p(u(\tau))|\le C_pe^{\sigma|\Iml\tau|}$. Also $|\Fh(\tau)|\le\|F\|_1e^{T_1|\Iml\tau|}$. So $\widehat{F_1}$ is entire
of exponential type at most $T_1+\sigma$ and square integrable on $\R$ ($|p(u)|\le1$ there), and the
Paley--Wiener theorem gives the support. \emph{(Computation:)} $T_1+\sigma=2.84993164043\ldots<U$.

(vii) \emph{(Computation:)} $M\operatorname{arcosh}(2/a-1)=120110016805.59\ldots<D$, and $2\log\mathsf J<1141$.
\end{proof}

The radius enters only through (vi): with these parameters the support condition is
$T<U-\sigma=\frac{74918843181703}{26735070867180}=2.80226835\ldots$, and $T_1$ is chosen below this value.

\begin{lemma}[the prime operator]\label{lem:pw-schur}
Let $\mathsf\Pi_U=Q_U\sum_{\log n<2U}\frac{\Lambda(n)}{\sqrt n}(V_{\log n}+V_{-\log n})Q_U$ on $L^2(I_U)$, so that the prime form at the window $U$ is
$\Pi_U(G)=-\langle\mathsf\Pi_UG,G\rangle$ for $G\in L^2(I_U)$. Then $\|\mathsf\Pi_U\|<19.486$.
\end{lemma}
\begin{proof}
Indeed $\langle V_{\log n}G,G\rangle+\langle V_{-\log n}G,G\rangle=2\Rel c_G(\log n)$, which gives the first
statement. Divide $[-U,U]$ into $2048$ equal cells, numbered $0,\dots,2047$ from the left. For each prime
power $q=\ell^k$ with $\log q<2U$ (so $q\le293$; there are $79$ of them) put
\[
o_q=\Big\lfloor\frac{2048\log q}{2U}\Big\rfloor,\qquad b_q=\Big\lceil10^8\frac{\log\ell}{\sqrt q}\Big\rceil .
\]
For positive weights $w_0,\dots,w_{2047}$ with $w_i=w_{2047-i}$, extended by $w_i=0$ for $i\ge2048$, put
\[
V_i=\sum_qb_q\max(w_{i+o_q},w_{i+o_q+1}),\qquad Y_i=V_i+V_{2047-i}.
\]
A point of cell $i$ translated by $\log q$ lies in cell $i+o_q$ or $i+o_q+1$, or outside $[-U,U]$;
translations by $-\log q$ are reduced to these by the reflection $t\mapsto-t$, which maps cell $i$ to cell
$2047-i$ and preserves the weights. Hence, for the step function $w$ equal to $w_i$ on cell $i$, the
positive operator $\mathsf\Pi'_U$ obtained from $\mathsf\Pi_U$ by replacing $\Lambda(q)/\sqrt q$ with $10^{-8}b_q$
satisfies $\mathsf\Pi'_Uw\le\max_i\frac{Y_i}{10^8w_i}\,w$ pointwise. By the Schur test, whose usual proof applies verbatim to a symmetric
positive combination of translations, $\|\mathsf\Pi_U\|\le\|\mathsf\Pi'_U\|\le\max_iY_i/(10^8w_i)$. The weights
are obtained from $w_i=10^7$ by seventy applications of
\[
v_i=18\cdot10^8w_i+Y_i,\qquad w_i\leftarrow1+\big\lfloor10^7v_i/\max_jv_j\big\rfloor,
\]
which preserve the symmetry. \emph{(Computation:)} with the final weights,
$1000\,Y_i<19486\cdot10^8w_i$ for $0\le i<2048$; the largest quotient is $19.4853834\ldots$.
\end{proof}

\begin{proof}[Proof of Theorem~\ref{thm:power-window}]
\emph{Reductions.} By Lemma~\ref{lem:monoT} it suffices to take $T=T_1$. For real $F_r,F_i\in\mathcal D_{T_1}$
the polarization of $W_{T_1}$ at $(F_r,F_i)$ is real and symmetric (the symbol $\sT$ is even, and the polar
part polarizes to $A(F_r)B(F_i)+B(F_r)A(F_i)$), so $W_{T_1}(F_r+iF_i)=W_{T_1}(F_r)+W_{T_1}(F_i)$, while
$\|F_r+iF_i\|^2=\|F_r\|^2+\|F_i\|^2$. By Lemma~\ref{lem:form-density} it therefore suffices to prove
$W_{T_1}(F)\ge e^{-D}$ for real $F\in C_c^\infty(I_{T_1})$ with $\|F\|=1$, i.e.\
$\frac1{2\pi}\int|\Fh|^2=1$. Define $F_1$ by $\widehat{F_1}=p(u)\Fh$. By Lemma~\ref{lem:pw-filter}(vi),(ii),
$F_1$ is supported in $I_U$ and $|\widehat{F_1}|\le|\Fh|$ on $\R$; since $p$ has real coefficients and $u$ is
real and even on $\R$, $\widehat{F_1}(-\tau)=\overline{\widehat{F_1}(\tau)}$; and $\Fh$ decays faster than any power on $\R$. So
$F_1\in C_c^\infty(I_U)$ is real, and $|\widehat{F_1}|^2=\Theta|\Fh|^2$.

\emph{The archimedean margin.} By Lemma~\ref{lem:mbounds}, $m(\tau)\ge m_{\rm lo}(R)$ for $|\tau|\ge R$ and
$m\ge C_0$ on $\R$; \emph{(computation)} $m_{\rm lo}(R)-19.486>0.10000099$ and $C_0-19.486>-26$. Hence
\begin{equation}\label{eq:pw-margin}
m(\tau)-19.486\ \ge\ \tfrac1{10}\ \ (|\tau|\ge R),\qquad m(\tau)-19.486\ \ge\ -26\ \ (\tau\in\R).
\end{equation}

\emph{A lower bound for $W_U(F_1)$.} By Lemma~\ref{lem:pw-schur}, $\Pi_U(F_1)\ge-19.486\|F_1\|^2$, so
\[
\mathfrak a_U(F_1)\ge\frac1{2\pi}\int(m-19.486)\Theta|\Fh|^2 .
\]
Split the margin $\frac1{10}$ on $|\tau|\ge R$ into two halves; use one half on $|\tau|>H$ and the other on
$|\tau|\ge2R$ (note $H>2R$), drop the nonnegative remainder on $R\le|\tau|<2R$, and use
\eqref{eq:pw-margin} and Lemma~\ref{lem:pw-filter}(ii) on $|\tau|<R$. With
Lemma~\ref{lem:pw-filter}(iii) and Lemma~\ref{lem:pw-leak} for $\Omega=2R$, $\eta=\eta(T_1,2R)$,
\[
\mathfrak a_U(F_1)\ge\frac1{40\pi}\int_{|\tau|>H}\Theta|\Fh|^2+\kappa\Big[\frac{\eta\,e^{M\operatorname{arcosh}6}}{80}-26\mathsf J^2\Big].
\]
The factor $80$ is $20\cdot4$: the half margin $\frac1{20}$ and the factor $\frac14$ of
Lemma~\ref{lem:pw-filter}(iii). For the polar term, $A(F_1)=\widehat{F_1}(-\frac i2)$ and $B(F_1)=\widehat{F_1}(\frac i2)$,
and $u$ is even, so $A(F_1)=p(u(\frac i2))A(F)$, $B(F_1)=p(u(\frac i2))B(F)$ with $p(u(\frac i2))$ real; and
$|A(F)|^2,|B(F)|^2\le\int_{-T_1}^{T_1}e^{\pm t}dt\le e^{T_1}$ by Cauchy--Schwarz. By
Lemma~\ref{lem:pw-filter}(v),
\[
2A(F_1)B(F_1)\ge-2e^{T_1}p(u(\tfrac i2))^2\ge-\kappa\,\mathsf K_{\rm pol},\qquad
\mathsf K_{\rm pol}=4e^{T_1}\mathsf J^2\exp\Big(M\operatorname{arcosh}\Big(1+\frac{b+2}{H^2a}\Big)\Big).
\]
\emph{(Computation:)} $\log(\eta e^{M\operatorname{arcosh}6}/80)>536525$ (it suffices to bound $-\log\eta$ by
$\log(2(1+2RT_1))+4RT_1+2$), $\log\mathsf K_{\rm pol}<1152$ and $2\log\mathsf J<1141$. Hence
\begin{equation}\label{eq:pw-FW}
W_U(F_1)\ge2\kappa+\frac1{40\pi}\int_{|\tau|>H}\Theta|\Fh|^2 .
\end{equation}

\emph{From $F_1$ to $F$.} Extension by zero changes neither the norm, the Tate moments nor the
autocorrelation, and at the new contacts the autocorrelation vanishes (proof of Lemma~\ref{lem:monoT});
hence $W_{T_1}(F)=W_U(F)$. Apply Lemma~\ref{lem:pw-zero}(iii) to $F$
and to $F_1$, with $|\widehat{F_1}|^2=\Theta|\Fh|^2$, and subtract:
\begin{equation}\label{eq:pw-identity}
W_{T_1}(F)=W_U(F_1)+\sum_\rho\big(1-\Theta(\gamma_\rho)\big)|\Fh(\gamma_\rho)|^2+\int_\R\big(1-\Theta(\tau)\big)|\Fh(\tau)|^2\,\mathcal C(\tau)\,d\tau .
\end{equation}
The zero sum is nonnegative by Lemma~\ref{lem:pw-filter}(ii). On $|\tau|\le H$, $0\le1-\Theta\le1$ and
Lemma~\ref{lem:pw-zero}(ii) bound the integrand below by $-e^{-D-2}|\Fh|^2$, so this part is at least
$-2\pi e^{-D-2}$. On $|\tau|>H$, Lemma~\ref{lem:pw-filter}(iv) and Lemma~\ref{lem:pw-zero}(i) give
\[
(1-\Theta)|\mathcal C|\le\frac{2z^{25}}{1-z}\Big(\frac H{|\tau|}\Big)^{50}10^9\log(3+|\tau|)
\le\frac{2z^{25}}{1-z}10^9\log(3+H_{\mathrm{PT}})<4.01\cdot10^{-8},
\]
since $t^{-50}\log(3+t)$ decreases for $t\ge H$ and \emph{(computation)} the last inequality holds. By
Lemma~\ref{lem:pw-filter}(iv) also $1-\Theta<\frac12$ there, so $\frac1{40\pi}\Theta\ge\frac1{80\pi}>\frac1{1000}>4.01\cdot10^{-8}$ there, and the
second term of \eqref{eq:pw-FW} absorbs this part of the integral. By \eqref{eq:pw-identity} and
\eqref{eq:pw-FW},
\[
W_{T_1}(F)\ge2\kappa-2\pi e^{-D-2}.
\]

\emph{The count.} $\kappa=\cosh^{-2}(d\operatorname{arcosh}(2/a-1))\ge e^{-M\operatorname{arcosh}(2/a-1)}>e^{-D}$ by
Lemma~\ref{lem:pw-filter}(vii), and $2\pi e^{-2}<1$. Hence $W_{T_1}(F)>2\kappa-e^{-D}>\kappa>e^{-D}$, which
is stronger than \eqref{eq:power-window}.
\end{proof}

\section{Reduction to piecewise exponentials}\label{sec:coer}

We impose both Tate conditions on each of $k$ blocks. Coercivity on
their common kernel leaves $2k-2$ piecewise exponential directions.

\begin{definition}\label{def:coer-V}
Let $k\ge1$ and let $B_1,\dots,B_k$ be the $k$ consecutive open intervals of length $2r$,
$r:=T/k$, that partition $\IT$ up to endpoints, $B_j=(-T+2(j-1)r,\,-T+2jr)$. Put
\[
\mathcal V_k:=\Big\{F\in L^2(\IT):\ \int_{B_j}F(t)e^{t/2}dt=\int_{B_j}F(t)e^{-t/2}dt=0\ \ (1\le j\le k)\Big\},
\qquad E_k:=\PT\ominus\mathcal V_k ,
\]
and $e_j^\pm:=e^{\pm t/2}\mathbf 1_{B_j}$. (Lemma~\ref{lem:coer-decomp} holds for any partition into
intervals; equal blocks are used from Lemma~\ref{lem:coer-lowpass} on.)
\end{definition}

\begin{lemma}[the decomposition $\PT=\mathcal V_k\oplus E_k$]\label{lem:coer-decomp}
\emph{(i)} $\mathcal V_k$ is a closed subspace of $\PT$. \emph{(ii)} The orthogonal complement of
$\mathcal V_k$ in $L^2(\IT)$ is $\operatorname{span}\{e_j^\pm\}$, of dimension $2k$.
\emph{(iii)} $E_k=\operatorname{span}\{e_j^\pm\}\cap\PT=\{\sum_jc_j^+e_j^++c_j^-e_j^-:\ A=B=0\}$ has
dimension $2k-2$, and $\PT=\mathcal V_k\oplus E_k$ (orthogonal in $L^2(\IT)$).
\emph{(iv)} $E_k\subset\mathcal D_T$, and $\PT\cap\mathcal D_T=(\mathcal V_k\cap\mathcal D_T)\oplus E_k$.
\end{lemma}
\begin{proof}
(i) The $2k$ conditions are continuous linear functionals; summing them over $j$ gives $A(F)=B(F)=0$.
(ii) $\mathcal V_k$ is the joint kernel of the functionals $F\mapsto\langle F,e_j^\pm\rangle$, and the $2k$
functions $e_j^\pm$ are linearly independent (disjoint supports; $e^{t/2},e^{-t/2}$ independent on an
interval). (iii) $E_k=\mathcal V_k^\perp\cap\PT$ since $\mathcal V_k\subset\PT$. The two functionals
$A,B$ restricted to $\operatorname{span}\{e_j^\pm\}$ are independent: on the two-dimensional space
$\operatorname{span}\{e_1^+,e_1^-\}$ their matrix is
$\begin{psmallmatrix}\int_{B_1}e^t&|B_1|\\|B_1|&\int_{B_1}e^{-t}\end{psmallmatrix}$, of determinant
$\int_{B_1}e^t\int_{B_1}e^{-t}-|B_1|^2>0$ (Cauchy--Schwarz, strict since $e^t$ is not constant). Hence
$\dim E_k=2k-2$. If $F\in\PT$ and $F=v+w$ with $v\in\mathcal V_k$, $w\in\mathcal V_k^\perp$, then
$w=F-v\in\PT$, so $w\in E_k$. (iv) $\widehat{e_j^\pm}(s)=\big[e^{(\pm\frac12+is)t}/(\pm\frac12+is)\big]_{\partial B_j}$,
so $|\widehat{e_j^\pm}(s)|\le2e^{T/2}(\frac14+s^2)^{-1/2}$ and $\int\log(2+|s|)|\widehat{e_j^\pm}|^2<\infty$.
Since $\mathcal D_T$ is a linear space containing $E_k$, the $\mathcal V_k$-component of an $F\in\PT\cap\mathcal D_T$
lies in $\mathcal D_T$.
\end{proof}

For a window $r>0$ we write $\lambda_{\mathrm A}(r)$ for the bottom of $\mathfrak a_r$ on the unit sphere of
$\mathcal P_r\cap\mathcal D_r$; it is attained (Lemma~\ref{lem:compact}) and continuous in $r$ (Lemma~\ref{lem:contA}).
For $2r\le\log2$ the contact set $\mathcal N_r$ is empty, so $\mathfrak a_r=\mT$ on $L^2(I_r)$ and $\lambda_{\mathrm A}(r)$ is a
purely archimedean quantity. We use $|\cF(u)|\le\cF(0)=\|F\|^2$ and $\cF(u)=0$ for $|u|\ge2T$.

\begin{lemma}[the reduced kernel]\label{lem:coer-K1}
Put $K_1(u):=K(u)-2e^{-u/2}=\dfrac{2e^{-5u/2}}{1-e^{-2u}}=2\sum_{j\ge1}e^{-(2j+\frac12)u}$ for $u>0$, with $K$ as in
Lemma~\ref{lem:K}. Then $0<K_1(u)\le u^{-1}$ for all $u>0$, and $K_1(u)=u^{-1}-\frac32+O(u)$ as $u\downarrow0$. Let
$F,G\in L^2(\IT)$ be real, $\supp F\subset\overline B$, $\supp G\subset\overline{B'}$, with $B,B'$ disjoint intervals and $B'$
to the right of $B$. Then
\begin{equation}\label{eq:coer-cross}
\mT(F+G)=\mT(F)+\mT(G)-\iint F(x)G(y)K(y-x)\,dx\,dy ,
\end{equation}
and if $\int F(x)e^{x/2}dx=0$ or $\int G(y)e^{-y/2}dy=0$, the kernel $K$ in \eqref{eq:coer-cross} may be replaced by $K_1$.
\end{lemma}
\begin{proof}
The three expressions for $K_1$ agree by Lemma~\ref{lem:K} ($K=2e^{-u/2}\sum_{j\ge0}e^{-2ju}$). $K_1\le u^{-1}$ is
$e^{5u/2}-e^{u/2}\ge2u$: the function $g(u)=e^{5u/2}-e^{u/2}-2u$ has $g(0)=g'(0)=0$ and
$g''=\frac{25}4e^{5u/2}-\frac14e^{u/2}>0$. The expansion follows from $\frac1{2\sinh(u/2)}=u^{-1}-\frac u{24}+O(u^3)$,
$\frac1{2\cosh(u/2)}=\frac12+O(u^2)$ and $2e^{-u/2}=2-u+O(u^2)$. For \eqref{eq:coer-cross}, apply \eqref{eq:Aform} with
$\NT=\emptyset$ to $F+G$: $c_{F+G}(u)=\cF(u)+c_G(u)+\int F(t)G(t-u)dt+\int G(t)F(t-u)dt$; for $u>0$ the first cross
integral vanishes ($t\in B$, $t-u\in B'$ is impossible), while the second equals $\iint F(x)G(y)\mathbf 1_{y-x=u}$; integrating
against $K(u)du$ over $u>0$ gives \eqref{eq:coer-cross}, the term $\|F+G\|^2=\|F\|^2+\|G\|^2$ having no cross part.
Finally $2e^{-(y-x)/2}=2e^{-y/2}e^{x/2}$, so $\iint F(x)G(y)\,2e^{-(y-x)/2}dxdy=2\int Fe^{x/2}\int Ge^{-y/2}$, which vanishes
under either Tate condition.
\end{proof}
The cross term in \eqref{eq:coer-cross} carries the factor $1$ for the unordered pair $\{B,B'\}$; the polarized form is
$\mT(F,G)=-\frac12\iint F(x)G(y)K(y-x)\,dxdy$.

\begin{lemma}[a block sector is a short window]\label{lem:coer-sector}
Let $B\subset\IT$ be an interval of length $2r$ with $2r\le\log2$. Then
\[
\inf\Big\{\mathfrak a_T(F):\ F\in\mathcal D_T,\ \supp F\subset\overline B,\ \int_BFe^{\pm t/2}=0,\ \|F\|=1\Big\}=\lambda_{\mathrm A}(r),
\]
independently of $T$ and of the position of $B$; the infimum is attained. In particular, for $F\in\mathcal V_k$ with
$r=T/k$, $2r\le\log2$, and $F_j:=F\mathbf 1_{B_j}$,
\begin{equation}\label{eq:coer-diag}
\sum_{j=1}^k\mT(F_j)\ \ge\ \lambda_{\mathrm A}(r)\,\|F\|^2 .
\end{equation}
\end{lemma}
\begin{proof}
Let $B=I_r+a$. For $F$ supported in $\overline B$ and $F_a(t):=F(t+a)$, supported in $\overline{I_r}$, one has $c_{F_a}=\cF$,
hence $\mT(F_a)=\mT(F)$ by \eqref{eq:Aform}, and $\int F_ae^{\pm t/2}=e^{\mp a/2}\int Fe^{\pm t/2}$, so the two Tate
conditions on $B$ are those of $\mathcal P_r$ for $F_a$. Since $\cF(\log n)=0$ for $\log n\ge2r$ and $\log2\ge2r$, also
$\PiT(F)=0$, so $\mathfrak a_T(F)=\mT(F)=\mathfrak a_r(F_a)$ and the infimum is $\lambda_{\mathrm A}(r)$, attained since
$\lambda_{\mathrm A}(r)$ is. \eqref{eq:coer-diag} is the sum over the blocks, each $F_j$ satisfying the Tate conditions of
$B_j$, together with $\sum_j\|F_j\|^2=\|F\|^2$.
\end{proof}


\subsection{A lower bound for the short window}\label{subsec:coer-short}

\begin{lemma}\label{lem:coer-K}
$K(u)\ge u^{-1}$ for $0<u\le2$.
\end{lemma}
\begin{proof}
By Lemma~\ref{lem:K}, $K(u)-u^{-1}=\frac{u-2\sinh(u/2)}{2u\sinh(u/2)}+\frac1{2\cosh(u/2)}$. From $\tanh x\le x$,
$2\sinh(u/2)\le u\cosh(u/2)$, so $u-2\sinh(u/2)\ge-u(\cosh(u/2)-1)=-2u\sinh^2(u/4)$, and the first term is
$\ge-\frac{\sinh^2(u/4)}{\sinh(u/2)}=-\frac12\tanh\frac u4\ge-\frac u8$. Hence $K(u)-u^{-1}\ge\frac1{2\cosh(u/2)}-\frac u8\ge0$
whenever $u\cosh(u/2)\le4$, which holds for $u\le2$ ($2\cosh1=3.09$).
\end{proof}

\begin{lemma}[the tail]\label{lem:coer-tail}
For $r>0$, $\int_{2r}^\infty K(u)\,du\ge\log\frac2r+\frac\pi2-r$.
\end{lemma}
\begin{proof}
By Theorem~\ref{thm:S}(iii) the integral is $\log\coth\frac r2+\pi-2\arctan e^r$; $\coth x\ge x^{-1}$ bounds the first
term, and $\frac{d}{dr}(\pi-2\arctan e^r)=-\frac1{\cosh r}\ge-1$ with value $\frac\pi2$ at $r=0$ bounds the second.
\end{proof}

\begin{lemma}[the truncated logarithm]\label{lem:coer-w}
Put $\mathsf w(s):=\int_0^2\frac{1-\cos sv}v\,dv=\log(2|s|)+\gamma-\operatorname{Ci}(2|s|)$ for $s\ne0$, $\mathsf w(0)=0$. Then $\mathsf w$ is
even, increasing in $|s|$, and for every $0<R\le\pi$ and all $s\in\R$, $\mathsf w(s)\ge\mathsf w(R)\min\{1,s^2/R^2\}$.
\end{lemma}
\begin{proof}
$\mathsf w'(s)=\int_0^2\sin(sv)\,dv=\frac{1-\cos2s}s\ge0$ for $s>0$. Next $\mathsf w(s)/s^2=\int_0^2v\,\frac{1-\cos sv}{(sv)^2}dv$, and
$x\mapsto\frac{1-\cos x}{x^2}=\frac12\big(\frac{\sin(x/2)}{x/2}\big)^2$ is decreasing on $(0,2\pi)$; for $0<s\le R\le\pi$ and $v\le2$
one has $sv\le2\pi$, so $\mathsf w(s)/s^2\ge\mathsf w(R)/R^2$. For $s\ge R$ use monotonicity.
\end{proof}

\begin{lemma}[a double Poincar\'e inequality]\label{lem:coer-poinc}
Let $h\in L^2(\R)$ with $\supp h\subset[-1,1]$ and $h''\in L^2(\R)$. Then $\|h''\|^2\ge\pi^2\|h'\|^2\ge\frac{\pi^4}4\|h\|^2$.
\end{lemma}
\begin{proof}
$h,h''\in L^2(\R)$ give $h\in H^2(\R)\subset C^1(\R)$, and the support condition gives $h(\pm1)=h'(\pm1)=0$. The second
inequality is the Dirichlet Poincar\'e inequality for $h\in H^1_0(-1,1)$ (first eigenvalue $(\pi/2)^2$). For the first,
$g:=h'\in H^1_0(-1,1)$ with $\int_{-1}^1g=h(1)-h(-1)=0$. Expand $g=\sum_{n\ge1}c_n\phi_n$ in the orthonormal Dirichlet
eigenbasis $\phi_n(x)=\sin\frac{n\pi(x+1)}2$ of $(-1,1)$, so that $\|g\|^2=\sum c_n^2$ and $\|g'\|^2=\sum(\frac{n\pi}2)^2c_n^2$. Since
$\int\phi_n=\frac4{n\pi}$ for odd $n$ and $0$ for even $n$, the constraint reads $c_1=-\sum_{n\ge3\ \mathrm{odd}}c_n/n$. The
even-indexed terms satisfy $(\frac{n\pi}2)^2\ge\pi^2$, and for the odd-indexed ones
\[
\sum_{n\ \mathrm{odd}}\Big(\frac{n^2\pi^2}4-\pi^2\Big)c_n^2=\frac{\pi^2}4\Big[\sum_{n\ge3\ \mathrm{odd}}(n^2-4)c_n^2-3c_1^2\Big]\ \ge\ 0,
\]
because by Cauchy--Schwarz $c_1^2\le\mathcal S\sum_{n\ge3\ \mathrm{odd}}(n^2-4)c_n^2$ with
$\mathcal S=\sum_{n\ge3\ \mathrm{odd}}\frac1{n^2(n^2-4)}=\frac14\big(\frac43-\frac{\pi^2}8\big)=0.0249<\frac13$ (partial fractions and
$\sum_{n\ \mathrm{odd}}n^{-2}=\pi^2/8$). Hence $\|g'\|^2\ge\pi^2\|g\|^2$.
\end{proof}

\begin{theorem}[the short window]\label{thm:coer-A}
Let $0<r\le1$ with $2r\le\log2$, and put $a:=r^2/4$. For every $R\in(0,\pi]$,
\begin{equation}\label{eq:coer-A}
\lambda_{\mathrm A}(r)\ \ge\ C_0+\log\frac2r+\frac\pi2-r+\mathsf w(R)\Big[1-\frac{16(R^2+a)^3}{27\pi^4R^2}\Big].
\end{equation}
In particular, with $R=2$ ($\mathsf w(2)=\log4+\gamma-\operatorname{Ci}(4)=2.1045$),
$\lambda_{\mathrm A}(r)\ge\log\frac1r-1.209-r$ for $0<r\le0.1$, so $\lambda_{\mathrm A}(r)\ge\log\frac1r-1.31$ on $(0,0.1]$.
\end{theorem}
\begin{proof}
Let $F\in\mathcal P_r\cap\mathcal D_r$, $\|F\|=1$. Since $\mathcal N_r=\emptyset$ and $\cF(u)=0$ for $u\ge2r$, \eqref{eq:Aform} gives
\[
\mathfrak a_r(F)=C_0+\int_{2r}^\infty K+\int_0^{2r}\big(1-\Rel\cF(u)\big)K(u)\,du .
\]
The tail is bounded by Lemma~\ref{lem:coer-tail}. Put $G(y):=r^{1/2}F(ry)$, so that $\supp G\subset[-1,1]$, $\|G\|=1$,
$\cF(u)=c_G(u/r)$, and $\int_0^{2r}(1-\Rel\cF)K=\int_0^2(1-\Rel c_G(v))\,rK(rv)\,dv\ge\int_0^2\frac{1-\Rel c_G(v)}v\,dv$, using
$1-\Rel c_G\ge0$ and Lemma~\ref{lem:coer-K} ($rv\le2r\le2$). By Plancherel, $1-\Rel c_G(v)=\frac1{2\pi}\int(1-\cos sv)|\widehat G(s)|^2ds$,
so by Tonelli the last integral equals $\frac1{2\pi}\int\mathsf w|\widehat G|^2$, and by Lemma~\ref{lem:coer-w}, for $0<R\le\pi$,
\[
\frac1{2\pi}\int\mathsf w|\widehat G|^2\ \ge\ \mathsf w(R)\Big[1-\frac1{2\pi}\int\Big(1-\frac{s^2}{R^2}\Big)^+|\widehat G(s)|^2ds\Big].
\]
The Tate conditions $\int Fe^{\pm t/2}=0$ read $\int Ge^{\pm ry/2}=0$, i.e.\ $\widehat G(\mp ir/2)=0$. $\widehat G$ is entire of
exponential type $\le1$ and in $L^2(\R)$, so $H(s):=\widehat G(s)/(s^2+a)$ is entire of exponential type $\le1$ and in $L^2(\R)$
($s^2+a\ge a>0$ on $\R$); by the Paley--Wiener theorem $H=\widehat h$ with $h\in L^2(\R)$, $\supp h\subset[-1,1]$, and
$G=-h''+ah$, so $h''\in L^2$. Then
$1=\|G\|^2=\frac1{2\pi}\int(s^2+a)^2|H|^2=\|h''\|^2+2a\|h'\|^2+a^2\|h\|^2\ge\|h''\|^2\ge\frac{\pi^4}4\|h\|^2$
(Lemma~\ref{lem:coer-poinc}). On the other hand $(1-\frac{s^2}{R^2})^+(s^2+a)^2\le M_R:=\frac{4(R^2+a)^3}{27R^2}$ (the
maximum of $(1-x/R^2)(x+a)^2$ on $0\le x\le R^2$ is at $x=(2R^2-a)/3$), so
$\frac1{2\pi}\int(1-s^2/R^2)^+|\widehat G|^2\le M_R\|h\|^2\le\frac{4M_R}{\pi^4}$. Collecting terms gives \eqref{eq:coer-A}. For
$R=2$ and $a\le\frac1{400}$ the bracket is $\ge0.9025$ and $\mathsf w(2)\cdot0.9025\ge1.899$; with
$C_0+\log2+\frac\pi2=-3.1082$ this gives $\log\frac1r-1.209-r$.
\end{proof}

Together with Corollary~\ref{cor:mono} and Lemma~\ref{lem:mbounds}(ii),
the lower estimate gives $\lambda_{\rm A}(r)=\log(1/r)+O(1)$ under the two Tate conditions.

Yoshida proves short-window coercivity after removing finitely many Fourier
modes \cite[Lemma~3]{Yoshida1992}; Connes and Consani obtain a finite-codimension
inequality from the archimedean operator \cite[Thm.~3.6]{ConnesConsani2021}.
Here the two Tate conditions give the explicit rate $\log(1/r)$.
For numerical and finite-section comparisons see \cite{Kim2026,CvS2025}.


\subsection{Coercivity of the archimedean form on $\mathcal V_k$}\label{subsec:coer-B}
For $R>0$ let $P_R$ be the orthogonal projection of $L^2(\R)$ onto the band $\{|s|\le R\}$,
$\widehat{P_RF}=\Fh\mathbf 1_{|s|\le R}$.

\begin{lemma}[low-pass leakage of $\mathcal V_k$]\label{lem:coer-lowpass}
Let $F\in\mathcal V_k$ with equal blocks of half-length $r$, and $R>0$. Then
\[
\|P_RF\|^2\le\varepsilon_{\mathrm{lp}}(Rr)\,\|F\|^2,\qquad \varepsilon_{\mathrm{lp}}(\theta):=\frac{\cosh^2r}{6}\Big(\theta^2+\frac{r^2}4\Big)^2 .
\]
\end{lemma}
\begin{proof}
$G:=P_RF$ is entire. Fix a block $B$ with centre $c$ and put, on $B$,
$u(x):=G(c)\cosh\frac{x-c}2+2G'(c)\sinh\frac{x-c}2\in\operatorname{span}\{e^{x/2},e^{-x/2}\}$, so that
$u(c)=G(c)$, $u'(c)=G'(c)$, $u''=u/4$. Then $\varrho:=G-u$ satisfies $\varrho(c)=\varrho'(c)=0$ and
$\varrho''-\frac\varrho4=G''-\frac G4=:H$, whence $\varrho(x)=\int_c^x2\sinh\frac{x-t}2\,H(t)\,dt$. By Cauchy--Schwarz and
$2\sinh\frac z2\le z\cosh\frac z2\le z\cosh r$ ($0\le z\le r$),
$|\varrho(x)|^2\le\cosh^2r\,\frac{|x-c|^3}3\,\|H\|_{L^2(B)}^2$, so $\int_B\varrho^2\le\cosh^2r\,\frac{r^4}6\|H\|_{L^2(B)}^2$.
Since $F_B:=F\mathbf 1_B$ is orthogonal to $e^{\pm x/2}\mathbf 1_B$, $\langle F_B,G\rangle=\langle F_B,\varrho\rangle\le\|F_B\|\,\|\varrho\|_{L^2(B)}$.
Summing over the blocks and using Cauchy--Schwarz once more,
\begin{align*}
\|G\|^2=\langle F,G\rangle&\le\|F\|\Big(\sum_B\|\varrho\|_{L^2(B)}^2\Big)^{1/2}\le\|F\|\cosh r\,\frac{r^2}{\sqrt6}\,\|G''-\tfrac14G\|_{L^2(\R)}\\
&\le\|F\|\cosh r\,\frac{r^2}{\sqrt6}\Big(R^2+\frac14\Big)\|G\|,
\end{align*}
the first equality because $P_R$ is an orthogonal projection, the last inequality because
$\widehat G$ is supported in $|s|\le R$. Squaring gives the claim, as $r^4(R^2+\frac14)^2=((Rr)^2+\frac{r^2}4)^2$.
\end{proof}

\begin{theorem}[coercivity on $\mathcal V_k$]\label{thm:coer-B}
Let $k\ge1$, $r=T/k\le1$, $\kappa:=\cosh^2r$, $a:=r^2/4$, $\varepsilon_{\mathrm{lp}}$ as in
Lemma~\ref{lem:coer-lowpass}, and $\theta_*:=(\sqrt{6/\kappa}-a)^{1/2}$ (the root of $\varepsilon_{\mathrm{lp}}(\theta_*)=1$); $m_{\mathrm{lo}}$ is
the lower bound of Lemma~\ref{lem:mbounds}(ii).
Then for every $F\in\mathcal V_k\cap\mathcal D_T$,
\begin{equation}\label{eq:coer-B}
\mT(F)\ \ge\ \Big(\log\frac1r-C_B(r)\Big)\|F\|^2,\qquad
C_B(r):=\log\frac1r-C_0\,\varepsilon_{\mathrm{lp}}(r)-\frac{2\kappa}3\big[\Phi_r(\theta_*)-\Phi_r(r)\big],
\end{equation}
where $\Phi_r$ is the elementary primitive
\[
\Phi_r(\theta):=\Big(\frac{\theta^4}4+\frac{a\theta^2}2\Big)\big(\log\theta-\log(2\pi r)\big)-\frac{\theta^4}{16}-\frac{a\theta^2}4
-r\Big(\frac{\theta^3}3+a\theta\Big)-\frac{r^2\theta^2}4-\frac{ar^2}2\log\theta
\]
of $(\theta^3+a\theta)\,m_{\mathrm{lo}}(\theta/r)$. Moreover $C_B(r)\to\frac14-\log\frac{6^{1/4}}{2\pi}=1.6399$ as
$r\to0$, and $C_B(0.02)=1.6573$, $C_B(0.0619)=1.6957$, $C_B(0.1)=1.7331$ (evaluations of the closed form).
\end{theorem}
\begin{proof}
Let $\|F\|=1$ and $\mu:=\frac1{2\pi}|\Fh|^2ds$, a probability measure, so that $\mT(F)=\int m\,d\mu$ by
\eqref{eq:form} with $\NT=\emptyset$. By Lemma~\ref{lem:mbounds}(i) and \eqref{eq:m}, $m$ is even, real-analytic,
increasing on $[0,\infty)$ with $m(0)=C_0$; hence $m(s)=C_0+\int_0^{|s|}m'(\rho)d\rho$ with $m'\ge0$ and, by Tonelli,
\[
\int m\,d\mu=C_0+\int_0^\infty m'(\rho)\,\mu(|s|\ge\rho)\,d\rho\ \ge\ C_0+\int_0^{\rho_*}m'(\rho)\big(1-\varepsilon_{\mathrm{lp}}(\rho r)\big)d\rho,
\qquad\rho_*:=\theta_*/r,
\]
since $\mu(|s|\ge\rho)=1-\|P_{\rho}F\|^2+\mu(|s|=\rho)\ge1-\varepsilon_{\mathrm{lp}}(\rho r)$ by Lemma~\ref{lem:coer-lowpass} and the
integrand is nonnegative on $[0,\rho_*]$, where $\varepsilon_{\mathrm{lp}}(\rho r)\le1$. Integrating by parts on $[0,\rho_*]$, with
vanishing boundary terms ($1-\varepsilon_{\mathrm{lp}}(\theta_*)=0$ and $m(0)-C_0=0$),
\[
\int_0^{\rho_*}m'(\rho)(1-\varepsilon_{\mathrm{lp}}(\rho r))d\rho=\int_0^{\theta_*}\big(m(\theta/r)-C_0\big)\varepsilon_{\mathrm{lp}}'(\theta)\,d\theta,
\]
so that
\[
\int m\,d\mu\ \ge\ C_0\,\varepsilon_{\mathrm{lp}}(0)+\int_0^{\theta_*}m(\theta/r)\,\varepsilon_{\mathrm{lp}}'(\theta)\,d\theta ,
\]
using $\int_0^{\theta_*}\varepsilon_{\mathrm{lp}}'=1-\varepsilon_{\mathrm{lp}}(0)$. Here $\varepsilon_{\mathrm{lp}}'(\theta)=\frac{2\kappa}3\theta(\theta^2+a)\ge0$. By
Lemma~\ref{lem:mbounds}(i),(ii), $m\ge\max\{C_0,m_{\mathrm{lo}}\}$ on $(0,\infty)$; we use $m\ge C_0$ on $[0,r]$ (i.e.\ $\rho\le1$)
and $m\ge m_{\mathrm{lo}}$ on $[r,\theta_*]$ (note $m_{\mathrm{lo}}(1)=-3.34>C_0$, and $m_{\mathrm{lo}}$ is increasing, so no further loss occurs
there). Hence
\[
\int m\,d\mu\ \ge\ C_0\,\varepsilon_{\mathrm{lp}}(0)+C_0\big(\varepsilon_{\mathrm{lp}}(r)-\varepsilon_{\mathrm{lp}}(0)\big)+\int_r^{\theta_*}m_{\mathrm{lo}}(\theta/r)\,\varepsilon_{\mathrm{lp}}'(\theta)\,d\theta
=C_0\,\varepsilon_{\mathrm{lp}}(r)+\frac{2\kappa}3\big[\Phi_r(\theta_*)-\Phi_r(r)\big],
\]
which is \eqref{eq:coer-B}; the primitive is checked by differentiation. As $r\to0$: $\kappa\to1$, $a\to0$,
$\theta_*\to6^{1/4}$, $\varepsilon_{\mathrm{lp}}(r)\to0$, $\frac{2\kappa}3\Phi_r(r)\to0$, and
$\frac{2\kappa}3\Phi_r(\theta_*)=\log\theta_*-\log(2\pi r)-\frac14+o(1)$ (the $a$- and $r$-terms are $o(1)$),
whence $C_B(r)\to\log2\pi-\log6^{1/4}+\frac14$. The three values are evaluations of the closed form (double precision).
\end{proof}

Corollary~\ref{cor:coer-L1} uses $C_B$ only at $r=T/k$. The bound
uses the full multiplier; the cost of separating adjacent blocks
is computed in Remark~\ref{rem:coer-nogo}.




\begin{proposition}[the norm of the prime form]\label{prop:coer-prime}
\emph{(i)} $|\PiT(F)|\le\Aconst\|F\|^2$ for all $F\in L^2(\IT)$ and all $T$.
\emph{(ii)} For $\ell>0$ let $L(T,\ell):=\lceil2T/\ell\rceil$. Then $|\Rel\cF(\ell)|\le\cos\frac{\pi}{L(T,\ell)+1}\,\|F\|^2$,
hence
\[
\|\PiT\|\ \le\ \sum_{n\in\NT}v_n\cos\frac\pi{L(T,\log n)+1}
=\begin{cases}
\Aconst/2, & T\le\log2,\\[2pt]
\dfrac{v_2}{\sqrt2}+\dfrac{v_3+v_4}2=1.6740, & \log2<T\le\tfrac12\log5 .
\end{cases}
\]
Here $\frac{v_2}{\sqrt2}+\frac{v_3+v_4}2=\log2+\frac{\log3}{\sqrt3}+\frac{\log2}2$; in the first case the bound is
$\frac{v_2}2=0.4901$ for $T\le\frac12\log3$ and $\frac{v_2+v_3}2=1.1244$ for $\frac12\log3<T\le\log2$. The bounds hold on all of $L^2(\IT)$, in particular on $\mathcal V_k$.
\end{proposition}
\begin{proof}
(i) is $|\cF(\log n)|\le\|F\|^2$. (ii) Let $V_\ell$ be the shift $F\mapsto F(\cdot-\ell)\mathbf 1_{\IT}$ on
$L^2(\IT)$, so that $\Rel\cF(\ell)=\frac12\langle(V_\ell+V_\ell^*)F,F\rangle$. The map
$F\mapsto\big(F(-T+y+j\ell)\big)_{0\le j<n(y)}$, $y\in[0,\ell)$, $n(y):=\#\{j\ge0:-T+y+j\ell<T\}\le L(T,\ell)$, is unitary
from $L^2(\IT)$ onto $\int^\oplus_{[0,\ell)}\C^{n(y)}dy$, and it carries $V_\ell$ to the fibrewise shift $j\mapsto j+1$.
Hence $V_\ell+V_\ell^*$ acts on the fibre $\C^{n}$ as the adjacency matrix of the path with $n$ vertices, whose
eigenvalues are $2\cos\frac{j\pi}{n+1}$, $1\le j\le n$; its norm is $2\cos\frac\pi{n+1}\le2\cos\frac\pi{L+1}$. For
$T\le\log2$ one has $2T/\log n\le2$ for all $n\ge2$, so $L=2$ and $\cos\frac\pi3=\frac12$. For $\log2<T\le\frac12\log5$:
$\NT=\{2,3,4\}$, $L(T,\log2)=3$ since $2<2T/\log2\le\log5/\log2=2.32$, and $L(T,\log3)=L(T,\log4)=2$ since
$2T\le\log5<2\log3$; $\cos\frac\pi4=\frac1{\sqrt2}$. The values use $v_2=\sqrt2\log2$, $v_3=2\log3/\sqrt3$, $v_4=\log2$.
\end{proof}
For $\frac12\log2<T<\frac12\log3$ the bound $\Aconst/2$ is the exact norm: on a set of fibres of positive
measure the fibre has length $2$ and the prime form acts there as $-v_2\begin{psmallmatrix}0&1/2\\1/2&0\end{psmallmatrix}$.
Once $T>\log2$ the orbit of $\log2$ has length $3$ and the path norm exceeds $1$, so the bound
$\|\PiT\|\le\Aconst/2$ is specific to $T\le\log2$.

\begin{corollary}[coercivity of the truncated form on $\mathcal V_k$]\label{cor:coer-L1}
Let $k\ge1$ with $r=T/k\le1$, and put $c(T,k):=\inf\{\mathfrak a_T(F):F\in\mathcal V_k\cap\mathcal D_T,\ \|F\|=1\}$. Then
\begin{equation}\label{eq:coer-L1}
c(T,k)\ \ge\ \log\frac kT-C_B\Big(\frac Tk\Big)-\|\PiT\| ,
\end{equation}
with $\|\PiT\|$ bounded by Proposition~\ref{prop:coer-prime}. In particular (evaluations of the closed form):
\begin{center}
\begin{tabular}{@{}lllll@{}}
\toprule
$T$ & $k$ & $\dim E_k$ & bound for $\|\PiT\|$ & lower bound for $c(T,k)$\\
\midrule
$\frac12\log3$ & $8$ & $14$ & $0.4901$ & $0.486$\\
$\log2$ & $12$ & $22$ & $1.1244$ & $0.035$\\
$0.69$ & $14$ & $26$ & $1.1244$ & $0.202$\\
$\frac12\log5$ & $25$ & $48$ & $1.6740$ & $0.094$\\
$\frac12\log5$ & $26$ & $50$ & $1.6740$ & $0.134$\\
$\frac12\log5$ & $30$ & $58$ & $1.6740$ & $0.281$\\
$\frac12\log5$ & $85$ & $168$ & $\Aconst=2.9420$ & $0.070$\\
\bottomrule
\end{tabular}
\end{center}
Moreover, for every $T>0$ there is $k_0(T)$ such that $c(T,k)>0$ for all $k\ge k_0(T)$; indeed the right-hand side
of \eqref{eq:coer-L1} tends to $+\infty$ with $k$, because $C_B(T/k)\to1.6399$ (Theorem~\ref{thm:coer-B}).
\end{corollary}
\begin{proof}
$\mathfrak a_T=\mT+\PiT$ on $\mathcal D_T$, Theorem~\ref{thm:coer-B} and Proposition~\ref{prop:coer-prime}.
\end{proof}

By \eqref{eq:coer-L1}, $k>Te^{\Aconst+C_B(T/k)}$ suffices.
As $C_B(T/k)\to1.6399$ and $\Aconst\asymp e^T$, this sufficient dimension
grows doubly exponentially. For $T\le\tfrac12\log5$ one can keep $2k-2<60$
by Corollary~\ref{cor:coer-L1}.




\subsection{The finite reduction}\label{subsec:coer-schur}

\begin{theorem}[reduction to the piecewise Tate exponentials]\label{thm:coer-schur}
Fix $T>0$ and $k\ge1$ with $c:=c(T,k)>0$ (Corollary~\ref{cor:coer-L1}: every $k\ge k_0(T)$). Put
$\mathcal V:=\mathcal V_k\cap\mathcal D_T$ and $M:=\Aconst-C_0$.
\emph{(i)} $\mathfrak a_T$ is an inner product on $\mathcal V$ for which $\mathcal V$ is complete, and for every
$e\in E_k$ the functional $v\mapsto\mathfrak a_T(e,v)$ is bounded on $(\mathcal V,\mathfrak a_T)$; hence there is a unique
$\mathsf Ke\in\mathcal V$ with $\mathfrak a_T(\mathsf Ke,v)=\mathfrak a_T(e,v)$ for all $v\in\mathcal V$. The map $\mathsf K:E_k\to\mathcal V$ is linear,
and $\|\mathsf K\|:=\sup\{\|\mathsf Ke\|:\|e\|=1\}<\infty$.
\emph{(ii)} For $v\in\mathcal V$ and $e\in E_k$,
\begin{equation}\label{eq:coer-schur}
\begin{split}
\mathfrak a_T(v+e)&=\mathfrak a_T(v+\mathsf Ke)+\mathfrak s_{T,k}(e),\\
\mathfrak s_{T,k}(e)&:=\mathfrak a_T(e)-\mathfrak a_T(\mathsf Ke)=\mathfrak a_T(e)-\sup_{u\in\mathcal V}\big[2\Rel\mathfrak a_T(e,u)-\mathfrak a_T(u)\big].
\end{split}
\end{equation}
Let $S_{T,k}$ be the matrix of the real quadratic form $\mathfrak s_{T,k}$ in an orthonormal basis of $E_k$ (a
real symmetric $(2k-2)\times(2k-2)$ matrix). For $k=1$, $E_k=0$ and $\lamA=c$.
For $k\ge2$, writing $\lambda_S=\min\sigma(S_{T,k})$, one has
\begin{equation}\label{eq:coer-sandwich}
\begin{aligned}
 \frac{\min\{c,\lambda_S\}}{(1+\|\mathsf K\|)^2}
 &\le\lamA\le\min\{c,\lambda_S\} &&(\lambda_S\ge0),\\
 \lambda_S &\le\lamA\le\frac{\lambda_S}{1+\|\mathsf K\|^2} &&(\lambda_S<0).
\end{aligned}
\end{equation}
\emph{(iii)} Consequently $\lamA>0$ if and only if $S_{T,k}\succ0$, and $\lamA\ge0$ if and only if $S_{T,k}\succeq0$.
\end{theorem}
\begin{proof}
(i) By Lemma~\ref{lem:mbounds}, $\sigma_T\ge C_0-\Aconst$, so $\mathfrak a_T+M\|\cdot\|^2\ge0$ on $\mathcal D_T$, and
by Lemma~\ref{lem:mbounds}(ii) and $m\ge C_0$ there are constants $c_1,c_2>0$ (depending on $T$) with
$c_1\int\log(2+|s|)|\Fh|^2\le\mathfrak a_T(F)+(M+1)\|F\|^2\le c_2\int\log(2+|s|)|\Fh|^2$ on $\mathcal D_T$. On $\mathcal V$,
$\mathfrak a_T\ge c\|\cdot\|^2$ gives $\mathfrak a_T(v)\ge\frac c{c+M+1}\big(\mathfrak a_T(v)+(M+1)\|v\|^2\big)$, so
$\mathfrak a_T$ is equivalent on $\mathcal V$ to the weighted norm $\int\log(2+|s|)|\widehat v|^2$, for which
$\mathcal D_T$ is complete and $\mathcal V_k\cap\mathcal D_T$ is closed ($\mathcal V_k$ is $L^2$-closed and the weighted norm
dominates the $L^2$ norm). For $e\in E_k\subset\mathcal D_T$ (Lemma~\ref{lem:coer-decomp}(iv)), the Cauchy--Schwarz
inequality for the nonnegative form $\mathfrak a_T+M\langle\cdot,\cdot\rangle$ gives
$|\mathfrak a_T(e,v)|\le(\mathfrak a_T(e)+M\|e\|^2)^{1/2}(\mathfrak a_T(v)+M\|v\|^2)^{1/2}+M\|e\|\|v\|\le C_e\,\mathfrak a_T(v)^{1/2}$,
so $\mathsf Ke$ exists by the Riesz representation theorem; $K$ is linear, and $\|\mathsf K\|<\infty$ since $\dim E_k<\infty$.
(ii) $\mathfrak a_T(v+\mathsf Ke)=\mathfrak a_T(v)+2\Rel\mathfrak a_T(\mathsf Ke,v)+\mathfrak a_T(\mathsf Ke)=\mathfrak a_T(v)+2\Rel\mathfrak a_T(e,v)+\mathfrak a_T(\mathsf Ke)$,
which rearranges to \eqref{eq:coer-schur}; the variational expression is $\mathfrak a_T(\mathsf Ke)=\sup_u[2\Rel\mathfrak a_T(\mathsf Ke,u)-\mathfrak a_T(u)]$
(maximum at $u=\mathsf Ke$). For \eqref{eq:coer-sandwich}: every $F\in\PT\cap\mathcal D_T$ is $v+e$ (Lemma~\ref{lem:coer-decomp}(iv)),
$\|F\|^2=\|v\|^2+\|e\|^2$, and $v\mapsto v+\mathsf Ke$ is a bijection of $\mathcal V$. Taking $e=0$ gives $\lamA\le c$. For $k\ge2$, take a unit least
eigenvector $e$ and $v=-\mathsf Ke$. The quotient is exactly
$\lambda_S/(1+\|\mathsf Ke\|^2)$, giving the upper bound in each case.
If $\lambda_S\ge0$, the square completion bounds the form below by
$\min(c,\lambda_S)(\|v+\mathsf Ke\|^2+\|e\|^2)$, and the inverse
change of variables on the orthogonal sum has norm at most $1+\|\mathsf K\|$.
If $\lambda_S<0$, drop the nonnegative square to get
$\mathfrak a_T(v+e)\ge\lambda_S\|e\|^2\ge\lambda_S\|v+e\|^2$.
These bounds also prove (iii).
\end{proof}

\begin{remark}[inertia of the two Schur complements]\label{rem:inertiatower}
If $S_{T,k}$ is nonsingular, the change of variables in
\eqref{eq:coer-schur} splits the primitive form into its coercive block
and $S_{T,k}$. Thus its negative index is $n_-(S_{T,k})$ and its radical
is zero. The spectrum of $H_T$ is discrete (Lemma~\ref{lem:compact}), so
zero could lie in it only as an eigenvalue, whose eigenvectors would lie in
the radical; hence zero is not in the spectrum of either parity sector,
and Theorem~\ref{thm:parityschur} gives
$n_-(W_T)=n_-(S_{T,k})+n_-(\mathfrak G_T)$ and
$n_0(W_T)=n_0(\mathfrak G_T)$. The size $2k-2$ counts blocks;
crossing a contact varies the entries continuously at fixed $k$.
\end{remark}



At $T=\tfrac12\log5$ one may take $k=25$ in \eqref{eq:coer-sandwich}.
Galerkin approximation bounds the Schur complement from above; a lower
certificate needs an error bound using $c>0$. For Legendre bases see
\cite[Thms.~1.1, 1.4]{Zhu2026} and \cite{Groskin2026}.

\begin{remark}[the route by operator norms]\label{rem:coer-nogo}
For adjacent blocks the Tate-orthogonal coupling has bilinear norm
$\pi$. In endpoint distances $x,y$ the kernel satisfies
$0<K_1(x+y)\le1/(x+y)$, so Hilbert's inequality gives the upper
bound \cite[Thm.~316]{HLP}. For sharpness take
$f_M(x)=x^{-1/2}\mathbf1_{[2re^{-M},2r]}(x)$ on each block.
Its squared norm is $M$; logarithmic coordinates give the Hilbert
pairing $\int_{[-M,0]^2}(2\cosh((u-v)/2))^{-1}du\,dv=\pi M-O(1)$.
The bounded kernel difference costs $O(r)$ since $\|f_M\|_1^2=O(r)$.
At fixed $r$ the projection onto the two Tate directions has bounded
norm, changing the pairing by $O(\sqrt M)$. Thus in
\eqref{eq:coer-cross} the coupling costs $\pi/2$ for two blocks
and at least this much for $k\ge2$. Subtracting the coupling and
prime norms separately in \eqref{eq:coer-A} requires
$\log(1/r)>1.21+r+\pi/2+1.674$ at $T=\tfrac12\log5$, hence
$k\ge70$ versus $25$ in Theorem~\ref{thm:coer-B}.
\end{remark}


\begin{lemma}[first-contact reduction]\label{lem:firstcontact}
If $\lamA<0$ for some $T$, then there is a window $T_*>\frac12\log5$ and a nonzero $F\in\mathcal P_{T_*}$ such
that $\lambda_{\mathrm A}(T)\ge0$ for $T\le T_*$, $\lambda_{\mathrm A}(T)<0$ for $T>T_*$, and
\[
 H_{T_*}\ge0,\qquad H_{T_*}F=0,\qquad JF=\varepsilon F,\quad \varepsilon\in\{\pm1\}.
\]
\end{lemma}
\begin{proof}
Set $T_*=\inf\{T:\lamA<0\}$. Monotonicity and continuity
(Lemmas~\ref{lem:monoT} and~\ref{lem:contA}) give the asserted signs
and $\lambda_{\mathrm A}(T_*)=0$; Theorem~\ref{thm:log5A} gives the
strict lower bound on $T_*$. Compact resolvent (Lemma~\ref{lem:compact})
gives a zero eigenvector, and reflection preserves its eigenspace.
\end{proof}
\section{A common family of tests for all windows}\label{sec:ladder}
There is one sequence of matrices which suffices to test the sign for every
window. Its entries can be computed before choosing any zeros. For $a>1$
write
\[
 L(a)=\frac1a+\frac1{a-1}-\frac12\log\pi+
 \frac12\psi(a/2)+\frac{\zeta'(a)}{\zeta(a)},\qquad
 K(a,b)=\frac{L(a)+L(b)}{a+b-1}.
\]
For real tests $h$ on the half-line, extended by zero, write $\mathcal W[h]$
for \eqref{eq:W} with all prime powers included. We use it only where
the integral, moments and prime series converge absolutely. Put
$h_a(y)=e^{-(a-1/2)y}\mathbf1_{y>0}$. The polarized form satisfies
$\mathcal W[h_a,h_b]=K(a,b)$. Indeed the positive shifts
in the symmetrized autocorrelation contribute the two Dirichlet series for
$\zeta'/\zeta$, divided by $a+b-1$. The moments contribute
$1/((a-1)b)+1/(a(b-1))$. For the archimedean integral the partial fractions
of the two Fourier transforms give Poisson kernels of widths $a-1/2$ and
$b-1/2$, divided by $a+b-1$. Poisson integration of $m$ supplies the
digamma and constant terms displayed above.

\begin{proposition}\label{ladder}
If $[K(32+i,32+j)]_{0\le i,j\le m}$ is positive semidefinite for every
$m\ge0$, then $W_T\ge0$ on $\mathcal D_T$ for every $T>0$.
\end{proposition}
\begin{proof}
Translate a real smooth compact test $h$ into $(0,\infty)$, preserving its
autocorrelation and moment product. The function $f(x)=x^{-63/2}h(-\log x)$
extends smoothly by zero at both ends of $[0,1]$. Approximate its derivative
by polynomials and integrate to obtain $P_j$ with $\epsilon_j=\|P_j-f\|_{C^1}\to0$.
The functions $e^{-63y/2}P_j(e^{-y})$ belong to the stated spans on $y>0$.
Their errors and first derivatives there are $O(\epsilon_j e^{-63y/2})$;
integration by parts after zero extension, including the boundary value,
bounds the Fourier error by $O(\epsilon_j/(1+|\tau|))$. Hence the
archimedean integral and both Tate moments converge. At $\log n$ the
autocorrelation error is $O(\epsilon_j n^{-63/2})$ by the common exponential
bound, summable against $v_n$. Thus the matrix inequalities pass to $h$,
to complex tests by polarization of the real symmetric form, and to the
form domain by core density (Lemma~\ref{lem:form-density}).
\end{proof}

The hypothesis of Proposition~\ref{ladder} concerns matrices of every
size. Conversely, the sign on all windows forces the sign of all these
matrices.

\begin{proposition}[the converse]\label{prop:ladder-converse}
If $W_T\ge0$ on $\mathcal D_T$ for every $T>0$, then for every $d\ge1$ and
all $a_1,\dots,a_d>1$ the matrix $[K(a_j,a_k)]_{j,k=1}^d$ is positive
semidefinite. In particular the matrices of Proposition~\ref{ladder} are
positive semidefinite in every degree.
\end{proposition}
\begin{proof}
The matrix is real and symmetric, so it suffices to show
$\sum_{j,k}x_jx_kK(a_j,a_k)=\mathcal W[h]\ge0$ for real $x_j$ and
$h=\sum_jx_jh_{a_j}$. Put $\alpha=\min_j a_j-\frac12>\frac12$, so
$|h(y)|\le Ce^{-\alpha y}$ for $y>0$. For $R>0$ let $h_R=h\mathbf 1_{(0,R)}$,
and let $F_R(t)=h_R(t+R/2)$, supported in $[-R/2,R/2]\subset I_T$ for
$T=R/2+1$. The function $F_R$ is bounded with two jumps, so
$|\widehat{F_R}(\tau)|=O(1/(1+|\tau|))$ and $F_R\in\mathcal D_T$. Translation
by $R/2$ does not change $|\widehat{F_R}|$ or the autocorrelation, and
multiplies $A$ and $B$ by $e^{-R/4}$ and $e^{R/4}$, so it does not change
$A\bar B$; and $c_{F_R}(\log n)=0$ for $\log n\ge R$, so the contacts of $W_T$
with $\log n\ge R$ and the prime powers outside $\NT$ contribute nothing.
Hence $\mathcal W[h_R]=W_T(F_R)\ge0$. We let $R\to\infty$.
(i) $h-h_R=\sum_jx_je^{-(a_j-1/2)R}h_{a_j}(\cdot-R)$, and each $h_{a}$ has
$\int\log(2+|\tau|)|\widehat{h_a}|^2<\infty$ since
$\widehat{h_a}(\tau)=(a-\frac12-i\tau)^{-1}$; as $|m(\tau)|\le C+\log(2+|\tau|)$,
the archimedean term of $h_R$ tends to that of $h$.
(ii) $A(h_{a}(\cdot-R))=e^{R/2}/(a-1)$ and $B(h_a(\cdot-R))=e^{-R/2}/a$, so
$A(h-h_R)=O(e^{-(\min_ja_j-1)R})$ and $B(h-h_R)=O(e^{-\alpha R})$; the moments
converge because every $a_j>1$.
(iii) $|c_{h_R}(u)|\le\int_{y>|u|}C^2e^{-\alpha y}e^{-\alpha(y-|u|)}dy\le
\frac{C^2}{2\alpha}e^{-\alpha|u|}$, and $c_{h_R}(u)\to c_h(u)$; since
$\sum_nv_ne^{-\alpha\log n}=2\sum_n\Lambda(n)n^{-\min_ja_j}<\infty$, dominated
convergence gives the convergence of the prime series.
Hence $\mathcal W[h]=\lim_R\mathcal W[h_R]\ge0$.
\end{proof}


\begin{proposition}[the matrices in every degree]\label{prop:all-degree-gram}
Assume \eqref{eq:dirichlet-input}, that is, that $\zeta$ has no zero with
$\Re s>\frac12$ (so that, by the functional equation, every nontrivial zero
lies on the critical line). Let $a_j>1$
and $u_j=a_j-\tfrac12$, $1\le j\le d$. Then, for $c_j\in\mathbb C$,
\[
 \sum_{j,k=1}^d c_j\overline{c_k}K(a_j,a_k)
 =\sum_{\gamma}\left|\sum_{j=1}^d\frac{c_j}{u_j-i\gamma}\right|^2 .
\]
Here $\tfrac12+i\gamma$ ranges over the nontrivial zeros with multiplicity.
In particular, under this assumption the hypothesis of Proposition~\ref{ladder} holds in every degree.
\end{proposition}
\begin{proof}
By evenness the order-one Hadamard product pairs the zeros at $\pm i\gamma$
of $\xi(\tfrac12+z)$; $N(H)=O(H\log H)$ gives locally uniform convergence
of the paired logarithmic derivatives away from the zeros. Thus
$L(\tfrac12+u)=\sum_\gamma(u-i\gamma)^{-1}$ in symmetric order.
Pair the terms at $\gamma,-\gamma$ in the two evaluations defining $K$:
$[(u-i\gamma)^{-1}+(v+i\gamma)^{-1}]/(u+v)
=[(u-i\gamma)(v+i\gamma)]^{-1}$. Summing gives the identity.
\end{proof}

\section{A prime sum on growing intervals}\label{sec:primesum}
\setlength{\abovedisplayskip}{1.8pt plus 2pt minus 1pt}
\setlength{\belowdisplayskip}{1.8pt plus 2pt minus 1pt}
The zero condition we need is
\begin{equation}\label{eq:dirichlet-input}
 \zeta(s)\ne0\qquad (\Re s>\tfrac12,\ s\ne1).
\end{equation}

For $x\ge1$ put $G(x)=\sum_{n\le x}\Lambda(n)(x/n-1)-x\log x+
(1+\gamma_E)x$, where $\gamma_E$ is Euler's constant.
There is also a pointwise comparison with a finite average near $x$.
Away from integers,
$G'(t)=\sum_{n\le t}\Lambda(n)/n-\log t+\gamma_E$.
With $m=\lfloor t\rfloor$, the identity
$\log(m!)=\sum_{n\le t}\Lambda(n)\lfloor t/n\rfloor$ and
Chebyshev's bound $\sum_{n\le t}\Lambda(n)=O(t)$ give
$\sum_{n\le t}\Lambda(n)/n=t^{-1}\log(m!)+O(1)=\log t+O(1)$.
Thus $G$ has bounded derivative between integers and is continuous
at each integer. For any $y_j\ge1$ with $|y_j-x|\le\sqrt{x}$ and
nonnegative weights of sum one, it follows that
$|G(x)-\sum_j w_jG(y_j)|\ll\sqrt{x}$ with an absolute constant.
This permits a bound for such an average to be used on the same
window, provided the bound holds for its chosen sample points.

An upper bound on this sum can be checked at integers. Indeed, on
$[n,n+1]$, $n\ge1$, the prime sum is affine (the new term at
$n+1$ vanishes there), so $G''(x)=-1/x$ in the interior. If $L_n$
interpolates the endpoint values linearly, concavity gives
$G\ge L_n$. Moreover $G-L_n-(x-n)(n+1-x)/(2n)$ is convex and
vanishes at both ends. Hence
\[
 0\le G(x)-L_n(x)\le (x-n)(n+1-x)/(2n)\le 1/(8n).
\]
In particular, if $G(k)\le C k^\theta$ at all integers of
$[Y,Y^J]$, with $C>0$ and $0<\theta<1$, concavity of $x^\theta$
gives $G(x)\le Cx^\theta+1/(8\lfloor x\rfloor)$ between the first
and last such integers. For fixed $1<J'<J$ this includes
$[\lceil Y\rceil,\lceil Y\rceil^{J'}]$ once $Y$ is large.
More generally, two sample points $1\le a<b$ suffice for an
upper bound between them. Let $L$ be the chord of $G$ on $[a,b]$.
The function $G(x)-L(x)-(x-a)(b-x)/(2a)$ is convex there:
its second derivative between integers is $1/a-1/x\ge0$,
and its derivative jumps upwards by $\Lambda(n)/n$ at each
interior integer $n$. Its endpoint values vanish. Consequently,
if $G(t)\le Ct^\theta$ at both endpoints, then
\[
 G(x)\le Cx^\theta+(b-a)^2/(8a)\qquad(a\le x\le b).
\]
Thus consecutive successful samples with $b-a\le D a^{3/4}$
cost at most $(D^2/8)\sqrt{x}$ throughout their interval.
Here $D$ must be uniform over the windows under consideration.

The Mellin transform of $G$ on $\Re s>1$ is
\[
 P(s)=-\frac{\zeta'(s)}{s(s-1)\zeta(s)}
       -\frac1{(s-1)^2}+\frac{1+\gamma_E}{s-1}.
\]
Indeed the integral of $(x/n-1)x^{-s-1}$ from $n$ to infinity
is $n^{-s}/[s(s-1)]$, and the Euler product converges absolutely.
Writing $t=s-1$, the expansions $-\zeta'/\zeta=t^{-1}-\gamma_E+O(t)$
and $1/[s(s-1)]=t^{-1}-1+O(t)$ give the singular terms
$t^{-2}-(1+\gamma_E)t^{-1}$, exactly canceled in $P$.
At a zero $\rho$ of multiplicity $m_\rho$ in $\Re s>0$, the residue
of $P$ is $-m_\rho/[\rho(\rho-1)]\ne0$.
For real $0<s<1$, grouping successive terms shows
$(1-2^{1-s})\zeta(s)>0$, while the Euler product gives
$\zeta(s)\ne0$ for real $s>1$. Thus all poles of $P$ in this
half-plane have nonzero ordinate.

\begin{proposition}\label{prop:growing-window}
Fix $1/2<\theta<1$ and $J>1$. If $G(x)\le C_\theta x^\theta$
throughout $[Y,Y^J]$ for arbitrarily large $Y$ with a common
$C_\theta$, then $\zeta$ has no zero in $\Re s>\theta$.
\end{proposition}
\begin{proof}
The function $g(u)=G(e^u)$ has Laplace transform $P$ and $|g(u)|\ll e^{2u}$
by $\Lambda(n)\le\log n$. Each pole just identified is simple, regardless of the zero multiplicity.

Put $U=\log Y$, let $\sigma>2$ be fixed, and take
$m=(1+J)/2$ and $k=\lfloor m\sigma U\rfloor$. Set $A_k(z)=(-1)^kP^{(k)}(z)/k!$.
With $d\nu_k=u^k e^{-\sigma u}\,du/k!$, the triangle inequality gives
$|A_k(\sigma+it)|+A_k(\sigma)\le
\int_0^\infty(|g|+g)\,d\nu_k=2\int_0^\infty\max(g,0)\,d\nu_k$.
On $[U,JU]$ the positive part is at most
$\max(C_\theta,0)e^{\theta u}$, whose integral over the whole half-line
is $\max(C_\theta,0)/(\sigma-\theta)^{k+1}$. Thus uniformly in $t$,
\[
 |A_k(\sigma+it)|+\Re A_k(\sigma)
 \le {2\max(C_\theta,0)\over(\sigma-\theta)^{k+1}}
 +O\!\left(\int_{u\notin[U,JU],\ u\ge0}
          {u^k e^{-(\sigma-2)u}\over k!}\,du\right).
\]
Put $a=\sigma-2$. The tail times $a^{k+1}$ is the probability that
a gamma variable $V$ of shape $k+1$ and rate $a$ lies outside
$[U,JU]$. Choose $1/m<v_-<1<v_+<J/m$. The normalized endpoints
$aU/(k+1),aJU/(k+1)$ tend to $(1-2/\sigma)/m,J(1-2/\sigma)/m$.
For large fixed $\sigma$ and large $U$ they enclose $[v_-,v_+]$.
Markov's inequality applied to $\exp((1-1/v)aV)$ bounds each
corresponding tail by $(ve^{1-v})^{k+1}$, $v=v_\pm$. Since both
bases are below one, choose $\sigma$ so that they remain below one
after multiplication by $(\sigma-\theta)/a$. The displayed bound is
then $O((\sigma-\theta)^{-k})$ along the chosen windows.

If $\Re\rho>\theta$ for some pole, let $D$ be the lesser of
$\sigma-\theta$ and the distance from $\sigma$ to the poles.
A centre $z=\sigma+it$ has a pole at distance $d<D$: when
$D=\sigma-\theta$ take $t=\Im\rho$, and otherwise move slightly
toward a nearest pole's nonzero ordinate. Perturb $t$ within the
strict inequality to make the nearest pole unique. Indeed the
squared distance to $\beta+i\gamma$ is
$(\sigma-\beta)^2+(t-\gamma)^2$. Poles of different ordinates
tie at most once, and poles of the same ordinate and different real
parts never tie since they lie to the left of $\sigma$. The excluded
ordinates are countable. Write $\rho_0$ for the nearest pole at the
chosen centre and $d=|z-\rho_0|$. By discreteness there is a radius
$d<r<D$ enclosing no other pole about $z$. With
$R_0=\operatorname{Res}_{\rho_0}P\ne0$, Cauchy's estimate on this
circle gives $A_k(z)=R_0/(z-\rho_0)^{k+1}+O(r^{-k})$.
The circle of the same radius about $\sigma$ encloses no pole,
so $A_k(\sigma)=O(r^{-k})$. Multiplying the window bound by
$d^{k+1}$ and letting $U\to\infty$ now gives $|R_0|\le0$,
since $d<r<D\le\sigma-\theta$.
\end{proof}

We specialize the character-family comparison of \cite{CharacterFamily}
to modulus one. The retained characters are the primitive characters
induced by powers of the chosen character. Every power of the
character modulo one is trivial, so the retained sum consists of
$G$ alone. The family comparison uses division and loss exponents $e_D,e_L$, which
are fixed before $\varepsilon$ is chosen; their values, the translation
factor $u^{e_D+e_L+3}$ and the resulting exponent are recorded once, in
Remark~\ref{rem:exponents}. A common family for the column and twist
estimates is given by Theorem~\ref{thm:alpoge-families}(a),(b). The
existence of members carrying the simultaneous estimates is the hypothesis
of Proposition~\ref{prop:simultaneous-slope}.
At modulus one the auxiliary subgroup has index one in $A$. The four
characters of $\{\pm1\}^2\times A$ trivial on it are generated by
$\chi_q=(\cdot/q)$ and $\chi_d=(\cdot/d)$ of coprime conductors; for the
hosts of Theorem~\ref{thm:alpoge-families}(b) this is
Lemma~\ref{lem:host-facts}(iv) below.
Removing the trivial character leaves the absolute error
\[
 R(x)=|G_{\chi_q}(x)|+|G_{\chi_d}(x)|+|G_{\chi_q\chi_d}(x)|,
 \qquad
 G_\chi(x)=\sum_{n\le x}\Lambda(n)\chi(n)(x/n-1)
             +x\frac{L'(1,\chi)}{L(1,\chi)} .
\]
The logarithmic derivative here is finite for each of the three
quadratic characters: for every nonprincipal Dirichlet character $\chi$ the
function $L(s,\chi)$ continues holomorphically to $\Re s>0$ and
$L(1,\chi)\ne0$ (Dirichlet's theorem; see \cite[Chapter~4]{Davenport}).

\begin{proposition}[division at one point]\label{prop:finite-to-window}
Fix $\varepsilon>0$, $C>0$ and $K\ge0$. Let $x\ge1$ and write
$u=x^\varepsilon$. Suppose that for a pair $(q,d)$ both of the following
estimates hold at $x$. First, $|G_\chi(x)|\le\sqrt{x}\,u$
for $\chi_q,\chi_d,\chi_q\chi_d$. Second, with
$D=G(x)-R(x)$, the comparison data satisfy
\begin{equation}\label{eq:finite-comparison}
 \frac{rmb\delta}{2}(D-4\sqrt{x}\,u)
 \le Cu^K\left[ra(1+m^2/k)
      +m\left(2nmbu+\frac{skbu^2}{n^2}\right)\right]
      +\frac{2rmb}{x}.
\end{equation}
Here $n=|A|$, $r,s$ are the sizes of the averaging sets,
$\delta=s/n\ge1/u$, $r\ge n/(2u^{1/8})$, $b>0$, $0\le a\le2xb$,
$\sqrt{x}/u\le m\le\sqrt{x}\,u^{1/2}$ and
$xm\le k\le n^2\sqrt{x}/u$.
Then
\[
 G(x)\le (20C+11)\sqrt{x}\,u^{K+21/8}.
\]
In particular, for $C=1$ and $K=5$ one has $G(x)\le 31\sqrt{x}\,u^{61/8}$,
and $G(x)\le x^{1/2+9\varepsilon}$ once $u^{11/8}\ge31$.
\end{proposition}
\begin{proof}
After division by $rmb\delta/2$, the three terms in brackets in
\eqref{eq:finite-comparison}, before multiplication by $Cu^K$, satisfy
\[
 \frac{2a}{b\delta}(m^{-1}+m/k)\le4\sqrt{x}\,u^2+4u,\qquad
 \frac{4nmu}{r\delta}\le8\sqrt{x}\,u^{21/8},\qquad
 \frac{2sku^2}{r\delta n^2}\le4\sqrt{x}\,u^{9/8}.
\]
The first uses $k\ge xm$, the second uses $n/r\le2u^{1/8}$,
and the third also uses $s/\delta=n$. Since $x,u\ge1$ the three bounds have
sum at most $20\sqrt{x}\,u^{21/8}$. The additive terms outside
the brackets are at most $4\sqrt{x}\,u+4u/x$.
Finally $G=D+R$ with $R\le3\sqrt{x}\,u$, so the terms outside the brackets,
including $R$, have sum at most $11\sqrt{x}\,u$. Summing gives the stated
bound. For the last assertion take $u^{11/8}\ge31$.
\end{proof}

The uniformity needed for zero exclusion concerns the exponent
of the loss. More precisely, take in Proposition~\ref{prop:finite-to-window}
$C=C_\varepsilon$ and a fixed $K\ge0$, with $C_\varepsilon>0$ independent
of the column, $x$ and the window. Then
\[
 G(x)\le (20C_\varepsilon+11)\sqrt{x}\,u^{K+21/8}.
\]
Given $\theta>1/2$, fix an admissible $\varepsilon$ with
$(K+21/8)\varepsilon<\theta-1/2$. Then $20C_\varepsilon+11$
is a single constant for all the selected windows, exactly as
required by Proposition~\ref{prop:growing-window}. In particular
no bound on the growth of $C_\varepsilon$ as $\varepsilon$
decreases is required.

The primes $q,d$, the host and the twist bounds at every $x$ of arbitrarily
long windows $[W,W^{3/2}]$ are supplied by
Theorem~\ref{thm:alpoge-families}(a),(b)(2), which is derived in
\cite[Section~\ref{A-i4:sec}]{Annex}.

The signed character sum entering the comparison admits a direct
prime-power formula. The character average over $A$ kills every
character nontrivial on $A$: translating the sum by an element on
which the character differs from one proves this assertion. On the
four surviving characters, evaluation at the sign element $(+1,-1)$
multiplies the $d$ and $qd$ terms by $-1$. Thus the quantity compared
to the arithmetic count is
\[
 H_{q,d}=G+G_{\chi_q}-G_{\chi_d}-G_{\chi_q\chi_d},
 \qquad G-R\le H_{q,d}.
\]
We give the arithmetic interpretation of the comparison. Use additive
coordinates $U=(\mathbb Z/2)^2\times A$, with the $d$ coordinate
first. Write $V=\{g:g_1=1\}$, so $|V|=2n$. For a prime power
$p^h$, let $\mu_{p^h}$ be uniform measure on the coset of $h$
times Frobenius modulo inertia in $U$. Its Fourier coefficient
$c_\xi(p^h)=\sum_g\xi(g)\mu_{p^h}(g)$ has modulus at most one.
Away from ramification this is just the character at Frobenius;
at a ramified prime the coset average is zero unless the character
is trivial on inertia. Define $G_\xi$ by the preceding prime-power
formula with these coefficients and the primitive logarithmic
derivative $B_\xi$, using $G$ for the trivial character. For $0<y<1$ the same
formulas apply with empty prime sums.

The column weights $\widetilde w$ are indexed by
$\Phi=\{(0,c,h)\}\subset U$, the CM type of the host. At $t\in A$ the weight attached
to a Frobenius class $g=(1,c,h)\in V$ is
\[
 w_g=\widetilde w(0,c+1,h+t).
\]
Indeed complex conjugation adds $(1,1,0)$ and the translation by $t$ adds
$(0,0,t)$. This translation bijects $V$ with $\Phi$, so that
$\sum_{g\in V}\log w_g=\sum_{\rho\in\Phi}\log\widetilde w(\rho)$ for every
$t$. In particular the logarithmic mean is independent of $t$ for
a fixed column, although its individual weights are permuted.
For $a>0$ put
\[
 \begin{aligned}
 P(a,w)&=2\sum_{g\in V}\sum_{p^h}(\log p)\mu_{p^h}(g)
                        (a/p^h-w_g)_+,\\
 M(a,w)&=a\left(\log a-1-\gamma_E+B_{\chi_d}
                   -\frac1{2n}\sum_{g\in V}\log w_g\right).
 \end{aligned}
\]
\begin{lemma}[the arithmetic of the host]\label{lem:host-facts}
Let $q,d,Q_S$, $N_E=Q_Sqd$, $E$, $\mathrm{Art}$, $I_p$ and $\mathrm{Fr}_p$ be
as in Theorem~\ref{thm:alpoge-families}(b), and use the additive
coordinates above, in which the first coordinate of $g\in U$ is $0$ or $1$
according as its $d$-sign is $1$ or $-1$, so that
$\Phi=\{g:g_1=0\}$.
\begin{enumerate}
\item[(i)] If $p\nmid N_E$, then $I_p=0$ and $\mathrm{Fr}_p=\mathrm{Art}(p)$.
\item[(ii)] If $p=2$, $p=q$, or $p$ is an odd prime dividing $Q_S$, then
$I_p\subset\Phi$ and $\mathrm{Fr}_p\in\Phi$.
\item[(iii)] $I_d=\{0,(1,0,0)\}$.
\item[(iv)] The characters of $U$ trivial on $A$ are the trivial character
and the $q$-sign, the $d$-sign and their product, whose composites with
$\mathrm{Art}$ are $\chi_q=(\cdot/q)$, $\chi_d=(\cdot/d)$ and $\chi_q\chi_d$.
\end{enumerate}
\end{lemma}
\begin{proof}
Write $N'_p$ for the prime-to-$p$ part of $N_E$. By definition $I_p$ is
the image under $\mathrm{Art}$ of the units $u\equiv1\pmod{N'_p}$, and
$\mathrm{Fr}_p=\mathrm{Art}(u_p)$ with $u_p\equiv p\pmod{N'_p}$; the $d$-sign
of $\mathrm{Art}(u)$ is $(u/d)$ and its $q$-sign is $(u/q)$. (i) If
$p\nmid N_E$ then $N'_p=N_E$, so $I_p=\mathrm{Art}(1)=0$ and $u_p\equiv p$.
(ii) Here $d\mid N'_p$. A unit $u\equiv1\pmod{N'_p}$ has $d$-sign
$(1/d)=1$, so $I_p\subset\Phi$, and $\mathrm{Fr}_p$ has $d$-sign $(p/d)$.
Now $(2/d)=1$ since $d\equiv7\pmod8$; and for an odd prime $l\ne d$,
quadratic reciprocity with $d\equiv3\pmod4$ gives $(l/d)=(-d/l)$, which is
$1$ for $l=q$ and for $l\mid Q_S$ by the host conditions. (iii) Here
$N'_d=Q_Sq$, so $I_d\subset\{0,\iota_d\}$ by the last condition on
$\mathrm{Art}$ in Theorem~\ref{thm:alpoge-families}(b), and $\iota_d=(1,0,0)$ in additive
coordinates. A unit $u\equiv1\pmod{Q_Sq}$ with $(u/d)=-1$, which exists by
the Chinese remainder theorem, has image $\iota_d$. (iv) A character of $U$
trivial on $A$ is a character of $(\Z/2)^2$, hence a product of powers of
the two signs; compose with $\mathrm{Art}$.
\end{proof}

For comparison with the column count, take an integer
$N\ge\max_{g\in V}a/w_g$. For the data of Theorem~\ref{thm:alpoge-families}(b)
such a cutoff is given by property~(3). A positive summand has $p^h<a_t/w_g$,
so both $p$ and $h$ are at most $N$. At $2$, at $q$ and at every odd prime $l\mid Q_S$ the
Frobenius and inertia lie in $\Phi$ (Lemma~\ref{lem:host-facts}(ii)),
that is, they have trivial $d$ coordinate, hence the coset
misses $V$. Away from the primes dividing $Q_Sqd$ the coset is a singleton
(Lemma~\ref{lem:host-facts}(i)), contributing
with coefficient $2\log p$ precisely when its first coordinate is one.
At $d$ the inertia group has order two (Lemma~\ref{lem:host-facts}(iii)),
so each coset element in $V$ contributes with coefficient $\log d$. These are exactly the unramified
and ramified summands of the truncated column count with the displayed
translation of the weights.
\begin{lemma}[the character identity]\label{lem:character-identity}
For $a>0$ and positive weights $(w_g)_{g\in V}$,
\[
 P(a,w)-M(a,w)=\frac1{2n}\sum_{g\in V}\sum_{\xi\in\widehat U}
                   \overline{\xi(g)}\,w_gG_\xi(a/w_g).
\]
\end{lemma}
\begin{proof}
Inversion gives $|U|\mu_{p^h}(g)$; the principal correction
is $-a(\log a-1-\gamma_E)+(a/2n)\sum_{g\in V}\log w_g$.
The only other character constant on $V$ is $\chi_d$, with value
$-1$, giving the remaining correction $-aB_{\chi_d}$.
\end{proof}

To identify the finite count at its actual weights, write $h_p$ for
the class of an odd prime in $A$ and $c_p\in\{0,1\}$ for the
additive value of $\chi_q(p)$ when $p\ne q$. At an unramified
$p$ only $\chi_d(p)=-1$ contributes, and then only odd powers:
the summand of $P$ at $p^i$ is
$2\log p\,(a/p^i-w_{(1,c_p,i h_p)})_+$ for odd $i$.
At $2$, at $q$ and at every odd prime $l\mid Q_S$ the $d$-character
splits and inertia lies in $\Phi$ (Lemma~\ref{lem:host-facts}(ii)), so their
inertia cosets miss $V$ at every power. At $d$ inertia varies
only the first coordinate, and $(d/q)=-1$. Exactly half of its coset
lies in $V$, giving the single term
$\log d\,(a/d^i-w_{(1,i\bmod2,i h_d)})_+$ for every $i\ge1$.
This term enters the count $P$ exactly (it is the $d$-summand of
$P(a,t,w)$ in \S\ref{subsec:assembly-dict}), so no estimate of its size is
needed. As noted above, restricting both $p$ and $i$ to $[1,N]$ is exact;
this applies separately to each numerator and translation before averaging.

The averaging sets $S,T\subset A$, of sizes $s$ and $r$, are part of the
data of Theorem~\ref{thm:alpoge-families}(b).

\begin{lemma}[the prime-power inequality]\label{lem:prime-power}
There is an absolute constant $x_*\ge3$ with the following property. Let
$0<\varepsilon<1/52$, $x\ge x_*$ and $u=x^\varepsilon$. Suppose that the
primes $q,d$, the host data, the sets $T,S\subset A$, the numerators
$a_t$, the weights $w_{v,j}$ and the numbers $a,b,b',\eta,X',\vartheta$
satisfy properties (1)--(4) of Theorem~\ref{thm:alpoge-families}(b) at $x$;
of (2) only the bounds $|G_\chi(x)|\le\sqrt x\,u$ are used, not the bound
on $|L'(1,\chi)/L(1,\chi)|$. Then, with $\delta=s/n$, with $P_*,M_*$ as in
\S\ref{subsec:assembly-dict}, and with $\mathbb E$ the mean over the
translations $v\in A$,
\begin{equation}\label{eq:prime-power}
 \frac{rmb\delta}{2}\bigl(G(x)-R(x)-4\sqrt{x}\,u\bigr)
 \le \mathbb E(P_*-M_*)+2rmb/x .
\end{equation}
\end{lemma}
\begin{proof}
The retained weights and the numerators below are those of (3); the bound
at discarded weights is the thin-scale bound of (4), which applies because
$a_t/(2b')\ge(x-1)X'/(2x)\ge X'/3$ for $x\ge3$ by (3); and the cost
condition used below is the sum of the two inequalities in (4)
(Remark~\ref{rem:alpoge-families}).
For $v\in A$, $t\in T$, use the subset
$V_{v,t}=\{(1,0,h):h\in S+v-t\}$. It is the inverse image of
$\{(0,1,h'):h'\in S+v\}$ under the preceding translation.
In the character sum above
retain only $g\in V_{v,t}$ and replace their weights by $b$.
At a fixed $t$ use a numerator $a_t$ independent of $v$. Each
$h\in A$ occurs for exactly $s$ choices of $v$, and summing over
$h$ kills the characters nontrivial on $A$. The average over $v$
of the resulting expression is exactly
$b\delta H_{q,d}(a_t/b)/2$.

Suppose the retained weights lie in $[b,b+\eta+1]$, $a_t/b\le x+1$,
and $|G_\xi(a_t/w_g)|\le(a_t/w_g)\vartheta$ at discarded weights.
Between breakpoints the derivative of $wG_\xi(a_t/w)$ is bounded by
$\psi(a_t/w)$ for nontrivial $\xi$ and equals $a_t/w-\psi(a_t/w)$
for the trivial character. Here $\psi(y)=\sum_{n\le y}\Lambda(n)$ is Chebyshev's
function (not the digamma function of the earlier sections), and
$\psi(y)\le4y\log2<3y$ by summing
$\psi(2k)-\psi(k)\le\log\binom{2k}k\le2k\log2$ up to the next power
of two. Continuity and $|\widehat U|=4n$ bound the error per numerator
by $4n(x+1)[3(\eta+1)+b\vartheta]$.

Summing over the $rm$ numerators costs at most $2rmb/x$ when
$4n(x+1)[3(\eta+1)+b\vartheta]\le2b/x$. For $|y_t-x|\le1$,
$y_t=a_t/b$, the character estimates at $x$ give
$|B_\chi|\le x^{-1/2+\varepsilon}+\sum_{n\le x}\Lambda(n)/n
\le c_0(1+\log x)$, with $c_0$ absolute, by Chebyshev's bound and partial
summation; only the twist bound of (2) enters here. Hence
$|H'_{q,d}|\le c(1+\log x)$ between breakpoints within distance one of $x$,
with $c$ absolute, while $H_{q,d}(x)\ge G(x)-R(x)$. By continuity
$H_{q,d}(y_t)\ge G(x)-R(x)-c(1+\log x)\ge G(x)-R(x)-4\sqrt{x}\,u$ as soon
as $c(1+\log x)\le4\sqrt x$; we take $x_*\ge3$ such that this holds for
every $x\ge x_*$. Thus $x_*$ is absolute. Summing over
$t\in T$ and $1\le j\le m$ bounds the left side of \eqref{eq:prime-power}
by the mean of $P-M$ plus $2rmb/x$.

For a translation $v$, the sum over $t\in T$ and $1\le j\le m$ of
$P(a_t,w_{v,j})-M(a_t,w_{v,j})$, with the weights translated as above, is
$P_*(v)-M_*(v)$ in the notation of \S\ref{subsec:assembly-dict}, so the mean
in question is $\mathbb E(P_*-M_*)$. (For the data of
Theorem~\ref{thm:alpoge-families}(b), property (5) gives one lift of the
package for each $v\in A$, whose count and principal term are $P_*(v)$ and
$M_*(v)$; a bound at every $v$ gives the same bound for
$\mathbb E(P_*-M_*)$. This is how the lemma is used in
Corollary~\ref{cor:obligation}.)
\end{proof}
Consequently a bound $\mathbb E(P_*-M_*)\le Cu^KZ$, with $Z$ the bracket in
\eqref{eq:finite-comparison}, gives \eqref{eq:finite-comparison} at $x$.

The arithmetic inputs behind Theorem~\ref{thm:alpoge-families} are proved in
\cite[Sections~\ref{A-i1:sec}--\ref{A-i4:sec}]{Annex} and summarised in
Section~\ref{sec:inputs}.

\begin{proposition}[the column criterion]\label{prop:simultaneous-slope}
Let $\varepsilon_j\downarrow0$ with $\varepsilon_j<1/52$, and let $C_j>0$
and $K_j\ge0$ with $\liminf_j\varepsilon_jK_j=0$. Suppose that for each $j$
there are a threshold $x_j\ge x_*$, with $x_*$ the absolute constant of
Lemma~\ref{lem:prime-power}, and arbitrarily large $W$ such that at every
$x\in[W,W^{3/2}]$ with $x\ge x_j$ there are data satisfying the hypotheses
of Lemma~\ref{lem:prime-power} with $\varepsilon=\varepsilon_j$, and, on
those same data,
\[
 \mathbb E(P_*-M_*)\le C_j u^{K_j} Z,\qquad u=x^{\varepsilon_j},
\]
where $Z$ is the bracket in \eqref{eq:finite-comparison}. Then for every
$1/2<\theta<1$ there are $C_\theta>0$ and arbitrarily large $Y$ with
$G(x)\le C_\theta x^\theta$ on $[Y,Y^{3/2}]$. Consequently $\zeta$ has no
zero in $\Re s>1/2$.
\end{proposition}
\begin{proof}
Fix $1/2<\theta<1$, and fix $j$ with
$1/2+(K_j+21/8)\varepsilon_j<\theta$; it exists since $\varepsilon_j\to0$
and $\liminf_j\varepsilon_jK_j=0$. Let $W\ge x_j$ be one of the given
windows and $x\in[W,W^{3/2}]$. Lemma~\ref{lem:prime-power} and the mean
bound give \eqref{eq:finite-comparison} at $x$ with $C=C_j$ and $K=K_j$.
The sizes and twist bounds required in
Proposition~\ref{prop:finite-to-window} are properties (1) and (2) of
Theorem~\ref{thm:alpoge-families}(b). Hence
$G(x)\le(20C_j+11)\sqrt{x}\,u^{K_j+21/8}\le(20C_j+11)x^\theta$ at every $x$
of $[W,W^{3/2}]$. The constant does not depend on $W$, so with $Y=W$ the
hypothesis of Proposition~\ref{prop:growing-window} holds with $J=3/2$, and
it excludes every zero to the right of $\theta$. Since $\theta>1/2$ was
arbitrary, the last assertion follows.
\end{proof}

\begin{remark}\label{rem:simultaneous-slope}
A multiplicative constant depending only on $\varepsilon_j$ can be included
in $C_j$. An extra factor $x^{h_j}$ in the final mean estimate adds $h_j$ to
$\varepsilon_jK_j$, since $x^{h_j}=u^{h_j/\varepsilon_j}$. The pair of
primes may depend on $x$. If the mean bound were known only at sample points
of the windows with successive gaps at most $Da^{3/4}$, with $D$ uniform
over the windows (for the fourth powers $k^4$ one has $D=15$), the chord
bound preceding Proposition~\ref{prop:growing-window} would add at most
$(D^2/8)\sqrt{x}$ between successive samples, and the same conclusion would
hold on windows $[Y,Y^{J}]$ with any fixed $1<J<3/2$. This is not needed:
Corollary~\ref{cor:obligation} gives the bound at every $x$ of the window.
\end{remark}


\section{The mean bound and the frame slope inequality}\label{sec:meanbound}
The criterion of the preceding section requires one estimate: the upper
bound for the column mean $\mathbb E(P_*-M_*)$ in
Proposition~\ref{prop:simultaneous-slope}, on the same fields and scales as
the three twist estimates, with an exponent loss that tends to zero. This
section proves it for column packages. The proof has three parts.
\begin{enumerate}
\item A slope inequality for Hermitian lattices carrying almost tight
frames (Theorem~\ref{thm:frameF}), proved in full in
\S\S\ref{subsec:HL}--\ref{subsec:frame-thm}. In the deduction below it
enters as condition (P8) of a column package; for the packages of
Theorem~\ref{thm:alpoge-families}(b)(5) that condition is supplied with the
package, and its verification for the columns is part of the formally
verified \cite[Proposition~\ref{A-i4:Lprop-fam}(b)]{Annex}.
\item A list of arithmetic estimates for the columns attached to one host
field, collected in a \emph{column package} (Definition~\ref{def:package}
in \S\ref{subsec:assembly-dict}). These are the loss bound, the divisor
bound and the level counts. The existence of packages with uniform
constants on arbitrarily large windows is Theorem~\ref{thm:alpoge-families};
it is due to L.~Alpoge \cite{CharacterFamily}, and we state it in
Section~\ref{sec:inputs}.
\item A short deduction (Theorem~\ref{thm:mean-bound}): for a column
package, the slope inequality gives
$P_*-M_*\le u^{e_D+e_L+3}Z$ at every translation, with constant one and an
exponent that does not depend on~$\varepsilon$. This is proved in full in
\S\ref{subsec:assembly-T5}.
\end{enumerate}
Corollary~\ref{cor:obligation} in Section~\ref{sec:inputs} then verifies the
hypothesis of Proposition~\ref{prop:simultaneous-slope}.


The first three subsections prove a lower bound for the mean slope of a ``column'' lattice in terms of a family of Hermitian lattices that map into it (Theorem~\ref{thm:frameF}). This is the only place where the geometry of numbers enters the mean bound. Its proof has three parts:
\begin{enumerate}
\item the Harder--Narasimhan formalism for Hermitian lattices over a totally complex field and the slope inequality for morphisms (\S\ref{subsec:HL}), both proved in full;
\item a linear-algebra statement about almost tight frames at a single archimedean place (Proposition~\ref{prop:frameLA});
\item an elementary counting argument and a layer-cake integration.
\end{enumerate}

\subsection{Hermitian lattices and slopes}\label{subsec:HL}

Throughout, $K$ is a totally complex number field of degree $n_K=[K:\Q]$ and discriminant $d_K$. We write $\Sigma_K$ for its set of infinite places, so $\#\Sigma_K=n_K/2$, and $\mathcal O_K$ for its ring of integers. For each $v\in\Sigma_K$ fix an embedding $\sigma_v\colon K\to\C$ representing $v$.

\begin{definition}[Metrised spaces]\label{def:metrised}
A \emph{metrised $K$-space} is a triple $(V,b,X)$ consisting of
\begin{itemize}
\item a finite-dimensional $K$-vector space $V$,
\item a $K$-basis $b=(b_j)_{j\in\iota}$ of $V$, and
\item invertible matrices $X_v\in\mathrm{GL}_\iota(\C)$, one for each $v\in\Sigma_K$.
\end{itemize}
For $x=\sum_jx_jb_j\in V$ and $v\in\Sigma_K$, the \emph{$v$-coordinates} of $x$ are $f_v(x)=X_v^{-1}\bigl(\sigma_v(x_j)\bigr)_j\in\C^\iota$. The \emph{realisation} of $V$ is the additive map
\[
\rho\colon V\longrightarrow E_V:=\bigoplus_{v\in\Sigma_K}\C^\iota\cong\R^{n_K\dim_KV},\qquad x\longmapsto(f_v(x))_v,
\]
where $E_V$ carries its standard Euclidean structure.
\end{definition}

\begin{lemma}\label{lem:realisation}
Let $L\subseteq V$ be a finitely generated $\mathcal O_K$-submodule and put $r=\operatorname{rk}_KL:=\dim_KKL$.
\begin{enumerate}
\item[(a)] $\rho$ maps $L$ injectively onto a discrete subgroup $\rho L$ of $E_V$ of $\Z$-rank $\operatorname{rk}L=n_Kr$.
\item[(b)] For each $v$, the $\C$-span of $f_v(L)$ in $\C^\iota$ has dimension $r$.
\end{enumerate}
\end{lemma}
\begin{proof}
Choose a $K$-basis $e_1,\dots,e_r$ of $KL$ inside $L$, and complete it to a basis of $V$. The coordinate matrix of $e_1,\dots,e_r$ has a nonzero $r\times r$ minor $\Delta\in K$. Applying $\sigma_v$ and the invertible matrix $X_v^{-1}$ preserves the nonvanishing of the corresponding exterior product, since $\sigma_v(\Delta)\neq0$. Hence $f_v(e_1),\dots,f_v(e_r)$ are $\C$-linearly independent. They span the $\C$-span of $f_v(L)$, because $L\subseteq\sum_aKe_a$ and $f_v$ is $\sigma_v$-semilinear. This proves (b).

For (a), $L$ is a finitely generated torsion-free $\mathcal O_K$-module, contained in $\sum_a K e_a$ and containing $\sum_a\mathcal O_Ke_a$. So it is a free $\Z$-module of rank $n_Kr$, commensurable with $\bigoplus_a\mathcal O_Ke_a$. The map $\mathcal O_K\to\prod_{v}\C$, $\alpha\mapsto(\sigma_v\alpha)_v$, is the Minkowski embedding, which maps $\mathcal O_K$ isomorphically onto a lattice in $\prod_v\C\cong\R^{n_K}$. By (b), $\rho$ restricted to $\bigoplus_a\mathcal O_Ke_a$ is the composition of $r$ copies of this embedding with the injective $\R$-linear map $(\C^{\Sigma_K})^r\to E_V$ given at each $v$ by the matrix $X_v^{-1}(\sigma_v e_a)_a$. Therefore $\rho$ is injective on $L$ and $\rho L$ is discrete of rank $n_Kr$.
\end{proof}

For a discrete subgroup $\Lambda\subset E$ of a Euclidean space with $\Z$-basis $e_1,\dots,e_m$, we define
\[
\operatorname{covol}\Lambda=\det(\langle e_a,e_c\rangle)^{1/2},\qquad \deg\Lambda=-\log\operatorname{covol}\Lambda,\qquad \operatorname{sl}\Lambda=\deg\Lambda/\operatorname{rk}\Lambda,
\]
with $\operatorname{covol}0=1$. For submodules $N\subseteq L$ we abbreviate $\deg N=\deg\rho N$ and $\operatorname{rk}N=\operatorname{rk}\rho N$.

We call $N\subseteq L$ \emph{saturated} if $N=L\cap KN$. For saturated $N$, let $\pi_N$ be the orthogonal projection of $E_V$ onto $(\R\rho N)^\perp$. The \emph{quotient lattice} $\rho L/\rho N:=\pi_N(\rho L)$ satisfies
\begin{equation}\label{eq:covol-mult}
\operatorname{covol}\rho L=\operatorname{covol}\rho N\cdot\operatorname{covol}(\rho L/\rho N),\qquad\text{so}\qquad \deg(\rho L/\rho N)=\deg L-\deg N .
\end{equation}
This holds because a $\Z$-basis of $\rho N$ extends to one of $\rho L$, and the Gram determinant factors after Gram--Schmidt.

\begin{lemma}[Elementary properties of the degree]\label{lem:deg-basic}
Let $A,B,N,N'\subseteq L$ be submodules.
\begin{enumerate}
\item[(a)] If $N'\subseteq N$ and $\operatorname{rk}N'=\operatorname{rk}N$, then $\deg N'=\deg N-\log[N:N']\le\deg N$.
\item[(b)] \emph{(Supermodularity)} $\operatorname{rk}(A+B)+\operatorname{rk}(A\cap B)=\operatorname{rk}A+\operatorname{rk}B$ and $\deg(A+B)+\deg(A\cap B)\ge\deg A+\deg B$.
\item[(c)] \emph{(Finiteness)} For every $c\in\R$, only finitely many submodules $N\subseteq L$ satisfy $\deg N\ge c$.
\end{enumerate}
\end{lemma}
\begin{proof}
(a) A finite-index sublattice has covolume multiplied by the index.

(b) The rank identity is the dimension formula for the $K$-spans. For the degree, by (a) we may replace $A\cap B$ by its saturation $S$ in $A$, which only increases $\deg(A\cap B)$; so assume $A\cap B$ is saturated in $A$. Projecting orthogonally to $\R\rho(A\cap B)$, \eqref{eq:covol-mult} gives
\[
\deg A-\deg(A\cap B)=\deg\pi(\rho A),\qquad \deg(A+B)-\deg B\ \ge\ \deg\bigl(\pi_B\rho(A+B)\bigr),
\]
where $\pi$ and $\pi_B$ denote the two projections. The second inequality is again by (a) applied to the saturation of $B$ in $A+B$. Now $\pi_B$ maps $\rho A$ onto $\pi_B\rho(A+B)$. Its restriction to $\pi(\rho A)$ is $\pi_B$ restricted to the orthogonal complement of $\R\rho(A\cap B)$ inside $\R\rho A$, and since $\R\rho(A\cap B)\subseteq\R\rho B$ this restriction is injective and has norm $\le1$. A linear map of norm $\le 1$ does not increase covolume, so $\deg\pi_B\rho(A+B)\ge\deg\pi(\rho A)$, which gives the claim.

(c) Let $N$ have rank $m$ and $\deg N\ge c$. Let $\bar N=L\cap KN$ be its saturation; then $\deg\bar N\ge\deg N\ge c$ and $[\bar N:N]\le e^{\deg\bar N-c}$. Write $\omega(\bar N)=u_1\wedge\dots\wedge u_m\in\bigwedge^m_\Z\rho L$ for a $\Z$-basis $u_1,\dots,u_m$ of $\rho\bar N$. Then $\|\omega(\bar N)\|=\operatorname{covol}\bar N\le e^{-c}$, and $\omega(\bar N)$ determines $\R\rho\bar N$, hence $\bar N$. The lattice $\bigwedge^m\rho L$ is discrete, so there are finitely many possibilities for $\bar N$. For each, the degree of $\bar N$ is bounded above by a constant depending only on $L$, because $\rho L$ has a nonzero shortest vector. Hence the index $[\bar N:N]$ is bounded, and there are finitely many possibilities for $N$.
\end{proof}

\begin{definition}[Extreme slopes]\label{def:extreme}
For a nonzero submodule $N\subseteq L$ put
\[
\operatorname{sl}_{\max}(N)=\sup\{\operatorname{sl}N':0\neq N'\subseteq N\},\qquad \operatorname{sl}_{\min}(N)=\inf\{\operatorname{sl}(\rho N/\rho N'):N'\subsetneq N\text{ saturated}\}.
\]
The supremum is finite and attained, by Lemma~\ref{lem:deg-basic}(c) applied to each rank.
\end{definition}

\begin{lemma}[Characterisation of the extreme slopes]\label{lem:muchar}
Let $0\ne N\subseteq L$ and $\theta\in\R$.
\begin{enumerate}
\item[(a)] $\operatorname{sl}_{\max}(N)\le\theta$ if and only if $\deg N'\le\theta\operatorname{rk}N'$ for all $N'\subseteq N$.
\item[(b)] $\theta\le\operatorname{sl}_{\min}(N)$ if and only if $\theta(\operatorname{rk}N-\operatorname{rk}N')\le\deg N-\deg N'$ for all $N'\subseteq N$.
\item[(c)] $\operatorname{sl}_{\min}(N)\le\operatorname{sl}N\le\operatorname{sl}_{\max}(N)$.
\item[(d)] There is a saturated $N_1\subseteq N$ with $\operatorname{sl}N_1=\operatorname{sl}_{\max}(N)$ and $\operatorname{sl}_{\min}(N_1)\ge\operatorname{sl}_{\max}(N)$.
\end{enumerate}
\end{lemma}
\begin{proof}
(a) is the definition. For (b), by \eqref{eq:covol-mult} the condition for saturated $N'$ is exactly $\operatorname{sl}(\rho N/\rho N')\ge\theta$. For a general $N'$, its saturation $\bar N'$ has the same rank and $\deg\bar N'\ge\deg N'$ by Lemma~\ref{lem:deg-basic}(a), so the saturated case implies the general one. (c) follows from (a) and (b) with $N'=N$, respectively $N'=0$.

(d) Among the submodules attaining $\operatorname{sl}_{\max}(N)$, choose $N_1$ of maximal rank; by Lemma~\ref{lem:deg-basic}(a) it is saturated. Let $N'\subsetneq N_1$ be saturated. Then $\deg N'\le\operatorname{sl}(N_1)\operatorname{rk}N'$, so
\[
\deg N_1-\deg N'\ge\operatorname{sl}(N_1)(\operatorname{rk}N_1-\operatorname{rk}N'),
\]
that is, $\operatorname{sl}(\rho N_1/\rho N')\ge\operatorname{sl}(N_1)$.
\end{proof}

\begin{proposition}[Harder--Narasimhan filtration]\label{prop:HN}
Let $L$ be a nonzero finitely generated $\mathcal O_K$-submodule of a metrised space, of $\Z$-rank $r$. For $\theta\in\R$, let $L_{\ge\theta}$ be the sum of all nonzero submodules $N\subseteq L$ with $\operatorname{sl}_{\min}(N)\ge\theta$ (and $0$ if there are none).
\begin{enumerate}
\item[(a)] $L_{\ge\theta}$ is saturated in $L$. If $L_{\ge\theta}\neq0$, then $\operatorname{sl}_{\min}(L_{\ge\theta})\ge\theta$. Every nonzero $N\subseteq L$ with $\operatorname{sl}_{\min}(N)\ge\theta$ satisfies $N\subseteq L_{\ge\theta}$.
\item[(b)] If $L_{\ge\theta}\subsetneq N\subseteq L$, then $\deg N-\deg L_{\ge\theta}<\theta(\operatorname{rk}N-\operatorname{rk}L_{\ge\theta})$.
\item[(c)] Put $\lambda_j=\sup\{\theta:j<\operatorname{rk}L_{\ge\theta}\}$ for $0\le j<r$. Then $\lambda_0\ge\dots\ge\lambda_{r-1}$, and
\[
\#\{j:\lambda_j\ge\theta\}=\operatorname{rk}L_{\ge\theta}\ \ \text{for all }\theta,\qquad \sum_j\lambda_j=\deg L,\qquad \lambda_0=\operatorname{sl}_{\max}(L),\qquad \lambda_{r-1}=\operatorname{sl}_{\min}(L).
\]
\end{enumerate}
In particular $\#\{j:\lambda_j<\theta\}=r-\operatorname{rk}L_{\ge\theta}$. The $\lambda_j$ are the \emph{Harder--Narasimhan slopes} of $L$.
\end{proposition}

\begin{proof}
\emph{Step 1: sums.} Let $A,B\subseteq L$ be nonzero with $\operatorname{sl}_{\min}\ge\theta$, and let $N'\subseteq A+B$. Apply Lemma~\ref{lem:deg-basic}(b) to the pair $(N'+A,B)$ and to the pair $(N',A)$, and then Lemma~\ref{lem:muchar}(b) to $B$ and to $A$:
\begin{align*}
\deg(A+B)-\deg(N'+A)&\ge\deg B-\deg((N'+A)\cap B)\ge\theta\bigl(\operatorname{rk}(A+B)-\operatorname{rk}(N'+A)\bigr),\\
\deg(N'+A)-\deg N'&\ge\deg A-\deg(N'\cap A)\ge\theta\bigl(\operatorname{rk}(N'+A)-\operatorname{rk}N'\bigr).
\end{align*}
The rank identities used here come from Lemma~\ref{lem:deg-basic}(b). Adding the two lines and applying Lemma~\ref{lem:muchar}(b) gives $\operatorname{sl}_{\min}(A+B)\ge\theta$.

\emph{Step 2: proof of (a).} $L$ is a Noetherian module, so $L_{\ge\theta}$ is a finite sum of submodules with $\operatorname{sl}_{\min}\ge\theta$. By Step~1, $\operatorname{sl}_{\min}(L_{\ge\theta})\ge\theta$ when $L_{\ge\theta}\neq0$. Maximality holds by construction. Let $S$ be the saturation of $L_{\ge\theta}$ and let $N'\subseteq S$. By Lemma~\ref{lem:deg-basic}(a),(b) and Lemma~\ref{lem:muchar}(b),
\[
\deg S-\deg N'\ge\deg(N'+L_{\ge\theta})-\deg N'\ge\deg L_{\ge\theta}-\deg(N'\cap L_{\ge\theta})\ge\theta\bigl(\operatorname{rk}L_{\ge\theta}-\operatorname{rk}(N'\cap L_{\ge\theta})\bigr).
\]
The right-hand side equals $\theta(\operatorname{rk}S-\operatorname{rk}N')$, because $\operatorname{rk}(N'+L_{\ge\theta})=\operatorname{rk}S$. Hence $\operatorname{sl}_{\min}(S)\ge\theta$, so $S\subseteq L_{\ge\theta}$, and $L_{\ge\theta}$ is saturated.

\emph{Step 3: proof of (b).} Write $F=L_{\ge\theta}$ and let $\bar Q=\pi_F(\rho L)$ be the quotient lattice. Submodules $F\subseteq N\subseteq L$ have images $\pi_F\rho N\subseteq\bar Q$, and $\deg N-\deg F=\deg\pi_F\rho N$ by \eqref{eq:covol-mult}, since $F$ is saturated in $N$. Suppose the claim fails for some $N$, so that $\bar Q$ has a nonzero subgroup of slope $\ge\theta$. The quotient lattice satisfies Lemmas~\ref{lem:deg-basic}--\ref{lem:muchar} verbatim, with subgroups of the form $\pi_F\rho N$ in place of submodules. So by Lemma~\ref{lem:muchar}(d) there is $F\subsetneq N_1\subseteq L$ with $\operatorname{sl}_{\min}(\pi_F\rho N_1)\ge\theta$. For $N'\subseteq N_1$, combining this with Lemma~\ref{lem:deg-basic}(b) for the pair $(N',F)$:
\begin{align*}
\deg N_1-\deg(N'+F)&=\deg\pi_F\rho N_1-\deg\pi_F\rho N'\ \ge\ \theta(\operatorname{rk}N_1-\operatorname{rk}(N'+F)),\\
\deg(N'+F)-\deg N'&\ge\deg F-\deg(N'\cap F)\ \ge\ \theta(\operatorname{rk}(N'+F)-\operatorname{rk}N').
\end{align*}
Adding gives $\operatorname{sl}_{\min}(N_1)\ge\theta$, so $N_1\subseteq F$ by (a), a contradiction.

\emph{Step 4: proof of (c).} The map $\theta\mapsto L_{\ge\theta}$ is decreasing. By Lemma~\ref{lem:muchar}(b),(c) and (a), it equals $L$ for $\theta\le\operatorname{sl}_{\min}(L)$ and $0$ for $\theta>\operatorname{sl}_{\max}(L)$.

We show that $\{\theta:\operatorname{rk}L_{\ge\theta}>j\}$ is closed. If $\theta_n\uparrow\lambda$, the saturated submodules $L_{\ge\theta_n}$ form a decreasing chain, which stabilises at some $N$. Then $\operatorname{sl}_{\min}(N)\ge\theta_n$ for all $n$, so $\operatorname{sl}_{\min}(N)\ge\lambda$ and $N\subseteq L_{\ge\lambda}$. This gives $\#\{j:\lambda_j\ge\theta\}=\operatorname{rk}L_{\ge\theta}$, and monotonicity of the $\lambda_j$ is clear.

Let $\theta_1>\dots>\theta_k$ be the values where the rank jumps, and put $F_i=L_{\ge\theta_i}$ and $F_0=0$. For $\operatorname{rk}F_{i-1}\le j<\operatorname{rk}F_i$ we have $\lambda_j=\theta_i$. By (a) and Lemma~\ref{lem:muchar}(b),
\[
\deg F_i-\deg F_{i-1}\ge\theta_i(\operatorname{rk}F_i-\operatorname{rk}F_{i-1}).
\]
For $\theta\in(\theta_i,\theta_{i-1})$ we have $L_{\ge\theta}=F_{i-1}$, so (b) gives $\deg F_i-\deg F_{i-1}<\theta(\operatorname{rk}F_i-\operatorname{rk}F_{i-1})$; letting $\theta\downarrow\theta_i$ gives the reverse inequality. Summing over $i$ gives $\sum_j\lambda_j=\deg L$.

Finally, $\lambda_{r-1}=\sup\{\theta:L_{\ge\theta}=L\}=\operatorname{sl}_{\min}(L)$. Also $\lambda_0=\sup\{\theta:L_{\ge\theta}\ne0\}$. If $L_{\ge\theta}\neq0$ then $\operatorname{sl}_{\max}(L)\ge\operatorname{sl}(L_{\ge\theta})\ge\theta$ by Lemma~\ref{lem:muchar}(c). Conversely, $L_{\ge\operatorname{sl}_{\max}(L)}\supseteq N_1\neq0$ by Lemma~\ref{lem:muchar}(d). Hence $\lambda_0=\operatorname{sl}_{\max}(L)$.
\end{proof}

This is the Harder--Narasimhan formalism for Hermitian lattices of
Stuhler \cite{Stuhler1976}, Grayson \cite{Grayson1984} and Bost \cite{Bost1996}; we give it here with proofs for
completeness.

\begin{definition}[Normalised slopes]\label{def:normslope}
Put $c_K=-\log2+\log|d_K|/n_K$. For a nonzero finitely generated $\mathcal O_K$-submodule $L$ of a metrised space, define
\[
\mu(L)=2\operatorname{sl}(L)+c_K,\qquad\mu_{\max}(L)=2\operatorname{sl}_{\max}(L)+c_K,\qquad\mu_{\min}(L)=2\operatorname{sl}_{\min}(L)+c_K .
\]
\end{definition}

\begin{remark}\label{rem:normalisation}
With this normalisation $\mu(L)=-2\log\operatorname{size}(L)$, where $\operatorname{size}(L)$ is the $(n_K\operatorname{rk}_KL)$-th root of the Arakelov height $2^{n_K\operatorname{rk}_KL/2}|d_K|^{-\operatorname{rk}_KL/2}\operatorname{covol}(\rho L)$. Multiplying all the matrices $X_v$ by scalars $c_v>0$ shifts $\mu$, $\mu_{\max}$ and $\mu_{\min}$ by the same amount, $\frac{2}{\#\Sigma_K}\sum_v\log c_v$. In the proof below the constant $c_K$ cancels, so these facts are not used; they explain the choice of normalisation.
\end{remark}

\begin{definition}[Height of a morphism]\label{def:htmap}
Let $(V,b,X)$ and $(W,b',X')$ be metrised spaces and let $\varphi\colon V\to W$ be $K$-linear. Let $\|\varphi\|_v$ be the operator norm of the matrix $X'^{-1}_v\sigma_v(\mathrm{Mat}_{b,b'}\varphi)X_v$, and put
\[
h(\varphi)=\frac{1}{\#\Sigma_K}\sum_{v\in\Sigma_K}\log\|\varphi\|_v^2 .
\]
\end{definition}

\begin{proposition}[Slope inequality]\label{prop:slopeineq}
In the setting of Definition~\ref{def:htmap}, let $A\subseteq V$ be a finitely generated $\mathcal O_K$-submodule with $\varphi(A)\ne0$. Then $\mu_{\min}(\varphi A)\ge\mu_{\min}(A)-h(\varphi)$.
\end{proposition}
\begin{proof}
Let $N\subsetneq\varphi A$ be saturated and put $A'=A\cap\varphi^{-1}(KN)$. Then $A'$ is saturated in $A$, and $\varphi$ induces a surjection $\bar\varphi\colon\rho A/\rho A'\to\rho\varphi A/\rho N$ between lattices of the same $\Z$-rank $m=n_Kr'$, where $r'=\operatorname{rk}_K\varphi A-\operatorname{rk}_KN$.

At the place $v$, the induced real-linear map is the restriction of $\varphi$ to the orthogonal complement of $\R\rho A'$, followed by an orthogonal projection, so its norm is at most $\|\varphi\|_v$. The $r'$-th complex exterior power of this map has norm at most $\|\varphi\|_v^{r'}$. As a real map on $\bigwedge^{2r'}$, the determinant of the induced map between the real spans is the squared modulus of the complex determinant, so it is at most $\|\varphi\|_v^{2r'}$.

Multiplying over $v$, and using that the image of a lattice under a linear map that is bijective on the real spans has covolume at most $|\det|$ times the covolume (with equality onto the image lattice, and $\bar\varphi$ surjective), we get
\[
\operatorname{covol}(\rho\varphi A/\rho N)\le\Bigl(\prod_v\|\varphi\|_v^{2r'}\Bigr)\operatorname{covol}(\rho A/\rho A').
\]
Take $-\log$ and divide by $m=n_Kr'=2\#\Sigma_Kr'$:
\[
\operatorname{sl}(\rho\varphi A/\rho N)\ge\operatorname{sl}(\rho A/\rho A')-\tfrac12h(\varphi)\ge\operatorname{sl}_{\min}(A)-\tfrac12h(\varphi).
\]
Take the infimum over $N$ and multiply by $2$.
\end{proof}

\subsection{Almost tight frames at one place}

For finite-dimensional complex inner-product spaces and a linear map $A$, put $\|A\|_{\mathrm{HS}}^2=\operatorname{tr}(A^*A)$. For linear maps $B_i\colon V_i\to H$ ($i\in I$, a finite set) and weights $c_i>0$, put
\[
K_i=c_iB_i^*B_i\in\operatorname{End}(V_i),\qquad Q=\sum_ic_iB_iB_i^*\in\operatorname{End}(H).
\]

\begin{lemma}\label{lem:fambound}
Let $B\colon V\to H$, $c>0$, and let $f_1,\dots,f_k\in V$ be orthonormal. Then
\[
\Bigl|\,c\sum_{m=1}^k\|Bf_m\|^2-k\,\Bigr|\le\sqrt k\,\|cB^*B-1\|_{\mathrm{HS}} .
\]
\end{lemma}
\begin{proof}
Put $A=cB^*B-1$. Then $c\sum_m\|Bf_m\|^2-k=\sum_m\langle f_m,Af_m\rangle$, and
\[
\Bigl|\sum_m\langle f_m,Af_m\rangle\Bigr|\le\sum_m\|Af_m\|\le\sqrt k\Bigl(\sum_m\|Af_m\|^2\Bigr)^{1/2}\le\sqrt k\,\|A\|_{\mathrm{HS}} .\qedhere
\]
\end{proof}

\begin{lemma}\label{lem:Wupper}
Let $B\colon V\to H$, let $N\subseteq V$ be a subspace, and let $g_1,\dots,g_w\in H$ be orthonormal with $\langle Bn,g_l\rangle=0$ for all $n\in N$ and all $l$. If $f_1,\dots,f_k$ is an orthonormal basis of $N^\perp$, then $\sum_l\|B^*g_l\|^2\le\sum_m\|Bf_m\|^2$.
\end{lemma}
\begin{proof}
The hypothesis says $B^*g_l\in N^\perp$. Let $P$ be the orthogonal projection onto $N^\perp$ and $R$ the orthogonal projection onto $\operatorname{span}(g_l)$. Then
\[
\sum_l\|B^*g_l\|^2=\|PB^*R\|_{\mathrm{HS}}^2\le\|PB^*\|_{\mathrm{HS}}^2=\|BP\|_{\mathrm{HS}}^2=\sum_m\|Bf_m\|^2 .\qedhere
\]
\end{proof}

\begin{lemma}\label{lem:tracebound}
Let $A\in\operatorname{End}(H)$ and let $g_1,\dots,g_w$ be orthonormal. Then $\bigl|\sum_l\Rel\langle g_l,Ag_l\rangle\bigr|\le\sqrt w\,\|A\|_{\mathrm{HS}}$.
\end{lemma}
\begin{proof}
Same as Lemma~\ref{lem:fambound}.
\end{proof}

\begin{proposition}[Frame linear algebra]\label{prop:frameLA}
Let $(V_i)_{i\in I}$ and $H$ be finite-dimensional complex inner-product spaces, let $B_i\colon V_i\to H$ be linear, and let $c_i>0$. Assume $R:=\sum_i\dim V_i>0$. Let $N_i\subseteq V_i$ be subspaces, and let $0\le e\le\frac14$ and $\eta\ge0$. Suppose:
\begin{align}
&\textstyle\sum_i\|K_i-1\|_{\mathrm{HS}}^2\le e\,R && \text{(almost tight)},\label{eq:LA-tight}\\
&\operatorname{tr}(Q^2)\le(1+\eta)(\operatorname{tr}Q)^2/\dim H && \text{(almost scalar)},\label{eq:LA-scalar}\\
&\textstyle\sum_i(\dim V_i-\dim N_i)\le\varepsilon R && \text{(small corank)}.\label{eq:LA-corank}
\end{align}
Then $\operatorname{codim}_H\bigl(\sum_iB_i(N_i)\bigr)\le(\eta+4\varepsilon+4\sqrt e)\dim H$.
\end{proposition}

\begin{proof}
Put $D=\dim H$, $t=\operatorname{tr}Q$, $\mathcal W=(\sum_iB_iN_i)^\perp$ and $w=\dim\mathcal W$. We may assume $D>0$. Note that \eqref{eq:LA-corank} forces $\varepsilon\ge0$. Write $h_i=\|K_i-1\|_{\mathrm{HS}}^2$ and $d_i=\dim V_i$.

\emph{Step 1.} Lemma~\ref{lem:fambound} on an orthonormal basis of $V_i$ gives $c_i\operatorname{tr}(B_i^*B_i)\ge d_i-\sqrt{d_ih_i}$. Summing over $i$ and using Cauchy--Schwarz together with \eqref{eq:LA-tight}:
\[
t=\sum_ic_i\operatorname{tr}(B_iB_i^*)\ge R-\Bigl(\sum_id_i\Bigr)^{1/2}\Bigl(\sum_ih_i\Bigr)^{1/2}\ge R(1-\sqrt e).
\]
Since $\sqrt e\le\frac12$, this gives $t>0$ and $R\le2t$.

\emph{Step 2.} Let $g_1,\dots,g_w$ be an orthonormal basis of $\mathcal W$ and put
\[
t_{\mathcal W}=\sum_l\Rel\langle g_l,Qg_l\rangle=\sum_ic_i\sum_l\|B_i^*g_l\|^2 .
\]
Each $g_l$ is orthogonal to $B_iN_i$. Lemma~\ref{lem:Wupper}, followed by Lemma~\ref{lem:fambound} on an orthonormal basis of $N_i^\perp$ (of size $k_i=d_i-\dim N_i$), gives $c_i\sum_l\|B_i^*g_l\|^2\le k_i+\sqrt{k_ih_i}$. Summing and using Cauchy--Schwarz with \eqref{eq:LA-tight} and \eqref{eq:LA-corank}:
\[
t_{\mathcal W}\le\varepsilon R+R\sqrt{\varepsilon e}.
\]

\emph{Step 3.} Let $s=t/D$ and $A=Q-s$. Then
\[
\|A\|_{\mathrm{HS}}^2=\operatorname{tr}Q^2-2st+Ds^2=\operatorname{tr}Q^2-t^2/D\le\eta t^2/D=:X_0
\]
by \eqref{eq:LA-scalar}. Also $t_{\mathcal W}=sw+\sum_l\Rel\langle g_l,Ag_l\rangle$, so Lemma~\ref{lem:tracebound} gives $\frac tDw-\sqrt w\sqrt{X_0}\le t_{\mathcal W}$.

\emph{Step 4.} If $\varepsilon\ge\frac14$, then $w\le D\le4\varepsilon D$ and we are done. Otherwise put $q=\sqrt{\varepsilon e}\le\sqrt e$ and $y=(D/t)\sqrt{X_0}$, so that $y^2\le\eta D$. Combining Steps 2 and 3, multiplying by $D/t$, and using $R\le2t$:
\[
w-\sqrt w\,y\le\frac{DR}{t}(\varepsilon+q)\le2D(\varepsilon+q).
\]
Since $\sqrt w\,y\le(w+y^2)/2$, this gives $w/2\le2D(\varepsilon+q)+\eta D/2$. Hence $w\le(\eta+4\varepsilon+4\sqrt e)D$.
\end{proof}

\subsection{The frame theorem}\label{subsec:frame-thm}

\begin{lemma}[Corank sharing]\label{lem:F-share}
Let $(x_{ij})_{i\in I,\,1\le j\le m_i}$ be real numbers, $s\in\R$ and $u,\tau\ge0$. Suppose $\bigl(\sum_im_i\bigr)s\le\sum_{i,j}x_{ij}$ and $x_{ij}\le s+u$ for all $i,j$. Then
\[
\#\{(i,j):x_{ij}<s-\tau\}\le\Bigl(\sum_im_i\Bigr)\frac{u}{u+\tau},
\]
where the right-hand side is read as $0$ if $u=\tau=0$.
\end{lemma}
\begin{proof}
Let $k$ be the number of pairs with $x_{ij}<s-\tau$ and let $m=\sum_im_i$. Then $ms\le\sum x_{ij}\le k(s-\tau)+(m-k)(s+u)$, so $k(u+\tau)\le mu$. If $u=0$, then every $x_{ij}\le s$ while the mean is $\ge s$, so all $x_{ij}=s$ and $k=0$.
\end{proof}

\begin{lemma}[Layer cake]\label{lem:layercake}
Let $x_1,\dots,x_D$ be real numbers with $D>0$, let $a\in\R$ and $\mathsf H\ge0$, and let $w$ be integrable on $[0,\mathsf H]$. Suppose $x_j\ge a-\mathsf H$ for all $j$, and $\#\{j:x_j<a-\tau\}\le w(\tau)D$ for all $\tau\in[0,\mathsf H]$. Then
\[
\frac1D\sum_jx_j\ge a-\int_0^{\mathsf H}\min(1,w(\tau))\,d\tau .
\]
\end{lemma}
\begin{proof}
Since $x_j\ge a-\mathsf H$, we have $a-x_j\le\int_0^{\mathsf H}\mathbf 1[x_j<a-\tau]\,d\tau$ (both sides vanish when $x_j\ge a$). Average over $j$; the averaged indicator is at most $\min(1,w(\tau))$.
\end{proof}

\smallskip
\noindent\emph{Setting.} The setting of the frame theorem (the ``column'' setting used in \S\ref{subsec:assembly-T5}) consists of:
\begin{itemize}
\item a metrised space $(W,b',X')$ of dimension $\kappa$, with a finitely generated $\mathcal O_K$-submodule $\mathrm{CC}\subseteq W$ of full rank, the \emph{column};
\item a nonempty finite family of metrised spaces $(V_i,b_i,X_i)_{i\in I}$, the \emph{words}, with $r_i=\dim_KV_i$ and $\sum_ir_i>0$, and for each $i$ a finitely generated $\mathcal O_K$-submodule $N_i\subseteq V_i$ of full rank;
\item $K$-linear maps $B_i\colon V_i\to W$ with $B_i(N_i)\subseteq\mathrm{CC}$, and real numbers $h_i$ with $h(B_i)\le h_i$ whenever $B_i\ne0$.
\end{itemize}

\begin{theorem}[Frame slope inequality]\label{thm:frameF}
In this setting, let $s\in\R$, $u\ge0$, $\mathsf H\ge0$, $0\le e\le\frac14$, $\eta\ge0$ and $0\le\delta<1$. Assume:
\begin{enumerate}
\item[(F1)] $s\sum_ir_i\le\sum_ir_i\bigl(\mu(N_i)-h_i\bigr)$;
\item[(F2)] $\mu_{\max}(N_i)-h_i\le s+u$ for every $i$;
\item[(F3)] $s-\mu_{\min}(\mathrm{CC})\le\mathsf H$;
\item[(F4)] there exist
  \begin{itemize}
  \item a place $v\in\Sigma_K$;
  \item invertible matrices $Y_i\in\mathrm{GL}_{r_i}(\C)$ and $Y'\in\mathrm{GL}_\kappa(\C)$;
  \item a finite set $\kappa_1$ with $(1-\delta)(\kappa+\#\kappa_1)\le\kappa$;
  \item weights $\gamma_i>0$;
  \item matrices $\mathsf B_i\in M_{(\kappa\sqcup\kappa_1)\times r_i}(\C)$ whose first $\kappa$ rows equal $Y'^{-1}\sigma_v(\mathrm{Mat}\,B_i)Y_i$,
  \end{itemize}
  such that the maps $\mathsf B_i$ satisfy \eqref{eq:LA-tight} with constant $e$ and \eqref{eq:LA-scalar} with constant $\eta$.
\end{enumerate}
Then
\[
\mu(\mathrm{CC})\ \ge\ s-\frac{1}{1-\delta}\Bigl(4u\log\Bigl(1+\frac{\mathsf H}{u}\Bigr)+(\eta+4\sqrt e)\,\mathsf H\Bigr),
\]
where $u\log(1+\mathsf H/u)$ is read as $0$ if $u=0$.
\end{theorem}

\begin{proof}
\emph{Step 0.} If $\kappa=0$, then (F4) forces $\kappa_1=\emptyset$. All $\mathsf B_i$ then vanish, and the left side of \eqref{eq:LA-tight} equals $\sum_ir_i>0$, contradicting $e\le\frac14$. Hence $\kappa>0$. By Lemma~\ref{lem:realisation}, $\operatorname{rk}N_i=n_Kr_i$ and $\operatorname{rk}\mathrm{CC}=n_K\kappa$.

\emph{Step 1.} Fix $\tau\ge0$. Let $\lambda_{i,j}$ ($0\le j<n_Kr_i$) be the Harder--Narasimhan slopes of $N_i$, and put $x_{ij}=2\lambda_{ij}+c_K-h_i$. By Proposition~\ref{prop:HN}(c), the mean of the $x_{ij}$ over $j$ is $\mu(N_i)-h_i$, and their maximum is $\mu_{\max}(N_i)-h_i$. So (F1) and (F2) allow us to apply Lemma~\ref{lem:F-share} with $m_i=n_Kr_i$.

Let $\theta_i=\frac12(s-\tau+h_i-c_K)$ and $N_i(\tau)=(N_i)_{\ge\theta_i}$. The condition $x_{ij}<s-\tau$ is equivalent to $\lambda_{ij}<\theta_i$. Proposition~\ref{prop:HN}(c) and $\operatorname{rk}=n_K\operatorname{rk}_K$ then give
\begin{equation}\label{eq:F-corank}
\sum_i\bigl(r_i-\operatorname{rk}_KN_i(\tau)\bigr)\le\frac{u}{u+\tau}\sum_ir_i .
\end{equation}

\emph{Step 2.} Put
\[
M(\tau)=\sum_iB_i(N_i(\tau)),\qquad \theta'=\tfrac12(s-\tau-c_K),\qquad \mathrm{CC}(\tau)=\mathrm{CC}_{\ge\theta'}.
\]
Suppose $B_iN_i(\tau)\neq0$. Then $\mu_{\min}(N_i(\tau))\ge2\theta_i+c_K$ by Proposition~\ref{prop:HN}(a), and Proposition~\ref{prop:slopeineq} gives
\[
\mu_{\min}(B_iN_i(\tau))\ge2\theta_i+c_K-h_i=2\theta'+c_K .
\]
Since $B_iN_i(\tau)\subseteq\mathrm{CC}$, Proposition~\ref{prop:HN}(a) gives $B_iN_i(\tau)\subseteq\mathrm{CC}(\tau)$. Hence $M(\tau)\subseteq\mathrm{CC}(\tau)$, and $\operatorname{rk}_KM(\tau)\le\operatorname{rk}_K\mathrm{CC}(\tau)$.

\emph{Step 3.} We use the following subspaces at the place $v$:
\begin{itemize}
\item $\mathcal N_i\subseteq\C^{r_i}$, the span of the $v$-coordinates of $N_i(\tau)$ computed with $Y_i$ in place of $X_{i,v}$;
\item $\mathcal T\subseteq\C^\kappa$, the span of the $v$-coordinates of $M(\tau)$ computed with $Y'$;
\item $\mathcal S=\sum_i\mathsf B_i\mathcal N_i$.
\end{itemize}
By Lemma~\ref{lem:realisation}(b), $\dim\mathcal N_i=\operatorname{rk}_KN_i(\tau)$ and $\dim\mathcal T=\operatorname{rk}_KM(\tau)$. Let $P\colon\C^{\kappa\sqcup\kappa_1}\to\C^\kappa$ be the coordinate projection, which is surjective with kernel of dimension $\#\kappa_1$. By (F4), $P\mathsf B_i=Y'^{-1}\sigma_v(B_i)Y_i$, so $P(\mathcal S)\subseteq\mathcal T$. Hence
\[
\dim\mathcal S\le\operatorname{rk}_KM(\tau)+\#\kappa_1 .
\]
By \eqref{eq:F-corank}, Proposition~\ref{prop:frameLA} applies with $\varepsilon=u/(u+\tau)$, giving
\[
\kappa+\#\kappa_1-\dim\mathcal S\le\Bigl(\eta+\frac{4u}{u+\tau}+4\sqrt e\Bigr)(\kappa+\#\kappa_1).
\]
Since $\kappa+\#\kappa_1\le(1-\delta)^{-1}\kappa$, this yields
\begin{equation}\label{eq:F-place}
\kappa-\operatorname{rk}_KM(\tau)\le w(\tau)\kappa,\qquad w(\tau):=\frac1{1-\delta}\Bigl(\eta+\frac{4u}{u+\tau}+4\sqrt e\Bigr).
\end{equation}

\emph{Step 4.} Let $\lambda'_j$ ($0\le j<n_K\kappa$) be the Harder--Narasimhan slopes of $\mathrm{CC}$, and put $x'_j=2\lambda'_j+c_K$. By Proposition~\ref{prop:HN}(c) and (F3), $x'_j\ge\mu_{\min}(\mathrm{CC})\ge s-\mathsf H$. For $\tau\in[0,\mathsf H]$, Proposition~\ref{prop:HN}(c), Step 2 and \eqref{eq:F-place} give
\[
\#\{j:x'_j<s-\tau\}=n_K\bigl(\kappa-\operatorname{rk}_K\mathrm{CC}(\tau)\bigr)\le n_K\bigl(\kappa-\operatorname{rk}_KM(\tau)\bigr)\le w(\tau)\,n_K\kappa .
\]

\emph{Step 5.} Lemma~\ref{lem:layercake} with $a=s$ and $D=n_K\kappa$, together with $\frac1D\sum_jx'_j=\mu(\mathrm{CC})$ (Proposition~\ref{prop:HN}(c)) and $\min(1,w)\le w$, gives
\[
\mu(\mathrm{CC})\ge s-\int_0^{\mathsf H}w(\tau)\,d\tau .
\]
For $u>0$, $\int_0^{\mathsf H}\frac{4u}{u+\tau}\,d\tau=4u\log(1+\mathsf H/u)$. For $u=0$, $w$ is constant.
\end{proof}

\begin{corollary}\label{cor:frameF-log}
Under the hypotheses of Theorem~\ref{thm:frameF},
\[
\mu(\mathrm{CC})\ge s-\log\#I-\frac{1}{1-\delta}\Bigl(4u\log\Bigl(1+\frac{\mathsf H}{u}\Bigr)+(\eta+4\sqrt e)\mathsf H\Bigr).
\]
\end{corollary}
\begin{proof}
$\log\#I\ge0$.
\end{proof}

\begin{remark}\label{rem:F-place-free}
The archimedean metrics enter Theorem~\ref{thm:frameF} only through the slopes and the heights $h(B_i)$. Hypothesis (F4) is a statement at a single place, in arbitrary auxiliary coordinates. Its role is purely dimensional: it forces the images of the large Harder--Narasimhan pieces of the words to fill almost all of the column. The constant $c_K$ cancels between the words and the column, because the thresholds $\theta_i$ and $\theta'$ differ exactly by $h_i/2$.
\end{remark}

\subsection{Column packages}\label{subsec:assembly-dict}
Throughout, $x\ge1$, $0<\varepsilon\le1$ and $u=x^\varepsilon$; $\gamma_E$ is Euler's constant and $G$ is the prime sum defined after \eqref{eq:dirichlet-input}.

Let $A$ be a finite abelian group of
order $n$, $E_0=(\Z/2)^2$, and $U=E_0\times A$. The CM type
$\Phi\subset U$ is the set of elements with first coordinate~$0$; $|\Phi|=2n$. Arithmetic data $D$ consist of primes
$q,d$, an inertia subgroup $I_p\subset U$ and a Frobenius element
$\mathrm{Fr}_p\in U$ for each prime~$p$.
For the hosts of Theorem~\ref{thm:alpoge-families}(b) they are defined there
through the Artin map, and Lemma~\ref{lem:host-facts} records the properties
used. For a nontrivial character $\xi$ of $U$,
$\ell(\xi)=L'(1,\chi_\xi)/L(1,\chi_\xi)$, where $\chi_\xi$ is the primitive
Dirichlet character with $\chi_\xi(p)=\xi(\mathrm{Fr}_p)$ at every prime $p$
with $I_p=0$. A prime $p$ with $p\nmid2qd$ and $\mathrm{Fr}_p\notin\Phi$ is
called \emph{inert}.
For $a>0$, $t\in A$, a weight $w:U\to\R_{>0}$ and a cutoff $N$, the
\emph{arithmetic count} $P(a,t,w)$ and the \emph{principal term} $M(a,w)$
are
\[
\begin{aligned}
 P(a,t,w)&=\sum_{\substack{p\le N,\ p\nmid2qd\\ \mathrm{Fr}_p\notin\Phi}}2\log p
   \sum_{\substack{1\le i\le N\\ i\mathrm{Fr}_p\notin\Phi}}
   \Bigl(\frac a{p^i}-w\bigl(\pi_t(i\mathrm{Fr}_p)\bigr)\Bigr)_+
   +\log d\sum_{i=1}^{N}\sum_{\substack{g\in i\mathrm{Fr}_d+I_d\\ g\notin\Phi}}
   \Bigl(\frac a{d^i}-w(\pi_t g)\Bigr)_+ ,\\
 M(a,w)&=a\Bigl(\log a-1-\gamma_E+B_{\chi_d}-\frac1{2n}\sum_{\rho\in\Phi}\log w(\rho)\Bigr),
\end{aligned}
\]
where $\pi_t(g)=c+g+(0,t)$ is the partner map ($c$ = complex conjugation)
and $B_{\chi_d}=\Rel\ell(\chi_d)$ is the logarithmic derivative attached to the
sign character.
These are $P(a,w)$ and $M(a,w)$ above (Lemma~\ref{lem:character-identity} and the paragraph preceding it) in
the coordinates used there; the identification of the summands at $p\nmid2qd$ and the
$d$-summands with the coset formula is the prime-power calculation of
the preceding section.

For a top set $T\subset A$ of size $r$, numerators $(a_t)_{t\in T}$ and
weights $w_{v,j}$ ($v\in A$ a translation, $1\le j\le m$),
\[
 P_*(v)=\sum_{t\in T}\sum_{j=1}^m P(a_t,t,w_{v,j}),\qquad
 M_*(v)=\sum_{t\in T}\sum_{j=1}^m M(a_t,w_{v,j}).
\]
The size bracket is
\begin{equation}\label{eq:assembly-Z}
 Z=ra\Bigl(1+\frac{m^2}k\Bigr)+m\Bigl(2nmbu+\frac{skbu^2}{n^2}\Bigr),
\end{equation}
with $s=|S|$ the size of the averaging set $S$ of
Theorem~\ref{thm:alpoge-families}(b), which defines the thick positions.
This is the bracket $Z$ of \eqref{eq:finite-comparison}.



\begin{definition}[column package]\label{def:package}
Fix finite sets $U_0$ (lifts) and $I_\omega\neq\emptyset$ ($\omega\in U_0$;
the columns, indexed by words), a finite top set of size $r$, integers
$m\ge0$, $f\ge0$. A \emph{record} consists of real numbers
$x,\varepsilon,a,b,b',k,L,L_a,B_0,N_c,\gamma,B,v_E,\log G_0,\log\Gamma,X''$,
numerators $a_t$, weights $w_{\omega,j,\rho}$ ($1\le j\le m$, $1\le\rho\le f$),
and, for each $\omega\in U_0$ and $i\in I_\omega$, numbers
$d_i,T_i,\mu(N_i),\mu_{\max}(N_i),h_i,\mu(M_i),\mu_{\max}(M_i),E_{i,j}$,
and per lift numbers
$\mu^\vee,\mu^\vee_{\min},F,g,\tau,P_*,I_O,H,\delta_0,\eta,e$, and
$\bar\delta,\bar\eta,\bar e$.
Constants are $c_{\rm div},e_D,c_E,c_H,c_{\rm count},c_{\rm slope},c_{\rm sl1},c_{\rm loss},e_L,c_{HH}$. Put $\Sigma_a=\sum_ta_t$,
$\log\Omega=\gamma+\log(\pi/v_E)-B$, $\phi(y)=y\log y-y$,
$S_\Omega=\sum_tm(\phi(a_t)-a_t\log\Omega)$, $E_i=\sum_jE_{i,j}$,
$[X]_\omega=\sum_id_iX_i/\sum_id_i$,
\[
 s_\omega=[\mu(N_i)-h_i]_\omega,\qquad
 \bar u_\omega=\max_{i\in I_\omega}(\mu_{\max}(N_i)-h_i)-s_\omega,\qquad
 \mathrm{LoP}=c_{\rm loss}\,r\,m\,b\,x^{e_L\varepsilon},
\]
\[
 \ell_i=c_{\rm sl1}I_O+c_{\rm sl1}rm\Bigl(\frac ak+b'\Bigr)\log T_i+c_{\rm sl1}X''+1,
 \qquad
 M_*(\omega)=\sum_{j=1}^m\sum_t a_t\Bigl(\log a_t-1-\gamma+B-\frac1f\sum_\rho\log w_{\omega,j,\rho}\Bigr).
\]
The record is a \emph{column package} if the following
eleven conditions hold for every $\omega\in U_0$ and $i,i'\in I_\omega$.
\begin{enumerate}
\item[(P0)] \emph{Positivity}: $x>0$, $b>0$,
$a\ge0$, $b'\ge0$, $k>0$, $v_E>0$, $f>0$, $a_t\ge0$, $w_{\omega,j,\rho}>0$,
$d_i>0$, $T_i>0$, $\mu(N_i)\le\mu_{\max}(N_i)$, $c_{\rm sl1}\ge0$,
$\bar\delta<1$, $\bar\eta\ge0$, $0\le\bar e\le1/4$.
\item[(P1)] \emph{Divisor bound}, see Input~C.
\item[(P2)] \emph{Exponent comparison}:
$|E_i-m\Sigma_a\log G_0|\le c_E\,m\,B_0$.
\item[(P3)] \emph{Height comparison}:
$|(h_i+\Sigma_a\log T_i)-(h_{i'}+\Sigma_a\log T_{i'})|\le c_H(raL_a+mB_0)$.
\item[(P4)] \emph{Counting},
see Input~D.
\item[(P5)] \emph{Reference slopes}:
$-c_{\rm slope}raL_a(1+m/k)\le\mu(M_i)$ and $\mu_{\max}(M_i)\le c_{\rm slope}raL_a$.
\item[(P6)] \emph{Source slopes}:
$\mu_{\max}(N_i)\le\mu_{\max}(M_i)+g+F-E_i+c_{\rm sl1}mB_0$ and
$\mu(M_i)+g+F-E_i-\ell_i\le\mu(N_i)$.
\item[(P7)] \emph{Losses},
see Input~B.
\item[(P8)] \emph{Frame slope inequality},
see Input~A.
\item[(P9)] \emph{Height of the frame}:
$s_\omega-\mu^\vee_{\min}\le H$ and $0<H\le x^{c_{HH}}b$.
\item[(P10)] \emph{Frame parameters}:
$0\le\delta_0\le\bar\delta$, $0\le\eta\le\bar\eta$, $0\le e\le\bar e$.
\end{enumerate}
The package satisfies the \emph{size windows with constants $(e_K,K)$} if
$1\le x$, $0<\varepsilon\le1$, $1\le b$, $xb\le a$, $1\le r$,
$m\le x^{1/2+\varepsilon/2}$, $xm\le k$, $0\le e_K$, $1\le k\le x^{e_K}$,
$k\le a$, $0\le L\le2u$, $0\le L_a\le2u$, $0\le b'$, $kb'\le a$,
$0\le X''\le a$, $0\le a_t\le2a$, $B_0,N_c\ge0$,
$\log|I_\omega|\le(k-3)\log2$, $\bar\delta\le1/2$, $\bar\eta,\bar e\ge0$,
$c_{\rm div},c_E,c_H,c_{\rm count},c_{\rm slope},c_{\rm sl1},c_{\rm loss}\in[0,K]$,
$e_D,e_L,c_{HH}\ge0$, and
$(\bar\eta+4\sqrt{\bar e})\,x^{c_{HH}}b\le K\,u\,ra$.
It satisfies the \emph{bracket bounds} if $B_0\le2nmbu$ and $N_c\le skbu^2/n^2$, and
$x$ is \emph{large} if
\begin{equation}\label{eq:assembly-largeX}
 \log x\le u,\qquad 200(1+K)^2(1+c_{HH})(1+e_K)\le u .
\end{equation}
Finally the package is \emph{read from the host} $(D,N,\ell,T,a,w)$ along $\operatorname{tr}:U_0\to A$
if the top index set is identified with $T$ (so $r=|T|$), $f=|\Phi|$ and a
bijection $\{1,\dots,f\}\simeq\Phi$ is fixed, $a_t$ and
$w_{\omega,j,\rho}=w_{\operatorname{tr}\omega,j}(\rho)$ (with $\rho\in\Phi$ through that
bijection) are the host values, $\gamma=\gamma_E$, $B=B_{\chi_d}$, and
$P_*(\omega)=\sum_{t\in T}\sum_jP(a_t,t,w_{\operatorname{tr}\omega,j})$.
\end{definition}

The four analytic-arithmetic inputs are the conditions (P1), (P4), (P7),
(P8) of a package. We record them separately because each rests on a
separate argument: Input~A is the frame slope inequality
(Theorem~\ref{thm:frameF}) applied to the columns, and Inputs~B--D are
proved in \cite{Annex}.

\begin{definition}[Input A: frame slope inequality]\label{inp:F}
For every $\omega$, with the package values $\delta=\delta_0(\omega)$,
$\eta=\eta(\omega)$, $e=e(\omega)$ and $H=H(\omega)$,
\[
 s_\omega-\log|I_\omega|-(1-\delta)^{-1}
 \Bigl(4\bar u_\omega\log\Bigl(1+\frac H{\bar u_\omega}\Bigr)+(\eta+4\sqrt e)H\Bigr)\le\mu^\vee ,
\]
where $\bar u_\omega\log(1+H/\bar u_\omega)$ is read as $0$ if $\bar u_\omega=0$.
These values are admissible for Theorem~\ref{thm:frameF}: by (P0) and (P10),
$0\le\delta<1$, $\eta\ge0$ and $0\le e\le1/4$, and by (P9), $H>0$ and
$s_\omega-\mu^\vee_{\min}\le H$.
\end{definition}

\begin{definition}[Input B: loss bound]\label{inp:loss}
For every $\omega$: $P_*(\omega)-g(\omega)\le c_{\rm loss}rmb\,x^{e_L\varepsilon}$
and $I_O(\omega)\le c_{\rm loss}rmb\,x^{e_L\varepsilon}$.
\end{definition}

\begin{definition}[Input C: divisor bound]\label{inp:div}
For every $\omega$ and $i\in I_\omega$,
\[
 \mu^\vee+h_i-F\le\sum_t\bigl(m(\phi(a_t)-a_t\log\Omega)-a_t\log T_i\bigr)
 +c_{\rm div}\,r\,a\,L_a+c_{\rm div}\,m(B_0+N_c)\,x^{e_D\varepsilon}.
\]
\end{definition}

\begin{definition}[Input D: counting]\label{inp:count}
For every $\omega$ and $i,i'\in I_\omega$, with
$\Lambda_\omega=\sum_{j=1}^m f^{-1}\sum_\rho\log(w_{\omega,j,\rho}v_E/\pi)+m\log G_0$,
\[
 \Lambda_\omega-c_{\rm count}\Bigl(L_a+\frac{m^2L_a}k\Bigr)\le\log T_i\le c_{\rm count}\,m\,L,
 \qquad |\log T_i-\log T_{i'}|\le c_{\rm count}\log k .
\]
\end{definition}




\subsection{The mean bound}\label{subsec:assembly-T5}

\begin{theorem}[mean bound]\label{thm:mean-bound}
Let a column package (Definition~\ref{def:package}) be read from a host along
$\operatorname{tr}:U_0\to A$, satisfy the size windows with constants
$(e_K,K)$ and the bracket bounds, and let $x$ be large in the sense
\eqref{eq:assembly-largeX}. Then for every lift $\omega\in U_0$, with
translation $v=\operatorname{tr}(\omega)$,
\[
 P_*(v)-M_*(v)\le u^{\,e_D+e_L+3}\,Z ,
\]
where $Z$ is \eqref{eq:assembly-Z}, $u=x^\varepsilon$, and $e_D$, $e_L$ are the
exponents of the package constants in the divisor bound (P1) and the loss bound (P7).
\end{theorem}

\begin{remark}\label{rem:mean-bound-paper}
Averaging over $v\in A$ (one lift per translation) gives, in the notation of
Proposition~\ref{prop:simultaneous-slope},
$\mathbb E(P_*-M_*)\le C_ju^{K_j}Z$ with $C_j=1$ and $K_j=e_D+e_L+3$
independent of $j$, so $\varepsilon_jK_j\to0$. No
multiplicative constant appears because all constants of the package are
absorbed by the second condition in \eqref{eq:assembly-largeX}, which only
enlarges the threshold in $x$.
\end{remark}

The proof occupies the rest of this subsection. Fix $\omega$ and drop it from
the notation; brackets $[\cdot]$ are rank-weighted means over $I_\omega$.

\begin{lemma}[identification]\label{lem:ident}
$M_*(\operatorname{tr}\omega)=\sum_{t\in T}\sum_j M(a_t,w_{\operatorname{tr}\omega,j})$
equals the package quantity $M_*(\omega)$ of Definition~\ref{def:package}, and
$P_*(\operatorname{tr}\omega)-M_*(\operatorname{tr}\omega)=P_*(\omega)-M_*(\omega)$.
\end{lemma}
\begin{proof}
Substitute the identifications of the package ($a_t$, $w$, $\gamma=\gamma_E$,
$B=\Rel\ell(\chi_d)$) into the definition of $M$, reindex the inner sum over
$\Phi$ by the bijection $\{1,\dots,f\}\simeq\Phi$ (so $f=|\Phi|$), and interchange
the finite sums over $t$ and $j$. The second assertion is the definition of
$P_*(\omega)$ in the package.
\end{proof}

\begin{lemma}[key identity]\label{lem:key-identity}
If $v_E>0$, $f>0$ and all weights are positive, then with
$\Lambda^{(5)}_\omega=\Lambda_\omega-c_{\rm count}(L_a+m^2L_a/k)$,
\[
 S_\Omega-\Sigma_a\Lambda^{(5)}_\omega+m\Sigma_a\log G_0=M_*(\omega)+\Sigma_a\,c_{\rm count}\Bigl(L_a+\frac{m^2L_a}k\Bigr).
\]
\end{lemma}
\begin{proof}
Expand $\log(w v_E/\pi)=\log w+\log v_E-\log\pi$ and
$\log\Omega=\gamma+\log\pi-\log v_E-B$; then
$S_\Omega=m\sum_ta_t(\log a_t-1-\gamma+B)-m(\log\pi-\log v_E)\Sigma_a$ and
$M_*=m\sum_ta_t(\log a_t-1-\gamma+B)-\Sigma_a\sum_jf^{-1}\sum_\rho\log w_{j\rho}$.
The identity is the linear combination of these two expansions.
\end{proof}

\begin{lemma}[lower bound for the source mean]\label{lem:s-lower}
Under (P0), (P1), (P2), (P4), (P5), (P6),
$g+\mu^\vee-K_5-[\ell_i]\le s$, where
$K_5=c_{\rm slope}raL_a(1+m/k)+m(\Sigma_a\log G_0+c_EB_0)+S_\Omega-\Sigma_a\Lambda^{(5)}+E_{\rm div}$ and
$E_{\rm div}=c_{\rm div}raL_a+c_{\rm div}m(B_0+N_c)x^{e_D\varepsilon}$.
\end{lemma}
\begin{proof}
For each column $i$: by (P6), $\mu(N_i)\ge\mu(M_i)+g+F-E_i-\ell_i$; by (P1),
$F-h_i\ge\mu^\vee-S_\Omega+\Sigma_a\log T_i-E_{\rm div}$; by (P4) and $\Sigma_a\ge0$,
$\Sigma_a\log T_i\ge \Sigma_a\Lambda^{(5)}$; by (P2), $E_i\le m(\Sigma_a\log G_0+c_EB_0)$; by (P5),
$\mu(M_i)\ge-c_{\rm slope}raL_a(1+m/k)$. Adding,
$\mu(N_i)-h_i\ge g+\mu^\vee-K_5-\ell_i$. Taking the rank-weighted mean (the
weights $d_i$ are positive by (P0)) gives the claim.
\end{proof}

\begin{lemma}[bound for the excess]\label{lem:uF}
Under (P0), (P2)--(P6), $\bar u\le K_3+[\ell_i]$, with
$K_3=c_{\rm slope}raL_a(2+m/k)+V+c_{\rm sl1}mB_0$ and
$V=2mc_EB_0+c_H(raL_a+mB_0)+\Sigma_a\,c_{\rm count}\log k$.
\end{lemma}
\begin{proof}
First, for $i,i'$: (P2) gives $E_{i'}-E_i\le2mc_EB_0$; (P3) gives
$h_{i'}-h_i\le c_H(raL_a+mB_0)+\Sigma_a(\log T_i-\log T_{i'})$, and (P4) bounds the last
term by $Ac_{\rm count}\log k$. Hence $(-E_i-h_i)-(-E_{i'}-h_{i'})\le V$. Next,
(P6) upper at $i$, (P6) lower at $i'$ and (P5) at both give
$\mu_{\max}(N_i)-h_i-(\mu(N_{i'})-h_{i'})\le K_3+\ell_{i'}$. Averaging in $i'$
with weights $d_{i'}$ gives $\mu_{\max}(N_i)-h_i-K_3-[\ell]\le s$ for every $i$;
take the maximum over $i$.
\end{proof}

\begin{lemma}\label{lem:phi-bound}
If $\bar u\ge0$, $0<H\le Xb$, $b>0$, $X>0$ and $\bar u\le U$, then
$\bar u\log(1+H/\bar u)\le\log(1+X)(U+b)$ (the left side read as $0$ when $\bar u=0$).
\end{lemma}
\begin{proof}
For $\bar u=0$ the right side is $\ge0$. For $\bar u>0$ apply
$\log y\le y-1$ to $y=(1+H/\bar u)/(1+X)$:
$\bar u\log(1+H/\bar u)\le\bar u\log(1+X)+\bar u(1+X)^{-1}+H(1+X)^{-1}-\bar u$.
Use $H(1+X)^{-1}\le b(1-(1+X)^{-1})$ (from $H\le Xb$),
$1-(1+X)^{-1}\le\log(1+X)$ (from $\log y\le y-1$ at $y=(1+X)^{-1}$),
$\bar u(1-(1+X)^{-1})\ge0$, and $\bar u\log(1+X)\le U\log(1+X)$.
\end{proof}

\begin{proposition}[pointwise bound]\label{prop:pointwise}
Under (P0)--(P10), for every $\omega$,
\[
 P_*(\omega)\le M_*(\omega)+\mathrm{Err}_\omega,
\]
\[
\begin{aligned}
 \mathrm{Err}_\omega={}&\Sigma_a\,c_{\rm count}\Bigl(L_a+\frac{m^2L_a}k\Bigr)+E_{\rm div}+mc_EB_0
 +c_{\rm slope}raL_a\Bigl(1+\frac mk\Bigr)+\log|I_\omega|+\mathrm{LoP}+M_{\rm Lo}\\
 &+(1-\bar\delta)^{-1}\Bigl(4\log(1+x^{c_{HH}})\bigl(K_3+M_{\rm Lo}+b\bigr)
   +(\bar\eta+4\sqrt{\bar e})\,x^{c_{HH}}b\Bigr),
\end{aligned}
\]
where $M_{\rm Lo}=c_{\rm sl1}\mathrm{LoP}+c_{\rm sl1}rm(a/k+b')c_{\rm count}mL+c_{\rm sl1}X''+1$.
\end{proposition}
\begin{proof}
\emph{Step 1 (the gain).} By Lemma~\ref{lem:key-identity},
$K_5=M_*+K_0$ with $K_0=Ac_{\rm count}(L_a+m^2L_a/k)+E_{\rm div}+mc_EB_0+c_{\rm slope}raL_a(1+m/k)$.
By (P8) (Input~A, at the package values of $\delta_0,\eta,e,H$),
$s\le\log|I|+\mu^\vee+(1-\delta_0)^{-1}(4\bar u\log(1+H/\bar u)+(\eta+4\sqrt e)H)$.
Combining with Lemma~\ref{lem:s-lower},
\[
 g\le M_*+K_0+[\ell]+\log|I|+(1-\delta_0)^{-1}\bigl(4\bar u\log(1+H/\bar u)+(\eta+4\sqrt e)H\bigr).
\]
By Lemma~\ref{lem:uF} and Lemma~\ref{lem:phi-bound} with $X=x^{c_{HH}}$ (P9) and
$U=K_3+[\ell]$, $4\bar u\log(1+H/\bar u)\le4\log(1+x^{c_{HH}})(K_3+[\ell]+b)$; by
(P10), $(\eta+4\sqrt e)H\le(\bar\eta+4\sqrt{\bar e})x^{c_{HH}}b$ and
$(1-\delta_0)^{-1}\le(1-\bar\delta)^{-1}$, all quantities being nonnegative.

\emph{Step 2 (the losses).} By (P4), $\log T_i\le c_{\rm count}mL$, hence
$[\ell]\le c_{\rm sl1}I_O+c_{\rm sl1}rm(a/k+b')c_{\rm count}mL+c_{\rm sl1}X''+1$,
and by Input~B (P7), $c_{\rm sl1}I_O\le c_{\rm sl1}\mathrm{LoP}$; thus $[\ell]\le M_{\rm Lo}$.
Again by (P7), $P_*\le g+\mathrm{LoP}$. Substituting $[\ell]\le M_{\rm Lo}$ into Step~1
(the coefficient $1+(1-\bar\delta)^{-1}4\log(1+x^{c_{HH}})$ of $[\ell]$ being
nonnegative) gives $P_*\le M_*+\mathrm{Err}_\omega$.
\end{proof}

\begin{lemma}[absorption of constants]\label{lem:absorb}
Under the size windows with $(e_K,K)$ and \eqref{eq:assembly-largeX},
$\mathrm{Err}_\omega\le x^{(e_D+e_L+3)\varepsilon}\,E_{\rm text}$, where
$E_{\rm text}=ra(1+m^2/k)+m(B_0+N_c)$.
\end{lemma}
\begin{proof}
Put $y=u$, $y_D=x^{e_D\varepsilon}$, $y_L=x^{e_L\varepsilon}$, all $\ge1$, $Y=y_Dy_Ly$
and $P=Y E_{\rm text}$. The size windows give $b\le a$, $m\le x$, $m\le k$,
$mb\le a$, $\Sigma_a\le2ra$, $\log|I_\omega|\le(k-3)\log2\le k\le a$,
$\log k\le(1+e_K)u$ (from $k\le x^{e_K}$ and $\log x\le u$),
$\log(1+x^{c_{HH}})\le(1+c_{HH})u$, and $(1-\bar\delta)^{-1}\le2$. With these,
each of the eight summands of $\mathrm{Err}_\omega$ except the last is at most
a polynomial in $K$ times $P$: for instance
$Ac_{\rm count}(L_a+m^2L_a/k)\le 4KP$, $E_{\rm div}\le3KP$,
$c_{\rm slope}raL_a(1+m/k)\le4KP$, $\log|I_\omega|\le P$, $\mathrm{LoP}\le KP$, and
$M_{\rm Lo}\le(5K^2+K+1)P$; the bracket $K_3+M_{\rm Lo}+b$ is at most
$(5K^2+13K+2K(1+e_K)+2)P$, and the last term is at most
$2\bigl(4(1+c_{HH})y\,(5K^2+13K+2K(1+e_K)+2)P+KP\bigr)$ using the last size-window
condition. Summing, $\mathrm{Err}_\omega\le\bigl[(5K^2+16K+2)+8(1+c_{HH})(5K^2+13K+2K(1+e_K)+2)\bigr]yP
\le200(1+K)^2(1+c_{HH})(1+e_K)\,yP\le y^2P$ by \eqref{eq:assembly-largeX}.
Finally $y^2P=x^{(e_D+e_L+3)\varepsilon}E_{\rm text}$.
\end{proof}

\begin{proof}[Proof of Theorem~\ref{thm:mean-bound}]
By Lemma~\ref{lem:ident}, $P_*(v)-M_*(v)=P_*(\omega)-M_*(\omega)$. By
Proposition~\ref{prop:pointwise} and Lemma~\ref{lem:absorb} this is at most
$x^{(e_D+e_L+3)\varepsilon}E_{\rm text}$. By the bracket bounds,
$E_{\rm text}\le ra(1+m^2/k)+m(2nmbu+skbu^2/n^2)=Z$ (the sizes $r=|T|$ agree by
the bijection of the top index set with $T$), and $x^{(e_D+e_L+3)\varepsilon}\ge0$.
\end{proof}

\begin{remark}[quantifiers and uniformity]\label{rem:mean-uniform}
Theorem~\ref{thm:mean-bound} is a deterministic statement at one $x$: no
quantifier over windows occurs. Uniformity in the window $W$ enters only
through Theorem~\ref{thm:alpoge-families}, which supplies, for each fixed $\varepsilon$, a single
family and constants $K,c_{HH},e_K,C_L$ and $x_0$ valid at every $x\ge x_0$ in
the window set of the family; the exponents $e_D,e_L$ are fixed before
$\varepsilon$. The threshold in \eqref{eq:assembly-largeX} then depends only on
$(\varepsilon,K,c_{HH},e_K)$.
Hence the constant $1$ and exponent $e_D+e_L+3$ are uniform in $x$, in the
translation, in the chosen columns and in $W$.
\end{remark}





\section{The theorem on families and the arithmetic inputs}\label{sec:inputs}

This section states the theorem on families of column packages, due to
L.~Alpoge, and deduces from it the column estimate required by
Proposition~\ref{prop:simultaneous-slope} (\S\ref{subsec:closing}). It then
states the arithmetic results from which the theorem is derived and shows
how they combine (\S\S\ref{inp:ssec-overview}--\ref{inp:ssec-status}).

\subsection{The theorem on families}\label{subsec:closing}
The following statement is the existence theorem for families of column
packages, specialised to the trivial character; it is due to L.~Alpoge
\cite{CharacterFamily}. It is derived in the companion paper \cite{Annex}
(\cite[Theorem~\ref{A-i4:thm-final}]{Annex}) from five arithmetic inputs, which
are stated in \S\S\ref{inp:ssec-overview}--\ref{inp:ssec-exist}. The
companion paper proves the reductions in full, quotes the published theorems
used in its written proofs with their statements, identifies which steps are
known theorems re-proved in the formalisation \cite{CharacterFamily}, and
states the new results taken from it as propositions marked \emph{formally
verified}, which are not reproved. The theorem rests on Alpoge's Lean
formalisation and has not yet been written out on paper in full. We have
compared the statement below clause by clause with the formal statement.

\begin{theorem}[Alpoge \cite{CharacterFamily}: packages with uniform constants]\label{thm:alpoge-families}
There are real numbers $e_D,e_L\ge0$ and $0<\varepsilon_1\le1/52$ with the
following property. For every $0<\varepsilon<\varepsilon_1$ there are real
numbers $K,c_{HH},e_K,C_L,x_0$ and a set $\mathcal X\subset\R$ such that:
\begin{enumerate}
\item[(a)] for every $y$ there is $W\ge y$ with $[W,W^{3/2}]\subset\mathcal X$;
\item[(b)] for every $x\in\mathcal X$ with $x\ge x_0$, writing $u=x^\varepsilon$,
there are primes $q,d$, a finite abelian group $A$ of order $n$, a host field
$E$ with group $U=(\Z/2)^2\times A$ and its arithmetic data, subsets
$T,S\subset A$ of sizes $r,s$, an integer $m\ge0$, a cutoff $N$, numerators
$(a_t)_{t\in T}$, weights $w_{v,j}$ for every $v\in A$ and $1\le j\le m$, and
real numbers $a,b,b',k,\eta,X',\vartheta$, with the following five properties.
The host is a host in the sense of \cite[Definition~\ref{A-i4:def-host}]{Annex};
in particular: $q\equiv3\pmod8$, $d\equiv7\pmod{16}$, $(-d/q)=1$, and
$(-d/l)=1$ for every odd prime $l\mid Q_S$, where $Q_S$ is an integer prime
to $qd$ and $N_E=Q_Sqd$; $E\subset\Q(\zeta_{N_E})$ contains $\sqrt{-d}$ and
$\sqrt{-q}$; the Artin map $\mathrm{Art}:(\Z/N_E)^\times\to U$ is such that
the $d$-sign and the $q$-sign of $\mathrm{Art}(u)$ (its actions on
$\sqrt{-d}$ and on $\sqrt{-q}$) are $(u/d)$ and $(u/q)$, and
$\mathrm{Art}(u)$ is trivial or equal to $\iota_d$ whenever $u\equiv1\pmod{Q_Sq}$, where
$\iota_d$ has $d$-sign $-1$, $q$-sign $1$ and trivial $A$-component; and
$I_p$, $\mathrm{Fr}_p$ are the inertia group and a Frobenius element of $p$
in $\mathrm{Gal}(E/\Q)\cong U$, namely the image under $\mathrm{Art}$ of the
units congruent to $1$ modulo the prime-to-$p$ part of $N_E$, and the image
of a unit congruent to $p$ modulo the prime-to-$p$ part and to $1$ modulo
the $p$-part. The consequences of these conditions that the argument uses
are collected in Lemma~\ref{lem:host-facts}.
\begin{enumerate}
\item[(1)] \emph{Sizes.} $b\ge1$, $xb\le a\le2xb$, $s/n\ge1/u$,
$r\ge n/(2u^{1/8})$, $\sqrt x/u\le m\le\sqrt x\,u^{1/2}$ and
$xm\le k\le n^2\sqrt x/u$.
\item[(2)] \emph{Twists.} For each of the three nontrivial characters $\chi$
of the host group that are trivial on $A$, $|G_\chi(x)|\le\sqrt x\,u$ and
$|L'(1,\chi)/L(1,\chi)|\le C_L\log x$.
\item[(3)] \emph{Weights.} $\eta,\vartheta\ge0$, $X'>0$, $b'>0$, $b'X'\le xb$ and
$|a_t/b-x|\le1$ for $t\in T$. For every $v\in A$ and every $j$ the weight
$w_{v,j}$ is positive on $\Phi$, lies in $[b,b+\eta+1]$ at the positions
$(0,1,h')$, $h'\in S+v$ (the \emph{thick positions} of $v$), is at most
$2b'$ at the other positions of $\Phi$, and satisfies
$a_t/w_{v,j}(\rho)\le N$ for $t\in T$, $\rho\in\Phi$.
\item[(4)] \emph{Thin scales and cost.} $|G_\xi(y)|\le y\vartheta$ for every
character $\xi$ of $U$ and every $y\ge X'/3$;
$4n(x+1)\vartheta\le1/x$; and $12n(x+1)(\eta+1)\le b/x$.
\item[(5)] \emph{Package.} There is a column package
(Definition~\ref{def:package}) read from this host along the identity map of
$A$, that is, with one lift for each element of $A$, with scale
$(x,\varepsilon)$, with these $a,b,b',k,m$, top set $T$ and set $S$ (so that
its bracket is the $Z$ of \eqref{eq:assembly-Z} for these values), which
satisfies the size windows with $(e_K,K)$ and the
bracket bounds, whose divisor exponent is at most $e_D$, whose loss exponent
is at most $e_L$, and whose constant $c_{HH}$ is at most the given value.
\end{enumerate}
\end{enumerate}
None of $e_D,e_L,\varepsilon_1$ depends on $\varepsilon$, and none of
$K,c_{HH},e_K,C_L,x_0$ depends on $x$ or on $W$.
\end{theorem}
\begin{remark}\label{rem:alpoge-families}
Properties (1) and (2) are the size and twist conditions of
Proposition~\ref{prop:finite-to-window}. Properties (3) and (4) are the
conditions under which the prime-power inequality,
Lemma~\ref{lem:prime-power}, is derived; the cost condition used
there is the sum of the two inequalities in (4). The bounds are pointwise in
the translation and hold at every $x$ of a window, as required in
Proposition~\ref{prop:simultaneous-slope}.
For the trivial character the translation subgroup is all of $A$, and $A$
is cyclic.
\end{remark}

\begin{remark}[the exponents]\label{rem:exponents}
The exponents of the argument are fixed here once. The windows are
$[W,W^{3/2}]$ (Theorem~\ref{thm:alpoge-families}(a)), and
$0<\varepsilon<\varepsilon_1\le1/52$. The divisor and loss exponents $e_D$,
$e_L$ do not depend on $\varepsilon$; in the derivation of
Theorem~\ref{thm:alpoge-families} one may take $e_D=2$ and $e_L=0$
(\cite[Theorem~\ref{A-i4:thm-final}, Step~2]{Annex};
Section~\ref{sec:inputs}). The exponent of the column estimate is
$K_j=e_D+e_L+3$ (Corollary~\ref{cor:obligation}), which with these values is
$5$; Proposition~\ref{prop:finite-to-window} with $C=1$ and $K=5$ then gives
$G(x)\le31\sqrt{x}\,u^{61/8}$. Of these values only the independence of
$e_D,e_L$ from $\varepsilon$ is used; the bound $\varepsilon_1\le1/52$ is
the range in which the families of \cite{Annex} are constructed.
\end{remark}

\begin{corollary}[the column estimate]\label{cor:obligation}
Assume Theorem~\ref{thm:alpoge-families}. Let $\varepsilon_j\downarrow0$
with $\varepsilon_j<\varepsilon_1$. Then the hypothesis of
Proposition~\ref{prop:simultaneous-slope} holds with $C_j=1$ and
$K_j=e_D+e_L+3$: for each $j$ there are arbitrarily large $W$ such that at
every $x\in[W,W^{3/2}]$ beyond a threshold depending only on
$\varepsilon_j$ the data of Theorem~\ref{thm:alpoge-families}(b)
satisfy the hypotheses of Lemma~\ref{lem:prime-power} and
\[
 \mathbb E(P_*-M_*)\le u^{e_D+e_L+3}Z,\qquad u=x^{\varepsilon_j}.
\]
In particular $\varepsilon_jK_j\to0$, and for every $1/2<\theta<1$ there
are $C_\theta>0$ and arbitrarily large $Y$ with $G(x)\le C_\theta x^\theta$
on $[Y,Y^{3/2}]$.
\end{corollary}
\begin{proof}
Fix $j$ and write $\varepsilon=\varepsilon_j$. Let $K,c_{HH},e_K,x_0$ be the
constants of Theorem~\ref{thm:alpoge-families}, and enlarge $x_0$ so that
\eqref{eq:assembly-largeX} holds for $x\ge x_0$; this threshold depends
only on $(\varepsilon,K,c_{HH},e_K)$ (Remark~\ref{rem:mean-uniform}); enlarge it
also beyond the absolute threshold $x_*$ of Lemma~\ref{lem:prime-power}.
Properties (1)--(4) of
Theorem~\ref{thm:alpoge-families}(b) are the hypotheses of that lemma. At
each $x\ge x_0$ of a window, Theorem~\ref{thm:mean-bound} applies to the
package of Theorem~\ref{thm:alpoge-families}(b)(5) and gives
$P_*(v)-M_*(v)\le u^{e_D'+e_L'+3}Z$ for every translation $v$, where
$e_D'\le e_D$ and $e_L'\le e_L$ are the exponents of the package. Since
$u\ge1$ and $Z\ge0$ the right side is at most $u^{e_D+e_L+3}Z$. Averaging
over the translations gives the displayed bound for the mean. The exponent
$K_j=e_D+e_L+3$ does not depend on $j$, so $\varepsilon_jK_j\to0$. The
last assertion is Proposition~\ref{prop:simultaneous-slope}.
\end{proof}

\subsection{Overview and dependency chain}\label{inp:ssec-overview}

The proofs of the arithmetic inputs are long, and they are given in the
companion paper \cite{Annex}, which has five sections:
\begin{center}
\begin{tabular}{lll}
\cite[Section~\ref{A-i1:sec}]{Annex} & the loss bound & (Annex Section A)\\
\cite[Section~\ref{A-i2:sec}]{Annex} & the divisor bound & (Annex Section B)\\
\cite[Section~\ref{A-i3:sec}]{Annex} & level counts and almost orthogonality & (Annex Section C)\\
\cite[Section~\ref{A-i5:sec}]{Annex} & the windowed Selmer bound & (Annex Section D)\\
\cite[Section~\ref{A-i4:sec}]{Annex} & existence of families with uniform constants & (Annex Section E)
\end{tabular}
\end{center}
We use the letters A--E only for the sections of \cite{Annex}. They are
not the letters of Definitions~\ref{inp:F}--\ref{inp:count} (Inputs A--D of
\S\ref{subsec:assembly-dict}), which name conditions of a column
package; the correspondence is recorded in Table~\ref{inp:tab-conditions}.

Nothing in the rest of this section is proved here except the final deduction
(Theorem~\ref{inp:thm-families}). For each input we state the main result,
outline its proof with references to \cite{Annex}, and say which clause of
Theorem~\ref{thm:alpoge-families}, or which condition (P0)--(P10) of
Definition~\ref{def:package}, it supplies. The last subsection
(\S\ref{inp:ssec-status}) says exactly what is proved in \cite{Annex}, what
is a published theorem, and what is taken from the formal verification
\cite{CharacterFamily} without a written proof. The reader should consult
it before relying on any statement of this section.


The proofs of the Riemann hypothesis (Theorem~\ref{thm:RH}) and of
$W_T\ge0$ (Corollary~\ref{cor:all-windows}) use
Theorem~\ref{thm:alpoge-families} only through
Corollary~\ref{cor:obligation}. The chain is
\[
\left.\begin{array}{l}
\text{(A) loss bound}\\
\text{(B) divisor bound}\\
\text{(C) level counts and almost orthogonality}\\
\text{(D) windowed Selmer bound}
\end{array}\right\}
\xrightarrow{\ \text{(E)}\ }
\text{Thm.~\ref{thm:alpoge-families}}
\Rightarrow
\text{Cor.~\ref{cor:obligation}}
\Rightarrow
\text{Thm.~\ref{thm:RH}}
\Rightarrow
W_T\ge0 ,
\]
where (E) is the existence argument for families.
In more detail:
\begin{enumerate}
\item The windowed Selmer bound (D) produces, at every scale $x$ of
arbitrarily long windows $[W,W^{3/2}]$, host fields whose Selmer groups of
Tate twists are uniformly small (the hypothesis $\mathrm{HypV}$). Through
the host supply of Section E it gives clause (a) of
Theorem~\ref{thm:alpoge-families} and the twist bounds (2).
\item The divisor bound (B) uses $\mathrm{HypV}$ to bound a weighted sum of
local congruence indices. It gives condition (P1) of a package, with
divisor exponent $e_D=2$.
\item The loss bound (A) gives condition (P7), with loss exponent
$e_L=0$.
\item The level counts (C) give condition (P4); the almost orthogonality
(Theorem S) and the irreducibility bound of (C) control the small frame
parameters that enter (P8) and (P10).
\item Section E fixes all constants in the right order, constructs the
families of automorphic origin, verifies that the four inputs hold on
them eventually and with uniform constants, and concludes with
\cite[Theorem~\ref{A-i4:thm-final}]{Annex}: the four inputs, together with
the statements listed in \cite[Remark~\ref{A-i4:rem-trust}]{Annex}, imply
Theorem~\ref{thm:alpoge-families}.
\item Corollary~\ref{cor:obligation} applies Theorem~\ref{thm:mean-bound} to
the packages of Theorem~\ref{thm:alpoge-families}(b)(5) and gives the
column estimate with exponent $K_j=e_D+e_L+3$, which with $e_D=2$, $e_L=0$
is $K_j=5$. The rest of the paper turns this into the Riemann hypothesis
(Theorem~\ref{thm:RH}) and then into $W_T\ge0$
(Corollary~\ref{cor:all-windows}).
\end{enumerate}

\begin{table}[ht]
\caption{Which input supplies which condition.}\label{inp:tab-conditions}
\small
\begin{tabular}{p{0.26\textwidth}p{0.27\textwidth}p{0.38\textwidth}}
\toprule
Condition & Content & Source\\
\midrule
(P1) & divisor bound (Def.~\ref{inp:div}) & Section B, \cite[Proposition~\ref{A-i2:prop-P1}]{Annex}\\
(P4) & counting (Def.~\ref{inp:count}) & Section C, \cite[Theorem~\ref{A-i3:thm-main}(d)]{Annex}\\
(P7) & losses (Def.~\ref{inp:loss}) & Section A, \cite[Theorem~\ref{A-i1:thm-loss}]{Annex}\\
(P8) & frame slope inequality (Def.~\ref{inp:F}) & Section E, the formally verified \cite[Proposition~\ref{A-i4:Lprop-fam}(b)]{Annex} (Theorem~\ref{thm:frameF} applied to the columns)\\
(P0), (P2), (P3), (P5), (P6), (P9), (P10) & bookkeeping of the package & Section E, \cite[Proposition~\ref{A-i4:Lprop-fam}(c)]{Annex}\\
clause (a) & windows of ratio $3/2$ & Sections D and E, \cite[Proposition~\ref{A-i4:prop-supply}]{Annex}\\
clause (b)(2) & twist bounds & Section E, \cite[Proposition~\ref{A-i4:Sprop-tw}]{Annex}\\
\bottomrule
\end{tabular}
\end{table}

Two remarks on this table. First, the package is \emph{assembled} in
Section E: the definition of ``chain-tied''
\cite[Definition~\ref{A-i4:def-chain}]{Annex} is that the data satisfy
(1)--(5) of Theorem~\ref{thm:alpoge-families}(b), and the statement that
the families constructed there are chain-tied is the core link
\cite[Proposition~\ref{A-i4:Lprop-fam}(c)]{Annex}, which is formally verified
and not proved in writing (\S\ref{inp:ssec-status}(c)). The conditions
(P0), (P2), (P3), (P5), (P6), (P9), (P10) are part of that statement. The
frame bound $\mathrm{FrameF}(C_H)$ for the columns, from which (P8) is obtained,
is one of the items of \cite[Proposition~\ref{A-i4:Lprop-fam}(b)]{Annex},
also formally verified, and enters the core link as a hypothesis. Second, the condition
labels (P1)--(P11), (M1)--(M7) and so on that occur inside individual
proofs of \cite{Annex} are local lists and have nothing to do with
Definition~\ref{def:package}.

\emph{The order of the constants.} Theorem~\ref{thm:alpoge-families}
quantifies in the order $e_D,e_L,\varepsilon_1$; then $\varepsilon$; then
$K,c_{HH},e_K,C_L,x_0$; then $x$. Corollary~\ref{cor:obligation} needs
exactly this order: the exponent $K_j=e_D+e_L+3$ must not depend on
$\varepsilon_j$, so that $\varepsilon_jK_j\to0$, while the threshold in $x$
may depend on everything chosen before it. In the derivation the two
exponents come out as absolute numbers: $e_L=0$ because the loss bound of
Section A has no power of $x$ in it (Theorem~\ref{inp:thm-loss}), and
$e_D=2$ because the division link is applied with $C_{\rm abs}=2$
(\cite[Proposition~\ref{A-i4:prop-links-formal}(a)]{Annex}). The
constants $K,c_{HH},e_K,C_L$ are produced by the core link
(Proposition~\ref{inp:prop-corelink} below) from the constants of the
family, which are fixed in stages by threshold data
(\cite[\S\ref{A-i4:ssec-closed}]{Annex}). Each input is used in a
\emph{closed} form: it holds for every family of a class that is large
with respect to a threshold datum chosen after $\varepsilon$ but before
$x$. This is why the statements below are formulated for families and
classes rather than for single hosts.

\emph{The level counts and (P4).} The counting condition $\mathrm{Count}(C)$
of Section C (Definition~\ref{inp:def-count}) has the same three
inequalities as Definition~\ref{inp:count}, indexed by words instead of
columns; that it gives (P4) for the packages of Section E is
\cite[Proposition~\ref{A-i4:hyp-S}(c)]{Annex}.

\subsection{The loss bound}\label{inp:ssec-loss}

\emph{Setting.} A family $\mathrm{Fm}$ at exponent $\varepsilon$ has, at
each large scale $x$ and index $i$, a datum $D_{x,i}$, a size datum
$L_{x,i}$ and a weight vector $X_{x,i}$; the scale quantities are
$b=e^{x^\varepsilon}$, $a=xb$, and $|\mathrm{TT}|$ is the number of tops,
$m$ the number of lifts. The \emph{ask} at a top $t$ and lift $j$ is a
truncated series $\mathrm{ASK}(N,a_t,t,w_j)$ in the prime powers up to the
cutoff $N$, and the \emph{gain} $\mathrm{Ga}(L_{x,i})$ is a sum over places
of local gains minus taxes. The local gains are defined as truncated,
shifted copies of the local ask series
\cite[Remark~\ref{A-i1:rem-gain}]{Annex}, so the loss bound is a lower bound
for the gain which holds essentially by construction; the real constraint
on the gain is the source-slope condition (P6) together with (P8) and
(P1). In the assembly the loss bound is used only in the form
$P_*\le g+\mathrm{LoP}$, $I_O\le\mathrm{LoP}$.

The objects in the next theorem are defined in
\cite[\S\ref{A-i1:ssec-family}]{Annex}: the family and its points
\cite[Definition~\ref{A-i1:def-family-data}]{Annex}, the conditions
$\mathrm{ChosenTT}$, $\mathrm{ChosenBadp}$, $\mathrm{ChW}$
\cite[Definition~\ref{A-i1:def-chosen}]{Annex}, the prime-by-prime
comparison $\mathrm{LinkG}$ \cite[Definition~\ref{A-i1:def-linkG}]{Annex},
and the gain $\mathrm{Ga}$, the index loss $\mathrm{IOsum}$ and the
readings \cite[Definition~\ref{A-i1:def-losses}]{Annex}.

\begin{theorem}[{\cite[Theorem~\ref{A-i1:thm-loss}]{Annex}}; loss bound for small-share weights]\label{inp:thm-loss}
Let $\mathrm{Fm}$ be a family at exponent $\varepsilon$ that satisfies
$\mathrm{ChosenTT}$, $\mathrm{ChosenBadp}$ and $\mathrm{LinkG}$, with constant
$C$ in $\mathrm{LinkG}$. Put $K=1+|C|+C_{\mathrm{ord}}$, where
$C_{\mathrm{ord}}=12$ is the absolute constant of
\cite[Lemma~\ref{A-i1:lem-orders}]{Annex}.
Then for all large $x$ and all indices $i$ at scale $x$ such that
$\mathrm{ChW}(X_{x,i})$ holds,
\begin{equation}\label{inp:eq-loss}
 \mathrm{IOsum}(L_{x,i})\le(3K+13)\,|\mathrm{TT}|\,m\,b,\qquad
 \sum_{t\in\mathrm{TT}}\sum_{j=1}^{m}\mathrm{ASK}(N,a_t,t,w_j)-\mathrm{Ga}(L_{x,i})
 \le(3K+13)\,|\mathrm{TT}|\,m\,b
\end{equation}
for every reading of $L_{x,i}$ and every admissible cutoff $N$.

Moreover the family satisfies $\mathrm{LossF}(0)$ with constant
$3K+13$: there is $x_0$ such that for every $x\ge x_0$ with
$x\in\mathrm{Win}$ there are an index $i_0$ at scale $x$ and integers
$(\bar a_t,\bar b_t)_{t\in\mathrm{TT}}$ with $\bar a_t$ in the amount window
of $L_{x,i_0}$ and $b'\le\bar b_t\le b_{\rm top}(x)$
\cite[Definition~\ref{A-i1:def-family-data}(8)]{Annex}, such
that~\eqref{inp:eq-loss} holds at every index $i$ at scale $x$ whose
amounts and top windows are $a_t=\bar a_t$, $b_t=\bar b_t$ for all tops
$t$.
\end{theorem}

\emph{Proof outline.} The proof in \cite[\S\ref{A-i1:ssec-overview}]{Annex}
has three layers.
\begin{enumerate}
\item An elementary, real-variable bound for an abstract ask--gain
difference (\cite[Theorem~\ref{A-i1:thm-elementary}]{Annex}). The primes are
sorted into four classes (served, first kind, second kind, high) and each
class is estimated separately. The only arithmetic inputs are Chebyshev's
and Mertens' bounds (\cite[Lemma~\ref{A-i1:lem-cheb}]{Annex}) and the
distribution of odd integers in the square classes modulo powers of $2$.
\item An instantiation (\cite[Theorem~\ref{A-i1:thm-setting}]{Annex}): the
column data of a family point satisfy the abstract hypotheses. The host
field enters only through elementary facts on Frobenius elements in a
subfield of $\Q(\zeta_N)$, with quadratic reciprocity
(\cite[Theorem~\ref{A-i1:thm-qr}]{Annex}). The only input that is not
bookkeeping is $\mathrm{LinkG}$. The second part of the theorem uses in
addition the existence of small-share weight vectors
(\cite[Proposition~\ref{A-i1:prop-chosenW}]{Annex}), which follows from
$\mathrm{LinkG}$. These two layers, and the main theorem itself, are
proved in full.
\item The source of $\mathrm{LinkG}$ (\cite[\S\ref{A-i1:ssec-linkG-source}]{Annex}).
It is reduced to group
compatibility of the families and the second admissibility condition
$\mathrm{Adm2F}$ \cite[Definition~\ref{A-i1:def-groupcompat}]{Annex}, through
a local statement about lattices in unramified quadratic Hida settings
(\cite[Proposition~\ref{A-i1:prop-loc}]{Annex}) and $37$ admission clauses
(\cite[Proposition~\ref{A-i1:prop-adm}]{Annex}). The admission clauses are
proved in \cite[\S\ref{A-i1:ssec-pf-adm}]{Annex}, using among other things
the Weyl form and sandwich statement
(\cite[Proposition~\ref{A-i1:prop-Fm}]{Annex}, proved in
\cite[\S\ref{A-i1:ssec-pf-Fm}]{Annex}) and
\cite[Propositions~\ref{A-i1:prop-A19c} and~\ref{A-i1:prop-A21}]{Annex}. What
is not proved in writing is the counting part of the clause
$\mathcal A_{27}$ (\cite[Proposition~\ref{A-i1:prop-segdeg-rest}]{Annex}),
the closed form (\cite[Proposition~\ref{A-i1:prop-closed}]{Annex}), the
first theorem at the family (\cite[Proposition~\ref{A-i1:prop-first}]{Annex}),
the consequences of group compatibility
(\cite[Proposition~\ref{A-i1:prop-gc}]{Annex}) and the remaining inputs of
$\mathrm{LinkG}$ (\cite[Proposition~\ref{A-i1:prop-linkG-res}]{Annex}); see
\S\ref{inp:ssec-status}(c).
\end{enumerate}
The first layer is the following statement, in the notation of
\cite[\S\S\ref{A-i1:ssec-params}--\ref{A-i1:ssec-elementary}]{Annex}: a
parameter record $P$ lies in the windows $\mathrm W(e,K)$ of
\cite[Definition~\ref{A-i1:def-windows}]{Annex}, $\sigma=100Kx^{4\varepsilon}/x$,
and $\mathrm{PCLH}$ is the list of prime-class hypotheses of
\cite[Definition~\ref{A-i1:def-PCLH}]{Annex}.

\begin{theorem}[{\cite[Theorem~\ref{A-i1:thm-elementary}]{Annex}}; elementary loss bound]\label{inp:thm-elementary}
Let $\varepsilon\in(0,\frac14)$ and $e,K\in\R$. There is $x_0$ with the
following property. Let $P\in\mathrm W(e,K)$ be a record with
$\varepsilon(P)=\varepsilon$, $C_{\mathrm{ord}}(P)\ge12$ and $x(P)\ge x_0$, and
let the data of \cite[Definition~\ref{A-i1:def-PCLH}]{Annex} satisfy $\mathrm{PCLH}$. Then
\[
\mathrm{ASK}_{\mathcal D}-\mathrm{Ga}_{\mathcal D}\le(2\log4\cdot C+13\sigma)\,|T|\,|J|\,b\le(3C+13)|T||J|b,\qquad
\sum_{p\in\mathrm{cand}}\log p\,\mathrm{io}(p)\le\sigma|T||J|b\le|T||J|b .
\]
\end{theorem}

\noindent Since $\sigma=100Kx^{4\varepsilon-1}\to0$, the proof gives the sharper bound
$(2\log4\cdot C+13\sigma)|\mathrm{TT}|mb$ in Theorem~\ref{inp:thm-loss}; the
constant $3K+13$ is not absolute, because it depends on the constant $C$ of
$\mathrm{LinkG}$.

The reduction of $\mathrm{LinkG}$ is the following statement. Here group compatibility
(\cite[Definition~\ref{A-i1:def-groupcompat}]{Annex}) asks for a record with
five fields (a hermitian level structure, form data at every word, a
comparison, ring isomorphisms at the local gatherings, and gathering
compatibilities; the last is a paraphrase of the formal predicate), and
$\mathrm{Adm2F}$ asks that the top-set and bad-prime conditions
$\mathrm{ChosenTT}$, $\mathrm{ChosenBadp}$ hold eventually. Closedness and
threshold data are as in \cite[Definition~\ref{A-i4:def-closed}]{Annex}.

\begin{proposition}[{\cite[Proposition~\ref{A-i1:prop-linkG}]{Annex}}; $\mathrm{LinkG}$ from group compatibility]\label{inp:prop-linkG}
Let $\mathcal C$ be a class of families. If $\mathcal C$ is closed for group
compatibility and closed for $\mathrm{Adm2F}$, then $\mathcal C$ is closed
for $\mathrm{LinkG}$. Explicitly: there is $\varepsilon_1>0$ such that for
every $\varepsilon\in(0,\varepsilon_1]$ there is a threshold datum $\kappa$
such that $\mathrm{LinkG}(\mathrm{Fm})$ holds for every family
$\mathrm{Fm}\in\mathcal C$ of exponent $\varepsilon$ which is $\kappa$-large
at some $C_H$.
\end{proposition}

\noindent The proof in \cite[\S\ref{A-i1:ssec-pf-linkG}]{Annex} is an
assembly of closedness statements: a service lattice condition proved from
an all-but-finitely-many statement for the lattices, zero chains, and the
gluing of the local comparisons, from the local statement
\cite[Proposition~\ref{A-i1:prop-loc}]{Annex} and the three formally verified
\cite[Propositions~\ref{A-i1:prop-first}, \ref{A-i1:prop-gc},
\ref{A-i1:prop-linkG-res}]{Annex}. The local statement is in turn reduced to
the $37$ admission clauses and the closed form
\cite[Proposition~\ref{A-i1:prop-closed}]{Annex}.

Group compatibility and $\mathrm{Adm2F}$ are not proved in Section A; they
are supplied for the families of Section E
(\cite[Proposition~\ref{A-i4:hyp-loss}(a)]{Annex}).

\emph{What it supplies.} Through the ties of the column package built from
the family at a large scale with these weights
(\cite[Lemma~\ref{A-i4:Llem-L8}]{Annex}), the second part of
Theorem~\ref{inp:thm-loss} is condition (P7) of
Definition~\ref{def:package} (Definition~\ref{inp:loss}), with
$c_{\rm loss}=3K+13$ and loss exponent $e_L=0$.

\subsection{The divisor bound}\label{inp:ssec-div}

\emph{Setting.} A datum $D$ has a CM host field $E=E_D$, abelian over $\Q$
of degree $2f$, generator primes $g_0,g_1$, a set of tops, and a coefficient
field with a finite set of places $w$, each with a residue characteristic
$p(w)$, ramification index $e_w$, weight $\lambda_w$ and twist label
$\mu_w$; a place is twisted if $\mu_w\ne\mu_0$. For each lift $j$ and place
$w$ there are local Hecke algebras with maximal ideals $i$ and local
congruence indices $\eta^{\rm loc}_i$, and $\eta_j(w)=\sum_i\eta^{\rm
loc}_i$. A size datum $L$ carries a scale $(x,\varepsilon)$, $x>1$,
$0<\varepsilon<1/52$, and the quantities $\mathsf L=x^\varepsilon$,
$b=e^{\mathsf L}$, $B_0=fmbx^\varepsilon$, $N_c$, and ranks
$d_{Z_0}(j)$. The full list is \cite[Definition~\ref{A-i2:def-data}]{Annex}.
For $l$ prime, $\ell_D(l,e)$ is the $\Z_l$-length of the
Tate--Shafarevich group of $\Q_l/\Z_l(\chi^{-e})$ (strict condition above
$l$, unramified elsewhere) \cite[Definition~\ref{A-i2:def-LData}]{Annex}. We
use the following definitions, copied from
\cite[Definitions~\ref{A-i2:def-vv}--\ref{A-i2:def-HypV}]{Annex}. For a real
sequence $(a_l)$, $\sum'_la_l$ is its sum if it is summable and $0$
otherwise.

\begin{definition}[{\cite[Definitions~\ref{A-i2:def-vv}, \ref{A-i2:def-LinkN}, \ref{A-i2:def-HypV}]{Annex}}]\label{inp:def-div}
For a datum $D$ and $l\in\N$ put $\mathrm{VV}_D(l)=\sum_{e=1}^{t-1}\ell_D(l,e)$ and
$\mathrm{vv}_D(l)=\mathrm{VV}_D(l)\log l$ if $l$ is prime, $\mathrm{vv}_D(l)=0$ otherwise.
$D$ is \emph{Sha-finite} if all the lengths $\ell_D(p,e)$, $0<e<t$, are
finite. For a size datum $L$ and $j\in\N$,
\[
\eta\Sigma_L(j)=\sum_{\substack{w\in\mathrm{supp}\,L\\ \mu_w\ne\mu_0}}\lambda_we_w\,\eta_j(w).
\]
For $C\in\R$, $\mathrm{LinkN}_{\rm at}(L,j,C)$ is the
statement: if $(\mathrm{vv}_D(l))_l$ is summable, then
\[
\eta\Sigma_L(j)\le d_{Z_0}(j)\Bigl(C\bigl(fmb\,\mathsf L^2+(\log x)^2N_c\bigr)+\textstyle\sum'_l\mathrm{vv}_D(l)\Bigr).
\]
$\mathrm{LinkN}(\mathcal F)$: there are $C_N,x_0$ such that $\mathrm{LinkN}_{\rm at}(L_{x,i},j,C_N)$
holds for all $x\ge x_0$, all $i\in I(x)$ with $D_{x,i}$ Sha-finite, and all lifts $j$.
$\mathrm{SmallLevelRankEq}(\mathcal F)$: eventually $d_{Z_0}(j)=D_0$ for all $i$ and lifts $j$,
where $D_0=D_0(D_{x,i})\in\N$ is a fixed rank attached to the datum.
For a number field $E$, $t\in\N$ and $B\in\R$, $\mathrm{VV}_\le(E,t,B)$ means: for every
algebraic closure $\Omega$ of $E$ and every choice of $l$-data $\mathrm{ld}(l)$ for all primes
$l$, (1) $\ell(\mathrm{Sha}(\mathrm{ld}(p),e))<\infty$ for all $p$ and $0<e<t$, and (2) for every
finite set $S$ of primes,
$\sum_{l\in S}\log l\sum_{e=1}^{t-1}\ell(\mathrm{Sha}(\mathrm{ld}(l),e))\le B$.
$\mathrm{HypV}(\mathcal F)$: eventually $\mathrm{VV}_\le(E_{D_{x,i}},t,\sqrt{b})$ for all $i\in I(x)$.
\end{definition}

\begin{theorem}[{\cite[Theorem~\ref{A-i2:thm-div}]{Annex}}; divisor bound]\label{inp:thm-div}
Let $\mathcal F$ be a family with exponent $\varepsilon>0$ satisfying
$\mathrm{HypV}(\mathcal F)$, $\mathrm{LinkN}(\mathcal F)$ and
$\mathrm{SmallLevelRankEq}(\mathcal F)$. Then there is $x_0$ such that for all
$x\ge x_0$, $i\in I(x)$ and lifts $j$,
\[
 \eta\Sigma_{L_{x,i}}(j)\le D_0\,(|C_N|+1)\,x^{\varepsilon}\,(B_0+N_c).
\]
\end{theorem}

The deduction of Theorem~\ref{inp:thm-div} from its three hypotheses is
short and is proved in full (\cite[\S\ref{A-i2:ssec-red}]{Annex}). The
substance is the hypothesis $\mathrm{LinkN}$, which is
\cite[Theorem~\ref{A-i2:thm-linkN}]{Annex}: under the published hypotheses
$\mathrm{Hypotheses}_2$ and Chebotarev's theorem, every family whose
members satisfy the twist condition $\mathrm{TwGood}$, the Galois
preconditions and the tie condition $\mathrm{Tie4}$ satisfies
$\mathrm{LinkN}$. Several terms of that statement ($\mathrm{TwGood}$, the
eighteen Galois preconditions, two clauses of $\mathrm{Tie4}$) are named
but not defined in \cite{Annex}, so \cite[Theorem~\ref{A-i2:thm-linkN}]{Annex}
is a paraphrase of the formal statement as far as these terms go, as is
said there.

\emph{Proof outline.} The proof (\cite[\S\S\ref{A-i2:ssec-perprime}--\ref{A-i2:ssec-local}]{Annex})
has three layers.
\begin{enumerate}
\item \emph{Summation over primes.} The weighted sum $\eta\Sigma_L(j)$ is
split by residue characteristic. At each prime $l$ an \emph{admissible
table} bounds the local index by local terms plus the $l$-part of the
Selmer length; the local terms are summed over $l$ by elementary
prime-sum estimates (\cite[Theorem~\ref{A-i2:thm-summed}]{Annex} and
\cite[Proposition~\ref{A-i2:prop-point}]{Annex}). This gives
\cite[Theorem~\ref{A-i2:thm-export}]{Annex}: admissible families satisfy
$\mathrm{LinkN}$ with $C_N=2000(|C_a|+1)^3$. This layer is proved in full.
\item \emph{Cuts.} A cut of the Hecke data
(\cite[Proposition~\ref{A-i2:prop-cut}]{Annex}) reduces the general case to
one where admissible tables exist, at a cost of $\log|G|$ for the level
group $G$; the bound is transferred back by
\cite[Theorem~\ref{A-i2:thm-transfer}]{Annex}. The existence of the tables
for the cut family is the paired-character condition
(\cite[Proposition~\ref{A-i2:prop-pcc}]{Annex}) and the commuting-endomorphism
condition (\cite[Proposition~\ref{A-i2:prop-cec}]{Annex}), which feed the
family conclusion \cite[Proposition~\ref{A-i2:prop-conclF}]{Annex}.
\item \emph{Local inputs} (\cite[\S\ref{A-i2:ssec-local}]{Annex}). At the
generator primes the local term comes from a crystalline lifting theorem
with explicit exponents (\cite[Proposition~\ref{A-i2:prop-lift}]{Annex}),
whose proof uses integral $p$-adic Hodge theory: Bhatt--Scholze
(\cite[Theorem~\ref{A-i2:thm-bs}]{Annex}), Fontaine, Herr, Wach--Berger, and
a chain of Breuil--Kisin steps. At all other primes it comes from a global
lattice-index argument of Ribet type
(\cite[Proposition~\ref{A-i2:prop-selmer}]{Annex}), and at large primes from
\cite[Proposition~\ref{A-i2:prop-large}]{Annex}. Here most of the steps are
taken from the formal verification (\S\ref{inp:ssec-status}(c)).
\end{enumerate}

The first layer is the following pair of statements. An \emph{admissible
table with constant $C$}, $\mathrm{AdmT}(L,\mathrm{tb},C)$, is a list of
local data (local lengths at the places of $E$, generator-prime terms, a
finite set $\mathrm{exc}$ of exceptional primes) satisfying the
clauses (adm1)--(exc) of \cite[Definition~\ref{A-i2:def-AdmT}]{Annex}; the
budgets $\lambda^B(l)$, $\lambda^G(l)$ are the explicit functions of
\cite[Definition~\ref{A-i2:def-budgets}]{Annex}; $\nu(D)=\log Q_S+\log(tr)$
bounds the cyclotomic order terms (\cite[Lemma~\ref{A-i2:lem-nu}]{Annex}); and $D$ is
\emph{admitted with constant $C$} if Sha-finiteness of $D$ implies the
existence of an admissible table with constant $C$.

\begin{theorem}[{\cite[Theorem~\ref{A-i2:thm-summed}]{Annex}}; summed local bound]\label{inp:thm-summed}
Assume $\mathrm{AdmT}(L,\mathrm{tb},C)$, $\mathrm{vv}_D$ summable, and let $j$ be a lift and
$\mathcal S=\{p(w):w\in\mathrm{supp}\,L\ \text{twisted}\}$. Then
\[
\eta\Sigma_L(j)\le d_{Z_0}(j)\Bigl(6+\nu(D)+\log h_E+\log|G|+\sum_{l\in\mathcal S}\log l\,\lambda^G(l)\Bigr)+d_{Z_0}(j)\textstyle\sum'_l\mathrm{vv}_D(l).
\]
\end{theorem}

\begin{theorem}[{\cite[Theorem~\ref{A-i2:thm-export}]{Annex}}; admissible families satisfy $\mathrm{LinkN}$]\label{inp:thm-export}
Let $\mathcal F$ be a family such that, for large $x$, every $D_{x,i}$ is admitted with constant
$C_a$. Then $\mathrm{LinkN}(\mathcal F)$ holds with $C_N=2000(|C_a|+1)^3$.
\end{theorem}

\noindent In Theorem~\ref{inp:thm-summed} the class number $h_E$ and the
level group $G$ enter only through $\log h_E\le Cx^C$ and
$\log|G|\le Cx^C\mathsf L$ (clause (Par)); at large scales the whole
bracket is at most $2000(C+1)^3fmb\,\mathsf L^2$
(\cite[Proposition~\ref{A-i2:prop-point}]{Annex}). The point of the cuts in layer 2 is that a
member need not itself be admitted: \cite[Theorem~\ref{A-i2:thm-transfer}]{Annex}
transfers $\mathrm{LinkN}$ from the cut family, at the cost of a
multiplicity bound and the rank clause (adm3).

\emph{What it supplies.} The divisor bound is consumed through division
kernel packages. \cite[Proposition~\ref{A-i2:prop-P1}]{Annex} shows that,
under the division-data clauses $\mathrm{DKCB}$ and the conclusion of
Theorem~\ref{inp:thm-div}, there are $C,x_0$ such that for all $x\ge x_0$,
$i\in I(x)$ and every word $u$,
\[
\mu^\vee+h(u)-F\le\mathrm{MAIN}(u)+C\,|T|\,a\,L_a+C\,m\,(B_0+N_c)\,x^{C_{\rm abs}\varepsilon},
\]
with $\mathrm{MAIN}(u)=\sum_{t\in T}\bigl(m(\phi(a_t)-a_t\log\Omega)-a_t\log T_u\bigr)$
and $\phi(y)=y\log y-y$. This is condition (P1) of
Definition~\ref{def:package} (Definition~\ref{inp:div}), with
$c_{\rm div}=C$ and divisor exponent $e_D=C_{\rm abs}$; the top-level
instance has $C_{\rm abs}=2$. The definition of $\mathrm{DKCB}$
(\cite[Definition~\ref{A-i2:def-DKCB}]{Annex}) names some division-data
predicates without defining them, so
\cite[Proposition~\ref{A-i2:prop-P1}]{Annex} is exact only in the displayed
inequality. The chain is: Section D gives $\mathrm{HypV}$;
\cite[Theorem~\ref{A-i2:thm-linkN}]{Annex} gives $\mathrm{LinkN}$;
Theorem~\ref{inp:thm-div} gives the divisor bound;
\cite[Proposition~\ref{A-i2:prop-P1}]{Annex} gives (P1).

\subsection{Level counts and almost orthogonality}\label{inp:ssec-counts}

\emph{Setting.} A family is a collection $(\mathcal D(x,i),\mathcal
L(x,i))$, $x>1$, $i\in I_x$. Each $\mathcal D$ is a column datum: a CM
extension $E/F$ with $f=[F:\Q]$, ranks $k,m$ and $r=k+m-1$, a finite set of
\emph{words} $u$, a prime level $T_u$ and a positive integer $d_u$ for each
word, an expected dimension $D_*>0$, a total dimension $D_{\rm all}$ and a
row space of dimension $D_{\rm rows}=(1-\delta)D_{\rm all}$, which defines
the irreducibility defect $\delta$. The datum has a kernel operator with
blocks $K(u,u')$. The record carries $L=x^\varepsilon$, $b=e^L$, an
auxiliary exponent $\varepsilon_z=\varepsilon/J$, and
$\Gamma=\exp(C_\Gamma m^5fL)$ with a fixed constant $C_\Gamma>0$. A
property $P$ is \emph{closed} on a class $\mathcal C$ if there is
$\varepsilon_1>0$ such that for every $0<\varepsilon\le\varepsilon_1$ there
is a threshold system $\mathcal K$ for the closure constants with
$P(\mathcal F)$ for every $\mathcal K$-large $\mathcal F\in\mathcal C$ of
exponent $\varepsilon$ (\cite[\S\ref{A-i3:ssec-setting}]{Annex}). We copy the
two definitions the theorem needs.

\begin{definition}[{\cite[Definition~\ref{A-i3:def-linkS}]{Annex}}; almost orthogonality]\label{inp:def-linkS}
Let $C_0\ge1$ and $\eta_0,\eta_1,e_0\in\R$.  The datum $\mathcal D$ is
\emph{$(C_0,\Gamma,\eta_0,\eta_1,e_0)$-almost orthogonal} if the following
implication holds: if $C_0^{-1}\le d_u\Gamma^2/D_*\le C_0$ for every word $u$
(the \emph{window}), then
\begin{enumerate}
\item[(i)] $\|K(u,u')\|_{HS}^2\le(1+\eta_0)\,d_ud_{u'}/D_*$ for all $u\ne u'$;
\item[(ii)] $\|K(u,u)-\mathrm{id}\|_{HS}^2\le e_0\,d_u$ for all $u$;
\item[(iii)] $D_{\rm all}\le(1+\eta_1)D_*$.
\end{enumerate}
The family $\mathcal F$ satisfies $\mathrm{LinkS}(C_0)$ if eventually there
are $\eta_0,\eta_1\in\R$, $e_0\ge0$ with
$\eta_0,\eta_1,\sqrt{e_0}\le\exp(-x^{\varepsilon/8})$ such that
$\mathcal D(x,i)$ is $(C_0,\Gamma(\mathcal L(x,i)),\eta_0,\eta_1,e_0)$-almost
orthogonal.  It satisfies $\mathrm{LinkIRR}$ if eventually
$0\le\delta(\mathcal D(x,i))\le\exp(-x^{\varepsilon_z/7})$.
\end{definition}

\begin{definition}[{\cite[Definition~\ref{A-i3:def-count}]{Annex}}; level counts]\label{inp:def-count}
Let $v=|d_E|^{1/(2f)}/2$ (the host parameter), $g_0,g_1$ the two auxiliary
primes of $\mathcal D$, $w_j(\tau)$ ($1\le j\le m$, $\tau$ an infinite place
of $F$) its archimedean weights, $a=xb$, $L_a=\log(2av)$.  For $C\in\R$ the
datum satisfies $\mathrm{Count}(C)$ if for every word $u$
\begin{enumerate}
\item $\displaystyle\sum_{j=1}^m\frac1f\sum_\tau\log\frac{w_j(\tau)v}{\pi}
 +m\log(g_0g_1)-C\Bigl(L_a+\frac{m^2L_a}{k}\Bigr)\le\log T_u$;
\item $\log T_u\le CmL$;
\item $|\log T_u-\log T_{u'}|\le C\log k$ for every word $u'$.
\end{enumerate}
The family satisfies $\mathrm{CountF}$ if there is $C\ge0$ with eventually
$\mathrm{Count}(C)$.
\end{definition}

\begin{theorem}[{\cite[Theorem~\ref{A-i3:thm-main}]{Annex}}; level counts and almost orthogonality]\label{inp:thm-counts}
Let $\mathcal C$ be the working class $D_3(\mathcal A)$ of
\cite[\S\ref{A-i4:ssec-links}]{Annex}, with its spectral selection.  (The
unitary Siegel--Weil formula (SW) of
\cite[Definition~\ref{A-i3:S-def-SW}]{Annex}, used in the proofs, is
Theorem~\ref{inp:thm-SW} below; it is not a hypothesis.)
\begin{enumerate}
\item[(a)] (Off-diagonal bound.)  If every selected member is locally tied and
the rank row and the loop row hold with constant $3$, then the conclusion of
\cite[Proposition~\ref{A-i3:prop-S-closed}]{Annex} is closed on $\mathcal C$: eventually, for
all words $u\ne u'$, if the standing hypotheses hold at
$(3,C_{k,\rm low},C_l^{\mathcal F})$ then
$\|K(u,u')\|_{HS}^2\le(1+b^{-6})\,d_ud_{u'}/D_*$.
\item[(b)] (Almost orthogonality.)  Under the slot hypotheses
(Sel), (Rk), (Lp), (Ty), (Mn), (Inst) of
\cite[Theorem~\ref{A-i3:thm-linkS-slots}]{Annex},
the property ``$\mathrm{TuSize}(\mathcal F)$ implies $\mathrm{LinkS}(3)$'' is
closed on $\mathcal C$.
\item[(c)] (Irreducibility.)  Assume also Green's parametrisation (G)
\cite[Definition~\ref{A-i3:def-green}]{Annex}, that every member of $\mathcal C$ satisfies
$\mathrm{TuSize}$, and the automorphic tie hypothesis (AT) of
\cite[Proposition~\ref{A-i3:prop-irr-row}]{Annex}.  Then $\mathrm{LinkIRR}$ is closed on
$\mathcal C$.
\item[(d)] (Level counts.)  Every family of $\mathcal C$ whose members
eventually have the volume tie (VT) of
\cite[Proposition~\ref{A-i3:prop-massratio}]{Annex}
satisfies the torus-pair and variation bounds of
\cite[Proposition~\ref{A-i3:prop-part5}]{Annex}, and hence $\mathrm{CountF}$.  (VT) holds
eventually on every family of a class on which the real-place second-layer
ties are closed \cite[Definition~\ref{A-i4:def-hyps}(vii)]{Annex}.
\end{enumerate}
\end{theorem}

The statement (SW) of \cite[Definition~\ref{A-i3:S-def-SW}]{Annex} is the
following published theorem, with part (i) of Borel and Harish-Chandra and
part (ii) the $T$-th Fourier coefficient of Weil's Siegel formula
(\cite[Proposition~\ref{A-i3:prop-SW}]{Annex}). Notation: $F$ totally real,
$E/F$ totally imaginary quadratic, $J\in GL_M(E)$ hermitian,
$h(x,y)=\bar x^{\,t}Jy$, $U=U_J$; for $n\ge0$, $V_n=(E^M)^n$ with the Gram
map $\mathrm{gr}(X)=(h(X_a,X_b))_{a,b}$; for nondegenerate
$T\in\mathrm{Herm}_n(F)$ and a test function $\Phi$ on $V_n(\mathbb A_F)$ put
$\theta_T(h)=\sum_{X\in V_n(F),\,\mathrm{gr}(X)=T}\Phi(h^{-1}X)$.

\begin{theorem}[Siegel--Weil formula in the convergent range;
\cite{BorelHC1962}, \cite{Weil1965}]\label{inp:thm-SW}
\begin{enumerate}
\item[(i)] $U(F)$ is discrete in $U(\mathbb A_F)$ and has finite covolume;
every measurable fundamental domain has the same finite Haar measure.
\item[(ii)] Let $2n<M$, $T\in\mathrm{Herm}_n(F)$ nondegenerate and
$\Phi=\Phi_\infty\otimes\bigotimes_v\Phi_v$ with $\Phi_\infty$ smooth and
compactly supported, $\Phi_v$ locally constant and compactly supported, and
$\Phi_v$ the characteristic function of the standard lattice for almost all
$v$. Give $V_n(F_v)$, $\mathrm{Herm}_n(F_v)$ (and the archimedean analogues)
the measures self-dual for a nontrivial character $\psi$ of $\mathbb
A_F/F$. Then the local densities $g_v(T)$ (values at $T$ of continuous
densities of the push-forward of $\Phi_v\,dX$ under $\mathrm{gr}$) exist,
$\prod_vg_v(T)$ converges, and for every Haar measure $\mu$ and fundamental
domain $\mathcal F$,
\[
\frac1{\mu(\mathcal F)}\int_{\mathcal F}\theta_T\,d\mu
 =g_\infty(T)\prod_{v<\infty}g_v(T).
\]
\end{enumerate}
\end{theorem}

\noindent \emph{Range and hypotheses.} Only $2n<M$ is used; then for every
Witt index $r\le M/2$ of $J$ one has $n+r<M$ and $M>2n$, which is Weil's
range of absolute convergence, where the Siegel--Weil identity holds without
regularisation (\cite[Theorem~\ref{A-i4:thm-sw}]{Annex}); (ii) is its $T$-th
Fourier coefficient at the identity, with the local Siegel integrals
identified with the densities by Fourier inversion. For (i), $U_J$ is
connected reductive without nontrivial $F$-rational characters (they factor
through the anisotropic norm-one torus). The statement (SW) of
\cite[Definition~\ref{A-i3:S-def-SW}]{Annex} is exactly
Theorem~\ref{inp:thm-SW}; it is also proved in the formal
companion \cite{CharacterFamily} by Siegel's inductive argument (rank one by
invariance and dilation, then orthogonal splitting), and in the main line of
\cite{CharacterFamily} the input (SW) is supplied by that proof.

\emph{Proof outline.}
\begin{enumerate}
\item[(a)] \emph{Theorem S off the diagonal}
(\cite[\S\S\ref{A-i3:ssec-bessel}--\ref{A-i3:ssec-S}]{Annex}). The archimedean
input is multiplicity one for Bessel models of discrete series of
$U(k+r,r+1)$ whose Harish-Chandra parameter is far from the walls
(\cite[Proposition~\ref{A-i3:prop-bessel}]{Annex}); the required spectral gap
is supplied by the archimedean weights of the family
(\cite[\S\ref{A-i3:ssec-gap}]{Annex}). Given multiplicity one, an orbital
integral of a matrix coefficient is a fixed multiple of the local Bessel
integral (\cite[Proposition~\ref{A-i3:prop-amp}]{Annex}), with the formal
degree from Schur orthogonality (Godement). The global off-diagonal bound
is then an abstract pair bound for positive kernels
(\cite[Theorems~\ref{A-i3:S-thm-abstract} and~\ref{A-i3:S-thm-pairbound}]{Annex}),
a mass formula (\cite[Theorem~\ref{A-i3:S-thm-massformula}]{Annex}), a
lattice-average recursion (\cite[Theorem~\ref{A-i3:S-thm-recursion}]{Annex})
and the verification of nineteen components
(\cite[Theorem~\ref{A-i3:S-thm-components}]{Annex}), which give Theorem S at
a pair (\cite[Theorem~\ref{A-i3:S-thm-pair}]{Annex}). What this step takes
on trust is listed in \cite[Remark~\ref{A-i3:S-rem-trust}]{Annex}.
\item[(b)] \emph{From Theorem S to $\mathrm{LinkS}$}
(\cite[\S\ref{A-i3:ssec-linkS}]{Annex}). The diagonal and Kloosterman
parts are bounded separately (\cite[Propositions~\ref{A-i3:prop-kpp-off}
and~\ref{A-i3:prop-kpp-diag}]{Annex}); the conclusion
(\cite[Theorems~\ref{A-i3:thm-linkS-standing} and~\ref{A-i3:thm-linkS-slots}]{Annex})
follows once the standing hypotheses hold along the family. The slot
hypotheses (Sel), (Rk), (Lp) are discharged from (SW) in
\cite[Proposition~\ref{A-i3:prop-slots-SW}]{Annex}.
\item[(c)] \emph{Irreducibility} (\cite[\S\ref{A-i3:ssec-irr}]{Annex}). The
defect $\delta$ is bounded by a share estimate for the weights of the
level-character idempotents (\cite[Proposition~\ref{A-i3:prop-W}]{Annex}),
using Green's parametrisation of the cuspidal representations of
$\mathrm{GL}_r(\mathbb F_g)$; the conclusion is
\cite[Theorem~\ref{A-i3:thm-linkIRR}]{Annex}.
\item[(d)] \emph{Level counts} (\cite[\S\ref{A-i3:ssec-count}]{Annex}). The
levels $T_u$ are primes in a window determined by the dimensions $d_u$
(\cite[Lemma~\ref{A-i3:lem-power}]{Annex}); such a prime exists by a prime
number theorem in short intervals for ray classes of $\Q(\sqrt{-d})$
(\cite[Lemma~\ref{A-i3:lem-levelprime} and
Theorem~\ref{A-i3:thm-shortprimes}]{Annex}), which is how the construction
fulfils this clause of the datum. The dimensions are
computed by a mass formula for a standard hermitian lattice
(\cite[Proposition~\ref{A-i3:prop-massratio}]{Annex}), the formal degree of
the discrete series (\cite[Proposition~\ref{A-i3:thm-HC-fdeg}]{Annex}), and
Landherr's theorem; the bounds follow in
\cite[Propositions~\ref{A-i3:prop-V}, \ref{A-i3:prop-part5}
and~\ref{A-i3:prop-countF}]{Annex}.
\end{enumerate}

Two statements show the shape of the argument. The first is the
elementary step that turns the level window into the dimension window used
in $\mathrm{LinkS}$.

\begin{lemma}[{\cite[Lemma~\ref{A-i3:lem-power}]{Annex}}]\label{inp:lem-power}
Let $f,k\ge1$ with $2f\le2^k$, $T\ge2$, $X>0$, $\Gamma\ge1$, $D_*,d^*>0$ with
$X^{fk}=D_*/d^*$, and $Y\le T\le Y(1+\frac1{fk})$ with $Y=X\Gamma^{-2/(fk)}$.
Then $\frac12D_*\Gamma^{-2}\le d^*(T^k-1)^f\le eD_*\Gamma^{-2}$,
$\log T\le\log X+\log2$ and $\log X-2\log\Gamma/(fk)\le\log T$.
\end{lemma}

\noindent The second is Theorem S at a single pair of words, stated in
the notation of \cite[\S\ref{A-i3:ssec-S}]{Annex}, where the size, standing,
gap, local tying and frame hypotheses, the loops
(RB$_{\varepsilon_L,d}$), (RS$_{\varepsilon_L,d}$), the hypothesis (FD) of
\cite[Lemma~\ref{A-i3:S-lem-FD}]{Annex} and the operator $K_0$ are defined.

\begin{theorem}[{\cite[Theorem~\ref{A-i3:S-thm-pair}]{Annex}}; Theorem S at a pair]\label{inp:thm-Spair}
Assume (SW), (FD) and the conclusion of \cite[Proposition~\ref{A-i3:S-prop-NG}]{Annex}.  For every $C_{\rm thr}$ there are constants
such that, for admissible size constants, there are $k_0$, $C_n\ge0$, $C'_0$, $C_{l,0}$ with: for $C'\ge C'_0$,
$C_l\ge C_{l,0}$, every datum $\mathcal D$, package $(X,Y)$ and words $u\ne u'$ satisfying the size hypotheses with
$k\ge k_0$, the standing hypotheses at $(C',C_l)$, the gap hypothesis, local tying, the frame condition and the
loops (RB$_{\varepsilon_L,d}$), (RS$_{\varepsilon_L,d}$) with $d=d_ud_{u'}/D_*$,
\[
\|K_0(u,u')\|^2\le(1+b^{-6})\frac{d_ud_{u'}}{D_*}\quad\text{and}\quad\|K(u,u')\|_{HS}^2\le(1+b^{-6})\frac{d_ud_{u'}}{D_*}.
\]
\end{theorem}

\noindent The bound $1+b^{-6}$ with $b=e^{x^\varepsilon}$ is what makes
$\eta_0\le\exp(-x^{\varepsilon/8})$ in $\mathrm{LinkS}$ attainable. Following \cite[\S\ref{A-i3:S-ssec-proof}]{Annex}, the proof
has four parts: (1) an exact Hilbert-space computation bounds the norm of an
operator written as an absolutely convergent sum of terms organised by
rows, through four row estimates (the main row, the loops, the
non-generic row and a Gram bound); (2) the main row is bounded by a
recursion on lattice averages over the levels, followed by the mass
formula, whose inputs come from Siegel--Weil; (3) the term system of a block
$(u,u')$ is constructed from the automorphic kernel
$\sum_\gamma\phi(x^{-1}\gamma y)$, and nineteen structural components say
that it has the properties required in (1) and (2); (4) the squared norm of
the term operator $K_0(u,u')$ equals $\|K(u,u')\|_{HS}^2$ off the diagonal.
Bessel uniqueness under the gap enters only through four of the
components.

\emph{What it supplies.} Part (d) is condition (P4) of
Definition~\ref{def:package} (Definition~\ref{inp:count}). Parts (b) and
(c) are the links $(\mathrm{LS}(3))$ and $(\mathrm{LIrr})$ of Section E
(\cite[Proposition~\ref{A-i4:hyp-S}]{Annex}); through the frame data they
control the small parameters $\eta$, $e$, $\delta$ that enter (P8) and
(P10).

\subsection{The windowed Selmer bound}\label{inp:ssec-selmer}

\emph{Setting.} For a number field $F$ and a prime $l$, an $l$-datum $D$
fixes valuation rings above the places $S$ of $F$ over $l$, and
\[
 \mathrm{Sha}^1_D(F,\Q_l/\Z_l(-e))=\ker\Bigl(H^1(G_{F,S},\Q_l/\Z_l(-e))\to\prod_{v\in S}H^1(D_v,\Q_l/\Z_l(-e))\Bigr),
\]
the Selmer group with the strict condition above $l$ and the unramified
condition elsewhere; $\operatorname{sha}_D(e)$ is its $\Z_l$-length
\cite[Definition~\ref{A-i5:def-ldata}]{Annex}. The fields in question are
the CM abelian fields attached to a Dirichlet character $\psi$ and a pair of
large primes $(q,d)$. For the trivial character, which is the case used,
\[
 E=\Q(\zeta_{2^{e+2}})^+\bigl(\sqrt{-q},\sqrt{-d}\bigr),\qquad [E:\Q]=2^{e+2}
\]
(\cite[Remark~\ref{A-i5:rem-trivial}]{Annex}). We copy the definition.

\begin{definition}[{\cite[Definition~\ref{A-i5:def-vwin}]{Annex}}; the windowed bound]\label{inp:def-vwin}
(1) For a number field $E$, $t\in\N$ and $B\in\R$, $\mathrm{VV}(E,t,B)$ means:
for every algebraic closure of $E$ and every family $(D_p)_p$ of $p$-data,
$\operatorname{sha}_{D_p}(e)<\infty$ for all primes $p$ and $0<e<t$, and for
every finite set $\mathcal S$ of primes
$\sum_{l\in\mathcal S}\log l\sum_{e=1}^{t-1}\operatorname{sha}_{D_l}(e)^\natural\le B$,
where $a^\natural=a$ if $a<\infty$ and $a^\natural=0$ otherwise.

(2) Let $\psi$ be a primitive Dirichlet character of conductor $N_\psi$ and
order $o_\psi$.  For $e,q,d\in\N$, a number field $E$ is a
\emph{$(\psi,e,q,d)$-field} if there are $N$ with $2^{e+2},N_\psi,q,d\mid N$,
an embedding $j\colon E\hookrightarrow\Q(\zeta_N)$ and a primitive $N$-th root
of unity $\zeta$ such that $\sigma_a\colon\zeta\mapsto\zeta^a$ fixes $j(E)$ if
and only if $a\equiv\pm1\bmod 2^{e+2}$, $\psi(a)=1$ and
$(\frac aq)=(\frac ad)=1$.

(3) For $\varepsilon,Y>0$, $(q,d)$ is a \emph{window pair}
$(q,d)\in\mathrm{WP}(\varepsilon,Y)$ if $q\equiv3\bmod8$ and $d\equiv7\bmod16$
are primes, $\tfrac12Y^{\varepsilon/5}<\log q\le Y^{\varepsilon/5}$,
$\tfrac12Y^{\varepsilon/5}(\log Y)^2<\log d\le Y^{\varepsilon/5}(\log Y)^2$,
$q\mid d+1$, and every odd prime $p\le Y^{\varepsilon/5}$ divides $d+1$.  An
integer $e$ is \emph{on page $x$}, $\mathrm{Page}_\varepsilon(x,e)$, if
$x^{2+8\varepsilon}\le2^e\le2x^{2+8\varepsilon}$.

(4) $\mathrm{VWinAt}(\psi,C_t)$: for all $0<\varepsilon<1/52$ and every
$e\colon\R\to\N$ with $\mathrm{Page}_\varepsilon(x,e(x))$ for $x\ge1$, there
are $J\ge1$ and $X_0$ such that for all $X\ge X_0$ and every
$(q,d)\in\mathrm{WP}(\varepsilon,X^{1/J})$ there is $Y$ with $X^{1/J}\le Y$,
$Y^2\le X$, such that for every $x\in[Y,Y^2]$ and every
$(\psi,e(x),q,d)$-field $E$,
$\mathrm{VV}\bigl(E,\lfloor x^{C_t}\rfloor,\sqrt{\exp(x^\varepsilon)}\bigr)$.
$\mathrm{VWin}$ is the statement that $\mathrm{VWinAt}(\psi,401)$ holds for every
primitive $\psi$.
\end{definition}

\begin{theorem}[{\cite[Theorem~\ref{A-i5:thm-vwin}]{Annex}}; windowed Selmer bound]\label{inp:thm-vwin}
$\mathrm{VWin}$ holds, given the published theorems and the formally
verified statements listed in
\cite[Theorem~\ref{A-i5:thm-vwin}]{Annex} \textup{(}reproduced in
\S\ref{inp:ssec-status}\textup{)}. In words: for the CM field $E$ attached to
$(\psi,q,d)$ and a page $e(x)$,
\[
 \sum_l\log l\sum_{1\le e<x^{401}}\operatorname{length}\mathrm{Sha}^1(E,\Q_l/\Z_l(-e))\le\sqrt{\exp(x^\varepsilon)} .
\]
\end{theorem}

\emph{Proof outline.} The proof (\cite[\S\S\ref{A-i5:ssec-assembly}--\ref{A-i5:ssec-proof}]{Annex})
is a reduction, written in full, to a small number of propositions.
\begin{enumerate}
\item \emph{Splitting.} For odd $l$ complex conjugation splits
$\mathrm{Sha}^1$ into plus and minus parts
(\cite[Lemma~\ref{A-i5:lem-split}]{Annex}); the prime $l=2$ is treated
separately. The window statements reduce the uniform bound to bounds at one
field and one prime.
\item \emph{Plus part.} After reduction to characters over $\Q$
(\cite[\S\ref{A-i5:ssec-red}]{Annex}) the primes $l$ fall into the cases I,
Z, II, III, IV, V, Vs of \cite[Definition~\ref{A-i5:def-cases}]{Annex},
which exhaust all possibilities
(\cite[Lemma~\ref{A-i5:lem-hostarith}]{Annex}). A Poitou--Tate argument
(\cite[Proposition~\ref{A-i5:prop-shaleV}]{Annex}) compares $\mathrm{Sha}^1$
with $H^2(G_{F,S},\Z_l(e+1))$, which is bounded through Soul\'e's theorem
(\cite[Theorem~\ref{A-i5:thm-soule}]{Annex}) and the bound
$\mathrm{ShaK}$ (\cite[Proposition~\ref{A-i5:prop-shak}]{Annex}), proved by
descent along cyclotomic layers, a counting step and a torsion bound for
the homology of $\mathrm{GL}_N$ of rings of integers. Cases Z, IV and Vs
are elementary (\cite[Propositions~\ref{A-i5:prop-countZ}, \ref{A-i5:prop-countIV},
\ref{A-i5:prop-countVs}]{Annex}). Cases I, II, III and V are the
\emph{automorphic counts}: an Eisenstein congruence for $\mathrm U(3)$ in
Case I (\cite[Theorem~\ref{A-i5:prop-corg}]{Annex} and
\cite[Proposition~\ref{A-i5:prop-countI}]{Annex}), a $\mathrm{GSp}_4$
argument in Cases II, III (\cite[Proposition~\ref{A-i5:prop-c23}]{Annex}), and
a $\mathrm U(3)$ lifting argument over $\Q(i)$ in Case V
(\cite[Proposition~\ref{A-i5:prop-c5}]{Annex}). Each is reduced in writing to
automorphy-lifting cores (\cite[\S\S\ref{A-i5:x5-ssec}, \ref{A-i5:x5-ssec-arch}]{Annex})
and to a count of admissible Galois representations
(\cite[Proposition~\ref{A-i5:x-prop:u3count}]{Annex}). In Case I Kato's
theorem (\cite[Theorem~\ref{A-i5:x-thm:kato}]{Annex}) enters.
\item \emph{Minus part.} An Euler-system bound with Stickelberger's
theorem (\cite[Theorem~\ref{A-i5:x-thm:thmbc}]{Annex},
\cite[Propositions~\ref{A-i5:prop-psibound} and~\ref{A-i5:prop-stick}]{Annex}),
proved in full.
\item \emph{Exceptional primes.} At the finite set $B_\psi$ of exceptional
primes ($B_{\mathbf 1}=\{3,7\}$) an Iwasawa-theoretic bound
(\cite[Proposition~\ref{A-i5:prop-thmbab}]{Annex}) using the
Ferrero--Washington theorem (\cite[Theorem~\ref{A-i5:thm-fw}]{Annex}).
\end{enumerate}

We quote three of the intermediate statements. The first two are proved
in full in \cite{Annex}. Write $V_D(e)$ for the length of
$H^2(G_{F,S},\Z_l(e+1))$ and $\mathrm{Sha}^2_D(F,\Z_l(e+1))$ for the kernel of
its localisation at $S$ (\cite[Definition~\ref{A-i5:def-ldata}]{Annex}), and
$\operatorname{sha}^\pm_D(e)$ for the lengths of the $\pm(-1)^e$
eigenspaces of complex conjugation on $\mathrm{Sha}^1_D$.

\begin{lemma}[{\cite[Lemma~\ref{A-i5:lem-split}]{Annex}}]\label{inp:lem-split}
For $E$ CM and $l$ odd, $\operatorname{sha}_D(e)=\operatorname{sha}^+_D(e)+\operatorname{sha}^-_D(e)$.
Consequently, $\mathrm{VV}_2(E,t,B_2)$, $\mathrm{VV}^-(E,t,B_-)$ and
$\mathrm{VV}^+(E,t,B_+)$ imply $\mathrm{VV}(E,t,B_2\log2+B_-+B_+)$.
\end{lemma}

\begin{proposition}[{\cite[Proposition~\ref{A-i5:prop-shaleV}]{Annex}}; Poitou--Tate step]\label{inp:prop-shaleV}
If $l\ne2$ or $F$ has no real place, then for every $l$-datum $D$ and $e\in\Z$,
$\operatorname{length}_{\Z_l}\mathrm{Sha}^2_D(F,\Z_l(e+1))=\operatorname{sha}_D(e)$; hence
$\operatorname{sha}_D(e)\le V_D(e)$, and
$V_D(e)\le\operatorname{sha}_D(e)+\sum_{v\in S}\operatorname{length} H^0(D_v,\Q_l/\Z_l(-e))$.
\end{proposition}

\noindent The third is the Case I count, in the notation of
\cite[\S\ref{A-i5:ssec-counts}]{Annex}, where the condition
$\mathrm{GD}(K,l,e,\psi)$ (\cite[Definition~\ref{A-i5:def-gd}]{Annex}) and the
Selmer lengths $\mathrm{Vn}(e,\psi)$, $V'(e,\psi)$ are defined. Its proof
(\cite[\S\ref{A-i5:ssec-pf-counts}]{Annex}) bounds the Selmer length of
$\psi\chi^e$ by the $l$-adic size of the congruences between an Eisenstein
trace and the Frobenius traces of irreducible conjugate self-dual
automorphic three-dimensional representations, and it rests on the
formally verified \cite[Propositions~\ref{A-i5:x5-prop-U1},
\ref{A-i5:x5-prop-U234}, \ref{A-i5:x5-prop-U4}, \ref{A-i5:x-prop:u3count}]{Annex}.

\begin{theorem}[{\cite[Theorem~\ref{A-i5:prop-corg}]{Annex}}; $\mathrm{CorGK}$, the Case~I count]\label{inp:thm-corg}
For every imaginary quadratic field $K$ there are $c_5,A\in\N$ such that for
all $e,N\ge1$, $H\le(\Z/N)^\times$, distinct primes $p,p'\nmid30N|d_K|$,
coefficient rings $(\mathcal O_l)_l$ and finite sets $\mathcal L$ of primes each
of which admits a character $\psi\colon(\Z/N)^\times\to\mathcal O_l^\times$ with
$\ker\psi=H$ and $\mathrm{GD}(K,l,e,\psi)$: all the $V'(e,\psi)$
are finite and
\[
\sum_{l\in\mathcal L}\log l\sum_{\ker\psi=H}\mathrm{Vn}(e,\psi)\le2c_5\,\varphi([(\Z/N)^\times:H])\,(N(e+1))^A(\log p+\log p'),
\]
the inner sum being over the $\mathcal O_l$-valued characters of $(\Z/N)^\times$.
\end{theorem}

\emph{What it supplies.} $\mathrm{VWin}$ is not a condition on a single
package (\cite[Remark~\ref{A-i5:rem-role}]{Annex}). For $\psi=\mathbf 1$ it
is the hypothesis \cite[Proposition~\ref{A-i4:hyp-selmer}]{Annex},
$\mathrm{VWinAt}(1,401)$. It is used twice: (i) through the host supply
(\cite[Proposition~\ref{A-i4:prop-supply}]{Annex}:
$\mathrm{VWinAt}(1,401)$ implies $\mathrm{HostExW}(3/2,\mathcal G_1)$) it
produces hosts on windows $[W,W^{3/2}]$, which is clause (a) of
Theorem~\ref{thm:alpoge-families}; (ii) it gives $\mathrm{HypV}$ for the
members of the working class (\cite[Lemma~\ref{A-i4:lem-hypV-cls}]{Annex}),
a hypothesis of Theorem~\ref{inp:thm-div}, and so bounds the Selmer sum
in condition (P1).

\subsection{Existence of families}\label{inp:ssec-exist}

\emph{Setting.} Section E works with the trivial character $\psi=1$. A
\emph{host} at scale $x$ (\cite[Definition~\ref{A-i4:def-host}]{Annex}) is a
CM field $E$, Galois over $\Q$ and contained in $\Q(\zeta_{N_E})$ with
$N_E=Q_Sqd$, with $\mathrm{Gal}(E/\Q)\cong U=(\Z/2)^2\times A$, $A$ finite
abelian of order $n$ (in fact cyclic of $2$-power order), together with its
arithmetic data and level primes and ranks; the primes satisfy
$(-d/q)=1$ and $(-d/l)=1$ for every odd prime $l\mid Q_S$, and $E$
contains $\sqrt{-d}$ and $\sqrt{-q}$ (see Theorem~\ref{thm:alpoge-families}(b)
and Lemma~\ref{lem:host-facts}). A \emph{family} of exponent
$\varepsilon$ attaches to each large $x$ and index $i$ a member $D(x,i)$
with a host; a family is \emph{chain-tied} at $x$
(\cite[Definition~\ref{A-i4:def-chain}]{Annex}) if there are a base member
and a map from $A$ to the members, all sharing frame and level data, from
which data satisfying (1)--(5) of Theorem~\ref{thm:alpoge-families}(b) can
be read. A family class is \emph{connected} if for every small
$\varepsilon$ it contains a family that is chain-tied, with constants
independent of $x$, at every large $x$ of its window set. The four
inputs above enter as hypotheses in the form of
\cite[Propositions~\ref{A-i4:hyp-div}, \ref{A-i4:hyp-loss}, \ref{A-i4:hyp-S}
and~\ref{A-i4:hyp-selmer}]{Annex}; the family estimates they must satisfy
are the \emph{links} of \cite[\S\ref{A-i4:ssec-hyps}]{Annex}.

\begin{theorem}[{\cite[Theorem~\ref{A-i4:thm-top}]{Annex}}; uniform existence of connected families]\label{inp:thm-top}
Assume the loss bound, the divisor bound, and the level counts and
Theorem~S, and the statements listed in
\cite[Remark~\ref{A-i4:rem-trust}]{Annex}. Let
$\lambda>1$, and let $\mathcal G$ be a requirement with
$\mathrm{HostExW}(\lambda,\mathcal G)$ all of whose class members satisfy
$\mathrm{HypV}$. Then $\mathcal G.\mathrm{cls}(\lambda)$ is connected, with
$e_D=2$ and $e_L=0$.
\end{theorem}

\begin{theorem}[{\cite[Theorem~\ref{A-i4:thm-final}]{Annex}}]\label{inp:thm-final}
Assume the loss bound \textup{(}\cite[Proposition~\ref{A-i4:hyp-loss}]{Annex}\textup{)}, the divisor bound
\textup{(}\cite[Proposition~\ref{A-i4:hyp-div}]{Annex}\textup{)}, the level counts and Theorem~S
\textup{(}\cite[Proposition~\ref{A-i4:hyp-S}]{Annex}\textup{)}, and the windowed Selmer bound
\textup{(}\cite[Proposition~\ref{A-i4:hyp-selmer}]{Annex}\textup{)}, and the statements listed in
\cite[Remark~\ref{A-i4:rem-trust}]{Annex}. Then Theorem~\ref{thm:alpoge-families}
holds.
\end{theorem}

The central definition, copied from \cite{Annex}, is the following.

\begin{definition}[{\cite[Definition~\ref{A-i4:def-chain}]{Annex}}; chain-tied, connected]\label{inp:def-chain}
Let $\mathcal F$ have exponent $\varepsilon$, and let
$e_D$, $e_L$, $K$, $c_{HH}$, $e_K$, $C_L$ and $x$ be real numbers. The family is \emph{chain-tied at $x$}
with these constants if there are a base index $i_0\in I(x)$ and a map
$\mathrm{pt}:A\to I(x)$ such that the following holds. Here $A$ is the
group of the host $H$ of $D(x,i_0)$. The columns at the translation
$v\in A$ are those of $D(x,\mathrm{pt}(v))$, which shares its frame and
level data with $D(x,i_0)$. From these there are data $q,d$, $T,S\subset A$
of sizes $r,s$, $m\ge0$, $N$, $(a_t)_{t\in T}$, $w_{v,j}$ for $v\in A$ and
$1\le j\le m$, and $a,b,b',k,\eta,X',\vartheta$, which satisfy properties
(1)--(5) of Theorem~\ref{thm:alpoge-families}(b) for $H$. The column
package in property (5) is read from $H$ along the identity of $A$, and its
constants satisfy $e_D^{\rm pkg}\le e_D$, $e_L^{\rm pkg}\le e_L$ and
$c_{HH}^{\rm pkg}\le c_{HH}$. A family class $\mathcal E$ is
\emph{connected} if the following holds. There are reals $e_D,e_L$ and
$\varepsilon_1>0$ such that for every $\varepsilon\in(0,\varepsilon_1)$
there are a family $\mathcal F\in\mathcal E$ of exponent $\varepsilon$ and
constants $K,c_{HH},e_K,C_L,x_0$ for which $\mathcal F$ is chain-tied at
every $x\in\mathrm{Win}(\mathcal F)$ with $x\ge x_0$.
\end{definition}

\noindent Thus Theorem~\ref{thm:alpoge-families} for the trivial character
is, up to the choice $\varepsilon_1\le1/52$ and the windows of ratio $3/2$,
the statement that one particular family class is connected. The engine
that produces connectedness is the following theorem, proved in full; its
inputs are named links, listed in \cite[\S\ref{A-i4:ssec-hyps}]{Annex}:
$(\mathrm{LN})$ is $\mathrm{LinkN}$, $(\mathrm{LEta})$ twist coverage,
$(\mathrm{LIrr})$ the irreducibility bound, $(\mathrm{SLR})$ the rank
equality, $(\mathrm{LG})$ the comparison $\mathrm{LinkG}$,
$(\mathrm{LDivt})$ the implication
$(\mathrm{LEta})\wedge(\mathrm{DivN})\Rightarrow(\mathrm{LDiv})$ to the
division bound $(\mathrm{LDiv})$, $(\mathrm{LS}(C_0))$ almost
orthogonality and $(\mathrm{LGl})$ the global link; $\mathrm{Des}$,
$\mathrm{WithTw}$, $\mathrm{WithSm}$ are the class operations of
\cite[\S\ref{A-i4:ssec-desWin}]{Annex}.

\begin{theorem}[{\cite[Theorem~\ref{A-i4:thm-desWin}]{Annex}}; connectedness from design existence]\label{inp:thm-desWin}
Let $\mathcal C,\mathcal C',\mathcal E$ be family classes and let
$C_{\rm abs},C_0,e_D,e_L\in\R$. Assume the divisor bound
\textup{(}\cite[Proposition~\ref{A-i4:hyp-div}]{Annex}\textup{)} and the following:
\begin{enumerate}
\item designs exist for $(\mathrm{WithTw}(\mathrm{WithSm}(\mathcal C')),C_0)$;
\item $\mathcal C'\subset\mathcal C\cap\mathcal E$;
\item every member of $\mathcal E$ satisfies $\mathrm{HypV}$;
\item $\mathrm{Des}(\mathcal C,C_0)$ is closed for each of the following
nine properties:
\begin{itemize}
\item $(\mathrm{LN})$; $(\mathrm{LEta})$; $(\mathrm{LIrr})$;
$(\mathrm{SLR})$; $(\mathrm{LG})$;
\item $\mathrm{TuSize}\wedge\mathrm{DGamSize}\Rightarrow(\mathrm{LDivt})$;
\item $\mathrm{ChosenE}(C_0,C_H)\Rightarrow$ (every $\nu_u$ is eventually
full and of positive rank);
\item $\mathrm{TuSize}\Rightarrow(\mathrm{LS}(C_0))$;
\item $\mathrm{TuSize}\wedge\mathrm{DGamSize}\Rightarrow(\mathrm{LGl})(C_{\rm abs},C_H,C_0,e_D,e_L)$.
\end{itemize}
\end{enumerate}
Then $\mathcal E$ is connected, with the exponents $e_D,e_L$.
\end{theorem}

In \cite[Theorem~\ref{A-i4:thm-final}]{Annex} the four inputs are used in
the following form; the first two parts of the list below are copied from
\cite[Propositions~\ref{A-i4:hyp-S}, \ref{A-i4:hyp-selmer}]{Annex}, and the
divisor and loss bounds are used as in
Theorems~\ref{inp:thm-div}, \ref{inp:thm-loss} and
Proposition~\ref{inp:prop-linkG}. Here $D_3(\mathcal A)$ is the working
class of \cite[\S\ref{A-i4:ssec-links}]{Annex}, and ``closed'' is in the
sense of \cite[Definition~\ref{A-i4:def-closed}]{Annex}.

\begin{proposition}[{\cite[Proposition~\ref{A-i4:hyp-S}]{Annex}}; the level counts and Theorem S, as a hypothesis]\label{inp:prop-hypS}
On the working class $D_3(\mathcal A)$:
\begin{enumerate}
\item[(a)] the implication $\mathrm{TuSize}\Rightarrow(\mathrm{LS}(3))$ is
closed;
\item[(b)] $(\mathrm{LIrr})$ is closed;
\item[(c)] for every member, condition (P4) holds at every point beyond a
threshold, with a constant $c_{\rm count}$.
\end{enumerate}
\end{proposition}

\begin{proposition}[{\cite[Proposition~\ref{A-i4:hyp-selmer}]{Annex}}; the windowed Selmer bound, as a hypothesis]\label{inp:prop-hypSel}
$\mathrm{VWinAt}(1,401)$ holds. Explicitly: for all
$0<\varepsilon<1/52$ and every $e:\R\to\N$ with
$x^{2+8\varepsilon}\le2^{e(x)}\le2x^{2+8\varepsilon}$, there are $J\ge1$
and $X_0$ with the following property. For all $X\ge X_0$ and every window
pair $(q,d)$ for $(\varepsilon,X^{1/J})$ there is $Y\ge X^{1/J}$ with
$Y^2\le X$ such that for every $x\in[Y,Y^2]$ and every $(1,e(x),q,d)$-field
$E$, $\mathrm{VV}(E,\lfloor x^{401}\rfloor,\sqrt{\exp(x^\varepsilon)})$
holds.
\end{proposition}

\begin{proposition}[{\cite[Proposition~\ref{A-i4:prop-supply}]{Annex}}; host supply]\label{inp:prop-supply}
$\mathrm{VWinAt}(1,401)$ implies $\mathrm{HostExW}(3/2,\mathcal G_1)$.
\end{proposition}

\noindent Here $\mathcal G_1$ is the requirement for the trivial character
(\cite[Definition~\ref{A-i4:def-G1}]{Annex}) and $\mathrm{HostExW}(\lambda,\mathcal G_1)$
says that for every $\varepsilon\in(0,1/52)$ there are constants, a
threshold $x_0$ and a set of scales whose $x$-coordinates contain intervals
$[W,W^\lambda]$ with $W$ arbitrarily large, such that every such scale with
$x\ge x_0$ carries valid ranks, an auxiliary prime $q_0$ and a good host
satisfying $\mathcal G_1$ (\cite[Definition~\ref{A-i4:def-G1}]{Annex}). The proof (\cite[\S\ref{A-i4:ssec-proof-supply}]{Annex}) is
an assembly of five formally verified constructions: the field of the
character, the window data, valid ranks, the host, and the twist bounds
(\cite[Propositions~\ref{A-i4:Sprop-field}--\ref{A-i4:Sprop-tw}]{Annex}). The
last of these gives, for the three nontrivial characters $\chi$ of $U$
trivial on $A$, $|G_\chi(x)|\le x^{1/2+\varepsilon}$ and
$|L'/L(1,\chi)|\le4\log x$.

The conclusion of the chain is carried by the following formally verified
statement, quoted exactly from \cite[Proposition~\ref{A-i4:Lprop-fam}(c)]{Annex}
in the notation of \cite[\S\ref{A-i4:ssec-links}]{Annex}.

\begin{proposition}[{\cite[Proposition~\ref{A-i4:Lprop-fam}(c)]{Annex}}; the core link; formally verified \cite{CharacterFamily}]\label{inp:prop-corelink}
Let $\mathcal F$ satisfy $\mathrm{TuSize}$ and the minimal global record.
Let $C_{\rm abs}\ge0$, $C_H>0$, $C_5$ and
$c_6>0$. If the family satisfies the count, the two slope conclusions
\textup{(}$\mathrm{SlopesM}$, $\mathrm{SlopesN}$\textup{)}, the lattice condition, the
nodes, $\mathrm{FrameF}(C_H)$, the losses $\mathrm{LossFS}(\cdot,0)$, the
character bound and the windows, and moreover
$\mathrm{ChosenTT}$, $\mathrm{ChosenQD}(C_5)$, $\mathrm{ChosenThin}(c_6)$,
$(\mathrm{LIrr})$ and $(\mathrm{LDiv})$ with constants $(C_{\rm abs},C_H)$,
then there are $K,c_{HH},e_K,C_L,x_0$ such that the family is chain-tied
with $(C_{\rm abs},0,K,c_{HH},e_K,C_L)$ at every $x\ge x_0$ in
$\mathrm{Win}$.
\end{proposition}

\noindent The hypotheses of Proposition~\ref{inp:prop-corelink} are the
items listed in \cite[Proposition~\ref{A-i4:Lprop-fam}(b)]{Annex}, also
formally verified, together with the links. Of the links,
$(\mathrm{LIrr})$ is Theorem~\ref{inp:thm-counts}(c), the losses come from
Theorem~\ref{inp:thm-loss}, and $(\mathrm{LDiv})$ contains condition (P1)
(\cite[\S\ref{A-i4:ssec-hyps}]{Annex}, clause (7) of $(\mathrm{LDiv})$).

\emph{Proof outline.} The logical layer of Section E is proved in full.
\begin{enumerate}
\item \emph{Order of constants.} Constants are fixed in stages: outer
constants of a family, then middle constants with lower bounds in terms of
the outer ones, then thresholds in $x$
(\cite[\S\ref{A-i4:ssec-closed}]{Annex}). A property is \emph{closed} on a
class if it holds, beyond a threshold system, for every large enough
family of the class (\cite[Definition~\ref{A-i4:def-closed}]{Annex}).
\item \emph{Connectedness from designs}
(\cite[Theorem~\ref{A-i4:thm-desWin}]{Annex}). If designs exist on a class
and nine link properties are closed on it, then the class is connected.
\item \emph{Design existence} (\cite[Theorems~\ref{A-i4:thm-choices}
and~\ref{A-i4:thm-design}]{Annex}), from Minkowski's bound and the thin-scale
input (\cite[Proposition~\ref{A-i4:thm-litthin}]{Annex}), which is derived
from the classical zero estimates and explicit formula.
\item \emph{Closedness of the links.} Five of them follow from the
inputs (\cite[Proposition~\ref{A-i4:prop-links-proved}]{Annex}); the
division and global links are reduced to their leaves
(\cite[Propositions~\ref{A-i4:prop-links-formal}
and~\ref{A-i4:prop-links-leaves}]{Annex}), and the leaves are the formally
verified \cite[Propositions~\ref{A-i4:Lprop-sel}, \ref{A-i4:Lprop-D},
\ref{A-i4:Lprop-pkg}, \ref{A-i4:Lprop-forms}, \ref{A-i4:Lprop-fam}]{Annex}.
\item \emph{The arithmetic construction}
(\cite[\S\ref{A-i4:ssec-arith}]{Annex}). Hosts are realised
(\cite[Theorem~\ref{A-i4:thm-realise}]{Annex}) using Hecke's theorem, the
least prime in a progression, and norm primes in short intervals
(\cite[Proposition~\ref{A-i4:thm-shortnorm}]{Annex}), which is proved from
the ray class family (\cite[Proposition~\ref{A-i4:prop-lfam-ray}]{Annex})
and zero statements for it (\cite[Proposition~\ref{A-i4:prop-lfam-zeros}]{Annex}).
The bundles of automorphic origin are built through the theta
correspondence for unitary groups, the Siegel--Weil formula and the mass
of a standard lattice (\cite[Theorem~\ref{A-i4:thm-genusmass}]{Annex}); the
$26$ rows they must satisfy are reduced to their leaves in
\cite[Theorem~\ref{A-i4:A-thm-allrows}]{Annex} and
\cite[Theorem~\ref{A-i4:thm-assembly}]{Annex}. The class selection and the
hypothesis structure are \cite[Theorem~\ref{A-i4:thm-classSel}]{Annex} and
\cite[Corollary~\ref{A-i4:cor-structure}]{Annex}.
\item \emph{Conclusion.} \cite[Theorem~\ref{A-i4:thm-top}]{Annex} combines
(2)--(5). In \cite[Theorem~\ref{A-i4:thm-final}]{Annex}: the Selmer bound
and the host supply give $\mathrm{HostExW}(3/2,\mathcal G_1)$;
\cite[Lemma~\ref{A-i4:lem-hypV-cls}]{Annex} gives $\mathrm{HypV}$ on the
class; \cite[Theorem~\ref{A-i4:thm-top}]{Annex} gives connectedness with
$e_D=2$, $e_L=0$; one puts $\varepsilon_1=\min(\varepsilon_1',1/52)$; the
windows of ratio $3/2$ give clause (a); and clause (b) is the chain-tied
property at the base member.
\end{enumerate}

\emph{What it supplies.} Theorem~\ref{thm:alpoge-families}, with $e_D=2$
and $e_L=0$. As stated in the proof of
\cite[Theorem~\ref{A-i4:thm-final}]{Annex}, clause (b) is obtained by
definition, and the passage from the family estimates to properties
(1)--(5) lies entirely in the core link
\cite[Proposition~\ref{A-i4:Lprop-fam}(c)]{Annex}, which is formally
verified; its hypotheses come from the other inputs and from
\cite[Theorem~\ref{A-i4:thm-top}]{Annex}.

\subsection{Status of the proofs}\label{inp:ssec-status}

We list, as precisely as \cite{Annex} allows, what the derivation of
Theorem~\ref{thm:alpoge-families} rests on. The lists follow the
statements and the trust remarks of \cite{Annex}
(\cite[Remarks~\ref{A-i4:rem-trust}, \ref{A-i3:S-rem-trust}]{Annex}, the
paragraph ``What is taken on trust'' at the end of
\cite[\S\ref{A-i2:ssec-P1}]{Annex}, and the dependency list of
\cite[Theorem~\ref{A-i5:thm-vwin}]{Annex}). Where these sources disagree
with each other, we list a statement under (c) if any of them treats it as
formally verified.

The formal development \cite{CharacterFamily} has a final theorem without
hypotheses. The published statements on which its argument is organised
(local representation theory, global automorphic forms and Galois
representations, Galois cohomology, base change, Siegel--Weil, complex
multiplication, and others) are isolated there as separate propositions and
each is proved inside the development; intermediate formal lemmas take
bundles of them as arguments, which is why they read as conditional, but
these bundles are discharged at the top. The published theorems listed in
(b) are therefore citations for the written proofs, not assumptions of the
formal theorem. The correspondence between a formal proposition and the
published theorem it is named after was read from the statements and has
not been audited item by item; the validity of the formal theorem does not
depend on these attributions.

\subsubsection*{(a) Proved in full in \cite{Annex}}
\begin{itemize}
\item \emph{Section A.} The elementary loss bound
(\cite[Theorem~\ref{A-i1:thm-elementary}]{Annex}), the instantiation
(\cite[Theorem~\ref{A-i1:thm-setting}]{Annex}), the small-share weights
(\cite[Proposition~\ref{A-i1:prop-chosenW}]{Annex}), the main theorem
\cite[Theorem~\ref{A-i1:thm-loss}]{Annex} from $\mathrm{LinkG}$; the
reduction of $\mathrm{LinkG}$ (\cite[Proposition~\ref{A-i1:prop-linkG}]{Annex})
and of the local statement (\cite[Proposition~\ref{A-i1:prop-loc}]{Annex}) to
the propositions listed in (c); the $37$ admission clauses
(\cite[Proposition~\ref{A-i1:prop-adm}]{Annex}) except the counting part of
$\mathcal A_{27}$; \cite[Propositions~\ref{A-i1:prop-Fm}, \ref{A-i1:prop-A19c},
\ref{A-i1:prop-A21}]{Annex}.
\item \emph{Section B.} The reduction to $\mathrm{LinkN}$ and
\cite[Theorem~\ref{A-i2:thm-div}]{Annex}; the summation over primes
(\cite[Theorems~\ref{A-i2:thm-summed}, \ref{A-i2:thm-export}]{Annex}); the cut
transfer (\cite[Theorem~\ref{A-i2:thm-transfer}]{Annex}) and parts (1), (2)
of \cite[Proposition~\ref{A-i2:prop-cut}]{Annex}; the generator-prime
summation, the lattice bound and the $\mathrm{Hom}$-term algebra; the
assembly of \cite[Theorem~\ref{A-i2:thm-linkN}]{Annex}; the reduction of the
crystalline lifting theorem \cite[Proposition~\ref{A-i2:prop-lift}]{Annex} to
the Breuil--Kisin steps; the assembly of
\cite[Proposition~\ref{A-i2:prop-P1}]{Annex}.
\item \emph{Section C.} The one-level inequality, the typed count, the Gram
bound, the level volumes, the identification of $D_*$, the identity cell
and the transport (\cite[\S\S\ref{A-i3:S-ssec-onelevel}--\ref{A-i3:S-ssec-pf-transport}]{Annex});
the abstract pair bound, the mass formula, the recursion and the assembly
of the nineteen components; the passage from Theorem S to $\mathrm{LinkS}$
modulo the Kloosterman inputs; the irreducibility bound
(\cite[Theorem~\ref{A-i3:thm-linkIRR}]{Annex}); the level counts
(\cite[Propositions~\ref{A-i3:prop-part5}, \ref{A-i3:prop-countF}]{Annex})
from \cite[Proposition~\ref{A-i3:prop-V}]{Annex}, itself proved modulo the
items in (c).
\item \emph{Section D.} The whole reduction of
\cite[Theorem~\ref{A-i5:thm-vwin}]{Annex}: the splitting, the window
statements, the reduction to characters over $\Q$, the case analysis, the
Poitou--Tate step, $\mathrm{ShaK}$ (with the counting step, the torsion
bound, the detection-and-construction statement
\cite[Proposition~\ref{A-i5:x5-prop-detect-construct}]{Annex} and the rank
bound \cite[Proposition~\ref{A-i5:x5-prop-bodd}]{Annex}), the counts in Cases
Z, IV, Vs, the minus part, the exceptional primes, the reductions of the
counts in Cases I, II, III, V to their cores, and the proofs of
\cite[Propositions~\ref{A-i5:x5-prop-U1} and~\ref{A-i5:x5-prop-Vcore}]{Annex}
from the smaller formally verified statements listed in (c).
\item \emph{Section E.} The logical layer
(\cite[Theorems~\ref{A-i4:thm-desWin}, \ref{A-i4:thm-top},
\ref{A-i4:thm-final}]{Annex}); design existence; the root-discriminant
estimate (\cite[Proposition~\ref{A-i4:prop-rd}]{Annex}); the kit parameters
(\cite[Proposition~\ref{A-i4:prop-kits}]{Annex}); the genus mass
(\cite[Theorem~\ref{A-i4:thm-genusmass}]{Annex}); the zero statements for
the ray class family (\cite[Proposition~\ref{A-i4:prop-lfam-zeros}]{Annex});
the reductions of the division and global links to their leaves;
realisation from its size leaves; the reduction of the $26$ rows and of the
open-origin rows to their leaves; the assembly of the host supply from the
five statements \cite[Propositions~\ref{A-i4:Sprop-field}--\ref{A-i4:Sprop-tw}]{Annex}.
\end{itemize}

\subsubsection*{(b) Published theorems cited}
References marked ``to locate'' or ``unconfirmed'' in \cite{Annex} are
still to be confirmed; we mark them with $^\dagger$.
\begin{itemize}
\item \emph{Section A.} Chebyshev's and Mertens' bounds (elementary loss
bound); quadratic reciprocity, Neukirch I (8.6)$^\dagger$
(\cite[Theorem~\ref{A-i1:thm-qr}]{Annex}; instantiation); the standard basis
of the dual Weyl module, cited in \cite{Annex} as J.~A.~Green,
D.~Eisenbud, C.~Procesi, Invent.\ Math.\ 56 (1980)$^\dagger$
(\cite[Theorem~\ref{A-i1:thm-stdbasis}]{Annex}; Weyl form); elementary
divisors, Bourbaki$^\dagger$ (\cite[Theorem~\ref{A-i1:thm-elemdiv}]{Annex});
Weyl's dimension formula$^\dagger$ (\cite[Theorem~\ref{A-i1:thm-weyl}]{Annex});
Kostant's commutation formula$^\dagger$
(\cite[Theorem~\ref{A-i1:thm-kostant}]{Annex}); basic facts on Witt
vectors$^\dagger$ (admission clauses).
\item \emph{Section B.} Neukirch I (8.2)
(\cite[Theorem~\ref{A-i2:thm-fund}]{Annex}; per-prime bound); Tate local
duality and the local Euler characteristic, NSW 7.2.6, 7.3.1$^\dagger$
(\cite[Theorems~\ref{A-i2:thm-tate}, \ref{A-i2:thm-euler}]{Annex}; cuts);
Bloch--Kato, local dimension and duality$^\dagger$
(\cite[Theorem~\ref{A-i2:thm-bk}]{Annex}); Chebotarev, Neukirch VII (13.4)
(\cite[Theorem~\ref{A-i2:thm-cheb}]{Annex}); global Artin
reciprocity$^\dagger$ (\cite[Theorem~\ref{A-i2:thm-artin}]{Annex}); the
integral crystalline comparison of Bhatt--Scholze, Thm.~1.2
(\cite[Theorem~\ref{A-i2:thm-bs}]{Annex}); Fontaine, Ast\'erisque 223$^\dagger$
(\cite[Theorem~\ref{A-i2:thm-fontaine}]{Annex}); Herr, Bull.\ Soc.\ Math.\
France 126$^\dagger$ (\cite[Theorem~\ref{A-i2:thm-herr}]{Annex}); Wach, and
Berger, Compositio Math.\ 140$^\dagger$ (\cite[Theorem~\ref{A-i2:thm-wach}]{Annex}). All are
used in the local inputs of the divisor bound.
\item \emph{Section C.} Harish-Chandra's discrete series
(\cite[Theorem~\ref{A-i3:thm-hc}]{Annex}); Schur orthogonality, Godement
(\cite[Theorem~\ref{A-i3:thm-godement}]{Annex}); the theorem of the highest
weight, Knapp (\cite[Theorem~\ref{A-i3:thm-knapp}]{Annex}); Landherr's
theorem (\cite[Theorem~\ref{A-i3:thm-landherr}]{Annex}; level counts);
integration in Cartan coordinates, Helgason
(\cite[Theorem~\ref{A-i3:thm-helgason}]{Annex}; Bessel uniqueness); the
inputs (F1)--(F3) of \cite[Lemma~\ref{A-i3:S-lem-FD}]{Annex} (Godement,
Borel, finite multiplicity in the discrete spectrum)$^\dagger$;
Gelfand--Graev multiplicity one$^\dagger$
(\cite[Lemma~\ref{A-i3:S-lem-GGmult}]{Annex}); Weyl's dimension formula;
and the hypotheses (SW) and (G), which are assumed.
\item \emph{Section D.} The local Langlands correspondence, Harris--Taylor
and Henniart$^\dagger$ (\cite[Theorem~\ref{A-i5:thm-llc}]{Annex});
Aubert--Baum--Plymen--Solleveld$^\dagger$
(\cite[Theorem~\ref{A-i5:thm-abps}]{Annex}); explicit level zero, Henniart and
Bushnell--Henniart$^\dagger$ (\cite[Theorem~\ref{A-i5:thm-lz}]{Annex});
local-global compatibility, Barnet-Lamb--Gee--Geraghty--Taylor with Shin,
Chenevier--Harris, Caraiani$^\dagger$
(\cite[Theorem~\ref{A-i5:thm-blggt}]{Annex}); Poitou--Tate and local
duality$^\dagger$ (\cite[Theorem~\ref{A-i5:thm-pt}]{Annex}); Soul\'e with
Borel$^\dagger$ (\cite[Theorem~\ref{A-i5:thm-soule}]{Annex}; $H^2$ bound);
Kato, Ast\'erisque 295, Thm.~14.2(1) (\cite[Theorem~\ref{A-i5:x-thm:kato}]{Annex};
Case I); Ferrero--Washington$^\dagger$
(\cite[Theorem~\ref{A-i5:thm-fw}]{Annex}; exceptional primes);
Stickelberger and Chebotarev (minus part); Quillen, Milnor--Moore and
Borel$^\dagger$ (rank bound); van der Kallen's homological stability with
Bass's stable-rank bound$^\dagger$ (torsion bound); and the classical
results of the counting steps (Dirichlet's unit theorem, finiteness of the
class number, Hilbert 90, Kronecker--Weber, Minkowski)$^\dagger$.
\item \emph{The $65$ published statements behind the count of admissible
representations.} The two counting statements $(\mathrm C_K)$,
$(\mathrm C_\Q)$ and the statement $(\mathrm N)$ behind
\cite[Proposition~\ref{A-i5:x-prop:u3count}]{Annex} are derived in
\cite{CharacterFamily} through the following known theorems, which the
formal development isolates as separate propositions and proves inside the
development; they are grouped by family as in \cite{Annex}; references for
each family are to be located, and it has not been checked whether every
one of them is used:
(H1) nine statements on automorphic representations (Clozel's
cohomological realisation and its compatibility with conjugation, Borel's
injectivity of cuspidal cohomology, strong multiplicity one, a
Gelfand--Kazhdan pairing, admissibility, Galois representations for totally
real $\mathrm{GL}_2$, quadratic base change, Galois representations of
newforms); (H2) eighteen statements of the local representation theory of
$\mathrm{GL}_n$ over $p$-adic fields and of finite groups of Lie type;
(H3) fifteen global statements (rapid decay, Harish-Chandra finiteness,
Jacquet--Shalika, the cuspidal spectrum, base change, local-global
compatibility, purity, Artin reciprocity, ray class fields, $l$-adic Hecke
characters, Weil--Deligne representations, Berger's theorem that de Rham
representations are potentially semistable, and others); (H4) twelve
statements of Galois cohomology (Poitou--Tate, finiteness, local duality,
Euler characteristics, the Greenberg--Wiles formula, Bloch--Kato); (H5)
eleven statements on CM abelian varieties and special values (including
Serre--Tate, the main theorem of complex multiplication, and
Siegel--Klingen--Deligne--Ribet). In total $9+18+15+12+11=65$.
\item \emph{Section E.} Stated there: Hecke
(\cite[Theorem~\ref{A-i4:thm-hecke}]{Annex}); the least prime in a
progression$^\dagger$ (\cite[Theorem~\ref{A-i4:thm-linnik}]{Annex}; host
realisation); Landherr and the Hasse norm theorem
(\cite[Theorem~\ref{A-i4:thm-landherr}]{Annex}); Green$^\dagger$
(\cite[Theorems~\ref{A-i4:thm-green}, \ref{A-i4:A-thm-green}]{Annex});
Bernstein--Zelevinsky (\cite[Theorems~\ref{A-i4:thm-bz},
\ref{A-i4:thm-bzfin}]{Annex}$^\dagger$); Serre
(\cite[Theorem~\ref{A-i4:thm-serre}]{Annex}); the two Siegel--Weil
formulas$^\dagger$ (\cite[Theorems~\ref{A-i4:thm-sw}, \ref{A-i4:thm-sw2}]{Annex});
Howe duality$^\dagger$ (\cite[Theorem~\ref{A-i4:thm-howe}]{Annex});
Minkowski's bound$^\dagger$ (\cite[Theorem~\ref{A-i4:thm-mink}]{Annex});
the classical zero estimates and truncated explicit formula$^\dagger$
(\cite[Theorem~\ref{A-i4:thm-classzeros}]{Annex}); Landau's lemma and the
truncated Perron formula$^\dagger$ (\cite[Lemmas~\ref{A-i4:lem-landau},
\ref{A-i4:lem-perron}]{Annex}); facts on cyclotomic fields and differents,
and non-vanishing of Hecke $L$-functions on $\Rel s\ge1$$^\dagger$. Not
stated in \cite{Annex}, isolated as separate propositions and proved inside
the formal development, all with references to locate: the published inputs of
\cite[Remark~\ref{A-i4:A-rem-pub}]{Annex} and the further list in
\cite[Remark~\ref{A-i4:rem-trust}(ii)]{Annex} (among them Clozel's
criterion, Galois representations attached to polarised cuspidal
representations with local-global compatibility, base change for definite
unitary groups, the doubling method, Jacquet--Shalika, the Weil
representation of unitary pairs, strong approximation, and the existence of
a local Langlands datum).
\end{itemize}

\subsubsection*{(c) Statements taken from the formal verification without a written proof}
The following statements are proved in \cite{CharacterFamily} and stated in
\cite{Annex}, as exactly as the definitions there allow, under the heading
``formally verified''; \cite{Annex} contains no written proof of them.
Several are marked in \cite{Annex} as paraphrases, not exact
transcriptions, because they use predicates of the formal development that
are named but not defined there; we mark these with $^\ast$. The whole of
Theorem~\ref{thm:alpoge-families} depends on this list.
\begin{itemize}
\item \emph{Automorphy-lifting cores} (Section D; none coincides exactly
with a published automorphy lifting theorem): the rigidity and type I
admissibility statements (U2), (U3)$^\ast$
(\cite[Proposition~\ref{A-i5:x5-prop-U234}]{Annex}); type II admissibility
and the length estimate (U3$'$), (U4)
(\cite[Proposition~\ref{A-i5:x5-prop-U4}]{Annex}); the $\mathrm{GSp}_4$
cores$^\ast$ (\cite[Proposition~\ref{A-i5:x5-prop-sp4core}]{Annex}); the
host inputs of Case V (\cite[Proposition~\ref{A-i5:x5-prop-Vhost}]{Annex});
irreducibility on open subgroups
(\cite[Lemma~\ref{A-i5:x5-lem-Vopen}]{Annex}); the cotangent estimates
(\cite[Proposition~\ref{A-i5:x5-prop-UM}]{Annex}); residually reducible
lifting at ordinary points (\cite[Proposition~\ref{A-i5:x5-prop-gdend}]{Annex});
\cite[Lemma~\ref{A-i5:x5-lem-U1steps}]{Annex}; host endomorphisms
(\cite[Proposition~\ref{A-i5:x5-prop-hostend}]{Annex});
\cite[Lemma~\ref{A-i5:x5-lem-Vsteps}]{Annex}; the algebraic Clifford lemma
quoted in the proof of \cite[Lemma~\ref{A-i5:x5-lem-irrcyc}]{Annex}; and the
count of admissible representations
(\cite[Proposition~\ref{A-i5:x-prop:u3count}]{Annex}), formally verified,
the $65$ published statements of (b) being themselves proved in the formal
development. The unitary lifting over $\Q(i)$
(\cite[Proposition~\ref{A-i5:x5-prop-Vcore}]{Annex}) and (U1)
(\cite[Proposition~\ref{A-i5:x5-prop-U1}]{Annex}) are proved in writing from
the items just listed.
\item \emph{Breuil--Kisin steps and the large-prime bound} (Section B):
packs, exactness, the class map, monodromy-kernel representatives, lifting
through exact triples, integral classes
(\cite[Propositions~\ref{A-i2:bk-b}, \ref{A-i2:bk-ex}, \ref{A-i2:bk-cl},
\ref{A-i2:bk-mon}, \ref{A-i2:bk-tri}, \ref{A-i2:bk-int}]{Annex}); fact (T4) of
\cite[Proposition~\ref{A-i2:prop-goodlat}]{Annex}, which is only sketched;
the large primes (\cite[Proposition~\ref{A-i2:prop-large}]{Annex}).
Further in Section B: the arithmetic of the column module$^\ast$
(\cite[Proposition~\ref{A-i2:prop-colconds}]{Annex}); point family, Selmer
signatures and thresholds$^\ast$ (\cite[Proposition~\ref{A-i2:prop-ptsel}]{Annex});
local admissibility data$^\ast$ (\cite[Proposition~\ref{A-i2:prop-locadm}]{Annex});
parts of the cut-cocycle bounds (\cite[Proposition~\ref{A-i2:prop-cutcoc}]{Annex}:
(a), the first assertion of (c), and (h)); the Galois representation of an
admissible list (\cite[Lemma~\ref{A-i2:lem-isorep}]{Annex}); parts of the
local inputs at a generator prime
(\cite[Proposition~\ref{A-i2:prop-pcclocal}]{Annex}); the division clauses
of a package and transport along ties
(\cite[Lemmas~\ref{A-i2:lem-divclauses}, \ref{A-i2:lem-tietransfer}]{Annex});
part (3) of \cite[Proposition~\ref{A-i2:prop-cut}]{Annex}, part (c) of
\cite[Proposition~\ref{A-i2:prop-potcount}]{Annex}, part (ii) of
\cite[Proposition~\ref{A-i2:prop-wordfacts}]{Annex}; and, as listed in the
trust paragraph of \cite[\S\ref{A-i2:ssec-P1}]{Annex},
\cite[Proposition~\ref{A-i2:prop-percut}]{Annex} and
\cite[Lemma~\ref{A-i2:lem-T3Z}]{Annex}.
\item \emph{The discrete-series model} (Section C):
\cite[Proposition~\ref{A-i3:prop-model}]{Annex}; with it the formal degree
under a gap (\cite[Proposition~\ref{A-i3:thm-HC-fdeg}]{Annex}), the local
data of the amplitude (\cite[Proposition~\ref{A-i3:prop-amp-in}]{Annex}),
the Gelfand--Kazhdan input (\cite[Proposition~\ref{A-i3:prop-GKin}]{Annex}),
the mass ratio (\cite[Proposition~\ref{A-i3:prop-massratio}]{Annex}),
\cite[Lemma~\ref{A-i3:lem-dataB}]{Annex}, the counting
(\cite[Proposition~\ref{A-i3:prop-counting}]{Annex}), the label weights
(\cite[Proposition~\ref{A-i3:prop-admLq-in}]{Annex}), the rows of the built
realisation (\cite[Proposition~\ref{A-i3:prop-allrows-built}]{Annex}), the
Kloosterman terms (\cite[Proposition~\ref{A-i3:prop-terms}]{Annex}), the
instance, the reduction $\mathrm{NG}$, the bridge facts and the leaves of
Theorem S (\cite[Propositions~\ref{A-i3:S-prop-instance}, \ref{A-i3:S-prop-NG},
\ref{A-i3:S-prop-bridge}, \ref{A-i3:S-prop-leaves}]{Annex}; the last is the
largest single item of Section C), the off-diagonal, covariance and
diagonal Kloosterman inputs (\cite[Propositions~\ref{A-i3:prop-kppoff-in},
\ref{A-i3:prop-cov-in}, \ref{A-i3:prop-kppdiag-in}]{Annex}), the local
estimates of the share bound (\cite[Proposition~\ref{A-i3:prop-W-loc}]{Annex}),
and the level character (\cite[Proposition~\ref{A-i3:prop-chi2}]{Annex}).
(\cite[Propositions~\ref{A-i3:S-prop-regions}, \ref{A-i3:S-prop-count}]{Annex}
are formally verified but not used for Theorem S.)
\item \emph{The extended closed form and the explicit minor map} (Section
A): these enter through the first theorem at the family
(\cite[Proposition~\ref{A-i1:prop-first}]{Annex}) and the remaining inputs
of $\mathrm{LinkG}$$^\ast$ (\cite[Proposition~\ref{A-i1:prop-linkG-res}]{Annex});
also the closed form (\cite[Proposition~\ref{A-i1:prop-closed}]{Annex}), the
counting part of $\mathcal A_{27}$
(\cite[Proposition~\ref{A-i1:prop-segdeg-rest}]{Annex}), and the
consequences of group compatibility$^\ast$
(\cite[Proposition~\ref{A-i1:prop-gc}]{Annex}). The written source of
\cite[Proposition~\ref{A-i1:prop-linkG-res}]{Annex} still has gaps, so no
complete written proof of it exists yet.
\item \emph{The automorphic row leaves} (Section E): the pointwise leaves
of the rows (\cite[Proposition~\ref{A-i4:A-prop-rowleaves}]{Annex}) and the
bundle-level leaves (\cite[Proposition~\ref{A-i4:A-prop-open}]{Annex}),
which are by far the largest items; the member inputs
(\cite[Proposition~\ref{A-i4:A-prop-inputs}(b)--(d)]{Annex}) and their
leaves (\cite[Proposition~\ref{A-i4:A-prop-inputs-leaves}]{Annex}); the
stages of the member statement
(\cite[Proposition~\ref{A-i4:A-prop-stages}]{Annex}); the Galois upgrade of
the norm row (\cite[Proposition~\ref{A-i4:A-prop-galois}]{Annex}); the
series input (\cite[Proposition~\ref{A-i4:A-prop-sern}]{Annex}); the
functional and covolume rows (\cite[Proposition~\ref{A-i4:A-prop-funcov}]{Annex});
the row bounds at $2$ and $3$ (\cite[Proposition~\ref{A-i4:A-prop-at23rows}]{Annex},
proved in \cite{Annex} modulo the split-case conclusion
\cite[Proposition~\ref{A-i4:A-prop-split}]{Annex}).
\item \emph{The division and global leaves} (Section E): the selection
clauses (\cite[Proposition~\ref{A-i4:Lprop-sel}(a)]{Annex}); the analytic
size estimates of the division side
(\cite[Proposition~\ref{A-i4:Lprop-D}]{Annex}, except the first claim of
(e)); the gap leaves (\cite[Proposition~\ref{A-i4:Lprop-gapleaves}]{Annex});
the division inputs at every word$^\ast$
(\cite[Proposition~\ref{A-i4:Lprop-pkg}]{Annex}); the level-zero leaves
(\cite[Proposition~\ref{A-i4:Lprop-forms}]{Annex}); and the family-level
items (\cite[Proposition~\ref{A-i4:Lprop-fam}]{Annex}). Part (c) of the last,
the core link, is the statement that the constructed families are
chain-tied, that is, that their data form column packages satisfying
(1)--(5) of Theorem~\ref{thm:alpoge-families}(b).
\item \emph{Smaller residues} (Section E): the five constructions of the
host supply (\cite[Propositions~\ref{A-i4:Sprop-field}, \ref{A-i4:Sprop-win},
\ref{A-i4:Sprop-ranks}, \ref{A-i4:Sprop-host}, \ref{A-i4:Sprop-tw}]{Annex}); the
leaves of the window count (\cite[Proposition~\ref{A-i4:Rprop-window-leaves}]{Annex},
in turn reduced in \cite{Annex} to \cite[Propositions~\ref{A-i4:Rprop-kappaI},
\ref{A-i4:Rprop-num}]{Annex}); the analytic and positivity leaves of the ray
class family (\cite[Proposition~\ref{A-i4:prop-ray-leaves}]{Annex}), from
which \cite[Proposition~\ref{A-i4:prop-lfam-ray}]{Annex} is proved.
\end{itemize}

We can now state the conclusion of the derivation.

\begin{theorem}\label{inp:thm-families}
Theorem~\ref{thm:alpoge-families} holds, under the published theorems and
the formally verified statements listed in \S\ref{inp:ssec-status}(b)
and~(c).
\end{theorem}

\begin{proof}
By \cite[Theorem~\ref{A-i4:thm-final}]{Annex}, whose four hypotheses are
Theorems~\ref{inp:thm-loss}, \ref{inp:thm-div}, \ref{inp:thm-counts} and
\ref{inp:thm-vwin} in the form of \cite[Propositions~\ref{A-i4:hyp-loss},
\ref{A-i4:hyp-div}, \ref{A-i4:hyp-S}, \ref{A-i4:hyp-selmer}]{Annex}.
\end{proof}



\section{Conclusion}\label{sec:conclusion}
\begin{theorem}[the Riemann hypothesis]\label{thm:RH}
Every nontrivial zero of the Riemann zeta function has real part $1/2$.
\end{theorem}
\begin{proof}
By Theorem~\ref{inp:thm-families}, Theorem~\ref{thm:alpoge-families}
holds. Let $\varepsilon_j\downarrow0$ with $\varepsilon_j<\varepsilon_1$. By
Corollary~\ref{cor:obligation} the hypothesis of
Proposition~\ref{prop:simultaneous-slope} holds with $C_j=1$ and
$K_j=e_D+e_L+3$. Proposition~\ref{prop:simultaneous-slope} then gives, for
every $1/2<\theta<1$, a constant $C_\theta>0$ and arbitrarily large $Y$ with $G(x)\le C_\theta x^\theta$ throughout
$[Y,Y^{3/2}]$, and Proposition~\ref{prop:growing-window} excludes the zeros
of $\zeta$ in $\Re s>\theta$. Since $\theta>1/2$ is arbitrary, $\zeta$ has no
zero in $\Re s>1/2$. If a nontrivial zero $\rho$ had $\Re\rho<1/2$, the
functional equation $\xi(s)=\xi(1-s)$ of the completed zeta function would
give the zero $1-\rho$, with $\Re(1-\rho)>1/2$. Hence every nontrivial zero
has real part $1/2$.
\end{proof}
The proof uses no result of \cite{P43}, \cite{P44} or \cite{Paper45}. In
the setting of the four rows described in the introduction, the theorem gives
the sign of the radial intersection form on every window: this is
Corollary~\ref{cor:all-windows} below, obtained through
Theorem~\ref{thm:all-support}, and it is there that \cite{Paper45} enters.

\subsection*{The sign of the form on every window}
\begin{theorem}\label{thm:all-support}
For $\zeta_{\mathbb Q}$ the Riemann hypothesis is equivalent both to
$W_T\ge0$ on every $\mathcal D_T$ and to $\lambda_{\rm A}(T)\ge0$ for all $T$.
Suppose that for each $1/2<\theta<1$ there exist $C_\theta>0$
and $J_\theta>1$ such that $G(x)\le C_\theta x^\theta$ on
$[Y,Y^{J_\theta}]$ for arbitrarily large $Y$. Then for every
$T>0$ and every $F\in\mathcal D_T$ one has
\[
 W_T(F)=\sum_{\rho}|\Phi_F(\rho)|^2\ge0 .
\]
Zeros are counted with multiplicities. Hence
$0\le\lambda^{\rm W}_T\le\lambda_{\rm A}(T)$ and
$I_\infty(\Gamma_F,\Gamma_F)\le0$ for smooth primitive tests.
\end{theorem}
\begin{proof}
Suppose first that a zero has real part $\beta>1/2$. Choose
$1/2<\theta<\min(1,\beta)$. Proposition~\ref{prop:growing-window}
shows that for each $C>0$ and each $J>1$, every sufficiently large
interval $[Y,Y^J]$ contains an $x$ with $G(x)>Cx^\theta$.
The hypothesis of the theorem rules this out with its own fixed
$C_\theta,J_\theta$, proving \eqref{eq:dirichlet-input}.

If a nontrivial zero $\rho$ lay to the left of the critical line,
the functional equation would give a zero $1-\rho$ to its right,
contrary to \eqref{eq:dirichlet-input}. Thus the prime-sum bound
implies the Riemann hypothesis.

For the forward implication of the equivalence, assume now that every
nontrivial zero lies on the critical line. For a smooth real test,
$1-\rho=\bar\rho$ and $\Phi_F(\bar\rho)=\overline{\Phi_F(\rho)}$.
The zero sum in \eqref{comparison} is therefore
$\sum_\rho |\Phi_F(\rho)|^2$.
For smooth real $U,V$, set $u_\rho=\Phi_U(\rho)$ and
$v_\rho=\Phi_V(\rho)$. Conjugation pairs the summands in
$\sum_\rho u_\rho\overline{v_\rho}$, so this absolutely convergent
sum is real. Expanding $|u_\rho+i v_\rho|^2$ and summing therefore
gives $W_T(U)+W_T(V)=W_T(U+iV)$ by complexification.
For $F\in\mathcal D_T$, take smooth $F_n$ converging in form norm
(Lemma~\ref{lem:form-density}(i)). Applying the identity to $F_n-F_m$,
for which $W_T(F_n-F_m)\to0$ by Lemma~\ref{lem:form-density}(ii),
makes $(\Phi_{F_n}(\rho))_\rho$ Cauchy in $\ell^2$, with one
coordinate for each occurrence of a zero. Since
$|e^{(\rho-1/2)t}|=1$, Cauchy--Schwarz gives
$|\Phi_{F_n-F}(\rho)|\le\sqrt{2T}\,\|F_n-F\|_2$.
This identifies every coordinate of the limit. Taking squared norms,
and using $W_T(F_n)\to W_T(F)$ (Lemma~\ref{lem:form-density}(ii)),
gives the formula for $F$. On $\mathcal P_T$ the polar term
vanishes, giving both inequalities. Conversely an off-line zero gives
a smooth primitive test with negative zero sum
\cite[Theorem~6.2]{Paper45}; both infima are negative on any window
containing its support.
\end{proof}
For the arithmetic criterion we take $J_\theta=3/2$:
Proposition~\ref{prop:simultaneous-slope} gives the bound at every point of
the windows $[W,W^{3/2}]$.

\begin{corollary}[the sign on every window]\label{cor:all-windows}
Assume Theorem~\ref{thm:alpoge-families}. Then for every $T>0$ and every
$F\in\mathcal D_T$,
\[
 W_T(F)=\sum_\rho|\Phi_F(\rho)|^2\ge0 ,
\]
so $0\le\lambda^{\rm W}_T\le\lambda_{\rm A}(T)$ for every $T>0$, and
$I_\infty(\Gamma_F,\Gamma_F)\le0$ for smooth primitive tests of arbitrary
support. In particular the positivity proved in Theorem~\ref{thm:log5W}
for $T\le\frac12\log5$, and in Theorem~\ref{thm:power-window} for
$T\le2.8022$ under (Z1)--(Z3), holds, in the form $\lambda^{\rm W}_T\ge0$,
on every window $(-T,T)$.
\end{corollary}
\begin{proof}
Corollary~\ref{cor:obligation} gives, for every $1/2<\theta<1$, a constant
$C_\theta$ and arbitrarily large $Y$ with $G(x)\le C_\theta x^\theta$ on
$[Y,Y^{3/2}]$. This is the hypothesis of Theorem~\ref{thm:all-support}
with $J_\theta=3/2$. The inequality itself comes from the identity \eqref{comparison}
\cite[Theorem~4.1]{Paper45}, through the third paragraph of the proof of
Theorem~\ref{thm:all-support}, applied once Theorem~\ref{thm:RH} is known.
\end{proof}

\subsection*{What is proved here and what is cited}
The column estimate required by Proposition~\ref{prop:simultaneous-slope} is
Corollary~\ref{cor:obligation}. Its proof in this paper is the deduction of
the mean bound from a column package (Theorem~\ref{thm:mean-bound}), proved
in full in \S\ref{subsec:assembly-T5}. The frame condition (P8) is a
condition of the package (Input~A): for the packages of
Theorem~\ref{thm:alpoge-families}(b)(5) it is supplied with the package, and
the verification of the hypotheses of the frame slope inequality for the
columns of a family is part of the formally verified
\cite[Proposition~\ref{A-i4:Lprop-fam}(b)]{Annex}. Section~\ref{sec:meanbound}
also gives a written proof of that inequality in general
(Theorem~\ref{thm:frameF}), including the Harder--Narasimhan formalism for
Hermitian lattices (\S\ref{subsec:HL}); it is not otherwise used in the proof
of Theorem~\ref{thm:RH}.

The existence of column packages with uniform constants on the same fields
and scales as the twist estimates (Theorem~\ref{thm:alpoge-families}) is due
to L.~Alpoge. His character-family argument and its formal verification
\cite{CharacterFamily} are the decisive input for global positivity: through
Theorem~\ref{thm:RH} and Corollary~\ref{cor:all-windows} they give the
Riemann hypothesis and $W_T\ge0$ on every window, and we owe this step to
him. In the companion paper \cite{Annex}, summarised in
Section~\ref{sec:inputs}, his theorem is derived from his loss bound,
divisor bound, level counts, existence theorem for families and windowed
bound for Selmer groups. There the reductions are proved in full, the
published theorems used in the written proofs are quoted with their
statements, and the steps that are known theorems re-proved in the
formalisation are identified. The new statements taken from the formal
development are stated exactly, as propositions marked \emph{formally
verified}; they are machine-checked in Lean~4 in \cite{CharacterFamily} and
are not reproved. The theorem on families therefore rests on that
formalisation. It has not yet been written out on paper in full:
\cite{Annex} is a first written account, and a complete written proof
remains to be given.

\paragraph{On what the proof rests.}
Theorem~\ref{thm:RH} and Corollary~\ref{cor:all-windows} rest on four
things. (a) The present paper and, for Corollary~\ref{cor:all-windows}
only, the author's earlier papers. Theorem~\ref{thm:RH} uses, from this
paper, the reduction to the prime sum
(Propositions~\ref{prop:growing-window}, \ref{prop:finite-to-window}
and~\ref{prop:simultaneous-slope}, Lemmas~\ref{lem:host-facts},
\ref{lem:character-identity} and~\ref{lem:prime-power}) and the mean bound
(Theorem~\ref{thm:mean-bound}, Corollary~\ref{cor:obligation}), together
with the functional equation of $\zeta$; it uses no result of \cite{P43},
\cite{P44} or \cite{Paper45}. These enter only in
Theorem~\ref{thm:all-support} and Corollary~\ref{cor:all-windows}: the
identity \eqref{comparison} \cite[Theorem~4.1]{Paper45}, which rests on the
constructions of \cite{P43} and \cite{P44}, turns the Riemann hypothesis
into $W_T(F)=\sum_\rho|\Phi_F(\rho)|^2\ge0$; and the criterion
\cite[Theorem~6.2]{Paper45} gives the converse implication of the
equivalence in Theorem~\ref{thm:all-support}, which neither
Theorem~\ref{thm:RH} nor Corollary~\ref{cor:all-windows} uses.
(b) Theorem~\ref{thm:alpoge-families}, which is not proved in this paper.
The existence of families with column packages is due to L.~Alpoge and
rests on his Lean formalisation \cite{CharacterFamily}; it has not yet been
written out on paper in full. The companion paper \cite{Annex}
(Theorem~\ref{inp:thm-families}) is a first written account: it derives
the theorem from the loss bound, the divisor bound, the level counts, the
windowed Selmer bound and the existence of families, proves the
reductions, identifies which steps are known theorems re-proved in the
formalisation, and states exactly the new statements taken from it.
(c) The published theorems used in the written proofs of this paper and of
\cite{Annex}, among them Dirichlet's theorem $L(1,\chi)\ne0$ (the
Harder--Narasimhan formalism for Hermitian lattices is not among them: it
is proved in \S\ref{subsec:HL}); they are citations for the written text.
Those of \cite{Annex} are listed in \S\ref{inp:ssec-status}(b), and
several of the exact references there are still to be confirmed. The
formal development does not assume them: it isolates each as a separate
proposition and proves it inside the development, and its final theorem
has no hypotheses. (d) The statements of \cite{Annex} that are taken from
L.~Alpoge's formal verification \cite{CharacterFamily} and are not
reproved in writing, listed in \S\ref{inp:ssec-status}(c). The formal
verification itself assumes nothing beyond Lean's standard axioms,
according to a reading of its sources; we have not yet completed an
independent compilation and axiom check. The passage from
the family estimates to properties (1)--(5) of
Theorem~\ref{thm:alpoge-families}(b) rests entirely on one of them, the
core link \cite[Proposition~\ref{A-i4:Lprop-fam}(c)]{Annex}, whose hypotheses
come from the other inputs and from \cite[Theorem~\ref{A-i4:thm-top}]{Annex};
the frame condition (P8) for the columns is part of the formally verified
\cite[Proposition~\ref{A-i4:Lprop-fam}(b)]{Annex}. Several of these
statements are paraphrases that use predicates of the formal development
which are named but not defined in writing, and for two of them
(\cite[Propositions~\ref{A-i1:prop-closed} and~\ref{A-i1:prop-linkG-res}]{Annex})
the written proofs in the source chapters are incomplete.

\begin{thebibliography}{99}
\bibitem{CharacterFamily} L.~Alpoge, Lean~4 formalisation of the character-family argument
for Dirichlet $L$-functions, source archive in fourteen parts, November 2026,
\url{https://alpo.ge/lean.tar.gz.part01} (first part).

\bibitem{BoasKac1945} R.~P.~Boas and M.~Kac, \emph{Inequalities for Fourier transforms of positive
functions}, Duke Math.~J. \textbf{12} (1945), 189--206.

\bibitem{Bombieri2000} E.~Bombieri, \emph{Remarks on Weil's quadratic functional in the theory of
prime numbers, I}, Atti Accad.~Naz.~Lincei Cl.~Sci.~Fis.~Mat.~Natur.~Rend.~Lincei (9) Mat.~Appl.
\textbf{11} (2000), no.~3, 183--233.

\bibitem{BorelHC1962} A.~Borel and Harish-Chandra, \emph{Arithmetic subgroups of algebraic groups},
Ann.\ of Math.\ (2) \textbf{75} (1962), 485--535.

\bibitem{Bost1996} J.-B.~Bost, \emph{P\'eriodes et isog\'enies des vari\'et\'es ab\'eliennes sur
les corps de nombres (d'apr\`es D.~Masser et G.~W\"ustholz)}, S\'eminaire Bourbaki, Exp.~795,
Ast\'erisque \textbf{237} (1996), 115--161.

\bibitem{ConnesConsani2021} A.~Connes and C.~Consani, \emph{Weil positivity and trace formula, the archimedean place}, Selecta Math.\ (N.S.) \textbf{27} (2021), no.~4, Paper No.~77; arXiv:2006.13771.

\bibitem{CvS2025} A.~Connes and W.~D.~van Suijlekom, \emph{Quadratic forms, real zeros and echoes of the spectral action}, arXiv:2511.23257 (2025).

\bibitem{Davenport} H.~Davenport, \emph{Multiplicative Number Theory}, 3rd ed., revised by
H.~L.~Montgomery, Graduate Texts in Mathematics \textbf{74}, Springer, New York, 2000.

\bibitem{DLMF} NIST Digital Library of Mathematical Functions, \url{https://dlmf.nist.gov}, \S5.9.

\bibitem{Grayson1984} D.~R.~Grayson, \emph{Reduction theory using semistability}, Comment.\
Math.\ Helv.\ \textbf{59} (1984), 600--634.

\bibitem{Groskin2026} A.~Groskin, \emph{A finite Guinand--Weil dictionary and archimedean tail
order for the truncated Weil quadratic form}, arXiv:2607.02828 (2026).

\bibitem{HLP} G.~H.~Hardy, J.~E.~Littlewood and G.~P\'olya, \emph{Inequalities}, 2nd ed., Cambridge Univ.~Press, 1952.

\bibitem{Haynsworth1968} E.~V.~Haynsworth, \emph{Determination of the inertia of a partitioned Hermitian
matrix}, Linear Algebra Appl.\ \textbf{1} (1968), 73--81.

\bibitem{IK} H.~Iwaniec and E.~Kowalski, \emph{Analytic Number Theory}, Amer.\ Math.\ Soc.
Colloq.\ Publ.\ \textbf{53}, Providence, RI, 2004.

\bibitem{Johansson2017} F.~Johansson, \emph{Arb: efficient arbitrary-precision midpoint-radius interval
arithmetic}, IEEE Trans.\ Comput.\ \textbf{66} (2017), 1281--1292.

\bibitem{Kim2026} T.~Kim, Y.~Hong, M.~Kim, S.~Choi, J.~Jang and M.~Kim, \emph{A numerical realization of Suzuki's Weil-quadratic-form operator: the archimedean spectral law, its universality, and an operator form of Weil's positivity criterion}, arXiv:2607.24830 (2026).

\bibitem{Krein1940} M.~G.~Kre\u{\i}n, \emph{Sur le probl\`eme du prolongement des fonctions
hermitiennes positives et continues}, Dokl.~Akad.~Nauk SSSR \textbf{26} (1940), 17--22.

\bibitem{Liu2026} V.~Liu, \emph{Certified Weil positivity beyond the unit window: source-exact
block-Schur and tail-compensation bounds for the Riemann zeta function}, manuscript dated 14 Sep 2026,
alphaXiv preprint.

\bibitem{MoriUpper} R.~Mori, \emph{Improved upper bounds for windowed Weil forms with exponential factor
$e^{-4\pi\mu}$}, Zenodo (2026), v3, doi:\nolinkurl{10.5281/zenodo.23363767}.

\bibitem{PlattTrudgian2021} D.~J.~Platt and T.~S.~Trudgian, \emph{The Riemann hypothesis is true up to
$3\cdot10^{12}$}, Bull.\ London Math.\ Soc.\ \textbf{53} (2021), 792--797.

\bibitem{Rockafellar} R.~T.~Rockafellar, \emph{Convex Analysis}, Princeton Univ.~Press, 1970.

\bibitem{Rosser1941} J.~B.~Rosser, \emph{Explicit bounds for some functions of prime numbers}, Amer.\ J.\
Math.\ \textbf{63} (1941), 211--232.

\bibitem{RosserYoheSchoenfeld1968} J.~B.~Rosser, J.~M.~Yohe and L.~Schoenfeld, \emph{Rigorous computation and
the zeros of the Riemann zeta-function}, Proc.\ IFIP Congress 1968, 70--76.

\bibitem{SSV2012} R.~L.~Schilling, R.~Song and Z.~Vondra\v{c}ek, \emph{Bernstein Functions:
Theory and Applications}, 2nd ed., de Gruyter, 2012.

\bibitem{Sasvari1994} Z.~Sasv\'ari, \emph{Positive Definite and Definitizable Functions},
Akademie Verlag, Berlin, 1994.

\bibitem{SongVondracek2003} R.~Song and Z.~Vondra\v{c}ek, \emph{Potential theory of subordinate
killed Brownian motion in a domain}, Probab.~Theory Related Fields \textbf{125} (2003), 578--592.

\bibitem{Stuhler1976} U.~Stuhler, \emph{Eine Bemerkung zur Reduktionstheorie quadratischer
Formen}, Arch.\ Math.\ (Basel) \textbf{27} (1976), 604--610.

\bibitem{Trefethen2013} L.~N.~Trefethen, \emph{Approximation Theory and Approximation Practice},
SIAM, Philadelphia, 2013.

\bibitem{P43} D.~A.~Trejo Pizzo,
\emph{A correspondence algebra and the Weil distribution on semilocal idele
class groups}, preprint, 2026.

\bibitem{P44} D.~A.~Trejo Pizzo,
\emph{An arithmetic Frobenius family on semilocal idele class groups},
preprint, 2026.

\bibitem{Paper45} D.~A.~Trejo Pizzo, \emph{An explicit formula for the trace on semilocal idele class groups}, preprint, 2026.

\bibitem{Annex} D.~A.~Trejo Pizzo, \emph{Column packages for character families: the
arithmetic inputs for the sign of the radial intersection form}, companion paper, preprint, 2026.

\bibitem{Weil52} A.~Weil, \emph{Sur les ``formules explicites'' de la th\'eorie des nombres
premiers}, Comm.~S\'em.~Math.~Univ.~Lund (1952), 252--265.

\bibitem{Weil1965} A.~Weil, \emph{Sur la formule de Siegel dans la th\'eorie des groupes
classiques}, Acta Math.\ \textbf{113} (1965), 1--87.

\bibitem{Yoshida1992} H.~Yoshida, \emph{On Hermitian forms attached to zeta functions}, in:
\emph{Zeta Functions in Geometry}, Adv.~Stud.~Pure Math.~\textbf{21}, Kinokuniya, Tokyo, 1992, 281--325.

\bibitem{Zhu2026} X.~Zhu, \emph{Weil positivity in compact windows: a finite reduction, certified
two-sided bounds, and a Landau--Widom decay law}, arXiv:2608.24827 (2026).
\end{thebibliography}
\end{document}
